\input{preamble}

% OK, start here.
%
\begin{document}

\title{Compléments sur les morphismes d'espaces}


\maketitle

\phantomsection
\label{section-phantom}

\tableofcontents

\section{Introduction}
\label{section-introduction}

\noindent
Dans ce chapitre, nous poursuivons l'étude des propriétés des morphismes
d'espaces algébriques. Une référence fondamentale est \cite{Kn}.





\section{Conventions}
\label{section-conventions}

\noindent
On suppose une fois pour toutes que tous les schémas appartiennent à
un grand site fppf $\Sch_{fppf}$. Tous les anneaux $A$ considérés
ont en outre la propriété que $\Spec(A)$ est (isomorphe) à un
objet de ce grand site.

\medskip\noindent
Soit $S$ un schéma et soit $X$ un espace algébrique sur $S$.
Dans ce chapitre et le suivant, nous noterons $X \times_S X$
le produit de $X$ par lui-même (dans la catégorie des espaces
algébriques sur $S$), au lieu de $X \times X$.








\section{Morphismes radiciels}
\label{section-radicial}

\noindent
Il s'avère qu'un morphisme radiciel n'est pas la même chose qu'un
morphisme universellement injectif, contrairement à ce qui se passe pour les
morphismes de schémas. C'est en fait une condition un peu plus forte.

\begin{definition}
\label{definition-radicial}
Soit $S$ un schéma. Soit $f : X \to Y$ un morphisme d'espaces algébriques
sur $S$. On dit que $f$ est {\it radiciel} si, pour tout morphisme
$\Spec(K) \to Y$, où $K$ est un corps, le réduit
$(\Spec(K) \times_Y X)_{red}$ est soit vide, soit
représentable par le spectre d'une extension de corps purement inséparable de $K$.
\end{definition}

\begin{lemma}
\label{lemma-radicial-implies-universally-injective}
Tout morphisme radiciel d'espaces algébriques est universellement injectif.
\end{lemma}

\begin{proof}
Soit $S$ un schéma. Soit $f : X \to Y$ un morphisme radiciel
d'espaces algébriques sur $S$.
Il résulte clairement de la définition que, pour un morphisme donné
$\Spec(K) \to Y$, il existe au plus un relèvement de ce morphisme
en un morphisme vers $X$. Nous concluons donc que $f$ est universellement
injectif par
Morphismes d'espaces,
Lemme \ref{spaces-morphisms-lemma-universally-injective}.
\end{proof}

\begin{example}
\label{example-universally-injective-not-radicial}
Les notions de morphisme universellement injectif et de morphisme radiciel ne sont plus équivalentes.
Par exemple, le morphisme
$$
X = [\Spec(\overline{\mathbf{Q}})/
\text{Gal}(\overline{\mathbf{Q}}/\mathbf{Q})]
\longrightarrow
S = \Spec(\mathbf{Q})
$$
de
Espaces, Exemple \ref{spaces-example-Qbar}
est universellement injectif, mais n'est pas radiciel au sens ci-dessus.
\end{example}

\noindent
Il arrive néanmoins souvent que la réciproque soit vraie.

\begin{lemma}
\label{lemma-when-universally-injective-radicial}
Soit $S$ un schéma. Soit $f : X \to Y$ un morphisme universellement injectif
d'espaces algébriques sur $S$.
\begin{enumerate}
\item Si $f$ est décent, alors $f$ est radiciel.
\item Si $f$ est quasi-séparé, alors $f$ est radiciel.
\item Si $f$ est localement séparé, alors $f$ est radiciel.
\end{enumerate}
\end{lemma}

\begin{proof}
Soit $\mathcal{P}$ une propriété des morphismes d'espaces algébriques
qui est stable par changement de base et par composition, et que possèdent les
immersions fermées. Supposons que $f : X \to Y$ possède $\mathcal{P}$ et
soit universellement injectif. Alors, dans la situation de la
Définition \ref{definition-radicial},
le morphisme $(\Spec(K) \times_Y X)_{red} \to \Spec(K)$
est universellement injectif et possède $\mathcal{P}$. Cela ramène le
problème consistant à démontrer
$$
\mathcal{P} + \text{universally injective}
\Rightarrow
\text{radicial}
$$
au problème consistant à démontrer que tout espace algébrique réduit non vide $X$
sur un corps, dont le morphisme structural $X \to \Spec(K)$ est universellement
injectif et possède $\mathcal{P}$, est représentable par le spectre d'un corps.
En effet, $X \to \Spec(K)$ sera alors un morphisme de schémas et
nous concluons par l'équivalence entre radiciel et universellement injectif pour les
morphismes de schémas, voir
Morphismes, Lemme \ref{morphisms-lemma-universally-injective}.

\medskip\noindent
Montrons (1). Supposons que $f$ soit décent et universellement injectif. Par
Espaces décents,
les Lemmes \ref{decent-spaces-lemma-base-change-relative-conditions},
\ref{decent-spaces-lemma-composition-relative-conditions}, et
\ref{decent-spaces-lemma-properties-trivial-implications}
(pour voir qu'une immersion est décente), l'argument du
premier paragraphe s'applique.
Soit $X$ un espace algébrique décent, réduit et non vide,
universellement injectif sur un corps $K$. En particulier, $|X|$
est réduit à un point. Par
Espaces décents, Lemme \ref{decent-spaces-lemma-when-field},
nous concluons que $X \cong \Spec(L)$ pour une certaine extension
$K \subset L$, comme voulu.

\medskip\noindent
Un morphisme quasi-séparé est décent, voir
Espaces décents,
Lemme \ref{decent-spaces-lemma-properties-trivial-implications}.
Ainsi, (1) implique (2).

\medskip\noindent
Montrons (3).
Rappelons que les axiomes de séparation sont stables par changement de base
et par composition, et que les immersions fermées sont séparées, voir
Morphismes d'espaces,
les Lemmes \ref{spaces-morphisms-lemma-base-change-separated},
\ref{spaces-morphisms-lemma-composition-separated}, et
\ref{spaces-morphisms-lemma-immersions-monomorphisms}.
L'argument du premier paragraphe de la démonstration s'applique donc.
Soit $X$ un espace algébrique réduit, universellement injectif et
localement séparé sur un corps $K$.
En particulier, $|X|$ est réduit à un point, donc $X$ est quasi-compact, voir
Propriétés des espaces, Lemme \ref{spaces-properties-lemma-quasi-compact-space}.
On peut trouver un morphisme étale surjectif $U \to X$ avec $U$ affine, voir
Propriétés des espaces,
Lemme \ref{spaces-properties-lemma-quasi-compact-affine-cover}.
Considérons le morphisme de schémas
$$
j :
U \times_X U
\longrightarrow
U \times_{\Spec(K)} U
$$
Comme $X \to \Spec(K)$ est universellement injectif, $j$ est surjectif,
et comme $X \to \Spec(K)$ est localement séparé, $j$ est une immersion.
Toute immersion surjective est une immersion fermée, voir
Schémas, Lemme \ref{schemes-lemma-immersion-when-closed}.
Ainsi, $R = U \times_X U$ est affine en tant que sous-schéma fermé d'un schéma affine.
En particulier, $R$ est quasi-compact.
Il s'ensuit que $X = U/R$ est quasi-séparé, et le résultat découle de (2).
\end{proof}

\begin{remark}
\label{remark-weakly-radicial}
Soit $X \to Y$ un morphisme d'espaces algébriques.
Pour certaines applications (des morphismes radiciels),
il suffit d'exiger que, pour tout
$\Spec(K) \to Y$, où $K$ est un corps,
\begin{enumerate}
\item l'espace $|\Spec(K) \times_Y X|$ soit réduit à un point,
\item il existe un monomorphisme
$\Spec(L) \to \Spec(K) \times_Y X$, et
\item $K \subset L$ soit purement inséparable.
\end{enumerate}
Au besoin, nous appellerons plus loin un tel morphisme {\it faiblement radiciel}.
Par exemple, si $X \to Y$ est un morphisme faiblement radiciel surjectif,
alors $X(k) \to Y(k)$ est surjectif pour tout corps algébriquement clos $k$.
Notons que le changement de base
$X_{\overline{\mathbf{Q}}} \to \Spec(\overline{\mathbf{Q}})$
du morphisme de
l'Exemple \ref{example-universally-injective-not-radicial}
est faiblement radiciel, mais non radiciel. L'analogue du
Lemme \ref{lemma-when-universally-injective-radicial}
énonce que, si $X \to Y$ possède la propriété ($\beta$) et est universellement
injectif, alors il est faiblement radiciel (démonstration omise).
\end{remark}

\begin{lemma}
\label{lemma-check-universally-injective}
Soit $S$ un schéma. Soit $f : X \to Y$ un morphisme d'espaces
algébriques sur $S$. Supposons que
\begin{enumerate}
\item $f$ soit localement de type fini,
\item pour tout morphisme étale $V \to Y$, l'application $|X \times_Y V| \to |V|$
soit injective.
\end{enumerate}
Alors $f$ est universellement injectif.
\end{lemma}

\begin{proof}
La question est locale pour la topologie étale sur $Y$ par
Morphismes d'espaces, Lemme
\ref{spaces-morphisms-lemma-universally-injective-local}.
On peut donc supposer que $Y$ est un schéma.
Alors $Y$ est en particulier décent et, par Espaces décents, Lemme
\ref{decent-spaces-lemma-conditions-on-point-in-fibre-and-qf}
nous voyons que $f$ est localement quasi-fini.
Soit $y \in Y$ un point et soit $X_y$ la fibre schématique.
Supposons $X_y$ non vide. Par Espaces sur les corps, Lemme
\ref{spaces-over-fields-lemma-locally-quasi-finite-over-field}
nous voyons que $X_y$ est un schéma localement quasi-fini sur
$\kappa(y)$. Comme $|X_y| \subset |X|$ est la fibre de $|X| \to |Y|$
au-dessus de $y$, nous voyons que $X_y$ possède un unique point $x$. Il en va de même
de $X_y \times_{\Spec(\kappa(y))} \Spec(k)$ pour toute
extension séparable finie $k/\kappa(y)$,
car on peut réaliser $k$ comme le corps résiduel en un point
situé au-dessus de $y$ dans un schéma étale sur $Y$,
voir Compléments sur les morphismes, Lemme
\ref{more-morphisms-lemma-realize-prescribed-residue-field-extension-etale}.
Ainsi, $X_y$ est géométriquement connexe, voir
Variétés, Lemme \ref{varieties-lemma-characterize-geometrically-disconnected}.
Cela implique que l'extension finie $\kappa(x)/\kappa(y)$
est purement inséparable.

\medskip\noindent
Nous concluons (dans le cas où $Y$ est un schéma)
que, pour tout $y \in Y$, soit la fibre $X_y$ est vide,
soit $(X_y)_{red} = \Spec(\kappa(x))$ avec
$\kappa(y) \subset \kappa(x)$ purement inséparable.
Ainsi, $f$ est radiciel (certains détails sont omis), donc universellement injectif par le
Lemme \ref{lemma-radicial-implies-universally-injective}.
\end{proof}




\section{Monomorphismes}
\label{section-monomorphisms}

\noindent
Cette section fait suite à
Morphismes d'espaces, Section \ref{spaces-morphisms-section-monomorphisms}.
Nous aimerions savoir si tout monomorphisme d'espaces
algébriques est représentable. Si vous pouvez le démontrer ou disposez d'un
contre-exemple, veuillez écrire à
\href{mailto:stacks.project@gmail.com}{stacks.project@gmail.com}.
Pour le moment, cela est établi dans les cas suivants :
\begin{enumerate}
\item pour les monomorphismes qui sont localement de type fini
(plus généralement, tout morphisme séparé et localement quasi-fini
est représentable d'après Morphismes d'espaces, Lemme
\ref{spaces-morphisms-lemma-locally-quasi-finite-separated-representable}
et un monomorphisme localement de type fini est
localement quasi-fini par Morphismes d'espaces, Lemme
\ref{spaces-morphisms-lemma-monomorphism-loc-finite-type-loc-quasi-finite}),
\item si le but est une somme disjointe de spectres d'anneaux locaux
de dimension zéro (Espaces décents, Lemme
\ref{decent-spaces-lemma-monomorphism-toward-disjoint-union-dim-0-rings}), et
\item pour les monomorphismes plats (voir ci-dessous).
\end{enumerate}

\begin{lemma}[David Rydh]
\label{lemma-flat-case}
Tout monomorphisme plat d'espaces algébriques est représentable par des schémas.
\end{lemma}

\begin{proof}
Soit $f : X \to Y$ un monomorphisme plat d'espaces algébriques.
Pour démontrer que $f$ est représentable, il faut montrer que
$X \times_Y V$ est un schéma pour tout schéma $V$ muni d'un morphisme vers $Y$.
Comme la propriété d'être un schéma est locale (Propriétés des espaces, 
Lemme \ref{spaces-properties-lemma-subscheme}), nous pouvons
supposer $V$ affine. Nous pouvons donc supposer que $Y = \Spec(B)$
est un schéma affine. Ensuite, nous pouvons supposer $X$ quasi-compact
en remplaçant $X$ par un ouvert quasi-compact. L'espace $X$ est
séparé, car $X \to X \times_{\Spec(B)} X$ est un isomorphisme.
En appliquant Limites d'espaces, Lemme \ref{spaces-limits-lemma-enough-local},
nous nous ramenons au cas où $B$ est local, où $X \to \Spec(B)$ est un
monomorphisme plat et où
il existe un point $x \in X$ dont l'image est le point fermé de $\Spec(B)$.
Alors $X \to \Spec(B)$ est surjectif, car les générisations
se relèvent le long des morphismes plats d'espaces algébriques séparés, voir
Espaces décents, Lemme \ref{decent-spaces-lemma-generalizations-lift-flat}.
Ainsi, $\{X \to \Spec(B)\}$ est un recouvrement fpqc.
Alors $X \to \Spec(B)$ est un morphisme qui devient un isomorphisme
après changement de base par $X \to \Spec(B)$. C'est donc un isomorphisme par
descente fpqc, voir Descente sur les espaces, Lemme
\ref{spaces-descent-lemma-descending-property-isomorphism}.
\end{proof}

\noindent
Le résultat suivant est (en un certain sens) une variante du lemme précédent.

\begin{lemma}
\label{lemma-ui-case}
Soit $S$ un schéma. Soit $f : X \to Y$ un monomorphisme quasi-compact
d'espaces algébriques tel que, pour tout $T \to Y$, le morphisme
$$
\mathcal{O}_T \to f_{T,*}\mathcal{O}_{X \times_Y T}
$$
soit injectif. Alors $f$ est un isomorphisme (et est donc représentable par des schémas).
\end{lemma}

\begin{proof}
La question est locale pour la topologie étale sur $Y$ ; nous pouvons donc supposer que $Y = \Spec(A)$
est affine. Alors $X$ est quasi-compact et nous pouvons choisir un schéma affine
$U = \Spec(B)$ et un morphisme étale surjectif $U \to X$
(Propriétés des espaces, Lemme
\ref{spaces-properties-lemma-quasi-compact-affine-cover}).
Notons que $U \times_X U = \Spec(B \otimes_A B)$. La catégorie des
$\mathcal{O}_X$-modules quasi-cohérents est donc équivalente à la
catégorie $DD_{B/A}$ des données de descente de modules pour $A \to B$.
Voir Propriétés des espaces, Proposition
\ref{spaces-properties-proposition-quasi-coherent},
Descente, Définition \ref{descent-definition-descent-datum-modules}, et
Descente, Sous-section \ref{descent-subsection-descent-modules-morphisms}.
D'autre part,
$$
A \to B
$$
est un homomorphisme d'anneaux universellement injectif. En effet, pour tout
$A$-module $M$, l'hypothèse du lemme montre que $A \oplus M \to B \otimes_A (A \oplus M)$
est injectif. Ainsi,
$DD_{B/A}$ est équivalente à la catégorie des $A$-modules par
Descente, Théorème \ref{descent-theorem-descent}. L'image inverse par
$f : X \to \Spec(A)$ détermine donc une équivalence entre les catégories de
modules quasi-cohérents. En particulier, $f^*$ est exact sur les
modules quasi-cohérents, et nous voyons que $f$ est plat
(un petit détail est omis). De plus, il est clair que $f$ est surjectif
(par exemple parce que $\Spec(B) \to \Spec(A)$ est surjectif).
Ainsi, $\{X \to \Spec(A)\}$ est un recouvrement fpqc.
Alors $X \to \Spec(A)$ est un morphisme qui devient un isomorphisme
après changement de base par $X \to \Spec(A)$. C'est donc un isomorphisme par
descente fpqc, voir Descente sur les espaces, Lemme
\ref{spaces-descent-lemma-descending-property-isomorphism}.
\end{proof}

\begin{lemma}
\label{lemma-flat-surjective-monomorphism}
Tout monomorphisme quasi-compact, plat et surjectif d'espaces algébriques
est un isomorphisme.
\end{lemma}

\begin{proof}
Un tel morphisme satisfait aux hypothèses du Lemme \ref{lemma-ui-case}.
\end{proof}







\section{Faisceau conormal d'une immersion}
\label{section-conormal-sheaf}

\noindent
Soit $S$ un schéma. Soit $i : Z \to X$ une immersion fermée d'espaces
algébriques sur $S$. Soit $\mathcal{I} \subset \mathcal{O}_X$ le faisceau
quasi-cohérent d'idéaux correspondant, voir
Morphismes d'espaces,
Lemme \ref{spaces-morphisms-lemma-closed-immersion-ideals}.
Considérons la suite exacte courte
$$
0 \to \mathcal{I}^2 \to \mathcal{I} \to \mathcal{I}/\mathcal{I}^2 \to 0
$$
de faisceaux quasi-cohérents sur $X$. Comme le faisceau $\mathcal{I}/\mathcal{I}^2$
est annulé par $\mathcal{I}$, il correspond à un faisceau sur $Z$ par
Morphismes d'espaces, Lemme \ref{spaces-morphisms-lemma-i-star-equivalence}.
Ce $\mathcal{O}_Z$-module quasi-cohérent est le
{\it faisceau conormal à $Z$ dans $X$} et se note souvent
$\mathcal{I}/\mathcal{I}^2$, par l'abus de notation signalé dans
Morphismes d'espaces,
Section \ref{spaces-morphisms-section-closed-immersions-quasi-coherent}.

\medskip\noindent
Lorsque $i : Z \to X$ est une immersion (localement fermée), on définit le
faisceau conormal de $i$ comme le faisceau conormal de
l'immersion fermée $i : Z \to X \setminus \partial Z$, voir
Morphismes d'espaces, Remarque \ref{spaces-morphisms-remark-immersion}.
On le note souvent
$\mathcal{I}/\mathcal{I}^2$, où $\mathcal{I}$ est le faisceau d'idéaux
de l'immersion fermée $i : Z \to X \setminus \partial Z$.

\begin{definition}
\label{definition-conormal-sheaf}
Soit $i : Z \to X$ une immersion. Le {\it faisceau conormal
$\mathcal{C}_{Z/X}$ à $Z$ dans $X$}, ou le {\it faisceau conormal de $i$},
est le $\mathcal{O}_Z$-module quasi-cohérent $\mathcal{I}/\mathcal{I}^2$
décrit ci-dessus.
\end{definition}

\noindent
Dans \cite[IV Définition 16.1.2]{EGA}, ce faisceau est noté
$\mathcal{N}_{Z/X}$. Nous ne suivrons pas cette convention, car nous souhaitons
réserver la notation $\mathcal{N}_{Z/X}$
au {\it faisceau normal de l'immersion}. Il est défini par
$$
\mathcal{N}_{Z/X} =
\SheafHom_{\mathcal{O}_Z}(\mathcal{C}_{Z/X}, \mathcal{O}_Z) =
\SheafHom_{\mathcal{O}_Z}(\mathcal{I}/\mathcal{I}^2, \mathcal{O}_Z)
$$
pourvu que le faisceau conormal soit de présentation finie (sinon le
faisceau normal pourrait même ne pas être quasi-cohérent). Nous reviendrons plus loin au
faisceau normal (insérer ici une référence future).

\begin{lemma}
\label{lemma-etale-conormal}
Soit $S$ un schéma. Soit $i : Z \to X$ une immersion.
Soit $\varphi : U \to X$ un morphisme étale, où $U$ est un schéma.
Posons $Z_U = U \times_X Z$, qui est un sous-schéma localement fermé de $U$.
Alors
$$
\mathcal{C}_{Z/X}|_{Z_U} = \mathcal{C}_{Z_U/U}
$$
canoniquement et fonctoriellement en $U$.
\end{lemma}

\begin{proof}
Soit $T \subset X$ un sous-espace fermé tel que $i$ définisse une immersion
fermée dans $X \setminus T$.
Soit $\mathcal{I}$ le faisceau quasi-cohérent d'idéaux sur
$X \setminus T$ définissant $Z$. Le lemme affirme alors simplement que
$\mathcal{I}|_{U \setminus \varphi^{-1}(T)}$ est le faisceau d'idéaux de
l'immersion $Z_U \to U \setminus \varphi^{-1}(T)$.
Cela résulte clairement de la construction de $\mathcal{I}$ dans
Morphismes d'espaces, Lemme \ref{spaces-morphisms-lemma-closed-immersion-ideals}.
\end{proof}

\begin{lemma}
\label{lemma-conormal-functorial}
Soit $S$ un schéma. Soit
$$
\xymatrix{
Z \ar[r]_i \ar[d]_f & X \ar[d]^g \\
Z' \ar[r]^{i'} & X'
}
$$
un diagramme commutatif d'espaces algébriques sur $S$.
Supposons que $i$ et $i'$ soient des immersions. Il existe un morphisme canonique
de $\mathcal{O}_Z$-modules
$$
f^*\mathcal{C}_{Z'/X'}
\longrightarrow
\mathcal{C}_{Z/X}
$$
\end{lemma}

\begin{proof}
Choisissons d'abord des sous-espaces ouverts $U' \subset X'$ et $U \subset X$ tels que
$g(U) \subset U'$ et que $i(Z) \subset U$ et $i(Z') \subset U'$
soient fermés (démonstration de l'existence omise). En remplaçant $X$ par $U$ et $X'$ par
$U'$, on peut supposer que $i$ et $i'$ sont des immersions fermées.
Soient $\mathcal{I}' \subset \mathcal{O}_{X'}$ et
$\mathcal{I} \subset \mathcal{O}_X$ les faisceaux quasi-cohérents
d'idéaux associés à $i'$ et $i$, voir
Morphismes d'espaces, Lemme \ref{spaces-morphisms-lemma-closed-immersion-ideals}.
Considérons la composée
$$
g^{-1}\mathcal{I}' \to g^{-1}\mathcal{O}_{X'}
\xrightarrow{g^\sharp} \mathcal{O}_X \to
\mathcal{O}_X/\mathcal{I} = i_*\mathcal{O}_Z
$$
Comme $g(i(Z)) \subset Z'$, nous concluons que cette composée est nulle (voir
l'énoncé sur les factorisations dans
Morphismes d'espaces,
Lemme \ref{spaces-morphisms-lemma-closed-immersion-ideals}).
On obtient ainsi un diagramme commutatif
$$
\xymatrix{
0 \ar[r] &
\mathcal{I} \ar[r] &
\mathcal{O}_X \ar[r] &
i_*\mathcal{O}_Z \ar[r] &
0 \\
0 \ar[r] &
g^{-1}\mathcal{I}' \ar[r] \ar[u] &
g^{-1}\mathcal{O}_{X'} \ar[r] \ar[u] &
g^{-1}i'_*\mathcal{O}_{Z'} \ar[r] \ar[u] &
0
}
$$
La ligne inférieure est exacte puisque $g^{-1}$ est un foncteur exact.
Par exactitude, on voit aussi que
$(g^{-1}\mathcal{I}')^2 = g^{-1}((\mathcal{I}')^2)$.
Le diagramme induit donc un morphisme
$g^{-1}(\mathcal{I}'/(\mathcal{I}')^2) \to \mathcal{I}/\mathcal{I}^2$.
En prenant l'image inverse (au moyen de $i^{-1}$, par exemple) sur $Z$, on obtient
$i^{-1}g^{-1}(\mathcal{I}'/(\mathcal{I}')^2) \to \mathcal{C}_{Z/X}$.
Comme $i^{-1}g^{-1} = f^{-1}(i')^{-1}$, cela donne un morphisme
$f^{-1}\mathcal{C}_{Z'/X'} \to \mathcal{C}_{Z/X}$, qui induit
le morphisme voulu.
\end{proof}

\begin{lemma}
\label{lemma-conormal-functorial-more}
Soit $S$ un schéma. Le faisceau conormal de la
Définition \ref{definition-conormal-sheaf} et sa fonctorialité donnée par le
Lemme \ref{lemma-conormal-functorial}
possèdent les propriétés suivantes :
\begin{enumerate}
\item Si $Z \to X$ est une immersion de schémas sur $S$, alors le faisceau
conormal coïncide avec celui de
Morphismes, Définition \ref{morphisms-definition-conormal-sheaf}.
\item Si, dans le
Lemme \ref{lemma-conormal-functorial},
tous les espaces sont des schémas, alors le morphisme
$f^*\mathcal{C}_{Z'/X'} \to \mathcal{C}_{Z/X}$ coïncide
avec celui construit dans
Morphismes, Lemme \ref{morphisms-lemma-conormal-functorial}.
\item Étant donné un diagramme commutatif
$$
\xymatrix{
Z \ar[r]_i \ar[d]_f & X \ar[d]^g \\
Z' \ar[r]^{i'} \ar[d]_{f'} & X' \ar[d]^{g'} \\
Z'' \ar[r]^{i''} & X''
}
$$
le morphisme $(f' \circ f)^*\mathcal{C}_{Z''/X''} \to \mathcal{C}_{Z/X}$
est la composée de
$f^*\mathcal{C}_{Z'/X'} \to \mathcal{C}_{Z/X}$
avec l'image inverse par $f$ de
$(f')^*\mathcal{C}_{Z''/X''} \to \mathcal{C}_{Z'/X'}$
\end{enumerate}
\end{lemma}

\begin{proof}
La démonstration est omise. Notons que l'assertion (1) est un cas particulier du
Lemme \ref{lemma-etale-conormal}.
\end{proof}

\begin{lemma}
\label{lemma-conormal-functorial-flat}
Soit $S$ un schéma. Soit
$$
\xymatrix{
Z \ar[r]_i \ar[d]_f & X \ar[d]^g \\
Z' \ar[r]^{i'} & X'
}
$$
un diagramme cartésien d'espaces algébriques sur $S$. Supposons
que $i$ et $i'$ soient des immersions. Alors le morphisme canonique
$f^*\mathcal{C}_{Z'/X'} \to \mathcal{C}_{Z/X}$ du
Lemme \ref{lemma-conormal-functorial}
est surjectif. Si $g$ est plat, alors c'est un isomorphisme.
\end{lemma}

\begin{proof}
Choisissons un diagramme commutatif
$$
\xymatrix{
U \ar[r] \ar[d] & X \ar[d] \\
U' \ar[r] & X'
}
$$
où $U$ et $U'$ sont des schémas et les flèches horizontales sont surjectives
et étales, voir
Espaces, Lemme \ref{spaces-lemma-lift-morphism-presentations}.
Alors, en utilisant les
Lemmes \ref{lemma-etale-conormal} et \ref{lemma-conormal-functorial-more},
nous voyons que la question se ramène au cas d'un morphisme de schémas.
Dans le cas des schémas, c'est
Morphismes, Lemme \ref{morphisms-lemma-conormal-functorial-flat}.
\end{proof}

\begin{lemma}
\label{lemma-transitivity-conormal}
Soit $S$ un schéma.
Soient $Z \to Y \to X$ des immersions d'espaces algébriques.
Alors il existe une suite exacte canonique
$$
i^*\mathcal{C}_{Y/X} \to
\mathcal{C}_{Z/X} \to
\mathcal{C}_{Z/Y} \to 0
$$
où les morphismes proviennent du
Lemme \ref{lemma-conormal-functorial}
et où $i : Z \to Y$ est le premier morphisme.
\end{lemma}

\begin{proof}
Soit $U$ un schéma et soit $U \to X$ un morphisme étale surjectif. Par les
Lemmes \ref{lemma-etale-conormal} et \ref{lemma-conormal-functorial-more},
l'exactitude de la suite se traduit immédiatement par
l'exactitude de la suite correspondante pour les immersions de schémas
$Z \times_X U \to Y \times_X U \to U$. Le lemme découle donc de
Morphismes, Lemme \ref{morphisms-lemma-transitivity-conormal}.
\end{proof}








\section{Le cône normal d'une immersion}
\label{section-normal-cone}

\noindent
Soit $S$ un schéma. Soit $i : Z \to X$ une immersion fermée d'espaces
algébriques sur $S$. Soit $\mathcal{I} \subset \mathcal{O}_X$ le
faisceau quasi-cohérent d'idéaux correspondant, voir
Morphismes d'espaces, Lemme
\ref{spaces-morphisms-lemma-closed-immersion-ideals}.
Considérons le faisceau quasi-cohérent de $\mathcal{O}_X$-algèbres graduées
$\bigoplus_{n \geq 0} \mathcal{I}^n/\mathcal{I}^{n + 1}$.
Comme chacun des faisceaux $\mathcal{I}^n/\mathcal{I}^{n + 1}$
est annulé par $\mathcal{I}$, cette algèbre graduée
correspond à un faisceau quasi-cohérent de $\mathcal{O}_Z$-algèbres graduées par
Morphismes d'espaces, Lemme \ref{spaces-morphisms-lemma-i-star-equivalence}.
Cette $\mathcal{O}_Z$-algèbre graduée quasi-cohérente est appelée
{\it algèbre conormale à $Z$ dans $X$} et se note souvent simplement
$\bigoplus_{n \geq 0} \mathcal{I}^n/\mathcal{I}^{n + 1}$
par l'abus de notation signalé dans
Morphismes d'espaces, Section
\ref{spaces-morphisms-section-closed-immersions-quasi-coherent}.

\medskip\noindent
Lorsque $i : Z \to X$ est une immersion (localement fermée), on définit l'algèbre
conormale de $i$ comme l'algèbre conormale de l'immersion fermée
$i : Z \to X \setminus \partial Z$, voir Morphismes d'espaces, Remarque
\ref{spaces-morphisms-remark-immersion}.
On la note souvent
$\bigoplus_{n \geq 0} \mathcal{I}^n/\mathcal{I}^{n + 1}$
où $\mathcal{I}$ est le faisceau d'idéaux
de l'immersion fermée $i : Z \to X \setminus \partial Z$.

\begin{definition}
\label{definition-conormal-algebra}
Soit $i : Z \to X$ une immersion. L'{\it algèbre conormale
$\mathcal{C}_{Z/X, *}$ à $Z$ dans $X$}, ou l'{\it algèbre conormale de $i$},
est le faisceau quasi-cohérent de $\mathcal{O}_Z$-algèbres graduées
$\bigoplus_{n \geq 0} \mathcal{I}^n/\mathcal{I}^{n + 1}$ décrit ci-dessus.
\end{definition}

\noindent
Ainsi, $\mathcal{C}_{Z/X, 1} = \mathcal{C}_{Z/X}$ est le faisceau conormal
de l'immersion. De plus, $\mathcal{C}_{Z/X, 0} = \mathcal{O}_Z$ et
$\mathcal{C}_{Z/X, n}$ est un $\mathcal{O}_Z$-module quasi-cohérent
caractérisé par la propriété
\begin{equation}
\label{equation-conormal-in-degree-n}
i_*\mathcal{C}_{Z/X, n} = \mathcal{I}^n/\mathcal{I}^{n + 1}
\end{equation}
où $i : Z \to X \setminus \partial Z$ et $\mathcal{I}$ est le faisceau d'idéaux
de $i$ comme ci-dessus. Enfin, notons qu'il existe un morphisme canonique surjectif
\begin{equation}
\label{equation-conormal-algebra-quotient}
\text{Sym}^*(\mathcal{C}_{Z/X}) \longrightarrow \mathcal{C}_{Z/X, *}
\end{equation}
de $\mathcal{O}_Z$-algèbres graduées quasi-cohérentes, qui est un isomorphisme
en degrés $0$ et $1$.

\begin{lemma}
\label{lemma-etale-conormal-algebra}
Soit $S$ un schéma. Soit $i : Z \to X$ une immersion d'espaces algébriques
sur $S$. Soit $\varphi : U \to X$ un morphisme étale, où $U$ est un
schéma. Posons $Z_U = U \times_X Z$, qui est un sous-schéma localement fermé de $U$.
Alors
$$
\mathcal{C}_{Z/X, *}|_{Z_U} = \mathcal{C}_{Z_U/U, *}
$$
canoniquement et fonctoriellement en $U$.
\end{lemma}

\begin{proof}
Soit $T \subset X$ un sous-espace fermé tel que $i$ définisse une immersion
fermée dans $X \setminus T$. Soit $\mathcal{I}$ le faisceau quasi-cohérent
d'idéaux sur $X \setminus T$ définissant $Z$. Le lemme découle alors
du fait que
$\mathcal{I}|_{U \setminus \varphi^{-1}(T)}$ est le faisceau d'idéaux de
l'immersion $Z_U \to U \setminus \varphi^{-1}(T)$.
Cela résulte clairement de la construction de $\mathcal{I}$ dans
Morphismes d'espaces, Lemme
\ref{spaces-morphisms-lemma-closed-immersion-ideals}.
\end{proof}

\begin{lemma}
\label{lemma-conormal-algebra-functorial}
Soit $S$ un schéma. Soit
$$
\xymatrix{
Z \ar[r]_i \ar[d]_f & X \ar[d]^g \\
Z' \ar[r]^{i'} & X'
}
$$
un diagramme commutatif d'espaces algébriques sur $S$.
Supposons que $i$ et $i'$ soient des immersions. Il existe un morphisme canonique
de $\mathcal{O}_Z$-algèbres graduées
$$
f^*\mathcal{C}_{Z'/X', *}
\longrightarrow
\mathcal{C}_{Z/X, *}
$$
\end{lemma}

\begin{proof}
Choisissons d'abord des sous-espaces ouverts $U' \subset X'$ et $U \subset X$ tels que
$g(U) \subset U'$ et que $i(Z) \subset U$ et $i(Z') \subset U'$
soient fermés (démonstration de l'existence omise). En remplaçant $X$ par $U$ et $X'$ par
$U'$, on peut supposer que $i$ et $i'$ sont des immersions fermées.
Soient $\mathcal{I}' \subset \mathcal{O}_{X'}$ et
$\mathcal{I} \subset \mathcal{O}_X$ les faisceaux quasi-cohérents
d'idéaux associés à $i'$ et $i$, voir
Morphismes d'espaces, Lemme \ref{spaces-morphisms-lemma-closed-immersion-ideals}.
Considérons la composée
$$
g^{-1}\mathcal{I}' \to g^{-1}\mathcal{O}_{X'}
\xrightarrow{g^\sharp} \mathcal{O}_X \to
\mathcal{O}_X/\mathcal{I} = i_*\mathcal{O}_Z
$$
Comme $g(i(Z)) \subset Z'$, nous concluons que cette composée est nulle (voir
l'énoncé sur les factorisations dans
Morphismes d'espaces,
Lemme \ref{spaces-morphisms-lemma-closed-immersion-ideals}).
On obtient ainsi un diagramme commutatif
$$
\xymatrix{
0 \ar[r] &
\mathcal{I} \ar[r] &
\mathcal{O}_X \ar[r] &
i_*\mathcal{O}_Z \ar[r] &
0 \\
0 \ar[r] &
g^{-1}\mathcal{I}' \ar[r] \ar[u] &
g^{-1}\mathcal{O}_{X'} \ar[r] \ar[u] &
g^{-1}i'_*\mathcal{O}_{Z'} \ar[r] \ar[u] &
0
}
$$
La ligne inférieure est exacte puisque $g^{-1}$ est un foncteur exact.
Par exactitude, on voit aussi que
$(g^{-1}\mathcal{I}')^n = g^{-1}((\mathcal{I}')^n)$ pour tout $n \geq 1$.
Le diagramme induit donc un morphisme
$g^{-1}((\mathcal{I}')^n/(\mathcal{I}')^{n + 1}) \to
\mathcal{I}^n/\mathcal{I}^{n + 1}$.
En prenant l'image inverse (au moyen de $i^{-1}$, par exemple) sur $Z$, on obtient
$i^{-1}g^{-1}((\mathcal{I}')^n/(\mathcal{I}')^{n + 1}) \to
\mathcal{C}_{Z/X, n}$.
Comme $i^{-1}g^{-1} = f^{-1}(i')^{-1}$, cela donne des morphismes
$f^{-1}\mathcal{C}_{Z'/X', n} \to \mathcal{C}_{Z/X, n}$, qui induisent
le morphisme voulu.
\end{proof}

\begin{lemma}
\label{lemma-conormal-algebra-functorial-flat}
Soit $S$ un schéma. Soit
$$
\xymatrix{
Z \ar[r]_i \ar[d]_f & X \ar[d]^g \\
Z' \ar[r]^{i'} & X'
}
$$
un carré cartésien d'espaces algébriques sur $S$ où
$i$ et $i'$ sont des immersions. Alors le morphisme canonique
$f^*\mathcal{C}_{Z'/X', *} \to \mathcal{C}_{Z/X, *}$ du
Lemme \ref{lemma-conormal-algebra-functorial}
est surjectif. Si $g$ est plat, alors c'est un isomorphisme.
\end{lemma}

\begin{proof}
On peut vérifier l'énoncé après localisation étale sur $X'$.
Dans ce cas, on peut supposer que $X' \to X$ est un morphisme de schémas ;
ainsi, $Z$ et $Z'$ sont des schémas et le résultat découle
du cas des schémas, voir
Diviseurs, Lemme \ref{divisors-lemma-conormal-algebra-functorial-flat}.
\end{proof}

\noindent
Nous adoptons pour les cônes et les fibrés vectoriels sur les
espaces algébriques les mêmes conventions que pour les schémas (pour lesquels nous suivons
les conventions de l'EGA), voir
Constructions, Sections \ref{constructions-section-cone} et
\ref{constructions-section-vector-bundle}.
En particulier, un fibré vectoriel est un objet très général
(et n'est pas localement isomorphe à un fibré en espaces affines).

\begin{definition}
\label{definition-normal-cone}
Soit $S$ un schéma. Soit $i : Z \to X$ une immersion d'espaces algébriques
sur $S$. Le {\it cône normal $C_ZX$} à $Z$ dans $X$ est
$$
C_ZX = \underline{\Spec}_Z(\mathcal{C}_{Z/X, *})
$$
voir Morphismes d'espaces,
Définition \ref{spaces-morphisms-definition-relative-spec}. Le
{\it fibré normal} à $Z$ dans $X$ est le fibré vectoriel
$$
N_ZX = \underline{\Spec}_Z(\text{Sym}(\mathcal{C}_{Z/X}))
$$
\end{definition}

\noindent
Ainsi, $C_ZX \to Z$ est un cône sur $Z$ et $N_ZX \to Z$ est un fibré vectoriel
sur $Z$. De plus, la surjection canonique
(\ref{equation-conormal-algebra-quotient}) d'algèbres graduées
définit une immersion fermée canonique
\begin{equation}
\label{equation-normal-cone-in-normal-bundle}
C_ZX \longrightarrow N_ZX
\end{equation}
de cônes sur $Z$.










\section{Faisceau des différentielles d'un morphisme}
\label{section-sheaf-differentials}

\noindent
Nous suggérons au lecteur de consulter la section correspondante
du chapitre sur l'algèbre commutative
(Algèbre, Section \ref{algebra-section-differentials}),
la section correspondante du chapitre sur les morphismes de schémas
(Morphismes, Section \ref{morphisms-section-sheaf-differentials}),
ainsi que
Modules sur les sites, Section \ref{sites-modules-section-differentials}.
Nous montrons d'abord que la notion de faisceau des différentielles d'un
morphisme de schémas coïncide avec celle du morphisme correspondant de
petits sites étales (annelés).

\medskip\noindent
Pour énoncer clairement le lemme suivant, nous revenons temporairement
à la notation $\mathcal{F}^a$ pour le faisceau de $\mathcal{O}_{X_\etale}$-modules
associé à un $\mathcal{O}_X$-module quasi-cohérent $\mathcal{F}$
sur le schéma $X$, voir
Descente, Définition \ref{descent-definition-structure-sheaf}.

\begin{lemma}
\label{lemma-match-modules-differentials}
Soit $f : X \to Y$ un morphisme de schémas. Soit
$f_{small} : X_\etale \to Y_\etale$ le morphisme associé
de petits sites étales, voir
Descente, Remarque \ref{descent-remark-change-topologies-ringed}.
Alors il existe un isomorphisme canonique
$$
(\Omega_{X/Y})^a = \Omega_{X_\etale/Y_\etale}
$$
compatible avec les dérivations universelles. Ici, le premier module
est le faisceau sur $X_\etale$ associé
au $\mathcal{O}_X$-module quasi-cohérent $\Omega_{X/Y}$, voir
Morphismes, Définition \ref{morphisms-definition-sheaf-differentials},
et le second module est celui de
Modules sur les sites,
Définition \ref{sites-modules-definition-module-differentials}.
\end{lemma}

\begin{proof}
Soit $h : U \to X$ un morphisme étale. Dans ce cas, le morphisme naturel
$h^*\Omega_{X/Y} \to \Omega_{U/Y}$ est un isomorphisme, voir
Compléments sur les morphismes,
Lemme \ref{more-morphisms-lemma-sheaf-differentials-etale-localization}.
Cela signifie qu'il existe une $\mathcal{O}_{Y_\etale}$-dérivation naturelle
$$
\text{d}^a : \mathcal{O}_{X_\etale} \longrightarrow (\Omega_{X/Y})^a
$$
puisque nous venons de voir que la valeur de $(\Omega_{X/Y})^a$ sur tout objet
$U$ de $X_\etale$ s'identifie canoniquement à
$\Gamma(U, \Omega_{U/Y})$. Par la propriété universelle de
$\text{d}_{X/Y} :
\mathcal{O}_{X_\etale}
\to
\Omega_{X_\etale/Y_\etale}$
il existe un unique morphisme $\mathcal{O}_{X_\etale}$-linéaire
$c : \Omega_{X_\etale/Y_\etale} \to (\Omega_{X/Y})^a$
tel que
$\text{d}^a = c \circ \text{d}_{X/Y}$.

\medskip\noindent
Réciproquement, supposons que $\mathcal{F}$ soit un
$\mathcal{O}_{X_\etale}$-module
et que $D : \mathcal{O}_{X_\etale} \to \mathcal{F}$ soit une
$\mathcal{O}_{Y_\etale}$-dérivation. On peut alors simplement restreindre
$D$ au petit site de Zariski $X_{Zar}$ de $X$. Comme les faisceaux sur $X_{Zar}$
s'identifient aux faisceaux sur $X$, voir
Descente, Remarque \ref{descent-remark-Zariski-site-space},
nous voyons que $D|_{X_{Zar}} : \mathcal{O}_X \to \mathcal{F}|_{X_{Zar}}$
est simplement une $Y$-dérivation « usuelle ». On obtient donc un morphisme
$\psi : \Omega_{X/Y} \longrightarrow \mathcal{F}|_{X_{Zar}}$
tel que $D|_{X_{Zar}} = \psi \circ \text{d}$. En particulier, si l'on
applique ceci avec $\mathcal{F} = \Omega_{X_\etale/Y_\etale}$,
on obtient un morphisme
$$
c' :
\Omega_{X/Y}
\longrightarrow
\Omega_{X_\etale/Y_\etale}|_{X_{Zar}}
$$
Considérons le morphisme de sites annelés
$\text{id}_{small, \etale, Zar} : X_\etale \to X_{Zar}$
étudié dans
Descente, Remarque \ref{descent-remark-change-topologies-ringed} et
Lemme \ref{descent-lemma-compare-sites}.
Comme le foncteur de restriction $\mathcal{F} \mapsto \mathcal{F}|_{X_{Zar}}$
est égal à $\text{id}_{small, \etale, Zar, *}$, que
$\text{id}_{small, \etale, Zar}^*$ est adjoint à gauche de
$\text{id}_{small, \etale, Zar, *}$ et que
$(\Omega_{X/Y})^a = \text{id}_{small, \etale, Zar}^*\Omega_{X/Y}$
nous voyons que $c'$ est adjoint à un morphisme
$$
c'' :
(\Omega_{X/Y})^a
\longrightarrow
\Omega_{X_\etale/Y_\etale}.
$$
Nous affirmons que $c''$ et $c'$ sont inverses l'un de l'autre.
Cette affirmation achève la démonstration du lemme.
Pour le voir, il suffit de montrer que $c''(\text{d}(f)) = \text{d}_{X/Y}(f)$
et $c(\text{d}_{X/Y}(f)) = \text{d}(f)$ si $f$ est une section locale de
$\mathcal{O}_X$ sur un ouvert de $X$. Nous omettons la vérification.
\end{proof}

\noindent
Cela permet de poser la définition suivante. Pour une autre construction, voir la
Remarque \ref{remark-alternative}.

\begin{definition}
\label{definition-sheaf-differentials}
Soit $S$ un schéma. Soit $f : X \to Y$ un morphisme d'espaces algébriques
sur $S$. Le {\it faisceau des différentielles $\Omega_{X/Y}$ de $X$ sur $Y$}
est le faisceau des différentielles
(Modules sur les sites,
Définition \ref{sites-modules-definition-sheaf-differentials})
du morphisme de topos annelés
$$
(f_{small}, f^\sharp) :
(X_\etale, \mathcal{O}_X)
\to
(Y_\etale, \mathcal{O}_Y)
$$
défini dans
Propriétés des espaces,
Lemme \ref{spaces-properties-lemma-morphism-ringed-topoi}.
La {\it $Y$-dérivation universelle} sera notée
$\text{d}_{X/Y} : \mathcal{O}_X \to \Omega_{X/Y}$.
\end{definition}

\noindent
D'après le
Lemme \ref{lemma-match-modules-differentials},
cela n'entre pas en conflit avec la notion déjà existante
lorsque $X$ et $Y$ sont représentables. Désormais, si $X$ et $Y$
sont représentables, nous ne distinguerons plus le faisceau des différentielles
défini ci-dessus de celui défini dans
Morphismes, Définition \ref{morphisms-definition-sheaf-differentials}.
Nous voulons relier ceci aux modules usuels des différentielles des
morphismes de schémas. Voici le lemme clé.

\begin{lemma}
\label{lemma-localize-differentials}
Soit $S$ un schéma. Soit $f : X \to Y$ un morphisme d'espaces algébriques
sur $S$. Considérons un diagramme commutatif quelconque
$$
\xymatrix{
U \ar[d]_a \ar[r]_\psi & V \ar[d]^b \\
X \ar[r]^f & Y
}
$$
où les flèches verticales sont des morphismes étales d'espaces algébriques. Alors
$$
\Omega_{X/Y}|_{U_\etale} = \Omega_{U/V}
$$
En particulier, si $U$ et $V$ sont des schémas, ceci est égal au faisceau usuel
des différentielles du morphisme de schémas $U \to V$.
\end{lemma}

\begin{proof}
D'après
Propriétés des espaces, Lemme \ref{spaces-properties-lemma-etale-morphism-topoi},
et l'Équation (\ref{spaces-properties-equation-restrict}),
on peut considérer la restriction d'un faisceau sur $X_\etale$ à
$U_\etale$ comme l'image inverse par $a_{small}$. De même pour $b$. Par
Modules sur les sites, Lemme \ref{sites-modules-lemma-localize-differentials},
on a
$$
\Omega_{X/Y}|_{U_\etale} =
\Omega_{\mathcal{O}_{U_\etale}/
a_{small}^{-1}f_{small}^{-1}\mathcal{O}_{Y_\etale}}
$$
Comme $a_{small}^{-1}f_{small}^{-1}\mathcal{O}_{Y_\etale}
= \psi_{small}^{-1}b_{small}^{-1}\mathcal{O}_{Y_\etale}
= \psi_{small}^{-1}\mathcal{O}_{V_\etale}$, le lemme en résulte.
\end{proof}

\begin{lemma}
\label{lemma-module-differentials-quasi-coherent}
Soit $S$ un schéma. Soit $f : X \to Y$ un morphisme d'espaces algébriques
sur $S$. Alors $\Omega_{X/Y}$ est un $\mathcal{O}_X$-module quasi-cohérent.
\end{lemma}

\begin{proof}
Choisissons un diagramme comme dans le
Lemme \ref{lemma-localize-differentials},
avec $a$ et $b$ surjectifs et $U$ et $V$ des schémas.
On voit alors que $\Omega_{X/Y}|_U = \Omega_{U/V}$, qui est
quasi-cohérent (par exemple, d'après
Morphismes, Lemme \ref{morphisms-lemma-differentials-diagonal}).
Nous concluons donc que $\Omega_{X/Y}$ est quasi-cohérent par
Propriétés des espaces,
Lemme \ref{spaces-properties-lemma-characterize-quasi-coherent}.
\end{proof}

\begin{remark}
\label{remark-alternative}
Maintenant que nous savons que $\Omega_{X/Y}$ est quasi-cohérent, nous pouvons tenter
de le construire d'une autre manière. On peut par exemple utiliser le résultat de
Propriétés des espaces,
Section \ref{spaces-properties-section-quasi-coherent-presentation},
pour construire le faisceau des différentielles par recollement.
Par exemple, si $Y$ est un schéma et si $U \to X$ est un morphisme étale surjectif
d'un schéma vers $X$, alors $\Omega_{U/Y}$ est
un $\mathcal{O}_U$-module quasi-cohérent et, comme $s, t : R \to U$
sont étales, on obtient un isomorphisme
$$
\alpha : s^*\Omega_{U/Y} \to \Omega_{R/Y} \to t^*\Omega_{U/Y}
$$
en utilisant
Morphismes, Lemme \ref{morphisms-lemma-triangle-differentials-smooth}.
On vérifie qu'il satisfait la condition de cocycle, et c'est terminé.
Si $Y$ n'est pas un schéma, on définit $\Omega_{U/Y}$ comme le conoyau
du morphisme $(U \to Y)^*\Omega_{Y/S} \to \Omega_{U/S}$, puis on procède comme
précédemment. Cette procédure en deux étapes est quelque peu disgracieuse. Une autre possibilité
consiste à recoller les faisceaux $\Omega_{U/V}$ pour tout diagramme comme dans le
Lemme \ref{lemma-localize-differentials},
mais cela n'est pas non plus très élégant. Les deux méthodes fonctionnent néanmoins et
donnent une construction légèrement plus élémentaire du faisceau des
différentielles.
\end{remark}

\begin{lemma}
\label{lemma-functoriality-differentials}
Soit $S$ un schéma. Soit
$$
\xymatrix{
X' \ar[d] \ar[r]_f & X \ar[d] \\
Y' \ar[r] & Y
}
$$
un diagramme commutatif d'espaces algébriques. Le morphisme
$f^\sharp : \mathcal{O}_X \to f_*\mathcal{O}_{X'}$, composé avec le morphisme
$f_*\text{d}_{X'/Y'} : f_*\mathcal{O}_{X'} \to f_*\Omega_{X'/Y'}$, est une
$Y$-dérivation. On obtient donc un morphisme canonique de $\mathcal{O}_X$-modules
$\Omega_{X/Y} \to f_*\Omega_{X'/Y'}$ et, par
l'adjonction entre $f_*$ et $f^*$, un
homomorphisme canonique de $\mathcal{O}_{X'}$-modules
$$
c_f : f^*\Omega_{X/Y} \longrightarrow \Omega_{X'/Y'}.
$$
Il est uniquement caractérisé par la propriété suivante :
$f^*\text{d}_{X/Y}(t)$ est envoyé sur $\text{d}_{X'/Y'}(f^* t)$
pour toute section locale $t$ de $\mathcal{O}_X$.
\end{lemma}

\begin{proof}
Il s'agit d'un cas particulier de
Modules sur les sites, Lemme
\ref{sites-modules-lemma-functoriality-differentials-ringed-topoi}.
\end{proof}

\begin{lemma}
\label{lemma-check-functoriality-differentials}
Soit $S$ un schéma. Soit
$$
\xymatrix{
X'' \ar[d] \ar[r]_g & X' \ar[d] \ar[r]_f & X \ar[d] \\
Y'' \ar[r] & Y' \ar[r] & Y
}
$$
un diagramme commutatif d'espaces algébriques sur $S$. Alors on a
$$
c_{f \circ g} = c_g \circ g^* c_f
$$
comme morphismes $(f \circ g)^*\Omega_{X/Y} \to \Omega_{X''/Y''}$.
\end{lemma}

\begin{proof}
Démonstration omise. Indication : utiliser la caractérisation de $c_f, c_g, c_{f \circ g}$
par l'effet de ces morphismes sur les sections locales.
\end{proof}

\begin{lemma}
\label{lemma-triangle-differentials}
Soit $S$ un schéma.
Soient $f : X \to Y$ et $g : Y \to B$ des morphismes d'espaces algébriques sur $S$.
Alors il existe une suite exacte canonique
$$
f^*\Omega_{Y/B} \to \Omega_{X/B} \to \Omega_{X/Y} \to 0
$$
où les morphismes proviennent d'applications du
Lemme \ref{lemma-functoriality-differentials}.
\end{lemma}

\begin{proof}
Cela découle de la version de ce résultat pour les schémas, voir
Morphismes, Lemme \ref{morphisms-lemma-triangle-differentials},
par localisation étale, voir le
Lemme \ref{lemma-localize-differentials}.
\end{proof}

\begin{lemma}
\label{lemma-immersion-differentials}
Soit $S$ un schéma. Si $X \to Y$ est une immersion
d'espaces algébriques sur $S$, alors $\Omega_{X/S}$ est nul.
\end{lemma}

\begin{proof}
Cela découle de la version de ce résultat pour les schémas, voir
Morphismes, Lemme \ref{morphisms-lemma-immersion-differentials},
par localisation étale, voir le
Lemme \ref{lemma-localize-differentials}.
\end{proof}

\begin{lemma}
\label{lemma-differentials-relative-immersion}
Soit $S$ un schéma. Soit $B$ un espace algébrique sur $S$.
Soit $i : Z \to X$ une immersion d'espaces algébriques sur $B$.
Il existe une suite exacte canonique
$$
\mathcal{C}_{Z/X} \to i^*\Omega_{X/B} \to \Omega_{Z/B} \to 0
$$
où la première flèche est induite par $\text{d}_{X/B}$
et où la seconde flèche provient du
Lemme \ref{lemma-functoriality-differentials}.
\end{lemma}

\begin{proof}
Il s'agit de la version pour les espaces algébriques de
Morphismes, Lemme \ref{morphisms-lemma-differentials-relative-immersion},
et elle découlera de ce lemme par
localisation étale, voir les
Lemmes \ref{lemma-localize-differentials} et
\ref{lemma-etale-conormal}.
Il faut toutefois s'assurer que l'on peut définir globalement la première flèche.
Expliquons donc ici le sens de « induite par $\text{d}_{X/B}$ ».
On peut en effet supposer que $i$ est une immersion fermée après avoir remplacé $X$
par un sous-espace ouvert. Soit $\mathcal{I} \subset \mathcal{O}_X$
le faisceau quasi-cohérent d'idéaux correspondant à $Z \subset X$.
Alors $\text{d}_{X/S} : \mathcal{I} \to \Omega_{X/S}$
envoie le sous-faisceau $\mathcal{I}^2 \subset \mathcal{I}$ dans
$\mathcal{I}\Omega_{X/S}$. Il induit donc un morphisme
$\mathcal{I}/\mathcal{I}^2 \to \Omega_{X/S}/\mathcal{I}\Omega_{X/S}$
qui est $\mathcal{O}_X/\mathcal{I}$-linéaire.
D'après
Morphismes d'espaces, Lemme \ref{spaces-morphisms-lemma-i-star-equivalence},
cela correspond au morphisme $\mathcal{C}_{Z/X} \to i^*\Omega_{X/S}$ voulu.
\end{proof}

\begin{lemma}
\label{lemma-differentials-relative-immersion-section}
Soit $S$ un schéma. Soit $B$ un espace algébrique sur $S$.
Soit $i : Z \to X$ une immersion d'espaces algébriques sur $B$, et
supposons que $i$ admette (localement pour la topologie étale) un inverse à gauche. Alors la
suite canonique
$$
0 \to \mathcal{C}_{Z/X} \to i^*\Omega_{X/B} \to \Omega_{Z/B} \to 0
$$
du
Lemme \ref{lemma-differentials-relative-immersion}
est (localement pour la topologie étale) exacte scindée.
\end{lemma}

\begin{proof}
Précision : nous affirmons que, si $g : X \to Z$ est un inverse à gauche de $i$
sur $B$, alors $i^*c_g$ est un inverse à droite du morphisme
$i^*\Omega_{X/B} \to \Omega_{Z/B}$.
Ceci étant dit, le résultat découle du résultat correspondant pour les
morphismes de schémas par localisation étale, voir les
Lemmes \ref{lemma-localize-differentials} et
\ref{lemma-etale-conormal}.
\end{proof}

\begin{lemma}
\label{lemma-base-change-differentials}
Soit $S$ un schéma.
Soit $X \to Y$ un morphisme d'espaces algébriques sur $S$.
Soit $g : Y' \to Y$ un morphisme d'espaces algébriques sur $S$.
Soit $X' = X_{Y'}$ le changement de base de $X$.
Notons $g' : X' \to X$ la projection.
Alors le morphisme
$$
(g')^*\Omega_{X/Y} \to \Omega_{X'/Y'}
$$
du
Lemme \ref{lemma-functoriality-differentials}
est un isomorphisme.
\end{lemma}

\begin{proof}
Cela découle de la version pour les schémas, voir
Morphismes, Lemme \ref{morphisms-lemma-base-change-differentials},
et de la localisation étale, voir le
Lemme \ref{lemma-localize-differentials}.
\end{proof}

\begin{lemma}
\label{lemma-differential-product}
Soit $S$ un schéma.
Soient $f : X \to B$ et $g : Y \to B$ des morphismes d'espaces algébriques
sur $S$ ayant le même but.
Soient $p : X \times_B Y \to X$ et $q : X \times_B Y \to Y$ les
morphismes de projection. Les morphismes du
Lemme \ref{lemma-functoriality-differentials}
$$
p^*\Omega_{X/B} \oplus q^*\Omega_{Y/B}
\longrightarrow
\Omega_{X \times_B Y/B}
$$
donnent un isomorphisme.
\end{lemma}

\begin{proof}
Cela découle de la version pour les schémas, voir
Morphismes, Lemme \ref{morphisms-lemma-differential-product},
et de la localisation étale, voir le
Lemme \ref{lemma-localize-differentials}.
\end{proof}

\begin{lemma}
\label{lemma-finite-type-differentials}
Soit $S$ un schéma.
Soit $f : X \to Y$ un morphisme d'espaces algébriques sur $S$.
Si $f$ est localement de type fini, alors $\Omega_{X/Y}$ est
un $\mathcal{O}_X$-module de type fini.
\end{lemma}

\begin{proof}
Cela découle de la version pour les schémas, voir
Morphismes, Lemme \ref{morphisms-lemma-finite-type-differentials},
et de la localisation étale, voir le
Lemme \ref{lemma-localize-differentials}.
\end{proof}

\begin{lemma}
\label{lemma-finite-presentation-differentials}
Soit $S$ un schéma.
Soit $f : X \to Y$ un morphisme d'espaces algébriques sur $S$.
Si $f$ est localement de présentation finie, alors $\Omega_{X/Y}$ est
un $\mathcal{O}_X$-module de présentation finie.
\end{lemma}

\begin{proof}
Cela découle de la version pour les schémas, voir
Morphismes, Lemme \ref{morphisms-lemma-finite-presentation-differentials},
et de la localisation étale, voir le
Lemme \ref{lemma-localize-differentials}.
\end{proof}

\begin{lemma}
\label{lemma-smooth-omega-finite-locally-free}
Soit $S$ un schéma.
Soit $f : X \to Y$ un morphisme lisse d'espaces algébriques sur $S$.
Alors le module des différentielles $\Omega_{X/Y}$
est localement libre de type fini.
\end{lemma}

\begin{proof}
L'énoncé est local pour la topologie étale sur $X$ et $Y$ par le
Lemme \ref{lemma-localize-differentials}.
Il découle donc du cas des schémas, voir
Morphismes, Lemme \ref{morphisms-lemma-smooth-omega-finite-locally-free}.
\end{proof}













\section{Invariance topologique du site étale}
\label{section-topological-invariance}

\noindent
Nous montrons que le site $X_{spaces, \etale}$ est un « invariant
topologique ». Il s'ensuit que $X_\etale$, qui
est formé des objets représentables de $X_{spaces, \etale}$,
est lui aussi un invariant topologique, voir le Lemme \ref{lemma-topological-invariance}.

\begin{theorem}
\label{theorem-topological-invariance}
Soit $S$ un schéma.
Soit $f : X \to Y$ un morphisme d'espaces algébriques sur $S$.
Supposons que $f$ soit entier, universellement injectif et surjectif.
Le foncteur
$$
V \longmapsto V_X = X \times_Y V
$$
définit une équivalence de catégories
$Y_{spaces, \etale} \to X_{spaces, \etale}$.
\end{theorem}

\begin{proof}
Le morphisme $f$ est représentable et constitue un homéomorphisme universel, voir
Morphismes d'espaces,
Section \ref{spaces-morphisms-section-universal-homeomorphisms}.

\medskip\noindent
Nous montrons d'abord que le foncteur est fidèle.
Supposons que $V', V$ soient des objets de $Y_{spaces, \etale}$ et
que $a, b : V' \to V$ soient des morphismes distincts sur $Y$.
Comme $V', V$ sont étales sur $Y$, l'égalisateur
$$
E =  V' \times_{(a, b), V \times_Y V, \Delta_{V/Y}} V
$$
de $a, b$ est lui aussi étale sur $Y$. Ainsi, $E \to V'$ est un monomorphisme étale
(c'est-à-dire une immersion ouverte) qui est un isomorphisme si et seulement s'il est
surjectif. Comme $X \to Y$ est un homéomorphisme universel, nous voyons que c'est
le cas si et seulement si $E_X = V'_X$, c'est-à-dire si et seulement si $a_X = b_X$.

\medskip\noindent
Nous montrons ensuite que le foncteur est pleinement fidèle.
Supposons que $V', V$ soient des objets de $Y_{spaces, \etale}$ et
que $c : V'_X \to V_X$ soit un morphisme sur $X$. Nous voulons construire
un morphisme $a : V' \to V$ sur $Y$ tel que $a_X = c$.
Soit $a' : V'' \to V'$ un morphisme étale surjectif tel que $V''$ soit
un espace algébrique séparé. Si l'on peut construire un morphisme
$a'' : V'' \to V$ tel que $a''_X = c \circ a'_X$, alors les deux composées
$$
V'' \times_{V'} V'' \xrightarrow{\text{pr}_i} V'' \xrightarrow{a''} V
$$
sont égales par la fidélité du foncteur établie au premier
paragraphe. Ainsi, $a''$ se factorise par un unique morphisme
$a : V' \to V$, puisque $V'$ est (en tant que faisceau) le quotient de $V''$ par
la relation d'équivalence $V'' \times_{V'} V''$. Nous pouvons donc supposer que
$V'$ est séparé. Dans ce cas, le graphe
$$
\Gamma_c \subset (V' \times_Y V)_X
$$
est ouvert et fermé (détails omis). Comme $X \to Y$ est un homéomorphisme
universel, il existe un sous-espace ouvert et fermé
$\Gamma \subset V' \times_Y V$ tel que $\Gamma_X = \Gamma_c$.
La projection $\Gamma \to V'$ est un morphisme étale dont le changement
de base à $X$ est un isomorphisme. Ainsi, $\Gamma \to V'$ est étale,
universellement injectif et surjectif, donc est un isomorphisme d'après
Morphismes d'espaces,
Lemme \ref{spaces-morphisms-lemma-etale-universally-injective-open}.
Ainsi, $\Gamma$ est le graphe d'un morphisme $a : V' \to V$, comme souhaité.

\medskip\noindent
Enfin, nous montrons que le foncteur est essentiellement surjectif.
Supposons que $U$ soit un objet de $X_{spaces, \etale}$.
Il faut trouver un objet $V$ de $Y_{spaces, \etale}$
tel que $V_X \cong U$. Soit $U' \to U$ un morphisme étale surjectif
tel que $U' \cong V'_X$ et $U' \times_U U' \cong V''_X$
pour certains objets $V'', V'$ de $Y_{spaces, \etale}$.
Alors, par la pleine fidélité du foncteur, on obtient des morphismes
$s, t : V'' \to V'$ tels que $t_X = \text{pr}_0$ et $s_X = \text{pr}_1$
comme morphismes $U' \times_U U' \to U'$. Puisque
$(\text{pr}_0, \text{pr}_1) : U' \times_U U' \to U' \times_S U'$
est une relation d'équivalence étale et que $U' \to V'$ et
$U' \times_U U' \to V''$ sont universellement injectifs et surjectifs,
on en déduit que
$(t, s) : V'' \to V' \times_S V'$ est une relation d'équivalence étale.
Alors le quotient $V = V'/V''$ (voir
Espaces, Théorème \ref{spaces-theorem-presentation})
est l'espace algébrique $V$ sur $Y$. Il existe un morphisme
$V' \to V$ tel que $V'' = V' \times_V V'$. On obtient ainsi un morphisme
$V \to Y$ (voir
Descente sur les espaces, Lemme
\ref{spaces-descent-lemma-fpqc-universal-effective-epimorphisms}).
Après changement de base à $X$, on voit que l'on a un morphisme $U' \to V_X$
et un isomorphisme compatible $U' \times_{V_X} U' = U' \times_U U'$, ce qui
implique que $V_X \cong U$ (en appliquant encore une fois le lemme qui vient d'être cité).

\medskip\noindent
Choisissons un schéma $W$ et un morphisme étale surjectif $W \to Y$.
Choisissons un schéma $U'$ et un morphisme étale surjectif $U' \to U \times_X W_X$.
Notons que $U'$ et $U' \times_U U'$ sont des schémas étales sur $X$ dont le
morphisme structural vers $X$ se factorise par le schéma $W_X$.
Ainsi, d'après
Cohomologie étale,
Théorème \ref{etale-cohomology-theorem-topological-invariance},
il existe des schémas $V', V''$ étales sur $W$ dont les changements de base à
$W_X$ sont respectivement isomorphes à $U'$ et $U' \times_U U'$.
Cela achève la démonstration.
\end{proof}

\begin{lemma}
\label{lemma-topological-invariance}
Sous les hypothèses et avec les notations du
Théorème \ref{theorem-topological-invariance},
l'équivalence de catégories
$Y_{spaces, \etale} \to X_{spaces, \etale}$
induit par restriction des équivalences de catégories
$Y_\etale \to X_\etale$ et $Y_{affine, \etale} \to X_{affine, \etale}$.
\end{lemma}

\begin{proof}
Il s'agit simplement de dire que, pour tout objet
$V \in Y_{spaces, \etale}$, $V$ est un schéma (respectivement un schéma affine) si et
seulement si $V \times_Y X$ est un schéma (respectivement un schéma affine). Comme $V \times_Y X \to V$
est entier, universellement injectif et surjectif (en tant que changement
de base de $X \to Y$), cela
découle du Lemme
\ref{spaces-limits-lemma-integral-universally-bijective-scheme} de Limites d'espaces et de la
Proposition \ref{spaces-limits-proposition-affine}.
\end{proof}

\begin{remark}
\label{remark-topological-invariance-etale-site}
\begin{reference}
Courriel de Lenny Taelman daté du 1er mai 2016.
\end{reference}
Un homéomorphisme universel d'espaces algébriques n'est pas nécessairement représentable, voir
Morphismes d'espaces,
Exemple \ref{spaces-morphisms-example-universal-homeomorphism}.
En fait, le Théorème \ref{theorem-topological-invariance} n'est
pas valable pour tous les homéomorphismes universels. Pour le voir, soit $k$ un
corps algébriquement clos de caractéristique $0$ et soit
$$
\mathbf{A}^1 \to X \to \mathbf{A}^1
$$
comme dans Morphismes d'espaces,
Exemple \ref{spaces-morphisms-example-universal-homeomorphism}.
Rappelons que le premier morphisme est étale et identifie
$t$ à $-t$ pour $t \in \mathbf{A}^1_k \setminus \{0\}$
et que le second morphisme est notre homéomorphisme universel.
Comme $\mathbf{A}^1_k$ n'admet aucun
revêtement fini étale connexe non trivial
(car $k$ est algébriquement clos de caractéristique zéro ; détails omis),
il suffit de construire un revêtement fini étale connexe non trivial
$Y \to X$. Pour cela, soit $Y$ la droite affine
à origine dédoublée
(Schémas, Exemple \ref{schemes-example-affine-space-zero-doubled}).
Alors $Y = Y_1 \cup Y_2$, avec $Y_i = \mathbf{A}^1_k$, est obtenu par recollement
le long de $\mathbf{A}^1_k \setminus \{0\}$.
Pour définir le morphisme $Y \to X$, nous utilisons les morphismes
$$
Y_1 \xrightarrow{1} \mathbf{A}^1_k \to X
\quad\text{and}\quad
Y_2 \xrightarrow{-1} \mathbf{A}^1_k \to X.
$$
Ces morphismes se recollent sur $Y_1 \cap Y_2$ par construction de $X$ et
définissent donc un morphisme $Y \to X$. En fait, nous affirmons que
$$
\xymatrix{
Y \ar[d] & Y_1 \amalg Y_2 \ar[l] \ar[d] \\
X & \mathbf{A}^1_k \ar[l]
}
$$
est un carré cartésien. Nous omettons les détails ; on peut utiliser par exemple
Groupoïdes, Lemme \ref{groupoids-lemma-criterion-fibre-product}.
Comme $\mathbf{A}^1_k \to X$ est étale et
surjectif, cela montre que $Y \to X$
est fini étale de degré $2$, ce qui fournit l'exemple voulu.

\medskip\noindent
Plus simplement, on peut raisonner comme suit. Le schéma $Y$ est muni de l'action du groupe
$G = \{+1, -1\}$, qui opère librement, l'élément $-1$ agissant en échangeant
$Y_1$ et $Y_2$ et en changeant le signe de la coordonnée. Alors
$X = Y/G$ (voir Espaces, Définition \ref{spaces-definition-quotient})
et donc $Y \to X$ est fini étale. On peut aussi montrer directement
qu'il existe un homéomorphisme universel $X \to \mathbf{A}^1_k$
en utilisant $t \mapsto t^2$ sur les espaces affines. En fait, ce $X$ est
le même que le $X$ ci-dessus.
\end{remark}








\section{Épaississements}
\label{section-thickenings}

\noindent
La terminologie suivante n'est peut-être pas tout à fait standard, mais elle est commode.

\begin{definition}
\label{definition-thickening}
Épaississements. Soit $S$ un schéma.
\begin{enumerate}
\item Nous disons qu'un espace algébrique $X'$ est un {\it épaississement} d'un espace
algébrique $X$ si $X$ est un sous-espace fermé de $X'$ et si les espaces topologiques
associés sont égaux.
\item Nous disons que $X'$ est un {\it épaississement du premier ordre} de $X$ si
$X$ est un sous-espace fermé de $X'$ et si le faisceau quasi-cohérent d'idéaux
$\mathcal{I} \subset \mathcal{O}_{X'}$ qui définit $X$ est de carré nul.
\item Nous disons que $X'$ est un {\it épaississement d'ordre fini} de $X$
si $X$ est un sous-espace fermé de $X$ et si le faisceau quasi-cohérent d'idéaux
$\mathcal{I} \subset \mathcal{O}_{X'}$ qui définit $X$ est nilpotent, c'est-à-dire s'il
existe un entier $n \geq 0$ tel que $\mathcal{I}^{n + 1} = 0$.
\item Nous disons que $X'$ est un {\it épaississement d'ordre $n$} de $X$
si $X$ est un sous-espace fermé de $X$ et si $\mathcal{I}^{n + 1} = 0$,
où $\mathcal{I} \subset \mathcal{O}_{X'}$ est le
faisceau quasi-cohérent d'idéaux qui définit $X$.
\item Étant donnés deux épaississements $X \subset X'$ et $Y \subset Y'$, un
{\it morphisme d'épaississements} est un morphisme $f' : X' \to Y'$ tel que
$f(X) \subset Y$, c'est-à-dire tel que $f'|_X$ se factorise par le sous-espace
fermé $Y$. Dans cette situation, nous posons $f = f'|_X : X \to Y$ et disons
que $(f, f') : (X \subset X') \to (Y \subset Y')$ est un morphisme
d'épaississements.
\item Soit $B$ un espace algébrique. Nous définissons de même les
{\it épaississements sur $B$} et les
{\it morphismes d'épaississements sur $B$}. Cela signifie que les espaces
$X, X', Y, Y'$ ci-dessus sont des espaces algébriques munis d'un morphisme
structural vers $B$ et que les morphismes
$X \to X'$, $Y \to Y'$ et $f' : X' \to Y'$ sont des morphismes sur $B$.
\end{enumerate}
\end{definition}

\noindent
L'équivalence fondamentale.
Notons que si $X \subset X'$ est un épaississement, alors $X \to X'$
est entier et universellement bijectif. Il en résulte que
\begin{equation}
\label{equation-equivalence-etale-spaces}
X_{spaces, \etale} = X'_{spaces, \etale}
\end{equation}
par le foncteur d'image inverse, voir le
Théorème \ref{theorem-topological-invariance}.
Nous pouvons donc considérer $\mathcal{O}_{X'}$ comme un faisceau sur
$X_{spaces, \etale}$. Nous obtenons ainsi une équivalence canonique
de topos localement annelés
\begin{equation}
\label{equation-fundamental-equivalence}
(\Sh(X'_{spaces, \etale}), \mathcal{O}_{X'})
\cong
(\Sh(X_{spaces, \etale}), \mathcal{O}_{X'})
\end{equation}
Dans la suite, nous la combinerons fréquemment avec le résultat de pleine fidélité de
Propriétés des espaces, Théorème \ref{spaces-properties-theorem-fully-faithful}.
Par exemple, à l'immersion fermée $i_X : X \to X'$ correspond
le morphisme surjectif $i_X^\sharp : \mathcal{O}_{X'} \to \mathcal{O}_X$.

\medskip\noindent
Soit $S$ un schéma et soit $B$ un espace algébrique sur $S$.
Soit $(f, f') : (X \subset X') \to (Y \subset Y')$ un morphisme
d'épaississements sur $B$. Notons que le diagramme de foncteurs continus
$$
\xymatrix{
X_{spaces, \etale} &
Y_{spaces, \etale} \ar[l] \\
X'_{spaces, \etale} \ar[u] &
Y'_{spaces, \etale} \ar[u] \ar[l]
}
$$
est commutatif et que les flèches verticales sont des équivalences. Ainsi,
$f_{spaces, \etale}$, $f_{small}$,
$f'_{spaces, \etale}$, and $f'_{small}$
définissent tous le même morphisme de topos. Nous pouvons donc considérer
$$
(f')^\sharp :
f_{spaces, \etale}^{-1}\mathcal{O}_{Y'}
\longrightarrow
\mathcal{O}_{X'}
$$
comme un morphisme de faisceaux de $\mathcal{O}_B$-algèbres qui s'insère dans le
diagramme commutatif
$$
\xymatrix{
f_{spaces, \etale}^{-1}\mathcal{O}_Y \ar[r]_-{f^\sharp} \ar[r] &
\mathcal{O}_X \\
f_{spaces, \etale}^{-1}\mathcal{O}_{Y'} \ar[r]^-{(f')^\sharp}
\ar[u]^{i_Y^\sharp} &
\mathcal{O}_{X'} \ar[u]_{i_X^\sharp}
}
$$
Ici, $i_X : X \to X'$ et $i_Y : Y \to Y'$ désignent les
immersions fermées données.

\begin{lemma}
\label{lemma-first-order-thickening-maps}
Soit $S$ un schéma. Soit $B$ un espace algébrique sur $S$.
Soient $X \subset X'$ et $Y \subset Y'$ des épaississements
d'espaces algébriques sur $B$. Soit $f : X \to Y$ un morphisme d'espaces
algébriques sur $B$. Étant donné un morphisme de $\mathcal{O}_B$-algèbres
$$
\alpha : f_{spaces, \etale}^{-1}\mathcal{O}_{Y'} \to \mathcal{O}_{X'}
$$
tel que le diagramme
$$
\xymatrix{
f_{spaces, \etale}^{-1}\mathcal{O}_Y \ar[r]_-{f^\sharp} \ar[r] &
\mathcal{O}_X \\
f_{spaces, \etale}^{-1}\mathcal{O}_{Y'} \ar[r]^-\alpha
\ar[u]^{i_Y^\sharp} &
\mathcal{O}_{X'} \ar[u]_{i_X^\sharp}
}
$$
soit commutatif, il existe un unique morphisme $(f, f')$
d'épaississements sur $B$ tel que $\alpha = (f')^\sharp$.
\end{lemma}

\begin{proof}
Pour trouver $f'$, d'après
Propriétés des espaces, Théorème \ref{spaces-properties-theorem-fully-faithful},
il suffit de montrer que le morphisme de topos annelés
$$
(f_{spaces, \etale}, \alpha) :
(\Sh(X_{spaces, \etale}), \mathcal{O}_{X'})
\longrightarrow
(\Sh(Y_{spaces, \etale}), \mathcal{O}_{Y'})
$$
est un morphisme de topos localement annelés. Cela découle directement
de la définition des morphismes de topos localement annelés
(Modules sur les sites,
Définition \ref{sites-modules-definition-morphism-locally-ringed-topoi}),
du fait que $(f, f^\sharp)$ est un morphisme de topos localement annelés
(Propriétés des espaces,
Lemme \ref{spaces-properties-lemma-morphism-locally-ringed}),
du fait que $\alpha$ s'insère dans le diagramme commutatif donné et
du fait que les noyaux de $i_X^\sharp$ et de $i_Y^\sharp$ sont
localement nilpotents. Enfin, l'égalité $f' \circ i_X = i_Y \circ f$
découle de la commutativité du diagramme et d'une nouvelle application de
Propriétés des espaces, Théorème \ref{spaces-properties-theorem-fully-faithful}.
Nous omettons la vérification que $f'$ est un morphisme sur $B$.
\end{proof}

\begin{lemma}
\label{lemma-open-subspace-thickening}
Soit $S$ un schéma. Soit $X \subset X'$ un épaississement
d'espaces algébriques sur $S$. Pour tout sous-espace ouvert $U \subset X$, il
existe un unique sous-espace ouvert $U' \subset X'$ tel que
$U = X \times_{X'} U'$.
\end{lemma}

\begin{proof}
Soit $U' \to X'$ l'objet de $X'_{spaces, \etale}$
qui correspond à l'objet $U \to X$ de $X_{spaces, \etale}$
par (\ref{equation-equivalence-etale-spaces}). Le morphisme
$U' \to X'$ est étale et universellement injectif, donc une immersion ouverte, voir
Morphismes d'espaces,
Lemme \ref{spaces-morphisms-lemma-etale-universally-injective-open}.
\end{proof}

\noindent
Soit $i_X : X \to X'$ un épaississement d'espaces algébriques. Toute section locale
du noyau $\mathcal{I} = \Ker(i_X^\sharp) \subset \mathcal{O}_{X'}$ est
localement nilpotente. Cela donne un certain contrôle sur le faisceau structural
$\mathcal{O}_{X'}$, mais, techniquement, la classe des épaississements d'ordre fini
$X \subset X'$ est beaucoup plus facile à manier.
En effet, si $X'$ est un épaississement d'ordre $n$, nous avons une filtration
$$
0 \subset \mathcal{I}^n \subset \mathcal{I}^{n - 1} \subset
\ldots \subset \mathcal{I} \subset \mathcal{O}_{X'}
$$
et nous voyons que $X'$ est muni d'une filtration par des sous-espaces fermés
$$
X = X_0 \subset X_1 \subset \ldots \subset X_{n - 1} \subset X_{n + 1} = X'
$$
telle que chaque paire $X_i \subset X_{i + 1}$ est un épaississement du premier ordre
sur $B$. De simples arguments de récurrence permettent de reformuler nombre de résultats démontrés pour les
épaississements du premier ordre en des résultats sur les épaississements d'ordre fini.

\begin{lemma}
\label{lemma-first-order-thickening-surjective}
Soit $S$ un schéma. Soit $X \subset X'$ un épaississement
d'espaces algébriques sur $S$. Soit $U$ un objet affine de
$X_{spaces, \etale}$. Alors
$$
\Gamma(U, \mathcal{O}_{X'}) \to \Gamma(U, \mathcal{O}_X)
$$
est surjectif, où l'on considère $\mathcal{O}_{X'}$ comme un faisceau sur
$X_{spaces, \etale}$ au moyen de (\ref{equation-fundamental-equivalence}).
\end{lemma}

\begin{proof}
Soit $U' \to X'$ le morphisme étale d'espaces algébriques tel que
$U = X \times_{X'} U'$, voir le Théorème \ref{theorem-topological-invariance}.
D'après Limites d'espaces, Lemme \ref{spaces-limits-lemma-affine}, nous voyons
que $U'$ est un schéma affine. Par conséquent,
$\Gamma(U, \mathcal{O}_{X'}) = \Gamma(U', \mathcal{O}_{U'}) \to
\Gamma(U, \mathcal{O}_U)$
est surjectif, puisque $U \to U'$ est une immersion fermée de schémas affines.
Nous donnons ci-dessous une démonstration directe pour les épaississements d'ordre fini,
cas le plus utilisé en pratique.
\end{proof}

\begin{proof}[Démonstration pour les épaississements d'ordre fini]
Nous pouvons supposer que $X \subset X'$ est un épaississement du premier ordre, par le
principe expliqué ci-dessus. Notons $\mathcal{I}$ le noyau de la surjection
$\mathcal{O}_{X'} \to \mathcal{O}_X$. Comme $\mathcal{I}$ est un
$\mathcal{O}_{X'}$-module quasi-cohérent et que $\mathcal{I}^2 = 0$ par définition
d'un épaississement du premier ordre, nous pouvons appliquer
Morphismes d'espaces, Lemme
\ref{spaces-morphisms-lemma-i-star-equivalence}
pour voir que $\mathcal{I}$ est un $\mathcal{O}_X$-module quasi-cohérent.
Le lemme découle donc de la suite exacte longue de cohomologie
associée à la suite exacte courte
$$
0 \to \mathcal{I} \to \mathcal{O}_{X'} \to \mathcal{O}_X \to 0
$$
et du fait que $H^1_\etale(U, \mathcal{I}) = 0$, puisque
$\mathcal{I}$ est quasi-cohérent, voir
Descente, Proposition \ref{descent-proposition-same-cohomology-quasi-coherent},
et Cohomologie des schémas, Lemme
\ref{coherent-lemma-quasi-coherent-affine-cohomology-zero}.
\end{proof}

\begin{lemma}
\label{lemma-thickening-scheme}
Soit $S$ un schéma. Soit $X \subset X'$ un épaississement d'espaces algébriques
sur $S$. Si $X$ est (représentable par) un schéma, alors $X'$ l'est aussi.
\end{lemma}

\begin{proof}
Notons que $X'_{red} = X_{red}$. Par conséquent, si $X$ est un schéma, alors
$X'_{red}$ est un schéma. Le résultat découle donc de
Limites d'espaces, Lemme
\ref{spaces-limits-lemma-reduction-scheme}.
Nous donnons ci-dessous une démonstration directe pour les épaississements d'ordre fini,
cas le plus souvent utilisé en pratique.
\end{proof}

\begin{proof}[Démonstration pour les épaississements d'ordre fini]
Il suffit de le démontrer lorsque $X'$ est un épaississement du premier ordre de $X$. D'après
Propriétés des espaces, Lemme \ref{spaces-properties-lemma-subscheme},
il existe un plus grand sous-espace ouvert de $X'$ qui soit un schéma. Nous devons donc
montrer que tout point $x$ de $|X'| = |X|$ est contenu dans un sous-espace ouvert de
$X'$ qui soit un schéma. En utilisant le
Lemme \ref{lemma-open-subspace-thickening},
nous pouvons remplacer $X \subset X'$ par $U \subset U'$, avec $x \in U$ et $U$
un schéma affine. Nous pouvons donc supposer que $X$ est affine.
Nous sommes ainsi ramenés au cas traité dans le paragraphe suivant.

\medskip\noindent
Supposons que $X \subset X'$ soit un épaississement du premier ordre, où $X$ est un schéma
affine. Posons $A = \Gamma(X, \mathcal{O}_X)$ et
$A' = \Gamma(X', \mathcal{O}_{X'})$. D'après le
Lemme \ref{lemma-first-order-thickening-surjective},
le morphisme $A' \to A$ est surjectif. Son noyau $I$ est un idéal de carré nul. D'après
Propriétés des espaces,
Lemme \ref{spaces-properties-lemma-morphism-to-affine-scheme},
nous obtenons un morphisme canonique $f : X' \to \Spec(A')$ qui s'insère
dans le diagramme commutatif suivant
$$
\xymatrix{
X \ar@{=}[d] \ar[r] &  X' \ar[d]^f \\
\Spec(A) \ar[r] & \Spec(A')
}
$$
Comme les flèches horizontales sont des épaississements, il est clair que $f$ est
universellement injectif et surjectif. Il suffit donc de montrer que
$f$ est étale, car alors
Morphismes d'espaces,
Lemme \ref{spaces-morphisms-lemma-etale-universally-injective-open},
entraînera que $f$ est un isomorphisme.

\medskip\noindent
Pour démontrer que $f$ est étale, choisissons un schéma affine $U'$ et un
morphisme étale $U' \to X'$. Il suffit de montrer que
$U' \to X' \to \Spec(A')$ est étale, voir
Propriétés des espaces, Définition \ref{spaces-properties-definition-etale}.
Écrivons $U' = \Spec(B')$. Posons $U = X \times_{X'} U'$. Comme $U$
est un sous-espace fermé de $U'$, c'est un sous-schéma fermé, et donc
$U = \Spec(B)$ avec $B' \to B$ surjectif. Notons
$J = \Ker(B' \to B)$ et remarquons que $J = \Gamma(U, \mathcal{I})$,
où $\mathcal{I} = \Ker(\mathcal{O}_{X'} \to \mathcal{O}_X)$
sur $X_{spaces, \etale}$, comme dans la démonstration du
Lemme \ref{lemma-first-order-thickening-surjective}.
Le morphisme $U' \to X' \to \Spec(A')$ induit un diagramme
commutatif
$$
\xymatrix{
0 \ar[r] &
J \ar[r] &
B' \ar[r] &
B \ar[r] & 0 \\
0 \ar[r] &
I \ar[r] \ar[u] &
A' \ar[r] \ar[u] &
A \ar[r] \ar[u] & 0
}
$$
Maintenant, comme $\mathcal{I}$ est un $\mathcal{O}_X$-module quasi-cohérent,
nous avons $\mathcal{I} = (\widetilde I)^a$, voir
Descente, Définition \ref{descent-definition-structure-sheaf},
pour les notations, et
Descente, Proposition \ref{descent-proposition-equivalence-quasi-coherent},
pour la justification. Nous voyons donc que $J = I \otimes_A B$.
Enfin, notons que $A \to B$ est étale, car $U \to X$ est étale en tant que
changement de base du morphisme étale $U' \to X'$.
Nous concluons que $A' \to B'$ est étale par
Algèbre, Lemme \ref{algebra-lemma-lift-etale-infinitesimal}.
\end{proof}

\begin{lemma}
\label{lemma-thickening-equivalence}
Soit $S$ un schéma. Soit $X \subset X'$ un épaississement d'espaces algébriques
sur $S$. Le foncteur
$$
V' \longmapsto V = X \times_{X'} V'
$$
définit une équivalence de catégories
$X'_\etale \to X_\etale$.
\end{lemma}

\begin{proof}
Le foncteur $V' \mapsto V$ définit une équivalence de catégories
$X'_{spaces, \etale} \to X_{spaces, \etale}$, voir le
Théorème \ref{theorem-topological-invariance}.
Il suffit donc de montrer que $V$ est un schéma si et seulement si $V'$ est
un schéma. C'est le contenu du
Lemme \ref{lemma-thickening-scheme}.
\end{proof}

\noindent
Les épaississements du premier ordre se décrivent comme suit.

\begin{lemma}
\label{lemma-first-order-thickening}
Soit $S$ un schéma.
Soit $f : X \to B$ un morphisme d'espaces algébriques sur $S$.
Considérons une suite exacte courte
$$
0 \to \mathcal{I} \to \mathcal{A} \to \mathcal{O}_X \to 0
$$
de faisceaux sur $X_\etale$, où $\mathcal{A}$ est un faisceau de
$f^{-1}\mathcal{O}_B$-algèbres, $\mathcal{A} \to \mathcal{O}_X$ est une surjection
de faisceaux de $f^{-1}\mathcal{O}_B$-algèbres, et $\mathcal{I}$ son noyau.
Si
\begin{enumerate}
\item $\mathcal{I}$ est un idéal de carré nul dans $\mathcal{A}$, et
\item $\mathcal{I}$ est quasi-cohérent en tant que $\mathcal{O}_X$-module,
\end{enumerate}
alors il existe un épaississement du premier ordre
$X \subset X'$ sur $B$ et un isomorphisme
$\mathcal{O}_{X'} \to \mathcal{A}$ de $f^{-1}\mathcal{O}_B$-algèbres
compatible aux surjections vers $\mathcal{O}_X$.
\end{lemma}

\begin{proof}
Dans cette démonstration, nous reprenons certains des arguments employés dans les
démonstrations des
Lemmes \ref{lemma-first-order-thickening-surjective} et
\ref{lemma-thickening-scheme}.
Nous traitons d'abord le cas $B = S = \Spec(\mathbf{Z})$.
Soit $U$ un schéma affine et soit $U \to X$ étale.
Alors
$$
0 \to \mathcal{I}(U) \to \mathcal{A}(U) \to \mathcal{O}_X(U) \to 0
$$
est exacte, car $H^1(U_\etale, \mathcal{I}) = 0$ puisque
$\mathcal{I}$ est quasi-cohérent, voir
Descente, Proposition \ref{descent-proposition-same-cohomology-quasi-coherent},
et Cohomologie des schémas, Lemme
\ref{coherent-lemma-quasi-coherent-affine-cohomology-zero}.
Si $V \to U$ est un morphisme d'objets affines de $X_{spaces, \etale}$,
alors
$$
\mathcal{I}(V) = \mathcal{I}(U) \otimes_{\mathcal{O}_X(U)} \mathcal{O}_X(V)
$$
puisque $\mathcal{I}$ est un $\mathcal{O}_X$-module quasi-cohérent, voir
Descente, Proposition \ref{descent-proposition-equivalence-quasi-coherent}.
Ainsi $\mathcal{A}(U) \to \mathcal{A}(V)$ est un homomorphisme d'anneaux
étale, voir
Algèbre, Lemme \ref{algebra-lemma-lift-etale-infinitesimal}.
Nous voyons donc que
$$
U \longmapsto U' = \Spec(\mathcal{A}(U))
$$
est un foncteur de $X_{affine, \etale}$ vers la catégorie des schémas
affines et des morphismes étales. En fait, nous affirmons que ce foncteur peut
s'étendre en un foncteur $U \mapsto U'$ sur tout $X_\etale$.
Pour le voir, si $U$ est un objet de $X_\etale$, notons que
$$
0 \to \mathcal{I}|_{U_{Zar}} \to \mathcal{A}|_{U_{Zar}} \to
\mathcal{O}_X|_{U_{Zar}} \to 0
$$
et que $\mathcal{I}|_{U_{Zar}}$ est un faisceau quasi-cohérent sur $U$, voir
Descente,
Proposition \ref{descent-proposition-equivalence-quasi-coherent-functorial}.
Par conséquent, d'après
Compléments sur les morphismes, Lemme \ref{more-morphisms-lemma-first-order-thickening},
nous obtenons un épaississement du premier ordre de schémas $U \subset U'$ tel que
$\mathcal{O}_{U'}$ soit isomorphe à $\mathcal{A}|_{U_{Zar}}$. Il est clair que
cette construction est compatible avec la construction affine ci-dessus.

\medskip\noindent
Choisissons une présentation $X = U/R$, voir
Espaces, Définition \ref{spaces-definition-presentation},
de sorte que $s, t : R \to U$ définissent une relation d'équivalence étale.
En appliquant le foncteur ci-dessus, nous obtenons une relation d'équivalence
étale $s', t' : R' \to U'$ dans les schémas. Considérons l'espace algébrique
$X' = U'/R'$ (voir
Espaces, Théorème \ref{spaces-theorem-presentation}).
Le morphisme $X = U/R \to U'/R' = X'$ est un épaississement du premier ordre.
Considérons $\mathcal{O}_{X'}$ comme un faisceau sur $X_\etale$.
Par construction, nous avons un isomorphisme
$$
\gamma :
\mathcal{O}_{X'}|_{U_\etale}
\longrightarrow
\mathcal{A}|_{U_\etale}
$$
tel que $s^{-1}\gamma$ coïncide avec $t^{-1}\gamma$ sur $R_\etale$.
Par conséquent, d'après
Propriétés des espaces, Lemme \ref{spaces-properties-lemma-descent-sheaf},
cela implique que $\gamma$ provient d'un unique isomorphisme
$\mathcal{O}_{X'} \to \mathcal{A}$, comme voulu.

\medskip\noindent
Pour traiter le cas d'un espace algébrique de base quelconque $B$, nous
construisons d'abord $X'$ comme espace algébrique sur $\mathbf{Z}$, ainsi que ci-dessus.
Nous utilisons ensuite l'isomorphisme $\mathcal{O}_{X'} \to \mathcal{A}$ pour
définir $f^{-1}\mathcal{O}_B \to \mathcal{O}_{X'}$. D'après le
Lemme \ref{lemma-first-order-thickening-maps},
cela définit un morphisme $X' \to B$ compatible avec le morphisme donné
$X \to B$, ce qui achève la démonstration.
\end{proof}

\begin{lemma}
\label{lemma-base-change-thickening}
Soit $S$ un schéma. Soit $Y \subset Y'$ un épaississement d'espaces algébriques
sur $S$. Soit $X' \to Y'$ un morphisme et posons $X = Y \times_{Y'} X'$.
Alors $(X \subset X') \to (Y \subset Y')$
est un morphisme d'épaississements. Si $Y \subset Y'$ est un épaississement du premier
ordre (respectivement d'ordre fini), alors $X \subset X'$ est un épaississement du premier
ordre (respectivement d'ordre fini).
\end{lemma}

\begin{proof}
Omis.
\end{proof}

\begin{lemma}
\label{lemma-composition-thickening}
Soit $S$ un schéma. Si $X \subset X'$ et $X' \subset X''$ sont des
épaississements d'espaces algébriques sur $S$, alors $X \subset X''$ en est un également.
\end{lemma}

\begin{proof}
Omis.
\end{proof}

\begin{lemma}
\label{lemma-descending-property-thickening}
La propriété d'être un épaississement est locale pour la topologie fpqc.
Il en va de même pour les épaississements du premier ordre.
\end{lemma}

\begin{proof}
L'énoncé signifie ceci : soit $S$ un schéma et soit
$X \to X'$ un morphisme d'espaces algébriques sur $S$.
Soit $\{g_i : X'_i \to X'\}$ un recouvrement fpqc d'espaces algébriques
tel que le changement de base $X_i \to X'_i$ soit un épaississement pour tout $i$.
Alors $X \to X'$ est un épaississement. Comme les morphismes $g_i$ sont conjointement
surjectifs, nous concluons que $X \to X'$ est surjectif. D'après
Descente sur les espaces, Lemme
\ref{spaces-descent-lemma-descending-property-closed-immersion},
nous concluons que $X \to X'$ est une immersion fermée.
Ainsi $X \to X'$ est un épaississement. Nous omettons la démonstration dans le
cas des épaississements du premier ordre.
\end{proof}








\section{Morphismes d'épaississements}
\label{section-morphisms-thickenings}

\noindent
Si $(f, f') : (X \subset X') \to (Y \subset Y')$ est un morphisme
d'épaississements d'espaces algébriques, les propriétés de $f$
se transmettent souvent à $f'$. Il existe plusieurs variantes.

\begin{lemma}
\label{lemma-thicken-property-morphisms}
Soit $S$ un schéma. Soit $(f, f') : (X \subset X') \to (Y \subset Y')$
un morphisme d'épaississements d'espaces algébriques sur $S$. Alors
\begin{enumerate}
\item $f$ est un morphisme affine si et seulement si $f'$ est un morphisme affine,
\item $f$ est un morphisme surjectif si et seulement si $f'$ est un morphisme surjectif,
\item $f$ est quasi-compact si et seulement si $f'$ est quasi-compact,
\item $f$ est universellement fermé si et seulement si $f'$ est universellement fermé,
\item $f$ est entier si et seulement si $f'$ est entier,
\item $f$ est (quasi-)séparé si et seulement si $f'$ est (quasi-)séparé,
\item $f$ est universellement injectif si et seulement si $f'$ est universellement injectif,
\item $f$ est universellement ouvert si et seulement si $f'$ est universellement ouvert,
\item $f$ est représentable si et seulement si $f'$ est représentable, et
\item ajouter d'autres propriétés ici.
\end{enumerate}
\end{lemma}

\begin{proof}
Observons que $Y \to Y'$ et $X \to X'$ sont des morphismes entiers et des
homéomorphismes universels. Cela entraîne immédiatement les assertions
(2), (3), (4), (7) et (8).
L'assertion (1) découle de
Limites d'espaces, Proposition \ref{spaces-limits-proposition-affine},
qui nous dit qu'il existe une correspondance bijective entre les
schémas affines étales sur $X$ et sur $X'$, et entre les schémas affines
étales sur $Y$ et sur $Y'$.
L'assertion (5) découle de (1) et (4) par
Morphismes d'espaces, Lemme
\ref{spaces-morphisms-lemma-integral-universally-closed}.
Enfin, notons que
$$
X \times_Y X = X \times_{Y'} X \to X \times_{Y'} X' \to X' \times_{Y'} X'
$$
est un épaississement (les deux flèches sont des épaississements par le
Lemme \ref{lemma-base-change-thickening}).
Ainsi, en appliquant (3) et (4) au morphisme
$(X \subset X') \to (X \times_Y X \to X' \times_{Y'} X')$
nous obtenons (6). Enfin, l'assertion (9) découle du fait qu'un espace
algébrique qui est un épaississement d'un schéma est encore un schéma, voir le
Lemme \ref{lemma-thickening-scheme}.
\end{proof}

\begin{lemma}
\label{lemma-thicken-property-morphisms-cartesian}
Soit $S$ un schéma. Soit $(f, f') : (X \subset X') \to (Y \subset Y')$ un
morphisme d'épaississements d'espaces algébriques sur $S$ tel que
$X = Y \times_{Y'} X'$. Si $X \subset X'$ est un épaississement d'ordre fini, alors
\begin{enumerate}
\item $f$ est une immersion fermée si et seulement si $f'$ est une immersion fermée,
\item $f$ est localement de type fini si et seulement si $f'$ est
localement de type fini,
\item $f$ est localement quasi-fini si et seulement si $f'$ est localement
quasi-fini,
\item $f$ est localement de type fini de dimension relative $d$ si et
seulement si $f'$ est localement de type fini de dimension relative $d$,
\item $\Omega_{X/Y} = 0$ si et seulement si $\Omega_{X'/Y'} = 0$,
\item $f$ est non ramifié si et seulement si $f'$ est non ramifié,
\item $f$ est propre si et seulement si $f'$ est propre,
\item $f$ est un morphisme fini si et seulement si $f'$ est un morphisme fini,
\item $f$ est un monomorphisme si et seulement si $f'$ est un monomorphisme,
\item $f$ est une immersion si et seulement si $f'$ est une immersion, et
\item ajouter d'autres propriétés ici.
\end{enumerate}
\end{lemma}

\begin{proof}
Choisissons un schéma $V'$ et un morphisme étale surjectif $V' \to Y'$.
Choisissons un schéma $U'$ et un morphisme étale surjectif
$U' \to X' \times_{Y'} V'$. Posons $V = Y \times_{Y'} V'$ et
$U = X \times_{X'} U'$. Pour les propriétés des morphismes qui sont locales
pour la topologie étale, nous pouvons alors nous ramener au morphisme d'épaississements de schémas
$(U \subset U') \to (V \subset V')$ et appliquer Compléments sur les morphismes, Lemme
\ref{more-morphisms-lemma-thicken-property-morphisms-cartesian}.
Cela démontre (2), (3), (4), (5) et (6).

\medskip\noindent
Les propriétés des morphismes figurant en (1), (7), (8), (9) et (10) sont stables
par changement de base ; par conséquent, si $f'$ possède la propriété $\mathcal{P}$, alors
$f$ la possède également. Voir
Espaces, Lemme \ref{spaces-lemma-base-change-immersions},
et
Morphismes d'espaces, Lemmes
\ref{spaces-morphisms-lemma-base-change-proper},
\ref{spaces-morphisms-lemma-base-change-integral}, et
\ref{spaces-morphisms-lemma-base-change-monomorphism}.

\medskip\noindent
Le sens intéressant de (1), (7), (8), (9) et (10) consiste à supposer
que $f$ possède la propriété et à en déduire que $f'$ la possède aussi.
Par récurrence sur l'ordre de l'épaississement, nous pouvons
supposer que $Y \subset Y'$ est un épaississement du premier ordre, voir la
discussion ci-dessus sur les épaississements d'ordre fini.

\medskip\noindent
Démonstration de (1). Choisissons un schéma $V'$ et un morphisme étale surjectif
$V' \to Y'$. Posons $V = Y \times_{Y'} V'$, $U' = X' \times_{Y'} V'$ et
$U = X \times_Y V$. Alors $U \to V$ est une immersion fermée, ce qui
implique que $U$ est un schéma, et entraîne à son tour que $U'$ est
un schéma (Lemme \ref{lemma-thickening-scheme}). Nous pouvons donc appliquer
le lemme dans le cas des schémas
(Compléments sur les morphismes, Lemme
\ref{more-morphisms-lemma-thicken-property-morphisms-cartesian})
à $(U \subset U') \to (V \subset V')$ pour conclure.

\medskip\noindent
Démonstration de (7). Elle résulte de la combinaison de (2) avec les
résultats du Lemme \ref{lemma-thicken-property-morphisms}
et du fait qu'être propre équivaut à être quasi-compact $+$
séparé $+$ localement de type fini $+$ universellement fermé.

\medskip\noindent
Démonstration de (8). Elle résulte de la combinaison de (2) avec les
résultats du Lemme \ref{lemma-thicken-property-morphisms}
et du fait qu'être fini équivaut à être entier $+$ localement
de type fini (Morphismes, Lemme \ref{morphisms-lemma-finite-integral}).

\medskip\noindent
Démonstration de (9). Comme $f$ est un monomorphisme, nous avons $X = X \times_Y X$.
Nous pouvons appliquer les résultats déjà démontrés au morphisme
d'épaississements $(X \subset X') \to (X \times_Y X \subset X' \times_{Y'} X')$.
Nous concluons par (1) que $X' \to X' \times_{Y'} X'$ est une immersion fermée.
En fait, c'est un épaississement du premier ordre, car l'idéal qui définit
l'immersion fermée $X' \to X' \times_{Y'} X'$ est contenu dans l'image inverse
de l'idéal $\mathcal{I} \subset \mathcal{O}_{Y'}$ qui définit $Y$ dans $Y'$.
En effet, $X = X \times_Y X = (X' \times_{Y'} X') \times_{Y'} Y$ est contenu
dans $X'$. Le faisceau conormal de l'immersion fermée
$\Delta : X' \to X' \times_{Y'} X'$ est égal à $\Omega_{X'/Y'}$
(c'est l'analogue de
Morphismes, Lemme \ref{morphisms-lemma-differentials-diagonal},
pour les espaces algébriques, et cela découle soit d'une localisation étale,
soit de la combinaison des
Lemmes \ref{lemma-differentials-relative-immersion-section} et
\ref{lemma-differential-product} ; certains détails sont omis).
Il suffit donc de montrer que $\Omega_{X'/Y'} = 0$, ce qui découle de (5)
et de l'énoncé correspondant pour $X/Y$.

\medskip\noindent
Démonstration de (10). Si $f : X \to Y$ est une immersion, elle se factorise sous la forme
$X \to V \to Y$, où $V \to Y$ est un sous-espace ouvert et $X \to V$ une
immersion fermée, voir
Morphismes d'espaces, Remarque \ref{spaces-morphisms-remark-immersion}.
Soit $V' \subset Y'$ le sous-espace ouvert dont
l'espace topologique sous-jacent $|V'|$ est égal à $|V| \subset |Y| = |Y'|$.
Alors $X' \to Y'$ se factorise par $V'$, et nous concluons par (1) que $X' \to V'$
est une immersion fermée. Cela achève la démonstration.
\end{proof}

\noindent
Le lemme suivant est une variante du précédent. Au lieu de supposer
que les épaississements en jeu sont d'ordre fini (ce qui permet de transférer
de $f$ à $f'$ la propriété d'être localement de type fini),
nous supposons plutôt que $f$ et $f'$ sont tous deux localement
de type fini.

\begin{lemma}
\label{lemma-properties-that-extend-over-thickenings}
Soit $S$ un schéma. Soit $(f, f') : (X \subset X') \to (Y \to Y')$ un
morphisme d'épaississements d'espaces algébriques sur $S$. Supposons que $f$ et $f'$
soient localement de type fini et que $X = Y \times_{Y'} X'$. Alors
\begin{enumerate}
\item $f$ est localement quasi-fini si et seulement si $f'$ est localement quasi-fini,
\item $f$ est fini si et seulement si $f'$ est fini,
\item $f$ est une immersion fermée si et seulement si $f'$ est une immersion fermée,
\item $\Omega_{X/Y} = 0$ si et seulement si $\Omega_{X'/Y'} = 0$,
\item $f$ est non ramifié si et seulement si $f'$ est non ramifié,
\item $f$ est un monomorphisme si et seulement si $f'$ est un monomorphisme,
\item $f$ est une immersion si et seulement si $f'$ est une immersion,
\item $f$ est propre si et seulement si $f'$ est propre, et
\item ajouter d'autres propriétés ici.
\end{enumerate}
\end{lemma}

\begin{proof}
Choisissons un schéma $V'$ et un morphisme étale surjectif $V' \to Y'$.
Choisissons un schéma $U'$ et un morphisme étale surjectif
$U' \to X' \times_{Y'} V'$. Posons $V = Y \times_{Y'} V'$ et
$U = X \times_{X'} U'$. Pour les propriétés des morphismes qui sont locales
pour la topologie étale, nous pouvons alors nous ramener au morphisme d'épaississements de schémas
$(U \subset U') \to (V \subset V')$ et appliquer
Compléments sur les morphismes, Lemme
\ref{more-morphisms-lemma-properties-that-extend-over-thickenings}.
Cela démontre (1), (4) et (5).

\medskip\noindent
Les propriétés figurant en (2), (3), (6), (7) et (8) sont stables
par changement de base ; par conséquent, si $f'$ possède la propriété $\mathcal{P}$, alors
$f$ la possède également. Voir Espaces, Lemme \ref{spaces-lemma-base-change-immersions},
et
Morphismes d'espaces, Lemmes
\ref{spaces-morphisms-lemma-base-change-proper},
\ref{spaces-morphisms-lemma-base-change-integral}, et
\ref{spaces-morphisms-lemma-base-change-monomorphism}.
Ainsi, dans chaque cas, il suffit de démontrer que, si $f$ possède
la propriété voulue, alors $f'$ la possède aussi.

\medskip\noindent
Le cas (2) découle du cas (5) du Lemme \ref{lemma-thicken-property-morphisms}
et du fait que les morphismes finis sont exactement
les morphismes entiers qui sont localement de type fini
(Morphismes d'espaces, Lemme \ref{spaces-morphisms-lemma-finite-integral}).

\medskip\noindent
Cas (3). Cela découle immédiatement de
Limites d'espaces, Lemme
\ref{spaces-limits-lemma-check-closed-infinitesimally}.

\medskip\noindent
Démonstration de (6). Comme $f$ est un monomorphisme, nous avons $X = X \times_Y X$.
Nous pouvons appliquer les résultats déjà démontrés au morphisme d'épaississements
$(X \subset X') \to (X \times_Y X \subset X' \times_{Y'} X')$.
Nous concluons par (3) que $\Delta_{X'/Y'} : X' \to X' \times_{Y'} X'$
est une immersion fermée. En fait, $\Delta_{X'/Y'}$ induit une bijection
$|X'| \to |X' \times_{Y'} X'|$, donc $\Delta_{X'/Y'}$ est un épaississement.
D'autre part, $\Delta_{X'/Y'}$ est localement de présentation finie d'après
Morphismes d'espaces, Lemme
\ref{spaces-morphisms-lemma-diagonal-morphism-finite-type}.
Autrement dit, $\Delta_{X'/Y'}(X')$ est défini par
un faisceau quasi-cohérent d'idéaux
$\mathcal{J} \subset \mathcal{O}_{X' \times_{Y'} X'}$ de type fini.
Comme $\Omega_{X'/Y'} = 0$ par (5), nous voyons que
le faisceau conormal de $X' \to X' \times_{Y'} X'$ est nul.
(Le faisceau conormal de l'immersion fermée $\Delta_{X'/Y'}$ est égal à
$\Omega_{X'/Y'}$ ; c'est l'analogue de
Morphismes, Lemme \ref{morphisms-lemma-differentials-diagonal},
pour les espaces algébriques, et cela découle soit d'une localisation étale,
soit de la combinaison des
Lemmes \ref{lemma-differentials-relative-immersion-section} et
\ref{lemma-differential-product} ; certains détails sont omis.)
Autrement dit, $\mathcal{J}/\mathcal{J}^2 = 0$.
Cela implique que $\Delta_{X'/Y'}$ est un isomorphisme, par exemple
par Algèbre, Lemme \ref{algebra-lemma-ideal-is-squared-union-connected}.

\medskip\noindent
Démonstration de (7). Si $f : X \to Y$ est une immersion, elle se factorise sous la forme
$X \to V \to Y$, où $V \to Y$ est un sous-espace ouvert et $X \to V$ une
immersion fermée, voir
Morphismes d'espaces, Remarque \ref{spaces-morphisms-remark-immersion}.
Soit $V' \subset Y'$ le sous-espace ouvert dont
l'espace topologique sous-jacent $|V'|$ est égal à $|V| \subset |Y| = |Y'|$.
Alors $X' \to Y'$ se factorise par $V'$, et nous concluons par (3) que $X' \to V'$
est une immersion fermée.

\medskip\noindent
Le cas (8) découle du Lemme \ref{lemma-thicken-property-morphisms}
et de la définition des morphismes propres comme morphismes quasi-compacts,
universellement fermés et séparés qui sont localement de type fini.
\end{proof}










\section{Groupes de Picard des épaississements}
\label{section-picard-group-thickening}

\noindent
Quelques éléments sur les groupes de Picard des épaississements.

\begin{lemma}
\label{lemma-picard-group-first-order-thickening}
Soit $S$ un schéma. Soit $X \subset X'$ un épaississement du premier ordre
d'espaces algébriques sur $S$, dont le faisceau d'idéaux est $\mathcal{I}$.
Il existe alors une suite exacte canonique
$$
\xymatrix{
0 \ar[r] &
H^0(X, \mathcal{I}) \ar[r] &
H^0(X', \mathcal{O}_{X'}^*) \ar[r] &
H^0(X, \mathcal{O}^*_X) \ar `r[d] `d[l] `l[llld] `d[dll] [dll] \\
& H^1(X, \mathcal{I}) \ar[r] &
\Pic(X') \ar[r] &
\Pic(X) \ar `r[d] `d[l] `l[llld] `d[dll] [dll] \\
& H^2(X, \mathcal{I}) \ar[r] & \ldots \ar[r] & \ldots
}
$$
de groupes abéliens.
\end{lemma}

\begin{proof}
Rappelons que $X_\etale = X'_\etale$, voir le
Lemme \ref{lemma-thickening-equivalence} et, plus généralement, la
discussion de la Section \ref{section-thickenings}.
La suite du lemme est la suite exacte longue de cohomologie
associée à la suite exacte courte de faisceaux de groupes abéliens
$$
0 \to \mathcal{I} \to \mathcal{O}_{X'}^* \to \mathcal{O}_X^* \to 0
$$
sur $X_\etale$, où le premier morphisme envoie une section locale $f$ de $\mathcal{I}$
sur la section inversible $1 + f$ de $\mathcal{O}_{X'}$.
Nous utilisons aussi l'identification du groupe de Picard d'un
site annelé au premier groupe de cohomologie du faisceau
des éléments inversibles, voir
Cohomologie sur les sites, Lemme \ref{sites-cohomology-lemma-h1-invertible}.
\end{proof}








\section{Voisinages infinitésimaux}
\label{section-first-order-infinitesimal-neighbourhood}

\noindent
Une construction naturelle d'épaississements d'ordre fini est la suivante.
Supposons que $i : Z \to X$ soit une immersion d'espaces algébriques. Choisissons un
sous-espace ouvert $U \subset X$ tel que $i$ identifie $Z$ au sous-espace
fermé $Z \subset U$ (voir
Morphismes d'espaces, Remarque \ref{spaces-morphisms-remark-immersion}).
Soit $\mathcal{I} \subset \mathcal{O}_U$ le
faisceau quasi-cohérent d'idéaux qui définit $Z$ dans $U$, voir
Morphismes d'espaces,
Lemme \ref{spaces-morphisms-lemma-closed-immersion-ideals}.
Pour $n \geq 1$, nous pouvons considérer le sous-espace fermé $Z_n \subset U$
défini par le faisceau quasi-cohérent d'idéaux $\mathcal{I}^{n + 1}$.

\begin{definition}
\label{definition-first-order-infinitesimal-neighbourhood}
Soit $i : Z \to X$ une immersion d'espaces algébriques.
\begin{enumerate}
\item Le {\it voisinage infinitésimal du premier ordre} de $Z$ dans $X$ est
l'épaississement du premier ordre $Z \subset Z_1$ sur $X$ décrit ci-dessus.
\item Le {\it voisinage infinitésimal d'ordre $n$} de $Z$ dans $X$ est
l'épaississement d'ordre $n$ $Z \subset Z_n$ sur $X$ décrit ci-dessus.
\end{enumerate}
\end{definition}

\noindent
Cet épaississement possède la propriété universelle suivante (ce qui dissipera
toute crainte que la construction ci-dessus ne dépende du choix de l'ouvert
$U$).

\begin{lemma}
\label{lemma-first-order-infinitesimal-neighbourhood}
Soit $i : Z \to X$ une immersion d'espaces algébriques.
\begin{enumerate}
\item Le voisinage infinitésimal du premier ordre $Z'$ de $Z$ dans $X$
possède la propriété universelle suivante :
pour tout diagramme commutatif
$$
\xymatrix{
Z \ar[d]_i & T \ar[l]^a \ar[d] \\
X & T' \ar[l]_b
}
$$
où $T \subset T'$ est un épaississement du premier ordre sur $X$, il existe
un unique morphisme $(a', a) : (T \subset T') \to (Z \subset Z')$
d'épaississements sur $X$.
\item Pour $n \geq 1$, le voisinage infinitésimal d'ordre $n$
$Z_n$ de $Z$ dans $X$ possède la propriété universelle suivante :
pour tout diagramme commutatif
$$
\xymatrix{
Z \ar[d]_i & T \ar[l]^a \ar[d] \\
X & T' \ar[l]_b
}
$$
où $T \subset T'$ est un épaississement d'ordre $n$ sur $X$, il existe
un unique morphisme $(a', a) : (T \subset T') \to (Z \subset Z_n)$
d'épaississements sur $X$.
\end{enumerate}
\end{lemma}

\begin{proof}
Nous ne démontrerons que (1).
Soit $U \subset X$ le sous-espace ouvert utilisé dans la construction de $Z'$,
c'est-à-dire un ouvert tel que $Z$ soit identifié à un sous-espace fermé de $U$
défini par le faisceau quasi-cohérent d'idéaux $\mathcal{I}$.
Comme $|T| = |T'|$, nous voyons que $|b|(|T'|) \subset |U|$. Nous pouvons donc
considérer $b$ comme un morphisme à valeurs dans $U$, voir
Propriétés des espaces,
Lemme \ref{spaces-properties-lemma-factor-through-open-subspace}.
Soit $\mathcal{J} \subset \mathcal{O}_{T'}$
le faisceau quasi-cohérent d'idéaux de carré nul qui définit $T$.
La commutativité du diagramme donne $b|_T = i \circ a$, où
$i : Z \to U$ est l'immersion fermée. Nous en déduisons que
$b^\sharp(b^{-1}\mathcal{I}) \subset \mathcal{J}$ d'après
Morphismes d'espaces,
Lemme \ref{spaces-morphisms-lemma-closed-immersion-ideals}.
Comme $T'$ est un épaississement du premier ordre de $T$, nous avons $\mathcal{J}^2 = 0$,
donc $b^\sharp(b^{-1}(\mathcal{I}^2)) = 0$. D'après
Morphismes d'espaces, Lemme \ref{spaces-morphisms-lemma-closed-immersion-ideals},
il s'ensuit que $b$ se factorise par $Z'$. En notant $a' : T' \to Z'$
cette factorisation, nous concluons.
\end{proof}

\begin{lemma}
\label{lemma-infinitesimal-neighbourhood-conormal}
Soit $i : Z \to X$ une immersion d'espaces algébriques.
Soit $Z \subset Z'$ le voisinage infinitésimal du premier ordre
de $Z$ dans $X$. Alors le diagramme
$$
\xymatrix{
Z \ar[r] \ar[d] & Z' \ar[d] \\
Z \ar[r] & X
}
$$
induit un morphisme de faisceaux conormaux
$\mathcal{C}_{Z/X} \to \mathcal{C}_{Z/Z'}$ d'après le
Lemme \ref{lemma-conormal-functorial}.
Ce morphisme est un isomorphisme.
\end{lemma}

\begin{proof}
Cela résulte immédiatement de la construction de $Z'$ ci-dessus.
\end{proof}





\section{Transformations formellement lisses, étales et non ramifiées}
\label{section-formally-smooth-etale-unramified}

\noindent
Rappelons qu'un homomorphisme d'anneaux $R \to A$ est dit
{\it formellement lisse}, resp.\ {\it formellement étale},
resp.\ {\it formellement non ramifié}
(voir Algèbre, Définition \ref{algebra-definition-formally-smooth},
resp.\ Définition \ref{algebra-definition-formally-etale},
resp.\ Définition \ref{algebra-definition-formally-unramified})
si, pour tout diagramme commutatif de flèches pleines
$$
\xymatrix{
A \ar[r] \ar@{-->}[rd] & B/I \\
R \ar[r] \ar[u] & B \ar[u]
}
$$
où $I \subset B$ est un idéal de carré nul, il
existe une flèche pointillée, resp.\ il en existe une unique, resp.\ il en existe
au plus une, qui rend le diagramme commutatif. Cela motive
l'analogue suivant pour les morphismes d'espaces algébriques et, plus
généralement, pour les foncteurs.

\begin{definition}
\label{definition-formally-smooth-etale-unramified}
Soit $S$ un schéma.
Soit $a : F \to G$ une transformation de foncteurs
$F, G : (\Sch/S)_{fppf}^{opp} \to \textit{Sets}$.
Considérons les diagrammes commutatifs de flèches pleines de la forme
$$
\xymatrix{
F \ar[d]_a & T \ar[d]^i \ar[l] \\
G & T' \ar[l] \ar@{-->}[lu]
}
$$
où $T$ et $T'$ sont des schémas affines et $i$ une immersion fermée
définie par un idéal de carré nul.
\begin{enumerate}
\item On dit que $a$ est {\it formellement lisse} si, pour tout
diagramme de flèches pleines comme ci-dessus, il existe une flèche pointillée qui rend le diagramme
commutatif\footnote{Ce n'est là qu'une des définitions possibles.
Une condition légèrement plus faible consisterait à exiger que
la flèche pointillée existe localement pour la topologie fppf sur $T'$. Cette notion plus faible
possède, en un certain sens, de meilleures propriétés formelles.}.
\item On dit que $a$ est {\it formellement étale} si, pour tout
diagramme de flèches pleines comme ci-dessus, il existe exactement une flèche pointillée qui rend le diagramme
commutatif.
\item On dit que $a$ est {\it formellement non ramifiée} si, pour tout
diagramme de flèches pleines comme ci-dessus, il existe au plus une flèche pointillée qui rend le diagramme
commutatif.
\end{enumerate}
\end{definition}

\begin{lemma}
\label{lemma-formally-etale-is-combination}
Soit $S$ un schéma.
Soit $a : F \to G$ une transformation de foncteurs
$F, G : (\Sch/S)_{fppf}^{opp} \to \textit{Sets}$.
Alors $a$ est formellement étale si et seulement si $a$ est à la fois formellement
lisse et formellement non ramifiée.
\end{lemma}

\begin{proof}
Cela résulte formellement de la définition.
\end{proof}

\begin{lemma}
\label{lemma-composition-formally-smooth-etale-unramified}
Composition.
\begin{enumerate}
\item Un composé de transformations de foncteurs formellement lisses est formellement
lisse.
\item Un composé de transformations de foncteurs formellement étales est formellement
étale.
\item Un composé de transformations de foncteurs formellement non ramifiées est
formellement non ramifié.
\end{enumerate}
\end{lemma}

\begin{proof}
C'est formel.
\end{proof}

\begin{lemma}
\label{lemma-base-change-formally-smooth-etale-unramified}
Soit $S$ un schéma appartenant à $\Sch_{fppf}$.
Soient $F, G, H : (\Sch/S)_{fppf}^{opp} \to \textit{Sets}$.
Soient $a : F \to G$, $b : H \to G$ des transformations de foncteurs.
Considérons le diagramme de produit fibré
$$
\xymatrix{
H \times_{b, G, a} F \ar[r]_-{b'} \ar[d]_{a'} & F \ar[d]^a \\
H \ar[r]^b & G
}
$$
\begin{enumerate}
\item Si $a$ est formellement lisse, alors le changement de base $a'$ est
formellement lisse.
\item Si $a$ est formellement étale, alors le changement de base $a'$ est
formellement étale.
\item Si $a$ est formellement non ramifiée, alors le changement de base $a'$ est
formellement non ramifiée.
\end{enumerate}
\end{lemma}

\begin{proof}
C'est formel.
\end{proof}

\begin{lemma}
\label{lemma-representable-property-formally-property}
Soit $S$ un schéma.
Soient $F, G : (\Sch/S)_{fppf}^{opp} \to \textit{Sets}$.
Soit $a : F \to G$ une transformation de foncteurs représentable.
\begin{enumerate}
\item Si $a$ est lisse, alors $a$ est formellement lisse.
\item Si $a$ est étale, alors $a$ est formellement étale.
\item Si $a$ est non ramifiée, alors $a$ est formellement non ramifiée.
\end{enumerate}
\end{lemma}

\begin{proof}
Considérons un diagramme commutatif de flèches pleines
$$
\xymatrix{
F \ar[d]_a & T \ar[d]^i \ar[l] \\
G & T' \ar[l] \ar@{-->}[lu]
}
$$
comme dans la
Définition \ref{definition-formally-smooth-etale-unramified}.
Alors $F \times_G T'$ est un schéma lisse (resp.\ étale, resp.\ non ramifié)
sur $T'$. Par conséquent, d'après
Compléments sur les morphismes, Lemme \ref{more-morphisms-lemma-smooth-formally-smooth}
(resp.\ Lemme \ref{more-morphisms-lemma-etale-formally-etale},
resp.\ Lemme \ref{more-morphisms-lemma-unramified-formally-unramified}),
nous pouvons tracer (resp.\ tracer de manière unique, resp.\ tracer d'au plus
une manière) la flèche pointillée du diagramme
$$
\xymatrix{
F \times_G T' \ar[d] & T \ar[d]^i \ar[l] \\
T' & T' \ar[l] \ar@{-->}[lu]
}
$$
et nous obtenons donc aussi l'assertion correspondante dans le premier diagramme.
\end{proof}

\begin{lemma}
\label{lemma-etale-on-top}
Soit $S$ un schéma appartenant à $\Sch_{fppf}$.
Soient $F, G, H : (\Sch/S)_{fppf}^{opp} \to \textit{Sets}$.
Soient $a : F \to G$, $b : G \to H$ des transformations de foncteurs.
Supposons que $a$ soit représentable, surjective et étale.
\begin{enumerate}
\item Si $b$ est formellement lisse, alors $b \circ a$ est formellement lisse.
\item Si $b$ est formellement étale, alors $b \circ a$ est formellement étale.
\item Si $b$ est formellement non ramifiée, alors $b \circ a$ est formellement non ramifiée.
\end{enumerate}
Réciproquement, considérons un diagramme commutatif de flèches pleines
$$
\xymatrix{
G \ar[d]_b & T \ar[d]^i \ar[l] \\
H & T' \ar[l] \ar@{-->}[lu]
}
$$
où $T'$ est un schéma affine sur $S$
et $i : T \to T'$ une immersion fermée définie par un idéal de carré nul.
\begin{enumerate}
\item[(4)] Si $b \circ a$ est formellement lisse, alors, pour tout $t \in T$,
il existe un morphisme étale de schémas affines $U' \to T'$ et un morphisme
$U' \to G$ tels que
$$
\xymatrix{
G \ar[d]_b & T \ar[l] & T \times_{T'} U' \ar[d] \ar[l]\\
H & T' \ar[l] & U' \ar[llu] \ar[l]
}
$$
le diagramme soit commutatif et $t$ appartienne à l'image de $U' \to T'$.
\item[(5)] Si $b \circ a$ est formellement non ramifiée, alors il existe au plus
une flèche pointillée dans le diagramme ci-dessus, c'est-à-dire que $b$ est formellement non ramifiée.
\item[(6)] Si $b \circ a$ est formellement étale, alors il existe exactement une
flèche pointillée dans le diagramme ci-dessus, c'est-à-dire que $b$ est formellement étale.
\end{enumerate}
\end{lemma}

\begin{proof}
Supposons que $b$ soit formellement lisse (resp.\ formellement étale,
resp.\ formellement non ramifiée). Puisqu'un morphisme étale est à la fois
lisse et non ramifié, nous voyons que $a$ est représentable et lisse
(resp.\ étale, resp. non ramifiée). Les assertions (1), (2) et (3)
résultent donc de la combinaison du
Lemme \ref{lemma-representable-property-formally-property}
et du
Lemme \ref{lemma-composition-formally-smooth-etale-unramified}.

\medskip\noindent
Supposons que $b \circ a$ soit formellement lisse. Considérons un diagramme
comme dans l'énoncé du lemme. Posons $W = F \times_G T$.
Par hypothèse, $W$ est un schéma étale surjectif sur $T$. D'après
Morphismes étales, Théorème \ref{etale-theorem-remarkable-equivalence},
il existe un schéma $W'$ étale sur $T'$ tel que $W = T \times_{T'} W'$.
Choisissons un sous-schéma ouvert affine $U' \subset W'$ tel que $t$ appartienne à
l'image de $U' \to T'$. Puisque $b \circ a$ est formellement
lisse, il existe un morphisme $U' \to F$ tel que
$$
\xymatrix{
F \ar[d]_{b \circ a} & W \ar[l] & T \times_{T'} U' \ar[d] \ar[l]\\
H & T' \ar[l] & U' \ar[llu] \ar[l]
}
$$
le diagramme soit commutatif. Le composé $U' \to F \to G$ fournit un
morphisme comme dans l'assertion (5) du lemme.

\medskip\noindent
Supposons que $f, g : T' \to G$ soient deux flèches pointillées qui complètent le
diagramme du lemme. Posons $W = F \times_G T$.
Par hypothèse, $W$ est un schéma étale surjectif sur $T$. D'après
Morphismes étales, Théorème \ref{etale-theorem-remarkable-equivalence},
il existe un schéma $W'$ étale sur $T'$ tel que $W = T \times_{T'} W'$.
Comme $a$ est formellement étale, les composés
$$
W' \to T' \xrightarrow{f} G
\quad\text{and}\quad
W' \to T' \xrightarrow{g} G
$$
se relèvent en des morphismes $f', g' : W' \to F$ (on relève sur les ouverts affines et l'on recolle par
unicité). Si maintenant $b \circ a : F \to H$ est formellement non ramifiée, alors
$f' = g'$ et donc $f = g$, puisque $W' \to T'$ est un recouvrement étale. Cela démontre
l'assertion (6) du lemme.

\medskip\noindent
Supposons que $b \circ a$ soit formellement étale. D'après l'assertion (4), nous
pouvons trouver, localement pour la topologie étale sur $T'$, une flèche pointillée qui complète le diagramme,
et cette flèche est unique d'après l'assertion (5). Nous pouvons donc recoller les
solutions locales pour obtenir l'assertion (6). Certains détails sont omis.
\end{proof}

\begin{remark}
\label{remark-tempting}
On pourrait être tenté de penser que, dans la situation du
Lemme \ref{lemma-etale-on-top},
nous avons
« $b$ formellement lisse » $\Leftrightarrow$ « $b \circ a$ formellement lisse ».
Cependant, cela est vraisemblablement faux en général.
\end{remark}

\begin{lemma}
\label{lemma-formally-permanence}
Soit $S$ un schéma.
Soient $F, G, H : (\Sch/S)_{fppf}^{opp} \to \textit{Sets}$.
Soient $a : F \to G$, $b : G \to H$ des transformations de foncteurs.
Supposons que $b$ soit formellement non ramifiée.
\begin{enumerate}
\item Si $b \circ a$ est formellement non ramifiée, alors $a$ est formellement non ramifiée.
\item Si $b \circ a$ est formellement étale, alors $a$ est formellement étale.
\item Si $b \circ a$ est formellement lisse, alors $a$ est formellement lisse.
\end{enumerate}
\end{lemma}

\begin{proof}
Soit $T \subset T'$ une immersion fermée de schémas affines définie par un idéal
de carré nul. Soient $g' : T' \to G$ et $f : T \to F$ tels que
$g'|_T = a \circ f$. Comme $b$ est formellement non ramifiée, il existe une correspondance
bijective entre
$$
\{f' : T' \to F \mid f = f'|_T\text{ and }a \circ f' = g'\}
$$
et
$$
\{f' : T' \to F \mid f = f'|_T\text{ and }b \circ a \circ f' = b \circ g'\}.
$$
Le lemme en résulte formellement.
\end{proof}








\section{Morphismes formellement non ramifiés}
\label{section-formally-unramified}

\noindent
Dans cette section, nous précisons ce que signifie, pour un morphisme d'espaces algébriques,
d'être formellement non ramifié.

\begin{definition}
\label{definition-formally-unramified}
Soit $S$ un schéma. Un morphisme $f : X \to Y$ d'espaces algébriques sur $S$
est dit {\it formellement non ramifié} s'il est formellement non ramifié en tant que
transformation de foncteurs au sens de la
Définition \ref{definition-formally-smooth-etale-unramified}.
\end{definition}

\noindent
Nous ne répéterons pas les résultats démontrés dans le cadre plus général des
transformations de foncteurs formellement non ramifiées à la
Section \ref{section-formally-smooth-etale-unramified}.
On peut caractériser cette propriété par l'annulation du
module des différentielles relatives, voir le
Lemme \ref{lemma-characterize-formally-unramified}.

\begin{lemma}
\label{lemma-formally-unramified}
Soit $S$ un schéma. Soit $f : X \to Y$ un morphisme d'espaces algébriques sur
$S$. Les assertions suivantes sont équivalentes :
\begin{enumerate}
\item $f$ est formellement non ramifié,
\item pour tout diagramme
$$
\xymatrix{
U \ar[d] \ar[r]_\psi & V \ar[d] \\
X \ar[r]^f & Y
}
$$
où $U$ et $V$ sont des schémas et les flèches verticales sont étales,
le morphisme de schémas $\psi$ est formellement non ramifié (au sens de
Compléments sur les morphismes,
Définition \ref{more-morphisms-definition-formally-unramified}), et
\item pour un tel diagramme dont les flèches verticales sont surjectives, le morphisme
$\psi$ est formellement non ramifié.
\end{enumerate}
\end{lemma}

\begin{proof}
Supposons $f$ formellement non ramifié. D'après le
Lemme \ref{lemma-representable-property-formally-property},
les morphismes $U \to X$ et $V \to Y$ sont formellement non ramifiés. D'après le
Lemme \ref{lemma-composition-formally-smooth-etale-unramified},
le composé $U \to Y$ est donc formellement non ramifié. Il résulte alors du
Lemme \ref{lemma-formally-permanence}
que $U \to V$ est formellement non ramifié. Ainsi, (1) implique (2), et (2)
implique trivialement (3).

\medskip\noindent
Supposons donné un diagramme comme en (3). D'après le
Lemme \ref{lemma-representable-property-formally-property},
le morphisme $V \to Y$ est formellement non ramifié. D'après le
Lemme \ref{lemma-composition-formally-smooth-etale-unramified},
le composé $U \to Y$ est donc formellement non ramifié. Il résulte alors du
Lemme \ref{lemma-etale-on-top}
que $X \to Y$ est formellement non ramifié, c'est-à-dire que (1) est vérifiée.
\end{proof}

\begin{lemma}
\label{lemma-formally-unramified-not-affine}
Soit $S$ un schéma.
Si $f : X \to Y$ est un morphisme formellement non ramifié d'espaces algébriques
sur $S$, alors, pour tout diagramme commutatif de flèches pleines
$$
\xymatrix{
X \ar[d]_f & T \ar[d]^i \ar[l] \\
Y & T' \ar[l] \ar@{-->}[lu]
}
$$
où $T \subset T'$ est un épaississement du premier ordre d'espaces algébriques
sur $S$, il existe au plus une flèche pointillée qui rend le diagramme commutatif.
Autrement dit, dans la
Définition \ref{definition-formally-unramified},
on peut supprimer la condition que $T$ soit un schéma affine.
\end{lemma}

\begin{proof}
Cela tient au fait qu'il existe un morphisme étale surjectif
$U' \to T'$, où $U'$ est une réunion disjointe de schémas affines (voir
Propriétés des espaces, Lemme
\ref{spaces-properties-lemma-cover-by-union-affines}),
et qu'un morphisme $T' \to X$ est déterminé par sa restriction à $U'$.
\end{proof}

\begin{lemma}
\label{lemma-composition-formally-unramified}
Un composé de morphismes formellement non ramifiés est formellement non ramifié.
\end{lemma}

\begin{proof}
C'est formel.
\end{proof}

\begin{lemma}
\label{lemma-base-change-formally-unramified}
Un changement de base d'un morphisme formellement non ramifié est formellement non ramifié.
\end{lemma}

\begin{proof}
C'est formel.
\end{proof}

\begin{lemma}
\label{lemma-characterize-formally-unramified}
Soit $S$ un schéma. Soit $f : X \to Y$ un morphisme d'espaces algébriques sur
$S$. Les assertions suivantes sont équivalentes :
\begin{enumerate}
\item $f$ est formellement non ramifié, et
\item $\Omega_{X/Y} = 0$.
\end{enumerate}
\end{lemma}

\begin{proof}
Cela résulte de la combinaison du
Lemme \ref{lemma-formally-unramified},
de Compléments sur les morphismes,
Lemme \ref{more-morphisms-lemma-formally-unramified-differentials},
et du
Lemme \ref{lemma-localize-differentials}.
\end{proof}

\begin{lemma}
\label{lemma-unramified-formally-unramified}
Soit $S$ un schéma.
Soit $f : X \to Y$ un morphisme d'espaces algébriques sur $S$.
Les assertions suivantes sont équivalentes :
\begin{enumerate}
\item le morphisme $f$ est non ramifié,
\item le morphisme $f$ est localement de type fini et $\Omega_{X/Y} = 0$, et
\item le morphisme $f$ est localement de type fini et formellement non ramifié.
\end{enumerate}
\end{lemma}

\begin{proof}
Choisissons un diagramme
$$
\xymatrix{
U \ar[d] \ar[r]_\psi & V \ar[d] \\
X \ar[r]^f & Y
}
$$
où $U$ et $V$ sont des schémas et les flèches verticales sont étales et
surjectives. Nous avons alors
\begin{align*}
f\text{ unramified}
& \Leftrightarrow
\psi\text{ unramified} \\
& \Leftrightarrow
\psi\text{ locally finite type and }\Omega_{U/V} = 0 \\
& \Leftrightarrow
f\text{ locally finite type and }\Omega_{X/Y} = 0 \\
& \Leftrightarrow
f\text{ locally finite type and formally unramified}
\end{align*}
Nous avons utilisé ici
Morphismes, Lemme \ref{morphisms-lemma-unramified-omega-zero}, et le
Lemme \ref{lemma-characterize-formally-unramified}.
\end{proof}

\begin{lemma}
\label{lemma-universally-injective-unramified}
Soit $S$ un schéma.
Soit $f : X \to Y$ un morphisme d'espaces algébriques sur $S$.
Les assertions suivantes sont équivalentes :
\begin{enumerate}
\item $f$ est non ramifié et est un monomorphisme,
\item $f$ est non ramifié et universellement injectif,
\item $f$ est localement de type fini et est un monomorphisme,
\item $f$ est universellement injectif, localement de type fini et
formellement non ramifié.
\end{enumerate}
De plus, dans ce cas, $f$ est aussi représentable, séparé et
localement quasi-fini.
\end{lemma}

\begin{proof}
Nous avons vu au
Lemme \ref{lemma-unramified-formally-unramified}
qu'être formellement non ramifié et localement de type fini équivaut
à être non ramifié.
Ainsi, (4) est équivalente à (2).
Un monomorphisme est certainement formellement non ramifié ; ainsi, (3) implique (4).
Il est clair que (1) implique (3). Enfin, si (2) est vérifiée, alors
$\Delta : X \to X \times_Y X$ est à la fois une immersion ouverte
(Morphismes d'espaces, Lemme
\ref{spaces-morphisms-lemma-diagonal-unramified-morphism})
et surjective
(Morphismes d'espaces, Lemme
\ref{spaces-morphisms-lemma-universally-injective}) ;
c'est donc un isomorphisme, c'est-à-dire que $f$ est un monomorphisme. Nous voyons ainsi que
(2) implique (1).
Enfin, nous voyons que $f$ est représentable, séparé et localement
quasi-fini d'après
Morphismes d'espaces, Lemmes
\ref{spaces-morphisms-lemma-monomorphism-loc-finite-type-loc-quasi-finite} et
\ref{spaces-morphisms-lemma-locally-quasi-finite-separated-representable}.
\end{proof}

\begin{lemma}
\label{lemma-characterize-closed-immersion}
Soit $S$ un schéma.
Soit $f : X \to Y$ un morphisme d'espaces algébriques sur $S$.
Les assertions suivantes sont équivalentes :
\begin{enumerate}
\item $f$ est une immersion fermée,
\item $f$ est universellement fermé, non ramifié et est un monomorphisme,
\item $f$ est universellement fermé, non ramifié et universellement injectif,
\item $f$ est universellement fermé, localement de type fini et est un monomorphisme,
\item $f$ est universellement fermé, universellement injectif, localement de
type fini et formellement non ramifié.
\end{enumerate}
\end{lemma}

\begin{proof}
L'équivalence de (2) -- (5) résulte immédiatement du
Lemme \ref{lemma-universally-injective-unramified}.
De plus, si (2) -- (5) sont vérifiées, alors $f$ est représentable.
De même, si (1) est vérifiée, alors $f$ est représentable.
Le résultat découle donc du cas des schémas, voir
Morphismes étales, Lemme \ref{etale-lemma-characterize-closed-immersion}.
\end{proof}





\section{Épaississements universels du premier ordre}
\label{section-universal-thickening}

\noindent
Soit $S$ un schéma.
Soit $h : Z \to X$ un morphisme d'espaces algébriques sur $S$.
Un {\it épaississement universel du premier ordre} de $Z$ sur $X$ est un
épaississement du premier ordre $Z \subset Z'$ sur $X$ tel que, pour
tout épaississement du premier ordre $T \subset T'$
sur $X$ et tout diagramme commutatif de flèches pleines
\begin{equation}
\label{equation-universal-first-order-thickening}
\vcenter{
\xymatrix{
& Z \ar[ld] & & T \ar[rd] \ar[ll]^a \\
Z' \ar[rrd] & & & & T' \ar@{..>}[llll]_{a'} \ar[lld]^b \\
 & & X
}
}
\end{equation}
il existe une unique flèche pointillée qui rend le diagramme commutatif.
Notons que, dans cette situation, $(a, a') : (T \subset T') \to (Z \subset Z')$
est un morphisme d'épaississements sur $X$. Ainsi, s'il existe un épaississement universel
du premier ordre, celui-ci est unique à isomorphisme unique près.
En général, un épaississement universel du premier ordre
n'existe pas, mais il existe lorsque $h$ est formellement non ramifié.
Avant de le démontrer, montrons qu'un épaississement universel du premier ordre
dans la catégorie des schémas l'est également dans la
catégorie des espaces algébriques.

\begin{lemma}
\label{lemma-check-universal-first-order-thickening}
Soit $S$ un schéma.
Soit $h : Z \to X$ un morphisme d'espaces algébriques sur $S$.
Soit $Z \subset Z'$ un épaississement du premier ordre sur $X$.
Les assertions suivantes sont équivalentes :
\begin{enumerate}
\item $Z \subset Z'$ est un épaississement universel du premier ordre,
\item pour tout diagramme (\ref{equation-universal-first-order-thickening})
où $T'$ est un schéma, il existe une unique flèche pointillée qui rend le diagramme commutatif, et
\item pour tout diagramme (\ref{equation-universal-first-order-thickening})
où $T'$ est un schéma affine, il existe une unique flèche pointillée qui rend le
diagramme commutatif.
\end{enumerate}
\end{lemma}

\begin{proof}
Les implications (1) $\Rightarrow$ (2) $\Rightarrow$ (3) sont formelles.
Supposons (3) et considérons un diagramme arbitraire
(\ref{equation-universal-first-order-thickening}).
Choisissons une présentation $T' = U'/R'$, voir
Espaces, Définition \ref{spaces-definition-presentation}.
Nous pouvons supposer que $U' = \coprod U'_i$ est une réunion disjointe
de schémas affines, de sorte que $R' = U' \times_{T'} U' = \coprod_{i, j} U'_i \times_T' U'_j$.
Pour chaque paire $(i, j)$, choisissons un recouvrement ouvert affine
$U'_i \times_T' U'_j = \bigcup_k R'_{ijk}$. Notons $U_i, R_{ijk}$
les produits fibrés avec $T$ au-dessus de $T'$. Alors chacun des couples
$U_i \subset U'_i$ et $R_{ijk} \subset R'_{ijk}$
est un épaississement du premier ordre de schémas affines.
Notons $a_i : U_i \to Z$, resp.\ $a_{ijk} : R_{ijk} \to Z$,
le composé de $a : T \to Z$ avec le morphisme
$U_i \to T$, resp.\ $R_{ijk} \to T$.
En appliquant (3) à $a_i : U_i \to Z$,
nous obtenons des morphismes uniques $a'_i : U'_i \to Z'$.
En appliquant (3) à $a_{ijk}$, nous voyons que les deux composés
$R'_{ijk} \to R'_i \to Z'$ et $R'_{ijk} \to R'_j \to Z'$
sont égaux. Ainsi, $a' = \coprod a'_i : U' = \coprod U'_i \to Z'$
descend au faisceau quotient $T' = U'/R'$, ce qui conclut.
\end{proof}

\begin{lemma}
\label{lemma-universal-thickening-over-formally-etale}
Soit $S$ un schéma.
Soient $Z \to Y \to X$ des morphismes d'espaces algébriques sur $S$.
Si $Z \subset Z'$ est un épaississement universel du premier ordre de
$Z$ sur $Y$ et si $Y \to X$ est formellement étale, alors $Z \subset Z'$ est
un épaississement universel du premier ordre de $Z$ sur $X$.
\end{lemma}

\begin{proof}
Cela est formel. En effet, d'après le
Lemme \ref{lemma-check-universal-first-order-thickening},
il suffit de considérer les diagrammes commutatifs de flèches pleines
(\ref{equation-universal-first-order-thickening})
où $T'$ est un schéma affine. Le composé
$T \to Z \to Y$ se relève de manière unique en $T' \to Y$, puisque $Y \to X$ est
supposé formellement étale. Par conséquent, le fait que
$Z \subset Z'$ soit un épaississement universel du premier ordre sur $Y$
fournit le morphisme voulu $a' : T' \to Z'$.
\end{proof}

\begin{lemma}
\label{lemma-etale-morphism-of-universal-thickenings}
Soit $S$ un schéma.
Soient $Z \to Y \to X$ des morphismes d'espaces algébriques sur $S$.
Supposons que $Z \to Y$ soit étale.
\begin{enumerate}
\item Si $Y \subset Y'$ est un épaississement universel du premier ordre de
$Y$ sur $X$, alors l'unique morphisme étale $Z' \to Y'$ tel
que $Z = Y \times_{Y'} Z'$ (voir le
Théorème \ref{theorem-topological-invariance})
est un épaississement universel du premier ordre de $Z$ sur $X$.
\item Si $Z \to Y$ est surjectif et si
$(Z \subset Z') \to (Y \subset Y')$ est un morphisme étale
d'épaississements du premier ordre sur $X$, tandis que $Z'$ est un épaississement universel du
premier ordre de $Z$ sur $X$, alors $Y'$ est un épaississement universel du
premier ordre de $Y$ sur $X$.
\end{enumerate}
\end{lemma}

\begin{proof}
Démonstration de (1). D'après le
Lemme \ref{lemma-check-universal-first-order-thickening},
il suffit de considérer les diagrammes commutatifs de flèches pleines
(\ref{equation-universal-first-order-thickening})
où $T'$ est un schéma affine. Le composé
$T \to Z \to Y$ se relève de manière unique en $T' \to Y'$, puisque $Y'$ est
l'épaississement universel du premier ordre. Le fait que
$Z' \to Y'$ soit étale implique alors (voir le
Lemme \ref{lemma-representable-property-formally-property})
que $T' \to Y'$ se relève en le
morphisme voulu $a' : T' \to Z'$.

\medskip\noindent
Démonstration de (2). Soit $T \subset T'$ un épaississement du premier ordre sur
$X$ et soit $a : T \to Y$ un morphisme. Posons $W = T \times_Y Z$
et notons $c : W \to Z$ la projection.
Soit $W' \to T'$ l'unique morphisme étale tel que
$W = T \times_{T'} W'$, voir le
Théorème \ref{theorem-topological-invariance}.
Notons que $W' \to T'$ est surjectif puisque $Z \to Y$ est surjectif.
Par hypothèse, nous obtenons un unique morphisme $c' : W' \to Z'$
sur $X$ et dont $c$ est la restriction à $W$. Par unicité, les deux restrictions
de $c'$ à $W' \times_{T'} W'$ sont égales (puisque les deux restrictions de
$c$ à $W \times_T W$ sont égales). Ainsi, $c'$ descend en un unique
morphisme $a' : T' \to Y'$, ce qui conclut.
\end{proof}

\begin{lemma}
\label{lemma-universal-thickening}
Soit $S$ un schéma.
Soit $h : Z \to X$ un morphisme formellement non ramifié d'espaces
algébriques sur $S$.
Il existe un épaississement universel du premier ordre $Z \subset Z'$ de
$Z$ sur $X$.
\end{lemma}

\begin{proof}
Choisissons un diagramme commutatif quelconque
$$
\xymatrix{
V \ar[d] \ar[r] & U \ar[d] \\
Z \ar[r] & X
}
$$
où $V$ et $U$ sont des schémas et les flèches verticales sont étales.
Notons que $V \to U$ est un morphisme formellement non ramifié de schémas, voir le
Lemme \ref{lemma-formally-unramified}.
En combinant le
Lemme \ref{lemma-check-universal-first-order-thickening}
et
Compléments sur les morphismes, Lemme \ref{more-morphisms-lemma-universal-thickening},
nous voyons qu'il existe un épaississement universel du premier ordre $V \subset V'$
de $V$ sur $U$. D'après le
Lemme \ref{lemma-universal-thickening-over-formally-etale}, partie (1),
$V'$ est un épaississement universel du premier ordre de $V$ sur $X$.

\medskip\noindent
Fixons un schéma $U$ et un morphisme étale surjectif $U \to X$.
L'argument ci-dessus montre que, pour tout $V \to Z$ étale, où $V$ est
un schéma tel que $V \to Z \to X$ se factorise par $U$, il existe un
épaississement universel du premier ordre $V \subset V'$ de $V$ sur $X$
(qui ne dépend toutefois pas de la factorisation choisie de $V \to X$
par $U$). Nous pouvons maintenant choisir $V$ tel que $V \to Z$ soit étale
surjectif (voir
Espaces, Lemme \ref{spaces-lemma-lift-morphism-presentations}).
Alors $R = V \times_Z V$ est un schéma étale sur $Z$ tel que
$R \to X$ se factorise lui aussi par $U$.
Nous obtenons donc des épaississements universels du premier ordre
$V \subset V'$ et $R \subset R'$ sur $X$.
Puisque $V \subset V'$ est un épaississement universel du premier ordre,
les deux projections $s, t : R \to V$ se relèvent en des morphismes
$s', t': R' \to V'$. D'après le
Lemme \ref{lemma-etale-morphism-of-universal-thickenings},
puisque $R'$ est l'épaississement universel du premier ordre de $R$ sur $X$,
ces morphismes sont étales.
Alors $(t', s') : R' \to V'$ est une relation d'équivalence étale,
et nous pouvons poser $Z' = V'/R'$. Puisque $V' \to Z'$ est étale surjectif
et que $v'$ est l'épaississement universel du premier ordre de $V$ sur $X$,
nous déduisons du
Lemme \ref{lemma-universal-thickening-over-formally-etale}, partie (2),
que $Z'$ est un épaississement universel du premier ordre de $Z$ sur $X$.
\end{proof}

\begin{definition}
\label{definition-universal-thickening}
Soit $S$ un schéma.
Soit $h : Z \to X$ un morphisme formellement non ramifié
d'espaces algébriques sur $S$.
\begin{enumerate}
\item L'{\it épaississement universel du premier ordre} de $Z$ sur $X$
est l'épaississement $Z \subset Z'$ construit au
Lemme \ref{lemma-universal-thickening}.
\item Le {\it faisceau conormal de $Z$ sur $X$} est le faisceau conormal
de $Z$ dans son épaississement universel du premier ordre $Z'$ sur $X$.
\end{enumerate}
Dans cette situation, nous notons souvent le faisceau conormal $\mathcal{C}_{Z/X}$.
\end{definition}

\noindent
Nous voyons ainsi qu'il existe une suite exacte courte de faisceaux
$$
0 \to \mathcal{C}_{Z/X} \to \mathcal{O}_{Z'} \to \mathcal{O}_Z \to 0
$$
sur $Z_\etale$, et $\mathcal{C}_{Z/X}$ est un $\mathcal{O}_Z$-module
quasi-cohérent. Le lemme suivant montre qu'il n'y a pas de
conflit entre cette définition et celle du cas où $Z \to X$
est une immersion.

\begin{lemma}
\label{lemma-immersion-universal-thickening}
Soit $S$ un schéma.
Soit $i : Z \to X$ une immersion d'espaces algébriques sur $S$. Alors
\begin{enumerate}
\item $i$ est formellement non ramifié,
\item l'épaississement universel du premier ordre de $Z$ sur $X$ est le voisinage
infinitésimal du premier ordre de $Z$ dans $X$ de la
Définition \ref{definition-first-order-infinitesimal-neighbourhood},
\item le faisceau conormal de $i$ au sens de la
Définition \ref{definition-conormal-sheaf}
coïncide avec le faisceau conormal de $i$ au sens de la
Définition \ref{definition-universal-thickening}.
\end{enumerate}
\end{lemma}

\begin{proof}
Une immersion d'espaces algébriques est, par définition, un morphisme représentable.
Par conséquent, d'après
Morphismes, Lemmes \ref{morphisms-lemma-open-immersion-unramified} et
\ref{morphisms-lemma-closed-immersion-unramified},
une immersion est non ramifiée (par le principe abstrait du chapitre
Espaces, Lemme
\ref{spaces-lemma-representable-transformations-property-implication}).
Elle est donc formellement non ramifiée d'après le
Lemme \ref{lemma-unramified-formally-unramified}.
Les autres assertions résultent de la combinaison des
Lemmes \ref{lemma-first-order-infinitesimal-neighbourhood} et
\ref{lemma-infinitesimal-neighbourhood-conormal}
avec les définitions.
\end{proof}

\begin{lemma}
\label{lemma-universal-thickening-unramified}
Soit $S$ un schéma.
Soit $Z \to X$ un morphisme formellement non ramifié d'espaces algébriques sur $S$.
Alors l'épaississement universel du premier ordre $Z'$ est formellement
non ramifié sur $X$.
\end{lemma}

\begin{proof}
Soit $T \subset T'$ un épaississement du premier ordre de schémas affines sur $X$.
Soit
$$
\xymatrix{
Z' \ar[d] & T \ar[l]^c \ar[d] \\
X & T' \ar[l] \ar[lu]^{a, b}
}
$$
un diagramme commutatif. Posons $T_0 = c^{-1}(Z) \subset T$ et
$T'_a = a^{-1}(Z)$ (au sens schématique).
Comme $Z'$ est un épaississement du premier ordre de $Z$, nous voyons que $T'$
est un épaississement du premier ordre de $T'_a$. De plus, puisque $c = a|_T$, nous avons
$T_0 = T \cap T'_a$ (au sens schématique). Comme $T'$ est un épaississement du premier ordre
de $T$, il s'ensuit que $T'_a$
est un épaississement du premier ordre de $T_0$. Or $a|_{T'_a}$ et $b|_{T'_a}$
sont des morphismes de $T'_a$ vers $Z'$ sur $X$ qui coïncident sur $T_0$ comme
morphismes vers $Z$. La propriété universelle de $Z'$ donne donc
$a|_{T'_a} = b|_{T'_a}$. Ainsi, $a$ et $b$ sont des morphismes depuis
l'épaississement du premier ordre $T'$ de $T'_a$ dont les restrictions à
$T'_a$ coïncident comme morphismes vers $Z$. En utilisant une nouvelle fois la propriété universelle de
$Z'$, nous obtenons $a = b$. Autrement dit, la propriété qui définit
un morphisme formellement non ramifié est vérifiée par $Z' \to X$, comme souhaité.
\end{proof}

\begin{lemma}
\label{lemma-universal-thickening-functorial}
Soit $S$ un schéma.
Considérons un diagramme commutatif d'espaces algébriques sur $S$
$$
\xymatrix{
Z \ar[r]_h \ar[d]_f & X \ar[d]^g \\
W \ar[r]^{h'} & Y
}
$$
où $h$ et $h'$ sont formellement non ramifiés. Soit $Z \subset Z'$ l'épaississement universel
du premier ordre de $Z$ sur $X$. Soit $W \subset W'$ l'épaississement universel
du premier ordre de $W$ sur $Y$. Il existe un morphisme canonique
$(f, f') : (Z, Z') \to (W, W')$ d'épaississements sur $Y$ qui s'insère dans
le diagramme commutatif suivant
$$
\xymatrix{
& & & Z' \ar[ld] \ar[d]^{f'} \\
Z \ar[rr] \ar[d]_f \ar[rrru] & & X \ar[d] & W' \ar[ld] \\
W \ar[rrru]|!{[rr];[rruu]}\hole \ar[rr] & & Y
}
$$
En particulier, le morphisme $(f, f')$ d'épaississements induit un morphisme
de faisceaux conormaux $f^*\mathcal{C}_{W/Y} \to \mathcal{C}_{Z/X}$.
\end{lemma}

\begin{proof}
La première assertion résulte immédiatement de la propriété universelle de $W'$.
Le morphisme induit sur les faisceaux conormaux est celui du
Lemme \ref{lemma-conormal-functorial}
appliqué à $(Z \subset Z') \to (W \subset W')$.
\end{proof}

\begin{lemma}
\label{lemma-universal-thickening-fibre-product}
Soit $S$ un schéma. Soit
$$
\xymatrix{
Z \ar[r]_h \ar[d]_f & X \ar[d]^g \\
W \ar[r]^{h'} & Y
}
$$
un diagramme de produit fibré d'espaces algébriques sur $S$, où
$h'$ est formellement non ramifié. Alors $h$ est formellement non ramifié et, si
$W \subset W'$ est l'épaississement universel du premier ordre de $W$ sur $Y$,
alors $Z = X \times_Y W \subset X \times_Y W'$ est l'épaississement universel
du premier ordre de $Z$ sur $X$. En particulier, le morphisme canonique
$f^*\mathcal{C}_{W/Y} \to \mathcal{C}_{Z/X}$ du
Lemme \ref{lemma-universal-thickening-functorial}
est surjectif.
\end{lemma}

\begin{proof}
Le morphisme $h$ est formellement non ramifié d'après le
Lemme \ref{lemma-base-change-formally-unramified}.
Il est clair que $X \times_Y W'$ est un épaississement du premier ordre.
On vérifie immédiatement qu'il possède la propriété universelle,
car $W'$ la possède (d'après les propriétés universelles des
produits fibrés). Voir le
Lemme \ref{lemma-conormal-functorial-flat}
pour comprendre pourquoi cela implique la surjectivité du morphisme de faisceaux conormaux.
\end{proof}

\begin{lemma}
\label{lemma-universal-thickening-fibre-product-flat}
Soit $S$ un schéma. Soit
$$
\xymatrix{
Z \ar[r]_h \ar[d]_f & X \ar[d]^g \\
W \ar[r]^{h'} & Y
}
$$
un diagramme de produit fibré d'espaces algébriques sur $S$, où
$h'$ est formellement non ramifié et $g$ plat. Dans ce cas, le morphisme correspondant
$Z' \to W'$ entre les épaississements universels du premier ordre est plat, et
$f^*\mathcal{C}_{W/Y} \to \mathcal{C}_{Z/X}$ est un isomorphisme.
\end{lemma}

\begin{proof}
La platitude est préservée par changement de base, voir
Morphismes d'espaces, Lemme \ref{spaces-morphisms-lemma-base-change-flat}.
La première assertion découle donc de la description de $W'$ au
Lemme \ref{lemma-universal-thickening-fibre-product}.
Il est clair que $X \times_Y W'$ est un épaississement du premier ordre.
On vérifie immédiatement qu'il possède la propriété universelle,
car $W'$ la possède (d'après les propriétés universelles des
produits fibrés). Voir le
Lemme \ref{lemma-conormal-functorial-flat}
pour comprendre pourquoi cela implique que le morphisme de faisceaux conormaux est un isomorphisme.
\end{proof}

\begin{lemma}
\label{lemma-universal-thickening-localize}
La formation des épaississements universels du premier ordre commute à la
localisation étale. Plus précisément, soit $h : Z \to X$ un morphisme formellement non ramifié
d'espaces algébriques sur un schéma de base $S$.
Considérons
$$
\xymatrix{
V \ar[d] \ar[r] & U \ar[d] \\
Z \ar[r] & X
}
$$
un diagramme commutatif dont les flèches verticales sont étales.
Soit $Z'$ l'épaississement universel du premier ordre de $Z$ sur $X$.
Alors $V \to U$ est formellement non ramifié, et l'épaississement universel du premier
ordre $V'$ de $V$ sur $U$ est étale sur $Z'$.
En particulier, $\mathcal{C}_{Z/X}|_V = \mathcal{C}_{V/U}$.
\end{lemma}

\begin{proof}
La première assertion résulte du
Lemme \ref{lemma-formally-unramified}.
La compatibilité des épaississements universels du premier ordre
résulte des
Lemmes \ref{lemma-universal-thickening-over-formally-etale} et
\ref{lemma-etale-morphism-of-universal-thickenings}.
\end{proof}

\begin{lemma}
\label{lemma-differentials-universally-unramified}
Soit $S$ un schéma. Soit $B$ un espace algébrique sur $S$.
Soit $h : Z \to X$ un morphisme formellement non ramifié d'espaces algébriques
sur $B$. Soit $Z \subset Z'$ l'épaississement universel du premier ordre de $Z$
sur $X$, de morphisme structural $h' : Z' \to X$. Le morphisme canonique
$$
\text{d}h' : (h')^*\Omega_{X/B} \to \Omega_{Z'/B}
$$
induit un isomorphisme
$h^*\Omega_{X/B} \to \Omega_{Z'/B} \otimes \mathcal{O}_Z$.
\end{lemma}

\begin{proof}
Le morphisme $c_{h'}$ est celui qui est défini au
Lemme \ref{lemma-functoriality-differentials}.
Si $i : Z \to Z'$ est l'immersion fermée donnée, alors
$i^*c_{h'}$ est un morphisme
$h^*\Omega_{X/S} \to \Omega_{Z'/S} \otimes \mathcal{O}_Z$.
Pour vérifier qu'il s'agit d'un isomorphisme, on se ramène au cas des schémas
par localisation étale, voir le
Lemme \ref{lemma-universal-thickening-localize}
et le
Lemme \ref{lemma-localize-differentials}.
Dans ce cas, le résultat est
Compléments sur les morphismes,
Lemme \ref{more-morphisms-lemma-differentials-universally-unramified}.
\end{proof}

\begin{lemma}
\label{lemma-universally-unramified-differentials-sequence}
Soit $S$ un schéma. Soit $B$ un espace algébrique sur $S$.
Soit $h : Z \to X$ un morphisme formellement non ramifié d'espaces
algébriques sur $B$.
Il existe une suite exacte canonique
$$
\mathcal{C}_{Z/X} \to h^*\Omega_{X/B} \to \Omega_{Z/B} \to 0.
$$
La première flèche est induite par $\text{d}_{Z'/B}$, où
$Z'$ est le voisinage universel du premier ordre de $Z$ sur $X$.
\end{lemma}

\begin{proof}
Nous savons qu'il existe une suite exacte canonique
$$
\mathcal{C}_{Z/Z'} \to
\Omega_{Z'/S} \otimes \mathcal{O}_Z \to
\Omega_{Z/S} \to 0.
$$
Cela résulte du
Lemme \ref{lemma-differentials-relative-immersion}.
Le résultat s'obtient donc en appliquant le
Lemme \ref{lemma-differentials-universally-unramified}.
\end{proof}

\begin{lemma}
\label{lemma-two-unramified-morphisms}
Soit $S$ un schéma. Soit
$$
\xymatrix{
Z \ar[r]_i \ar[rd]_j & X \ar[d] \\
& Y
}
$$
un diagramme commutatif d'espaces algébriques sur $S$,
où $i$ et $j$ sont formellement non ramifiés. Il existe alors une
suite exacte canonique
$$
\mathcal{C}_{Z/Y} \to
\mathcal{C}_{Z/X} \to
i^*\Omega_{X/Y} \to 0
$$
où la première flèche provient du
Lemme \ref{lemma-universal-thickening-functorial}
et la seconde du
Lemme \ref{lemma-universally-unramified-differentials-sequence}.
\end{lemma}

\begin{proof}
Les morphismes étant définis, la vérification de l'exactitude de la suite
se ramène au cas des schémas par localisation étale, voir le
Lemme \ref{lemma-universal-thickening-localize}
et le
Lemme \ref{lemma-localize-differentials}.
Dans ce cas, le résultat est
Compléments sur les morphismes,
Lemme \ref{more-morphisms-lemma-two-unramified-morphisms}.
\end{proof}

\begin{lemma}
\label{lemma-transitivity-conormal-unramified}
Soit $S$ un schéma.
Soient $Z \to Y \to X$ des morphismes formellement non ramifiés
d'espaces algébriques sur $S$.
\begin{enumerate}
\item Si $Z \subset Z'$ est l'épaississement universel du premier ordre de $Z$
sur $X$ et si $Y \subset Y'$ est l'épaississement universel du premier ordre de $Y$
sur $X$, alors il existe un morphisme $Z' \to Y'$ et $Y \times_{Y'} Z'$ est
l'épaississement universel du premier ordre de $Z$ sur $Y$.
\item Il existe une suite exacte canonique
$$
i^*\mathcal{C}_{Y/X} \to
\mathcal{C}_{Z/X} \to
\mathcal{C}_{Z/Y} \to 0
$$
où les morphismes proviennent du
Lemme \ref{lemma-universal-thickening-functorial}
et $i : Z \to Y$ est le premier morphisme.
\end{enumerate}
\end{lemma}

\begin{proof}
Le morphisme $h : Z' \to Y'$ de (1) provient du
Lemme \ref{lemma-universal-thickening-functorial}.
L'assertion selon laquelle $Y \times_{Y'} Z'$ est l'épaississement universel du premier ordre
de $Z$ sur $Y$ résulte immédiatement des propriétés universelles
de $Z'$ et de $Y'$. D'après le
Lemme \ref{lemma-transitivity-conormal},
nous avons une suite exacte
$$
(i')^*\mathcal{C}_{Y \times_{Y'} Z'/Z'} \to
\mathcal{C}_{Z/Z'} \to
\mathcal{C}_{Z/Y \times_{Y'} Z'} \to 0
$$
où $i' : Z \to Y \times_{Y'} Z'$ est le morphisme donné. D'après le
Lemme \ref{lemma-conormal-functorial-flat},
il existe une surjection
$h^*\mathcal{C}_{Y/Y'} \to \mathcal{C}_{Y \times_{Y'} Z'/Z'}$.
Combinée aux égalités
$\mathcal{C}_{Y/Y'} = \mathcal{C}_{Y/X}$,
$\mathcal{C}_{Z/Z'} = \mathcal{C}_{Z/X}$, et
$\mathcal{C}_{Z/Y \times_{Y'} Z'} = \mathcal{C}_{Z/Y}$,
cette surjection démontre le lemme.
\end{proof}













\section{Morphismes formellement étales}
\label{section-formally-etale}

\noindent
Dans cette section, nous explicitons ce que signifie pour un morphisme d'espaces algébriques
d'être formellement étale.

\begin{definition}
\label{definition-formally-etale}
Soit $S$ un schéma. Un morphisme $f : X \to Y$ d'espaces algébriques sur $S$
est dit {\it formellement étale} s'il est formellement étale en tant que
transformation de foncteurs au sens de la
Définition \ref{definition-formally-smooth-etale-unramified}.
\end{definition}

\noindent
Nous ne redonnerons pas les résultats démontrés dans le cadre plus général des
transformations de foncteurs formellement étales dans la
Section \ref{section-formally-smooth-etale-unramified}.

\begin{lemma}
\label{lemma-formally-etale}
Soit $S$ un schéma. Soit $f : X \to Y$ un morphisme d'espaces algébriques sur
$S$. Les assertions suivantes sont équivalentes :
\begin{enumerate}
\item $f$ est formellement étale,
\item pour tout diagramme
$$
\xymatrix{
U \ar[d] \ar[r]_\psi & V \ar[d] \\
X \ar[r]^f & Y
}
$$
où $U$ et $V$ sont des schémas et où les flèches verticales sont étales,
le morphisme de schémas $\psi$ est formellement étale (au sens de
Compléments sur les morphismes,
Définition \ref{more-morphisms-definition-formally-etale}), et
\item il existe un tel diagramme dont les flèches verticales sont surjectives et dans lequel
$\psi$ est formellement étale.
\end{enumerate}
\end{lemma}

\begin{proof}
Supposons $f$ formellement étale. D'après le
Lemme \ref{lemma-representable-property-formally-property},
les morphismes $U \to X$ et $V \to Y$ sont formellement étales. Ainsi, d'après le
Lemme \ref{lemma-composition-formally-smooth-etale-unramified},
la composée $U \to Y$ est formellement étale. Il résulte alors du
Lemme \ref{lemma-formally-permanence}
que $U \to V$ est formellement étale. Ainsi, (1) implique (2), et (2)
implique trivialement (3).

\medskip\noindent
Supposons donné un diagramme comme en (3). D'après le
Lemme \ref{lemma-representable-property-formally-property},
le morphisme $V \to Y$ est formellement étale. Ainsi, d'après le
Lemme \ref{lemma-composition-formally-smooth-etale-unramified},
la composée $U \to Y$ est formellement étale. Il résulte alors du
Lemme \ref{lemma-etale-on-top}
que $X \to Y$ est formellement étale, c'est-à-dire que (1) est vérifiée.
\end{proof}

\begin{lemma}
\label{lemma-formally-etale-not-affine}
Soit $S$ un schéma.
Soit $f : X \to Y$ un morphisme formellement étale d'espaces algébriques sur $S$.
Alors, pour tout diagramme commutatif de flèches pleines
$$
\xymatrix{
X \ar[d]_f & T \ar[d]^i \ar[l]_a \\
Y & T' \ar[l] \ar@{-->}[lu]
}
$$
où $T \subset T'$ est un épaississement du premier ordre d'espaces algébriques
sur $Y$, il existe une unique flèche pointillée qui rend le diagramme commutatif.
Autrement dit, dans la
Définition \ref{definition-formally-etale},
on peut supprimer la condition que $T$ soit affine.
\end{lemma}

\begin{proof}
Soit $U' \to T'$ un morphisme étale surjectif, où $U' = \coprod U'_i$
est une somme disjointe de schémas affines. Posons
$U_i = T \times_{T'} U'_i$. On obtient alors des morphismes
$a'_i : U'_i \to X$ tels que $a'_i|_{U_i}$ soit égal à la composée
$U_i \to T \to X$. Par unicité (voir le
Lemme \ref{lemma-formally-unramified-not-affine}),
on voit que $a'_i$ et $a'_j$ coïncident sur le produit fibré
$U'_i \times_{T'} U'_j$. Par suite, $\coprod a'_i : U' \to X$
descend en un unique morphisme $a' : T' \to X$.
\end{proof}

\begin{lemma}
\label{lemma-composition-formally-etale}
La composée de morphismes formellement étales est formellement étale.
\end{lemma}

\begin{proof}
C'est formel.
\end{proof}

\begin{lemma}
\label{lemma-base-change-formally-etale}
Un changement de base d'un morphisme formellement étale est formellement étale.
\end{lemma}

\begin{proof}
C'est formel.
\end{proof}

\begin{lemma}
\label{lemma-characterize-formally-etale}
Soit $S$ un schéma.
Soit $f : X \to Y$ un morphisme d'espaces algébriques sur $S$.
Les assertions suivantes sont équivalentes :
\begin{enumerate}
\item $f$ est formellement étale,
\item $f$ est formellement non ramifié et l'épaississement universel du premier ordre
de $X$ sur $Y$ est égal à $X$,
\item $f$ est formellement non ramifié et $\mathcal{C}_{X/Y} = 0$, et
\item $\Omega_{X/Y} = 0$ et $\mathcal{C}_{X/Y} = 0$.
\end{enumerate}
\end{lemma}

\begin{proof}
En réalité, la dernière assertion n'a de sens que parce que $\Omega_{X/Y} = 0$
implique que $\mathcal{C}_{X/Y}$ est défini par le
Lemme \ref{lemma-characterize-formally-unramified}
et la
Définition \ref{definition-universal-thickening}.
Cela montre aussi que (3) et (4) sont équivalentes.

\medskip\noindent
Chacune des hypothèses (1), (2) et (3) implique que $f$ est formellement
non ramifié. On peut donc supposer $f$ formellement non ramifié. L'équivalence
de (1), (2) et (3) résulte de la propriété universelle de l'épaississement
universel du premier ordre $X'$ de $X$ sur $S$ et du fait que
$X = X' \Leftrightarrow \mathcal{C}_{X/Y} = 0$, puisque,
par définition, $\mathcal{C}_{X/Y} = \mathcal{C}_{X/X'}$
est le faisceau d'idéaux de $X$ dans $X'$.
\end{proof}

\begin{lemma}
\label{lemma-unramified-flat-formally-etale}
Un morphisme plat et non ramifié est formellement étale.
\end{lemma}

\begin{proof}
Cela résulte du cas des schémas, voir
Compléments sur les morphismes,
Lemme \ref{more-morphisms-lemma-unramified-flat-formally-etale},
et de la localisation étale, voir les
Lemmes \ref{lemma-formally-unramified} et \ref{lemma-formally-etale},
ainsi que
Morphismes d'espaces, Lemme \ref{spaces-morphisms-lemma-flat-local}.
\end{proof}

\begin{lemma}
\label{lemma-etale-formally-etale}
Soit $S$ un schéma.
Soit $f : X \to Y$ un morphisme d'espaces algébriques sur $S$.
Les assertions suivantes sont équivalentes :
\begin{enumerate}
\item le morphisme $f$ est étale, et
\item le morphisme $f$ est localement de présentation finie et
formellement étale.
\end{enumerate}
\end{lemma}

\begin{proof}
Cela résulte du cas des schémas, voir
Compléments sur les morphismes,
Lemme \ref{more-morphisms-lemma-etale-formally-etale},
et de la localisation étale, voir le
Lemme \ref{lemma-formally-etale},
ainsi que
Morphismes d'espaces,
Lemmes \ref{spaces-morphisms-lemma-finite-presentation-local} et
\ref{spaces-morphisms-lemma-etale-local}.
\end{proof}















\section{Déformations infinitésimales de morphismes}
\label{section-action-by-derivations}

\noindent
Dans cette section, nous expliquons comment une dérivation peut servir à
déformer infinitésimalement un morphisme. Dans toute la section, nous utilisons le fait qu'un
faisceau sur un épaississement $X'$ de $X$ peut être considéré comme un faisceau sur $X$, voir les
Équations (\ref{equation-equivalence-etale-spaces}) et
(\ref{equation-fundamental-equivalence}).

\begin{lemma}
\label{lemma-difference-derivation}
Soit $S$ un schéma. Soit $B$ un espace algébrique sur $S$.
Soient $X \subset X'$ et $Y \subset Y'$ deux épaississements du premier ordre
d'espaces algébriques sur $B$.
Soient $(a, a'), (b, b') : (X \subset X') \to (Y \subset Y')$
deux morphismes d'épaississements sur $B$. Supposons que
\begin{enumerate}
\item $a = b$, et
\item les deux morphismes $a^*\mathcal{C}_{Y/Y'} \to \mathcal{C}_{X/X'}$
(Lemme \ref{lemma-conormal-functorial})
soient égaux.
\end{enumerate}
Alors le morphisme $(a')^\sharp - (b')^\sharp$ se factorise sous la forme
$$
\mathcal{O}_{Y'} \to \mathcal{O}_Y \xrightarrow{D}
a_*\mathcal{C}_{X/X'} \to a_*\mathcal{O}_{X'}
$$
où $D$ est une $\mathcal{O}_B$-dérivation.
\end{lemma}

\begin{proof}
Au lieu de travailler sur $Y$, travaillons sur $X$. L'avantage est que le foncteur d'image inverse
$a^{-1}$ est exact. En utilisant (1) et (2), on obtient un diagramme commutatif
à lignes exactes
$$
\xymatrix{
0 \ar[r] &
\mathcal{C}_{X/X'} \ar[r] &
\mathcal{O}_{X'} \ar[r] &
\mathcal{O}_X \ar[r] & 0 \\
0 \ar[r] &
a^{-1}\mathcal{C}_{Y/Y'} \ar[r] \ar[u] &
a^{-1}\mathcal{O}_{Y'}
\ar[r] \ar@<1ex>[u]^{(a')^\sharp} \ar@<-1ex>[u]_{(b')^\sharp} &
a^{-1}\mathcal{O}_Y \ar[r] \ar[u] & 0
}
$$
Or il est général que, dans une telle situation, la différence des
morphismes de $\mathcal{O}_B$-algèbres $(a')^\sharp$ et $(b')^\sharp$ est une
$\mathcal{O}_B$-dérivation de $a^{-1}\mathcal{O}_Y$ dans $\mathcal{C}_{X/X'}$.
Par adjonction entre les foncteurs $a^{-1}$ et $a_*$, cela revient
à une $\mathcal{O}_B$-dérivation de
$\mathcal{O}_Y$ dans $a_*\mathcal{C}_{X/X'}$. Nous omettons certains détails.
\end{proof}

\noindent
Notons que, dans la situation du lemme ci-dessus, on peut écrire
$D$ sous la forme
\begin{equation}
\label{equation-D}
D = \text{d}_{Y/B} \circ \theta
\end{equation}
où $\theta$ est une application $\mathcal{O}_Y$-linéaire
$\theta : \Omega_{Y/B} \to a_*\mathcal{C}_{X/X'}$.
Bien entendu, par adjonction encore, on peut alors considérer $\theta$ comme une
application $\mathcal{O}_X$-linéaire
$\theta : a^*\Omega_{Y/B} \to \mathcal{C}_{X/X'}$.

\begin{lemma}
\label{lemma-action-by-derivations}
Soit $S$ un schéma. Soit $B$ un espace algébrique sur $S$.
Soit $(a, a') : (X \subset X') \to (Y \subset Y')$
un morphisme d'épaississements du premier ordre sur $B$.
Soit
$$
\theta : a^*\Omega_{Y/B} \to \mathcal{C}_{X/X'}
$$
une application $\mathcal{O}_X$-linéaire. Il existe alors un unique morphisme d'épaississements
$(b, b') : (X \subset X') \to (Y \subset Y')$ tel que
(1) et (2) du
Lemme \ref{lemma-difference-derivation}
soient vérifiées et que la dérivation $D$ et $\theta$ soient reliées par
l'Équation (\ref{equation-D}).
\end{lemma}

\begin{proof}
Considérons le morphisme
$$
\alpha = (a')^\sharp + D : a^{-1}\mathcal{O}_{Y'} \to \mathcal{O}_{X'}
$$
où $D$ est comme dans l'Équation (\ref{equation-D}). Puisque $D$ est une
$\mathcal{O}_B$-dérivation, il en résulte que $\alpha$ est un morphisme de
faisceaux de $\mathcal{O}_B$-algèbres. Par construction, on a
$i_X^\sharp \circ \alpha = a^\sharp \circ i_Y^\sharp$, où
$i_X : X \to X'$ et $i_Y : Y \to Y'$ sont les immersions fermées données. D'après le
Lemme \ref{lemma-first-order-thickening-maps},
on obtient un unique morphisme
$(a, b') : (X  \subset X') \to (Y \subset Y')$ d'épaississements
sur $B$ tel que $\alpha = (b')^\sharp$. En posant $b = a$,
le résultat est démontré.
\end{proof}

\begin{remark}
\label{remark-action-by-derivations}
Mêmes hypothèses et notations qu'au Lemme \ref{lemma-action-by-derivations}.
L'action d'une section locale $\theta$ sur $a'$ est parfois notée
$\theta \cdot a'$. Notons que cela signifie seulement que
$(a')^\sharp$ et $(\theta \cdot a')^\sharp$ diffèrent d'une dérivation
$D$ reliée à $\theta$ par l'Équation (\ref{equation-D}).
\end{remark}

\begin{lemma}
\label{lemma-sheaf}
Soit $S$ un schéma. Soit $B$ un espace algébrique sur $S$.
Soient $X \subset X'$ et $Y \subset Y'$ des épaississements du premier ordre
sur $B$. Supposons donnés un morphisme $a : X \to Y$ et un morphisme
$A : a^*\mathcal{C}_{Y/Y'} \to \mathcal{C}_{X/X'}$ de
$\mathcal{O}_X$-modules. Pour un objet $U'$ de
$(X')_{spaces, \etale}$ avec $U = X \times_{X'} U'$
considérons les morphismes $a' : U' \to Y'$ tels que
\begin{enumerate}
\item $a'$ soit un morphisme sur $B$,
\item $a'|_U = a|_U$, et
\item le morphisme induit
$a^*\mathcal{C}_{Y/Y'}|_U \to \mathcal{C}_{X/X'}|_U$
soit la restriction de $A$ à $U$.
\end{enumerate}
Alors la règle
\begin{equation}
\label{equation-sheaf}
U' \mapsto
\{a' : U' \to Y'\text{ such that (1), (2), (3) hold.}\}
\end{equation}
définit un faisceau d'ensembles sur $(X')_{spaces, \etale}$.
\end{lemma}

\begin{proof}
Notons $\mathcal{F}$ la règle du lemme.
Le morphisme de restriction $\mathcal{F}(U') \to \mathcal{F}(V')$ associé à
$V' \subset U' \subset X'$
pour la règle $\mathcal{F}$ est bien le morphisme de restriction $a' \mapsto a'|_{V'}$.
Avec cette définition, il est clair que $\mathcal{F}$ est un
faisceau, car les morphismes d'espaces algébriques vérifient la descente étale, voir
Descente sur les espaces,
Lemme \ref{spaces-descent-lemma-fpqc-universal-effective-epimorphisms}.
\end{proof}

\begin{lemma}
\label{lemma-action-sheaf}
Mêmes notations et hypothèses qu'au Lemme \ref{lemma-sheaf}.
Nous identifions les faisceaux sur $X$ et $X'$ au moyen de
(\ref{equation-equivalence-etale-spaces}).
Le faisceau
$$
\SheafHom_{\mathcal{O}_X}(a^*\Omega_{Y/B}, \mathcal{C}_{X/X'})
$$
opère sur le faisceau (\ref{equation-sheaf}). En outre, cette action
est simplement transitive pour tout objet $U'$ de $(X')_{spaces, \etale}$
au-dessus duquel le faisceau (\ref{equation-sheaf}) possède une section.
\end{lemma}

\begin{proof}
C'est la combinaison des
Lemmes \ref{lemma-difference-derivation},
\ref{lemma-action-by-derivations},
et \ref{lemma-sheaf}.
\end{proof}

\begin{remark}
\label{remark-special-case}
Un cas particulier des
Lemmes \ref{lemma-difference-derivation},
\ref{lemma-action-by-derivations},
\ref{lemma-sheaf}, et
\ref{lemma-action-sheaf}
est celui où $Y = Y'$. Dans ce cas, le morphisme $A$ est toujours nul.
Le faisceau du
Lemme \ref{lemma-sheaf}
est simplement donné par la règle
$$
U' \mapsto
\{a' : U' \to Y\text{ over }B\text{ with } a'|_U = a|_U\}
$$
sur lequel opère le faisceau
$\SheafHom_{\mathcal{O}_X}(a^*\Omega_{Y/B}, \mathcal{C}_{X/X'})$.
\end{remark}

\begin{remark}
\label{remark-another-special-case}
Un autre cas particulier des
Lemmes \ref{lemma-difference-derivation},
\ref{lemma-action-by-derivations},
\ref{lemma-sheaf}, et
\ref{lemma-action-sheaf}
est celui où $B$ est lui-même l'espace ambiant d'un épaississement $Z \subset Z' = B$
et où $Y = Z \times_{Z'} Y'$. Représentons la situation par
$$
\xymatrix{
(X \subset X') \ar@{..>}[rr]_{(a, ?)} \ar[rd]_{(g, g')} & &
(Y \subset Y') \ar[ld]^{(h, h')} \\
& (Z \subset Z')
}
$$
Dans ce cas, le morphisme $A : a^*\mathcal{C}_{Y/Y'} \to \mathcal{C}_{X/X'}$
est déterminé par $a$ : le morphisme
$h^*\mathcal{C}_{Z/Z'} \to \mathcal{C}_{Y/Y'}$ est surjectif (car nous avons
supposé $Y = Z \times_{Z'} Y'$) ; par conséquent, son image inverse
$g^*\mathcal{C}_{Z/Z'} = a^*h^*\mathcal{C}_{Z/Z'} \to
a^*\mathcal{C}_{Y/Y'}$ est surjective, et la composée
$g^*\mathcal{C}_{Z/Z'} \to a^*\mathcal{C}_{Y/Y'} \to \mathcal{C}_{X/X'}$
doit être le morphisme canonique induit par $g'$. Ainsi, le faisceau du
Lemme \ref{lemma-sheaf}
est simplement donné par la règle
$$
U' \mapsto
\{a' : U' \to Y'\text{ over }Z'\text{ with } a'|_U = a|_U\}
$$
sur lequel opère le faisceau
$\SheafHom_{\mathcal{O}_X}(a^*\Omega_{Y/Z}, \mathcal{C}_{X/X'})$.
\end{remark}

\begin{lemma}
\label{lemma-action-by-derivations-etale-localization}
Soit $S$ un schéma. Considérons un diagramme commutatif d'épaississements
du premier ordre
$$
\vcenter{
\xymatrix{
(T_2 \subset T_2') \ar[d]_{(h, h')} \ar[rr]_{(a_2, a_2')} & &
(X_2 \subset X_2') \ar[d]^{(f, f')} \\
(T_1 \subset T_1') \ar[rr]^{(a_1, a_1')} & &
(X_1 \subset X_1')
}
}
\quad
\begin{matrix}
\text{and a commutative} \\
\text{diagram}
\end{matrix}
\quad
\vcenter{
\xymatrix{
X_2' \ar[r] \ar[d] & B_2 \ar[d] \\
X_1' \ar[r] & B_1
}
}
$$
d'espaces algébriques sur $S$,
où $X_2 \to X_1$ et $B_2 \to B_1$ sont étales.
Pour toute application $\mathcal{O}_{T_1}$-linéaire
$\theta_1 : a_1^*\Omega_{X_1/B_1} \to \mathcal{C}_{T_1/T'_1}$, soit
$\theta_2$ la composée
$$
\xymatrix{
a_2^*\Omega_{X_2/B_2} \ar@{=}[r] &
h^*a_1^*\Omega_{X_1/B_1} \ar[r]^-{h^*\theta_1} &
h^*\mathcal{C}_{T_1/T'_1} \ar[r] &
\mathcal{C}_{T_2/T'_2}
}
$$
(le signe d'égalité est expliqué dans la démonstration). Alors le diagramme
$$
\xymatrix{
T_2' \ar[rr]_{\theta_2 \cdot a_2'} \ar[d] & & X'_2 \ar[d] \\
T_1' \ar[rr]^{\theta_1 \cdot a_1'} & & X'_1
}
$$
est commutatif, où les actions $\theta_2 \cdot a_2'$ et $\theta_1 \cdot a_1'$
sont celles de la Remarque \ref{remark-action-by-derivations}.
\end{lemma}

\begin{proof}
Le signe d'égalité provient de l'identification
$f^*\Omega_{X_1/S_1} = \Omega_{X_2/S_2}$ que l'on obtient
parce que la construction du faisceau des différentielles est
compatible avec la localisation étale (tant à la source qu'au but), voir le
Lemme \ref{lemma-localize-differentials}.
En effet, on a ainsi
$a_2^*\Omega_{X_2/S_2} = a_2^*f^*\Omega_{X_1/S_1} =
h^*a_1^*\Omega_{X_1/S_1}$ puisque $f \circ a_2 = a_1 \circ h$.
Cela étant, la commutativité du diagramme se vérifie
localement pour la topologie étale. On peut donc supposer que $T'_i$, $X'_i$,
$B_2$ et $B_1$ sont des schémas ; dans ce cas, le lemme
résulte de
Compléments sur les morphismes, Lemme
\ref{more-morphisms-lemma-action-by-derivations-etale-localization}.
Autre démonstration : d'après le Lemme \ref{lemma-first-order-thickening-maps},
il suffit de montrer qu'un certain diagramme de faisceaux
d'anneaux sur $X_1'$ est commutatif ; on raisonne alors exactement
comme dans la démonstration du lemme précité de
Compléments sur les morphismes, Lemme
\ref{more-morphisms-lemma-action-by-derivations-etale-localization}
pour voir que c'est bien le cas.
\end{proof}










\section{Déformations infinitésimales d'espaces algébriques}
\label{section-deform}

\noindent
Le lemme élémentaire suivant est souvent un outil commode pour vérifier si
une déformation infinitésimale d'un morphisme est plate.

\begin{lemma}
\label{lemma-deform}
Soit $S$ un schéma. Soit $(f, f') : (X \subset X') \to (Y \subset Y')$ un
morphisme d'épaississements du premier ordre d'espaces algébriques sur $S$. Supposons que
$f$ soit plat. Alors les assertions suivantes sont équivalentes :
\begin{enumerate}
\item $f'$ est plat et $X = Y \times_{Y'} X'$, et
\item le morphisme canonique $f^*\mathcal{C}_{Y/Y'} \to \mathcal{C}_{X/X'}$
est un isomorphisme.
\end{enumerate}
\end{lemma}

\begin{proof}
Choisissons un schéma $V'$ et un morphisme étale surjectif $V' \to Y'$.
Choisissons un schéma $U'$ et un morphisme étale surjectif
$U' \to X' \times_{Y'} V'$. Posons $U = X \times_{X'} U'$ et
$V = Y \times_{Y'} V'$. D'après notre définition d'un morphisme plat
d'espaces algébriques, le morphisme induit $g : U \to V$ est un morphisme plat
de schémas, et $f'$ est plat si et seulement si le morphisme correspondant
$g' : U' \to V'$ est plat. De plus, $X = Y \times_{Y'} X'$ si et seulement
si $U = V \times_{V'} V'$. Enfin, le morphisme
$f^*\mathcal{C}_{Y/Y'} \to \mathcal{C}_{X/X'}$
est un isomorphisme si et seulement si
$g^*\mathcal{C}_{V/V'} \to \mathcal{C}_{U/U'}$ est un isomorphisme.
Le lemme résulte donc de son analogue pour les morphismes de schémas, voir
Compléments sur les morphismes, Lemme \ref{more-morphisms-lemma-deform}.
\end{proof}

\noindent
Le lemme suivant est la version « nilpotente » du
« critère de platitude par fibres », voir la
Section \ref{section-criterion-flat-fibres}.

\begin{lemma}
\label{lemma-flatness-morphism-thickenings}
Soit $S$ un schéma. Considérons un diagramme commutatif
$$
\xymatrix{
(X \subset X') \ar[rr]_{(f, f')} \ar[rd] & & (Y \subset Y') \ar[ld] \\
& (B \subset B')
}
$$
d'épaississements d'espaces algébriques sur $S$. Supposons que
\begin{enumerate}
\item $X'$ soit plat sur $B'$,
\item $f$ soit plat,
\item $B \subset B'$ soit un épaississement d'ordre fini, et
\item $X = B \times_{B'} X'$ et $Y = B \times_{B'} Y'$.
\end{enumerate}
Alors $f'$ est plat et $Y'$ est plat sur $B'$ en tout point de
l'image de $f'$.
\end{lemma}

\begin{proof}
Choisissons un schéma $U'$ et un morphisme étale surjectif $U' \to B'$.
Choisissons un schéma $V'$ et un morphisme étale surjectif
$V' \to U' \times_{B'} Y'$.
Choisissons un schéma $W'$ et un morphisme étale surjectif
$W' \to V' \times_{Y'} X'$. Soient $U, V, W$ les changements de base
de $U', V', W'$ par $B \to B'$. Alors la platitude de $f'$ est
équivalente à celle de $W' \to V'$, et nous savons
que $W \to V$ est plat. Nous pouvons donc appliquer le lemme
dans le cas des schémas au diagramme
$$
\xymatrix{
(W \subset W') \ar[rr] \ar[rd] & & (V \subset V') \ar[ld] \\
& (U \subset U')
}
$$
d'épaississements de schémas. Voir
Compléments sur les morphismes, Lemme
\ref{more-morphisms-lemma-flatness-morphism-thickenings}.
L'assertion relative à la platitude de $Y'/B'$ aux points de
l'image de $f'$ se démontre de la même manière.
\end{proof}

\noindent
De nombreuses propriétés des morphismes de schémas sont préservées par les
déformations plates.

\begin{lemma}
\label{lemma-deform-property}
Soit $S$ un schéma. Considérons un diagramme commutatif
$$
\xymatrix{
(X \subset X') \ar[rr]_{(f, f')} \ar[rd] & & (Y \subset Y') \ar[ld] \\
& (B \subset B')
}
$$
d'épaississements d'espaces algébriques sur $S$. Supposons que $B \subset B'$
soit un épaississement d'ordre fini, que $X'$ soit plat sur $B'$, que $X = B \times_{B'} X'$,
et que $Y = B \times_{B'} Y'$. Alors
\begin{enumerate}
\item $f$ est représentable si et seulement si $f'$ est représentable,
\label{item-representable}
\item $f$ est plat si et seulement si $f'$ est plat,
\label{item-flat}
\item $f$ est un isomorphisme si et seulement si $f'$ est un isomorphisme,
\label{item-isomorphism}
\item $f$ est une immersion ouverte si et seulement si $f'$ est une immersion ouverte,
\label{item-open-immersion}
\item $f$ est quasi-compact si et seulement si $f'$ est quasi-compact,
\label{item-quasi-compact}
\item $f$ est universellement fermé si et seulement si $f'$ est universellement fermé,
\label{item-universally-closed}
\item $f$ est (quasi-)séparé si et seulement si $f'$ est (quasi-)séparé,
\label{item-separated}
\item $f$ est un monomorphisme si et seulement si $f'$ est un monomorphisme,
\label{item-monomorphism}
\item $f$ est surjectif si et seulement si $f'$ est surjectif,
\label{item-surjective}
\item $f$ est universellement injectif si et seulement si $f'$ est universellement injectif,
\label{item-universally-injective}
\item $f$ est affine si et seulement si $f'$ est affine,
\label{item-affine}
\item
\label{item-finite-type}
$f$ est localement de type fini si et seulement si $f'$ est localement de type fini,
\item $f$ est localement quasi-fini si et seulement si $f'$ est localement quasi-fini,
\label{item-quasi-finite}
\item
\label{item-finite-presentation}
$f$ est localement de présentation finie si et seulement si $f'$ est localement de
présentation finie,
\item
\label{item-relative-dimension-d}
$f$ est localement de type fini de dimension relative $d$ si et seulement si
$f'$ est localement de type fini de dimension relative $d$,
\item $f$ est universellement ouvert si et seulement si $f'$ est universellement ouvert,
\label{item-universally-open}
\item $f$ est syntomique si et seulement si $f'$ est syntomique,
\label{item-syntomic}
\item $f$ est lisse si et seulement si $f'$ est lisse,
\label{item-smooth}
\item $f$ est non ramifié si et seulement si $f'$ est non ramifié,
\label{item-unramified}
\item $f$ est étale si et seulement si $f'$ est étale,
\label{item-etale}
\item $f$ est propre si et seulement si $f'$ est propre,
\label{item-proper}
\item $f$ est entier si et seulement si $f'$ est entier,
\label{item-integral}
\item $f$ est fini si et seulement si $f'$ est fini,
\label{item-finite}
\item
\label{item-finite-locally-free}
$f$ est fini localement libre (de rang $d$) si et seulement si $f'$
est fini localement libre (de rang $d$), et
\item ajouter d'autres cas ici.
\end{enumerate}
\end{lemma}

\begin{proof}
Le cas (\ref{item-representable}) résulte du
Lemme \ref{lemma-thicken-property-morphisms}.

\medskip\noindent
Choisissons un schéma $U'$ et un morphisme étale surjectif $U' \to B'$.
Choisissons un schéma $V'$ et un morphisme étale surjectif
$V' \to U' \times_{B'} Y'$.
Choisissons un schéma $W'$ et un morphisme étale surjectif
$W' \to V' \times_{Y'} X'$. Soient $U, V, W$ les changements de base
de $U', V', W'$ par $B \to B'$. Considérons le diagramme
$$
\xymatrix{
(W \subset W') \ar[rr] \ar[rd] & & (V \subset V') \ar[ld] \\
& (U \subset U')
}
$$
d'épaississements de schémas. Pour chacune des propriétés qui sont
locales pour la topologie étale à la source et au but, le résultat découle immédiatement
du résultat correspondant pour les morphismes d'épaississements de schémas,
appliqué au diagramme ci-dessus. Ainsi, les cas
(\ref{item-flat}), (\ref{item-finite-type}),
(\ref{item-quasi-finite}), (\ref{item-finite-presentation}),
(\ref{item-relative-dimension-d}), (\ref{item-syntomic}),
(\ref{item-smooth}), (\ref{item-unramified}), (\ref{item-etale})
résultent des cas correspondants de
Compléments sur les morphismes, Lemme \ref{more-morphisms-lemma-deform-property}.

\medskip\noindent
Puisque $X \to X'$ et $Y \to Y'$ sont des homéomorphismes universels,
toute question portant sur la topologie des morphismes
$X \to Y$ et $X' \to Y'$ reçoit la même réponse. Ainsi,
les cas (\ref{item-quasi-compact}), (\ref{item-universally-closed}),
(\ref{item-surjective}), (\ref{item-universally-injective}), et
(\ref{item-universally-open}) sont vérifiés.

\medskip\noindent
Dans chacun des cas restants, nous démontrons seulement l'implication
$f\text{ has }P \Rightarrow f'\text{ has }P$, car l'autre
implication résulte du fait que $P$ est stable par
changement de base, voir
Espaces, Lemme \ref{spaces-lemma-base-change-immersions} et
Morphismes d'espaces, Lemmes
\ref{spaces-morphisms-lemma-base-change-separated},
\ref{spaces-morphisms-lemma-base-change-monomorphism},
\ref{spaces-morphisms-lemma-base-change-affine},
\ref{spaces-morphisms-lemma-base-change-proper},
\ref{spaces-morphisms-lemma-base-change-integral}, et
\ref{spaces-morphisms-lemma-base-change-finite-locally-free}.

\medskip\noindent
Le cas (\ref{item-open-immersion}). Supposons que $f$ soit une immersion ouverte.
Alors $f'$ est étale d'après (\ref{item-etale}) et universellement injectif
d'après (\ref{item-universally-injective}) ;
ainsi, $f'$ est une immersion ouverte, voir
Morphismes d'espaces, Lemme
\ref{spaces-morphisms-lemma-etale-universally-injective-open}.
On peut éviter d'utiliser ce lemme en
utilisant d'abord (\ref{item-representable}) pour se ramener au cas des schémas.

\medskip\noindent
Le cas (\ref{item-isomorphism}) résulte des cas
(\ref{item-open-immersion}) et (\ref{item-surjective}).

\medskip\noindent
Le cas (\ref{item-separated}). Voir le
Lemme \ref{lemma-thicken-property-morphisms}.

\medskip\noindent
Le cas (\ref{item-monomorphism}). Supposons que $f$ soit un monomorphisme.
Considérons le morphisme diagonal $\Delta_{X'/Y'} : X' \to X' \times_{Y'} X'$.
Le changement de base de $\Delta_{X'/Y'}$ par $B \to B'$ est $\Delta_{X/Y}$,
qui est un isomorphisme par hypothèse. D'après (\ref{item-isomorphism}),
nous concluons que $\Delta_{X'/Y'}$ est un isomorphisme et donc que
$f'$ est un monomorphisme.

\medskip\noindent
Le cas (\ref{item-affine}). Voir le Lemme \ref{lemma-thicken-property-morphisms}.

\medskip\noindent
Le cas (\ref{item-proper}). Voir le
Lemme \ref{lemma-thicken-property-morphisms-cartesian}.

\medskip\noindent
Le cas (\ref{item-integral}). Voir le
Lemme \ref{lemma-thicken-property-morphisms}.

\medskip\noindent
Le cas (\ref{item-finite}). Voir le
Lemme \ref{lemma-thicken-property-morphisms-cartesian}.

\medskip\noindent
Le cas (\ref{item-finite-locally-free}). Supposons $f$ fini localement libre.
D'après (\ref{item-finite}), $f'$ est fini.
D'après (\ref{item-flat}), $f'$ est plat.
D'après (\ref{item-finite-presentation}), $f'$ est localement de
présentation finie. Ainsi, $f'$ est fini localement libre d'après
Morphismes d'espaces, Lemme \ref{spaces-morphisms-lemma-finite-flat}.
\end{proof}

\noindent
Le lemme suivant est la version « localement nilpotente » du
« critère de platitude par fibres », voir la
Section \ref{section-criterion-flat-fibres}.

\begin{lemma}
\label{lemma-flatness-morphism-thickenings-fp-over-ft}
Soit $S$ un schéma. Considérons un diagramme commutatif
$$
\xymatrix{
(X \subset X') \ar[rr]_{(f, f')} \ar[rd] & & (Y \subset Y') \ar[ld] \\
& (B \subset B')
}
$$
d'épaississements d'espaces algébriques sur $S$. Supposons que
\begin{enumerate}
\item $Y' \to B'$ soit localement de type fini,
\item $X' \to B'$ soit plat et localement de présentation finie,
\item $f$ soit plat, et
\item $X = B \times_{B'} X'$ et $Y = B \times_{B'} Y'$.
\end{enumerate}
Alors $f'$ est plat et, pour tout $y' \in |Y'|$ dans l'image de $|f'|$,
le morphisme $Y' \to B'$ est plat en $y'$.
\end{lemma}

\begin{proof}
Choisissons un schéma $U'$ et un morphisme étale surjectif $U' \to B'$.
Choisissons un schéma $V'$ et un morphisme étale surjectif
$V' \to U' \times_{B'} Y'$.
Choisissons un schéma $W'$ et un morphisme étale surjectif
$W' \to V' \times_{Y'} X'$. Soient $U, V, W$ les changements de base
de $U', V', W'$ par $B \to B'$. Alors la platitude de $f'$ est
équivalente à celle de $W' \to V'$, et nous savons
que $W \to V$ est plat. Nous pouvons donc appliquer le lemme
dans le cas des schémas au diagramme
$$
\xymatrix{
(W \subset W') \ar[rr] \ar[rd] & & (V \subset V') \ar[ld] \\
& (U \subset U')
}
$$
d'épaississements de schémas. Voir
Compléments sur les morphismes, Lemme
\ref{more-morphisms-lemma-flatness-morphism-thickenings-fp-over-ft}.
L'assertion relative à la platitude de $Y'/B'$ aux points de
l'image de $f'$ se démontre de la même manière.
\end{proof}

\noindent
De nombreuses propriétés des morphismes de schémas sont préservées par les
déformations plates comme dans le lemme ci-dessus.

\begin{lemma}
\label{lemma-deform-property-fp-over-ft}
Soit $S$ un schéma. Considérons un diagramme commutatif
$$
\xymatrix{
(X \subset X') \ar[rr]_{(f, f')} \ar[rd] & & (Y \subset Y') \ar[ld] \\
& (B \subset B')
}
$$
d'épaississements d'espaces algébriques sur $S$.
Supposons que $Y' \to B'$ soit localement de type fini,
que $X' \to B'$ soit plat et localement de présentation finie,
que $X = B \times_{B'} X'$ et que $Y = B \times_{B'} Y'$. Alors
\begin{enumerate}
\item $f$ est représentable si et seulement si $f'$ est représentable,
\label{item-representable-fp-over-ft}
\item $f$ est plat si et seulement si $f'$ est plat,
\label{item-flat-fp-over-ft}
\item $f$ est un isomorphisme si et seulement si $f'$ est un isomorphisme,
\label{item-isomorphism-fp-over-ft}
\item $f$ est une immersion ouverte si et seulement si $f'$ est une immersion ouverte,
\label{item-open-immersion-fp-over-ft}
\item $f$ est quasi-compact si et seulement si $f'$ est quasi-compact,
\label{item-quasi-compact-fp-over-ft}
\item $f$ est universellement fermé si et seulement si $f'$ est universellement fermé,
\label{item-universally-closed-fp-over-ft}
\item $f$ est (quasi-)séparé si et seulement si $f'$ est (quasi-)séparé,
\label{item-separated-fp-over-ft}
\item $f$ est un monomorphisme si et seulement si $f'$ est un monomorphisme,
\label{item-monomorphism-fp-over-ft}
\item $f$ est surjectif si et seulement si $f'$ est surjectif,
\label{item-surjective-fp-over-ft}
\item $f$ est universellement injectif si et seulement si $f'$ est universellement injectif,
\label{item-universally-injective-fp-over-ft}
\item $f$ est affine si et seulement si $f'$ est affine,
\label{item-affine-fp-over-ft}
\item $f$ est localement quasi-fini si et seulement si $f'$ est localement quasi-fini,
\label{item-quasi-finite-fp-over-ft}
\item
\label{item-relative-dimension-d-fp-over-ft}
$f$ est localement de type fini de dimension relative $d$ si et seulement si
$f'$ est localement de type fini de dimension relative $d$,
\item $f$ est universellement ouvert si et seulement si $f'$ est universellement ouvert,
\label{item-universally-open-fp-over-ft}
\item $f$ est syntomique si et seulement si $f'$ est syntomique,
\label{item-syntomic-fp-over-ft}
\item $f$ est lisse si et seulement si $f'$ est lisse,
\label{item-smooth-fp-over-ft}
\item $f$ est non ramifié si et seulement si $f'$ est non ramifié,
\label{item-unramified-fp-over-ft}
\item $f$ est étale si et seulement si $f'$ est étale,
\label{item-etale-fp-over-ft}
\item $f$ est propre si et seulement si $f'$ est propre,
\label{item-proper-fp-over-ft}
\item $f$ est fini si et seulement si $f'$ est fini,
\label{item-finite-fp-over-ft}
\item
\label{item-finite-locally-free-fp-over-ft}
$f$ est fini localement libre (de rang $d$) si et seulement si $f'$
est fini localement libre (de rang $d$), et
\item ajouter d'autres cas ici.
\end{enumerate}
\end{lemma}

\begin{proof}
Le cas (\ref{item-representable-fp-over-ft}) résulte du
Lemme \ref{lemma-thicken-property-morphisms}.

\medskip\noindent
Choisissons un schéma $U'$ et un morphisme étale surjectif $U' \to B'$.
Choisissons un schéma $V'$ et un morphisme étale surjectif
$V' \to U' \times_{B'} Y'$.
Choisissons un schéma $W'$ et un morphisme étale surjectif
$W' \to V' \times_{Y'} X'$. Soient $U, V, W$ les changements de base
de $U', V', W'$ par $B \to B'$. Considérons le diagramme
$$
\xymatrix{
(W \subset W') \ar[rr] \ar[rd] & & (V \subset V') \ar[ld] \\
& (U \subset U')
}
$$
d'épaississements de schémas. Pour chacune des propriétés qui sont
locales pour la topologie étale à la source et au but, le résultat découle immédiatement
du résultat correspondant pour les morphismes d'épaississements de schémas,
appliqué au diagramme ci-dessus. Ainsi, les cas
(\ref{item-flat-fp-over-ft}),
(\ref{item-quasi-finite-fp-over-ft}),
(\ref{item-relative-dimension-d-fp-over-ft}),
(\ref{item-syntomic-fp-over-ft}),
(\ref{item-smooth-fp-over-ft}),
(\ref{item-unramified-fp-over-ft}),
(\ref{item-etale-fp-over-ft})
résultent des cas correspondants de
Compléments sur les morphismes, Lemme \ref{more-morphisms-lemma-deform-property-fp-over-ft}.

\medskip\noindent
Puisque $X \to X'$ et $Y \to Y'$ sont des homéomorphismes universels,
toute question portant sur la topologie des morphismes
$X \to Y$ et $X' \to Y'$ reçoit la même réponse. Ainsi,
les cas (\ref{item-quasi-compact-fp-over-ft}),
(\ref{item-universally-closed-fp-over-ft}),
(\ref{item-surjective-fp-over-ft}),
(\ref{item-universally-injective-fp-over-ft}), et
(\ref{item-universally-open-fp-over-ft}) sont vérifiés.

\medskip\noindent
Dans chacun des cas restants, nous démontrons seulement l'implication
$f\text{ has }P \Rightarrow f'\text{ has }P$, car l'autre
implication résulte du fait que $P$ est stable par
changement de base, voir
Espaces, Lemme \ref{spaces-lemma-base-change-immersions} et
Morphismes d'espaces, Lemmes
\ref{spaces-morphisms-lemma-base-change-separated},
\ref{spaces-morphisms-lemma-base-change-monomorphism},
\ref{spaces-morphisms-lemma-base-change-affine},
\ref{spaces-morphisms-lemma-base-change-proper},
\ref{spaces-morphisms-lemma-base-change-integral}, et
\ref{spaces-morphisms-lemma-base-change-finite-locally-free}.

\medskip\noindent
Le cas (\ref{item-open-immersion-fp-over-ft}).
Supposons que $f$ soit une immersion ouverte.
Alors $f'$ est étale d'après (\ref{item-etale-fp-over-ft}) et universellement injectif
d'après (\ref{item-universally-injective-fp-over-ft}) ;
ainsi, $f'$ est une immersion ouverte, voir
Morphismes d'espaces, Lemme
\ref{spaces-morphisms-lemma-etale-universally-injective-open}.
On peut éviter d'utiliser ce lemme en
utilisant d'abord (\ref{item-representable-fp-over-ft}) pour se ramener au cas des schémas.

\medskip\noindent
Le cas (\ref{item-isomorphism-fp-over-ft}) résulte des cas
(\ref{item-open-immersion-fp-over-ft}) et (\ref{item-surjective-fp-over-ft}).

\medskip\noindent
Le cas (\ref{item-separated-fp-over-ft}). Voir le
Lemme \ref{lemma-thicken-property-morphisms}.

\medskip\noindent
Le cas (\ref{item-monomorphism-fp-over-ft}). Supposons que $f$ soit un monomorphisme.
Considérons le morphisme diagonal $\Delta_{X'/Y'} : X' \to X' \times_{Y'} X'$.
Le changement de base de $\Delta_{X'/Y'}$ par $B \to B'$ est $\Delta_{X/Y}$,
qui est un isomorphisme par hypothèse. D'après (\ref{item-isomorphism-fp-over-ft}),
nous concluons que $\Delta_{X'/Y'}$ est un isomorphisme et donc que
$f'$ est un monomorphisme.

\medskip\noindent
Le cas (\ref{item-affine-fp-over-ft}).
Voir le Lemme \ref{lemma-thicken-property-morphisms}.

\medskip\noindent
Le cas (\ref{item-proper-fp-over-ft}). Voir le
Lemme \ref{lemma-properties-that-extend-over-thickenings}.

\medskip\noindent
Le cas (\ref{item-finite-fp-over-ft}). Voir le
Lemme \ref{lemma-properties-that-extend-over-thickenings}.

\medskip\noindent
Le cas (\ref{item-finite-locally-free-fp-over-ft}).
Supposons $f$ fini localement libre.
D'après (\ref{item-finite-fp-over-ft}), $f'$ est fini.
D'après (\ref{item-flat-fp-over-ft}), $f'$ est plat.
De plus, $f'$ est localement de présentation finie d'après
Morphismes d'espaces, Lemme
\ref{spaces-morphisms-lemma-finite-presentation-permanence}.
Ainsi, $f'$ est fini localement libre d'après
Morphismes d'espaces, Lemme \ref{spaces-morphisms-lemma-finite-flat}.
\end{proof}













\section{Morphismes formellement lisses}
\label{section-formally-smooth}

\noindent
Dans cette section, nous introduisons la notion de morphisme formellement lisse
$X \to Y$ d'espaces algébriques. Un tel morphisme est
caractérisé par la propriété que, pour $T$, les points de $X$ à valeurs dans ce schéma se relèvent
aux épaississements infinitésimaux de $T$, pourvu que $T$ soit affine.
Le résultat principal affirme qu'un morphisme formellement lisse et
localement de présentation finie est lisse, voir le
Lemme \ref{lemma-smooth-formally-smooth}.
Il se trouve que ce critère est souvent plus facile à utiliser que le
critère jacobien.

\begin{definition}
\label{definition-formally-smooth}
Soit $S$ un schéma. Un morphisme $f : X \to Y$ d'espaces algébriques sur $S$
est dit {\it formellement lisse} s'il est formellement lisse comme
transformation de foncteurs au sens de la
Définition \ref{definition-formally-smooth-etale-unramified}.
\end{definition}

\noindent
Dans le cas des morphismes formellement non ramifiés et formellement étales,
on pourrait omettre l'hypothèse que $T'$ soit affine, voir les
Lemmes \ref{lemma-formally-unramified-not-affine} et
\ref{lemma-formally-etale-not-affine}.
Ce n'est plus vrai dans le cas des morphismes formellement lisses.
En fait, une condition un peu plus naturelle consisterait à pouvoir
compléter la flèche pointillée localement pour la topologie étale sur $T'$.
En effet, l'analyse de la démonstration du
Lemme \ref{lemma-smooth-formally-smooth}
montre que cela serait équivalent à la définition sous sa forme
actuelle. Il est aussi vrai qu'exiger l'existence de la flèche
pointillée localement pour la topologie fppf sur $T'$ suffirait, mais cela est un peu
plus difficile à démontrer.

\medskip\noindent
Nous ne reformulerons pas les résultats démontrés dans le cadre plus général des
transformations de foncteurs formellement lisses à la
Section \ref{section-formally-smooth-etale-unramified}.

\begin{lemma}
\label{lemma-composition-formally-smooth}
Une composée de morphismes formellement lisses est formellement lisse.
\end{lemma}

\begin{proof}
Démonstration omise.
\end{proof}

\begin{lemma}
\label{lemma-base-change-formally-smooth}
Un changement de base d'un morphisme formellement lisse est formellement lisse.
\end{lemma}

\begin{proof}
Démonstration omise, mais voir
Algèbre, Lemme \ref{algebra-lemma-base-change-fs}
pour la version algébrique.
\end{proof}

\begin{lemma}
\label{lemma-formally-etale-unramified-smooth}
Soit $f : X \to S$ un morphisme de schémas.
Alors $f$ est formellement étale si et seulement si
$f$ est formellement lisse et formellement non ramifié.
\end{lemma}

\begin{proof}
Démonstration omise.
\end{proof}

\noindent
Voici un lemme auxiliaire qui sera englobé par le
Lemme \ref{lemma-formally-smooth}.

\begin{lemma}
\label{lemma-helper-formally-smooth}
Soit $S$ un schéma. Soit
$$
\xymatrix{
U \ar[d] \ar[r]_\psi & V \ar[d] \\
X \ar[r]^f & Y
}
$$
un diagramme commutatif de morphismes d'espaces algébriques sur $S$.
Si les flèches verticales sont étales et si $f$ est formellement lisse, alors
$\psi$ est formellement lisse.
\end{lemma}

\begin{proof}
D'après le
Lemme \ref{lemma-representable-property-formally-property},
les morphismes $U \to X$ et $V \to Y$ sont formellement étales. D'après le
Lemme \ref{lemma-composition-formally-smooth-etale-unramified},
la composée $U \to Y$ est formellement lisse. D'après le
Lemme \ref{lemma-formally-permanence},
nous voyons que $\psi : U \to V$ est formellement lisse.
\end{proof}

\noindent
Le lemme suivant est le résultat principal de cette section.
Combiné avec la
Proposition de Limites d'espaces
\ref{spaces-limits-proposition-characterize-locally-finite-presentation},
il implique que l'on peut reconnaître si un morphisme d'espaces algébriques
$f : X \to Y$ est lisse au moyen de propriétés « simples » de la
transformation de foncteurs $X \to Y$.

\begin{lemma}[Critère infinitésimal de relèvement]
\label{lemma-smooth-formally-smooth}
Soit $S$ un schéma.
Soit $f : X \to Y$ un morphisme d'espaces algébriques sur $S$.
Les assertions suivantes sont équivalentes :
\begin{enumerate}
\item Le morphisme $f$ est lisse.
\item Le morphisme $f$ est localement de présentation finie et
formellement lisse.
\end{enumerate}
\end{lemma}

\begin{proof}
Supposons que $f : X \to S$ soit localement de présentation finie et formellement lisse.
Considérons un diagramme commutatif
$$
\xymatrix{
U \ar[d] \ar[r]_\psi & V \ar[d] \\
X \ar[r]^f & Y
}
$$
où $U$ et $V$ sont des schémas et où les flèches verticales sont étales et
surjectives. D'après le
Lemme \ref{lemma-helper-formally-smooth},
nous voyons que $\psi : U \to V$ est formellement lisse. D'après
Morphismes d'espaces, Lemme
\ref{spaces-morphisms-lemma-finite-presentation-local}
le morphisme $\psi$ est localement de présentation finie.
Par conséquent, dans le cas des schémas, le morphisme
$\psi$ est lisse, voir
Compléments sur les morphismes, Lemme \ref{more-morphisms-lemma-smooth-formally-smooth}.
Ainsi, $f$ est lisse, voir
Morphismes d'espaces, Lemme
\ref{spaces-morphisms-lemma-smooth-local}.

\medskip\noindent
Réciproquement, supposons que $f : X \to Y$ soit lisse.
Considérons un diagramme commutatif de flèches pleines
$$
\xymatrix{
X \ar[d]_f & T \ar[d]^i \ar[l]^a \\
Y & T' \ar[l] \ar@{-->}[lu]
}
$$
comme dans la Définition \ref{definition-formally-smooth}. Nous allons montrer que la
flèche pointillée existe, ce qui démontrera que $f$ est formellement lisse.
Soit $\mathcal{F}$ le faisceau d'ensembles sur $(T')_{spaces, \etale}$ du
Lemme \ref{lemma-sheaf}, dans le cas particulier examiné à la
Remarque \ref{remark-special-case}. Soit
$$
\mathcal{H} =
\SheafHom_{\mathcal{O}_T}(a^*\Omega_{X/Y}, \mathcal{C}_{T/T'})
$$
le faisceau de $\mathcal{O}_T$-modules sur $T_{spaces, \etale}$
muni de l'action $\mathcal{H} \times \mathcal{F} \to \mathcal{F}$
du Lemme \ref{lemma-action-sheaf}.
L'action $\mathcal{H} \times \mathcal{F} \to \mathcal{F}$
fait de $\mathcal{F}$ un pseudo-torseur sous $\mathcal{H}$, voir
Cohomologie sur les sites, Définition \ref{sites-cohomology-definition-torsor}.
Notre but est de montrer que $\mathcal{F}$ est un torseur trivial sous $\mathcal{H}$.
Il y a deux étapes : (I) Pour montrer que $\mathcal{F}$ est un torseur,
nous devons montrer que $\mathcal{F}$ possède localement pour la topologie étale
une section. (II) Pour montrer que $\mathcal{F}$ est le torseur trivial,
il suffit de montrer que $H^1(T_\etale, \mathcal{H}) = 0$, voir
Cohomologie sur les sites, Lemme \ref{sites-cohomology-lemma-torsors-h1}.

\medskip\noindent
Démontrons d'abord (I). Pour cela, choisissons un diagramme commutatif
$$
\xymatrix{
U \ar[d] \ar[r]_\psi & V \ar[d] \\
X \ar[r]^f & Y
}
$$
où $U$ et $V$ sont des schémas et où les flèches verticales sont étales et
surjectives. Comme $f$ est supposé lisse, nous voyons que $\psi$ est lisse et
donc formellement lisse d'après le
Lemme \ref{lemma-representable-property-formally-property}.
D'après le même lemme, le morphisme $V \to Y$ est formellement étale. Ainsi, d'après le
Lemme \ref{lemma-composition-formally-smooth-etale-unramified},
la composée $U \to Y$ est formellement lisse. L'étape (I) résulte alors du
Lemme \ref{lemma-etale-on-top}, partie (4).

\medskip\noindent
Démontrons enfin (II). D'après le
Lemme \ref{lemma-finite-presentation-differentials},
nous voyons que $\Omega_{X/S}$ est de présentation finie.
Ainsi, $a^*\Omega_{X/S}$ est de présentation finie (voir
Propriétés d'espaces,
Section \ref{spaces-properties-section-properties-modules}).
Par conséquent, le faisceau
$\mathcal{H} =
\SheafHom_{\mathcal{O}_T}(a^*\Omega_{X/Y}, \mathcal{C}_{T/T'})$
est quasi-cohérent d'après
Propriétés d'espaces,
Lemme \ref{spaces-properties-lemma-properties-quasi-coherent}.
Ainsi, d'après
Descente, Proposition \ref{descent-proposition-same-cohomology-quasi-coherent}
et Cohomologie des schémas, Lemme
\ref{coherent-lemma-quasi-coherent-affine-cohomology-zero}
nous avons
$$
H^1(T_{spaces, \etale}, \mathcal{H}) =
H^1(T_\etale, \mathcal{H}) =
H^1(T, \mathcal{H}) = 0
$$
comme voulu.
\end{proof}

\noindent
Les morphismes lisses satisfont une forte propriété locale de relèvement, voir le
Lemme \ref{lemma-smooth-strong-lift}. Si, dans le lemme, nous
supposons que $T'$ est affine, alors
nous ignorons s'il est nécessaire de prendre un recouvrement étale.
Plus précisément, si nous avons un diagramme commutatif
$$
\xymatrix{
X \ar[d] & T \ar[l] \ar[d] \\
Y & T' \ar[l] \ar@{..>}[lu]
}
$$
d'espaces algébriques où $X \to Y$ est lisse
et où $T \to T'$ est un épaississement de schémas affines,
existe-t-il toujours une flèche pointillée rendant le diagramme
commutatif ? Si vous connaissez la réponse ou une référence, veuillez écrire à
\href{mailto:stacks.project@gmail.com}{stacks.project@gmail.com}.

\begin{lemma}
\label{lemma-smooth-strong-lift}
Soit $S$ un schéma. Considérons un diagramme commutatif
$$
\xymatrix{
X \ar[d] & T \ar[l] \ar[d] \\
Y & T' \ar[l]
}
$$
d'espaces algébriques sur $S$, où $X \to Y$ est lisse
et où $T \to T'$ est un épaississement. Alors il existe un
recouvrement étale $\{T'_i \to T'\}$ tel que nous
puissions trouver la flèche pointillée dans
$$
\xymatrix{
X \ar[d] & T \ar[l] \ar[d] & T \times_{T'} T'_i \ar[l] \ar[d] \\
Y & T' \ar[l] & T'_i \ar[l] \ar@{..>}[llu]
}
$$
qui rend le diagramme commutatif (pour tout $i$).
\end{lemma}

\begin{proof}
Choisissons un recouvrement étale $\{Y_i \to Y\}$ où chaque $Y_i$ est affine.
Après avoir remplacé $T'$ par le recouvrement étale induit, nous pouvons supposer
que $Y$ est affine.

\medskip\noindent
Supposons que $Y$ soit affine. Choisissons un recouvrement étale $\{X_i \to X\}$.
Il donne naissance à un recouvrement étale de $T$. Ce recouvrement étale de $T$
provient d'un recouvrement étale de $T'$
(d'après le Théorème \ref{theorem-topological-invariance}, voir la
discussion de la Section \ref{section-thickenings}).
Nous pouvons donc supposer que $X$ est affine.

\medskip\noindent
Supposons que $X$ et $Y$ soient affines. Nous pouvons encore remplacer
$T'$ par un recouvrement étale et supposer que $T'$ est affine. Dans ce cas, le lemme résulte de
Algèbre, Lemme \ref{algebra-lemma-smooth-strong-lift}.
\end{proof}







\noindent
Nous poursuivons un peu pour montrer que la propriété d'être formellement lisse est locale
pour la topologie étale à la source. Pour commencer, nous montrons qu'un morphisme formellement lisse possède un
faisceau de différentielles bien comporté. La notion de module quasi-cohérent
localement projectif est définie dans
Propriétés d'espaces,
Section \ref{spaces-properties-section-locally-projective}.

\begin{lemma}
\label{lemma-formally-smooth-sheaf-differentials}
Soit $S$ un schéma.
Soit $f : X \to Y$ un morphisme formellement lisse d'espaces algébriques sur $S$.
Alors $\Omega_{X/Y}$ est localement projectif sur $X$.
\end{lemma}

\begin{proof}
Choisissons un diagramme
$$
\xymatrix{
U \ar[d] \ar[r]_\psi & V \ar[d] \\
X \ar[r]^f & Y
}
$$
où $U$ et $V$ sont des schémas affines (!) et où les flèches verticales sont étales.
D'après le
Lemme \ref{lemma-helper-formally-smooth},
nous voyons que $\psi : U \to V$ est formellement lisse. Ainsi,
$\Gamma(V, \mathcal{O}_V) \to \Gamma(U, \mathcal{O}_U)$ est
un homomorphisme d'anneaux formellement lisse, voir
Compléments sur les morphismes, Lemme \ref{more-morphisms-lemma-affine-formally-smooth}.
Par conséquent, d'après
Algèbre, Lemme \ref{algebra-lemma-characterize-formally-smooth-again},
le $\Gamma(U, \mathcal{O}_U)$-module
$\Omega_{\Gamma(U, \mathcal{O}_U)/\Gamma(V, \mathcal{O}_V)}$
est projectif. Ainsi, $\Omega_{U/V}$ est localement projectif, voir
Propriétés, Section \ref{properties-section-locally-projective}.
Puisque $\Omega_{X/Y}|_U = \Omega_{U/V}$, nous voyons que $\Omega_{X/Y}$ est
lui aussi localement projectif. (En effet, nous pouvons trouver un recouvrement étale de
$X$ par des $U$ affines s'insérant dans des diagrammes comme ci-dessus --- les détails sont
omis.)
\end{proof}

\begin{lemma}
\label{lemma-h1-is-zero}
Soit $T$ un schéma affine.
Soient $\mathcal{F}$ et $\mathcal{G}$ des
$\mathcal{O}_T$-modules quasi-cohérents sur $T_\etale$.
Considérons le faisceau Hom interne
$\mathcal{H} = \SheafHom_{\mathcal{O}_T}(\mathcal{F}, \mathcal{G})$
sur $T_\etale$.
Si $\mathcal{F}$ est localement projectif, alors
$H^1(T_\etale, \mathcal{H}) = 0$.
\end{lemma}

\begin{proof}
D'après la définition d'un faisceau localement projectif sur un espace algébrique (voir
Propriétés d'espaces,
Définition \ref{spaces-properties-definition-locally-projective}),
nous voyons que $\mathcal{F}_{Zar} = \mathcal{F}|_{T_{Zar}}$ est un faisceau
localement projectif sur le schéma $T$. Ainsi, $\mathcal{F}_{Zar}$ est un
facteur direct d'un $\mathcal{O}_{T_{Zar}}$-module libre. Nous
concluons alors (puisque $\mathcal{F} = (\mathcal{F}_{Zar})^a$, voir
Descente, Proposition \ref{descent-proposition-equivalence-quasi-coherent})
que $\mathcal{F}$ est un facteur direct d'un $\mathcal{O}_T$-module libre
sur $T_\etale$. Nous pouvons donc supposer que
$\mathcal{F} = \bigoplus_{i \in I} \mathcal{O}_T$ est un module libre.
Dans ce cas, $\mathcal{H} = \prod_{i \in I} \mathcal{G}$ est
un produit de modules quasi-cohérents. D'après
Cohomologie sur les sites, Lemme \ref{sites-cohomology-lemma-cohomology-products},
nous concluons que $H^1 = 0$, car la cohomologie d'un faisceau quasi-cohérent
sur un schéma affine est nulle, voir
Descente, Proposition \ref{descent-proposition-same-cohomology-quasi-coherent}
et Cohomologie des schémas, Lemme
\ref{coherent-lemma-quasi-coherent-affine-cohomology-zero}.
\end{proof}

\begin{lemma}
\label{lemma-formally-smooth}
Soit $S$ un schéma.
Soit $f : X \to Y$ un morphisme d'espaces algébriques sur
$S$. Les assertions suivantes sont équivalentes :
\begin{enumerate}
\item $f$ est formellement lisse,
\item pour tout diagramme
$$
\xymatrix{
U \ar[d] \ar[r]_\psi & V \ar[d] \\
X \ar[r]^f & Y
}
$$
où $U$ et $V$ sont des schémas et où les flèches verticales sont étales,
le morphisme de schémas $\psi$ est formellement lisse (au sens de
Compléments sur les morphismes,
Définition \ref{more-morphisms-definition-formally-unramified}), et
\item pour un tel diagramme à flèches verticales surjectives, le morphisme
$\psi$ est formellement lisse.
\end{enumerate}
\end{lemma}

\begin{proof}
Nous avons vu que (1) implique (2) et (3) dans le
Lemme \ref{lemma-helper-formally-smooth}.
Supposons (3).
La démonstration du fait que $f$ est formellement lisse est entièrement semblable à
la démonstration de (1) $\Rightarrow$ (2) du
Lemme \ref{lemma-smooth-formally-smooth}.

\medskip\noindent
Considérons un diagramme commutatif de flèches pleines
$$
\xymatrix{
X \ar[d]_f & T \ar[d]^i \ar[l]^a \\
Y & T' \ar[l] \ar@{-->}[lu]
}
$$
comme dans la Définition \ref{definition-formally-smooth}.
Nous allons montrer que la flèche pointillée existe, ce qui
démontrera que $f$ est formellement lisse.
Soit $\mathcal{F}$ le faisceau d'ensembles sur $(T')_{spaces, \etale}$ du
Lemme \ref{lemma-sheaf}, dans le cas particulier examiné à la
Remarque \ref{remark-special-case}.
Soit
$$
\mathcal{H} =
\SheafHom_{\mathcal{O}_T}(a^*\Omega_{X/Y}, \mathcal{C}_{T/T'})
$$
le faisceau de $\mathcal{O}_T$-modules sur $T_{spaces, \etale}$
muni de l'action $\mathcal{H} \times \mathcal{F} \to \mathcal{F}$ comme dans le
Lemme \ref{lemma-action-sheaf}.
L'action $\mathcal{H} \times \mathcal{F} \to \mathcal{F}$
fait de $\mathcal{F}$ un pseudo-torseur sous $\mathcal{H}$, voir
Cohomologie sur les sites, Définition \ref{sites-cohomology-definition-torsor}.
Notre but est de montrer que $\mathcal{F}$ est un torseur trivial sous $\mathcal{H}$.
Il y a deux étapes : (I) Pour montrer que $\mathcal{F}$ est un torseur,
nous devons montrer que $\mathcal{F}$ possède localement pour la topologie étale
une section. (II) Pour montrer que $\mathcal{F}$ est le torseur trivial,
il suffit de montrer que $H^1(T_\etale, \mathcal{H}) = 0$, voir
Cohomologie sur les sites, Lemme \ref{sites-cohomology-lemma-torsors-h1}.

\medskip\noindent
Démontrons d'abord (I). Pour cela, considérons un diagramme
(qui existe puisque nous supposons (3))
$$
\xymatrix{
U \ar[d] \ar[r]_\psi & V \ar[d] \\
X \ar[r]^f & Y
}
$$
où $U$ et $V$ sont des schémas, où les flèches verticales sont étales et
surjectives, et où $\psi$ est formellement lisse. D'après le
Lemme \ref{lemma-representable-property-formally-property},
le morphisme $V \to Y$ est formellement étale. Ainsi, d'après le
Lemme \ref{lemma-composition-formally-smooth-etale-unramified},
la composée $U \to Y$ est formellement lisse. L'étape (I) résulte alors du
Lemme \ref{lemma-etale-on-top}, partie (4).

\medskip\noindent
Démontrons enfin (II). D'après le
Lemme \ref{lemma-formally-smooth-sheaf-differentials},
nous voyons que $\Omega_{U/V}$ est localement projectif.
Ainsi, $\Omega_{X/Y}$ est localement projectif, voir
Descente sur les espaces,
Lemme \ref{spaces-descent-lemma-locally-projective-descends}.
Par conséquent, $a^*\Omega_{X/Y}$ est localement projectif, voir
Propriétés d'espaces, Lemme
\ref{spaces-properties-lemma-locally-projective-pullback}.
Ainsi,
$$
H^1(T_\etale, \mathcal{H}) =
H^1(T_\etale,
\SheafHom_{\mathcal{O}_T}(a^*\Omega_{X/Y}, \mathcal{C}_{T/T'}) = 0
$$
d'après le
Lemme \ref{lemma-h1-is-zero},
comme voulu.
\end{proof}

\begin{lemma}
\label{lemma-descending-property-formally-smooth}
La propriété $\mathcal{P}(f) =$ « $f$ est formellement lisse »
est locale pour la topologie fpqc sur la base.
\end{lemma}

\begin{proof}
Soit $f : X \to Y$ un morphisme d'espaces algébriques sur un schéma $S$.
Choisissons un ensemble d'indices $I$ et des diagrammes
$$
\xymatrix{
U_i \ar[d] \ar[r]_{\psi_i} & V_i \ar[d] \\
X \ar[r]^f & Y
}
$$
à flèches verticales étales, où $U_i$ et $V_i$ sont des schémas affines. De plus,
supposons que $\coprod U_i \to X$ et $\coprod V_i \to Y$ soient surjectifs, voir
Propriétés d'espaces,
Lemme \ref{spaces-properties-lemma-cover-by-union-affines}.
D'après le
Lemme \ref{lemma-formally-smooth},
nous voyons que $f$ est formellement lisse si et seulement si chacun des morphismes
$\psi_i$ est formellement lisse. Nous nous ramenons donc au cas d'un morphisme
de schémas affines. Dans ce cas, le résultat résulte de
Algèbre, Lemme \ref{algebra-lemma-descent-formally-smooth}.
Certains détails sont omis.
\end{proof}

\begin{lemma}
\label{lemma-triangle-differentials-formally-smooth}
Soit $S$ un schéma.
Soient $f : X \to Y$ et $g : Y \to Z$ des morphismes d'espaces algébriques sur $S$.
Supposons que $f$ soit formellement lisse. Alors la suite
$$
0 \to f^*\Omega_{Y/Z} \to \Omega_{X/Z} \to \Omega_{X/Y} \to 0
$$
du Lemme \ref{lemma-triangle-differentials}
est exacte courte.
\end{lemma}

\begin{proof}
Cela résulte du cas des schémas, voir
Compléments sur les morphismes,
Lemme \ref{more-morphisms-lemma-triangle-differentials-formally-smooth},
par localisation étale, voir les
Lemmes \ref{lemma-formally-smooth} et \ref{lemma-localize-differentials}.
\end{proof}

\begin{lemma}
\label{lemma-differentials-formally-unramified-formally-smooth}
Soit $S$ un schéma. Soit $B$ un espace algébrique sur $S$.
Soit $h : Z \to X$ un morphisme formellement non ramifié d'espaces algébriques
sur $B$.
Supposons que $Z$ soit formellement lisse sur $B$. Alors la
suite exacte canonique
$$
0 \to \mathcal{C}_{Z/X} \to h^*\Omega_{X/B} \to \Omega_{Z/B} \to 0
$$
du
Lemme \ref{lemma-universally-unramified-differentials-sequence}
est exacte courte.
\end{lemma}

\begin{proof}
Soit $Z \to Z'$ l'épaississement universel du premier ordre de $Z$ sur $X$.
La démonstration du
Lemme \ref{lemma-universally-unramified-differentials-sequence}
montre que notre suite s'identifie à la suite
$$
\mathcal{C}_{Z/Z'} \to \Omega_{Z'/B} \otimes \mathcal{O}_Z \to
\Omega_{Z/B} \to 0.
$$
Puisque $Z \to S$ est formellement lisse, nous pouvons, localement pour la topologie étale sur $Z'$, trouver
un inverse à gauche $Z' \to Z$ sur $B$ du morphisme d'inclusion $Z \to Z'$.
Ainsi, la suite est localement scindée pour la topologie étale, voir le
Lemme \ref{lemma-differentials-relative-immersion-section}.
\end{proof}

\begin{lemma}
\label{lemma-two-unramified-morphisms-formally-smooth}
Soit $S$ un schéma. Soit
$$
\xymatrix{
Z \ar[r]_i \ar[rd]_j & X \ar[d]^f \\
& Y
}
$$
un diagramme commutatif d'espaces algébriques sur $S$,
où $i$ et $j$ sont formellement non ramifiés et où $f$ est formellement lisse.
Alors la suite exacte canonique
$$
0 \to
\mathcal{C}_{Z/Y} \to
\mathcal{C}_{Z/X} \to
i^*\Omega_{X/Y} \to 0
$$
du
Lemme \ref{lemma-two-unramified-morphisms}
est exacte et localement scindée.
\end{lemma}

\begin{proof}
Notons $Z \to Z'$ l'épaississement universel du premier ordre de $Z$ sur $X$.
Notons $Z \to Z''$ l'épaississement universel du premier ordre de $Z$ sur $Y$.
D'après le
Lemme \ref{lemma-universally-unramified-differentials-sequence},
il existe un morphisme canonique $Z' \to Z''$ tel que nous ayons un diagramme
commutatif
$$
\xymatrix{
Z \ar[r]_{i'} \ar[rd]_{j'} & Z' \ar[r]_a \ar[d]^k & X \ar[d]^f \\
& Z'' \ar[r]^b & Y
}
$$
La suite ci-dessus s'identifie à la suite
$$
\mathcal{C}_{Z/Z''} \to
\mathcal{C}_{Z/Z'} \to
(i')^*\Omega_{Z'/Z''} \to 0
$$
par nos définitions des faisceaux conormaux aux morphismes formellement non ramifiés.
Soit $U'' \to Z''$ un morphisme étale avec $U''$ affine.
Notons $U \to Z$ et $U' \to Z'$ les schémas affines correspondants,
étales sur $Z$ et $Z'$.
Comme $f$ est formellement lisse, il existe un morphisme $h : U'' \to X$
qui coïncide avec $i$ sur $U$ et tel que $f \circ h$ soit égal à $b|_{U''}$.
Puisque $Z'$ est l'épaississement universel du premier ordre, nous obtenons un unique
morphisme $g : U'' \to Z'$ tel que $g = a \circ h$. La propriété
universelle de $Z''$ implique que $k \circ g$ est le morphisme d'inclusion
$U'' \to Z''$. Ainsi, $g$ est un inverse à gauche de $k$. La situation est
$$
\xymatrix{
U \ar[d] \ar[r] & Z' \ar[d]^k \\
U'' \ar[r] \ar[ru]^g & Z''
}
$$
Ainsi, $g$ induit un morphisme $\mathcal{C}_{Z/Z'}|_U \to \mathcal{C}_{Z/Z''}|_U$
qui est un inverse à gauche du morphisme
$\mathcal{C}_{Z/Z''} \to \mathcal{C}_{Z/Z'}$ sur $U$.
\end{proof}











\section{Lissité sur une base noethérienne}
\label{section-smooth-Noetherian}

\noindent
Cette section est l'analogue de
Compléments sur les morphismes, Section \ref{more-morphisms-section-smooth-Noetherian}.

\begin{lemma}
\label{lemma-lifting-along-artinian-at-point}
Soit $S$ un schéma. Soit $f : X \to Y$ un morphisme d'espaces algébriques
sur $S$. Soit $x \in |X|$.
Supposons que $Y$ soit localement noethérien et $f$ localement de type fini.
Les assertions suivantes sont équivalentes :
\begin{enumerate}
\item $f$ est lisse en $x$,
\item pour tout diagramme commutatif de flèches pleines
$$
\xymatrix{
X \ar[d]_f & \Spec(B) \ar[d]^i \ar[l]^-\alpha \\
Y & \Spec(B') \ar[l]_-{\beta} \ar@{-->}[lu]
}
$$
où $B' \to B$ est une surjection d'anneaux locaux telle que
$\Ker(B' \to B)$ soit de carré nul, et où $\alpha$ envoie le
point fermé de $\Spec(B)$ sur $x$, il existe
une flèche pointillée qui rend le diagramme commutatif, et
\item la même condition qu'en (2), mais où $B' \to B$ parcourt les petites
extensions (voir Algèbre, Définition \ref{algebra-definition-small-extension}).
\end{enumerate}
\end{lemma}

\begin{proof}
La condition (1) signifie qu'il existe un sous-espace ouvert $X' \subset X$
tel que $X' \to Y$ soit lisse. Ainsi, (1) implique les conditions (2) et (3) d'après le
Lemme \ref{lemma-smooth-formally-smooth}. La condition (2) implique
trivialement la condition (3). Supposons (3). Choisissons un diagramme commutatif
$$
\xymatrix{
X \ar[d] & U \ar[l] \ar[d] \\
Y & V \ar[l]
}
$$
où $U$ et $V$ sont affines, où les flèches horizontales sont étales, et
tel qu'il existe un point $u \in U$ envoyé sur $x$. Considérons ensuite
un diagramme
$$
\xymatrix{
X \ar[d] & U \ar[l] \ar[d] & \Spec(B) \ar[d]^i \ar[l]^-\alpha  \\
Y & V \ar[l] & \Spec(B') \ar[l]_-{\beta}
}
$$
comme en (3), mais pour $u \in U \to V$. Soit $\gamma : \Spec(B') \to X$
la flèche fournie par notre hypothèse que (3) vaut pour $X$.
Puisque $U \to X$ est étale et
donc formellement étale (Lemme \ref{lemma-etale-formally-etale}),
le morphisme $\gamma$
possède un unique relèvement à $U$ compatible avec $\alpha$. Puisque
$V \to Y$ est étale, donc formellement étale, ce relèvement est compatible
avec $\beta$. Ainsi, (3) vaut pour $u \in U \to V$, et nous concluons
que $U \to V$ est lisse en $u$ d'après
Compléments sur les morphismes, Lemme
\ref{more-morphisms-lemma-lifting-along-artinian-at-point}.
Cela montre que $X \to Y$ est lisse en $x$, et
achève la démonstration.
\end{proof}

\noindent
Il est parfois utile de savoir qu'il suffit de vérifier le
critère de relèvement pour les petites extensions « centrées » en des points
de type fini (voir
Morphismes d'espaces, Section \ref{spaces-morphisms-section-points-finite-type}).

\begin{lemma}
\label{lemma-lifting-along-artinian}
Soit $S$ un schéma. Soit $f : X \to Y$ un morphisme d'espaces algébriques
sur $S$. Supposons que $Y$ soit localement noethérien et $f$ localement de type fini.
Les assertions suivantes sont équivalentes :
\begin{enumerate}
\item $f$ est lisse,
\item pour tout diagramme commutatif de flèches pleines
$$
\xymatrix{
X \ar[d]_f & \Spec(B) \ar[d]^i \ar[l]^-\alpha \\
Y & \Spec(B') \ar[l]_-{\beta} \ar@{-->}[lu]
}
$$
où $B' \to B$ est une petite extension d'anneaux locaux artiniens
et où $\beta$ est de type fini (!), il existe une flèche pointillée qui rend
le diagramme commutatif.
\end{enumerate}
\end{lemma}

\begin{proof}
Si $f$ est lisse, alors le critère infinitésimal de relèvement
(Lemme \ref{lemma-smooth-formally-smooth}) affirme que
$f$ est formellement lisse, et (2) est vérifiée.

\medskip\noindent
Supposons que $f$ ne soit pas lisse. L'ensemble des points $x \in X$ où $f$ n'est pas lisse
forme une partie fermée $T$ de $|X|$. D'après
Morphismes d'espaces, Lemme 
\ref{spaces-morphisms-lemma-enough-finite-type-points}, il existe
un point $x \in T \subset X$ tel que $x \in X_{\text{ft-pts}}$. Choisissons un
diagramme commutatif
$$
\xymatrix{
X \ar[d] & U \ar[l] \ar[d] & u \ar@{|->}[d] \\
Y & V \ar[l] & v
}
$$
où $U$ et $V$ sont affines, où les flèches horizontales sont étales, et
tel qu'il existe un point $u \in U$ envoyé sur $x$. Alors $u$
est un point de type fini de $U$. Puisque $U \to V$ n'est pas lisse au
point $u$, d'après
Compléments sur les morphismes,
Lemme \ref{more-morphisms-lemma-lifting-along-artinian-at-point},
il existe un diagramme
$$
\xymatrix{
X \ar[d] & U \ar[l] \ar[d] & \Spec(B) \ar[d]^i \ar[l]^-\alpha  \\
Y & V \ar[l] & \Spec(B') \ar[l]_-{\beta} \ar@{-->}[lu]
}
$$
où $B' \to B$ est une petite extension d'anneaux locaux (artiniens),
telle que le corps résiduel de $B$ soit égal à $\kappa(v)$ et
que la flèche pointillée n'existe pas. Puisque $U \to V$ est
de type fini, nous voyons que $v$ est un point de type fini de $V$. D'après
Morphismes, Lemme \ref{morphisms-lemma-artinian-finite-type},
le morphisme $\beta$ est de type fini ; par conséquent, la composée
$\Spec(B) \to Y$ est elle aussi de type fini.
En raisonnant exactement comme dans la démonstration du
Lemme \ref{lemma-lifting-along-artinian-at-point}
(en utilisant que $U \to X$ et $V \to Y$ sont étales, donc
formellement étales),
nous voyons qu'il ne peut exister de flèche $\Spec(B) \to X$
s'insérant dans le rectangle extérieur du dernier diagramme affiché.
Autrement dit, (2) n'est pas vérifiée, et la démonstration est terminée.
\end{proof}

\noindent
Voici une application utile.

\begin{lemma}
\label{lemma-check-smoothness-on-infinitesimal-nbhds}
Soit $S$ un schéma. Soit $f : X \to Y$ un morphisme
d'espaces algébriques sur $S$. Supposons que $f$ soit
localement de type fini et $Y$ localement noethérien.
Soit $Z \subset Y$ un sous-espace fermé dont le voisinage infinitésimal d'ordre $n$ est
$Z_n \subset Y$. Posons $X_n = Z_n \times_Y X$.
\begin{enumerate}
\item Si $X_n \to Z_n$ est lisse pour tout $n$, alors $f$
est lisse en tout point de $f^{-1}(Z)$.
\item Si $X_n \to Z_n$ est étale pour tout $n$, alors $f$
est étale en tout point de $f^{-1}(Z)$.
\end{enumerate}
\end{lemma}

\begin{proof}
Supposons que $X_n \to Z_n$ soit lisse pour tout $n$.
Soit $x \in X$ un point situé au-dessus d'un point de $Z$.
Étant donnés une petite extension $B' \to B$ et des morphismes $\alpha$, $\beta$ comme dans le
Lemme \ref{lemma-lifting-along-artinian-at-point}, partie (3),
l'idéal maximal de $B'$ est nilpotent (puisque $B'$ est artinien),
et le morphisme $\beta$ se factorise donc par $Z_n$, tandis que $\alpha$ se factorise
par $X_n$ pour un $n$ convenable. Ainsi, la propriété de relèvement pour
$X_n \to Z_n$ s'applique et fournit la flèche pointillée souhaitée dans le diagramme.
Cela démontre (1). La partie (2) résulte de (1) et du fait qu'un morphisme
est étale si et seulement s'il est lisse de dimension relative $0$.
\end{proof}











\section{Le complexe cotangent naïf}
\label{section-netherlander}

\noindent
Cette section prolonge
Modules sur les sites, Section \ref{sites-modules-section-netherlander},
qui prolonge à son tour la discussion menée dans
Algèbre, Section \ref{algebra-section-netherlander}.

\begin{definition}
\label{definition-netherlander}
Soit $S$ un schéma.
Soit $f : X \to Y$ un morphisme d'espaces algébriques sur $S$.
Le {\it complexe cotangent naïf de $f$}
est le complexe défini dans Modules sur les sites, Définition
\ref{sites-modules-definition-cotangent-complex-morphism-ringed-topoi}
pour le morphisme de topos annelés $f_{small}$ entre les
petits sites étales de $X$ et $Y$, voir
Propriétés d'espaces, Lemme
\ref{spaces-properties-lemma-morphism-ringed-topoi}.
Notation : $\NL_f$ ou $\NL_{X/Y}$.
\end{definition}

\noindent
Les lemmes suivants montrent que cette définition est compatible avec celle
pour les homomorphismes d'anneaux et pour les schémas, et que $\NL_{X/Y}$ est un
objet de $D_\QCoh(\mathcal{O}_X)$.

\begin{lemma}
\label{lemma-NL-etale-localization}
Soit $S$ un schéma. Considérons un diagramme commutatif
$$
\xymatrix{
U \ar[d]_p \ar[r]_g & V \ar[d]^q \\
X \ar[r]^f & Y
}
$$
d'espaces algébriques sur $S$ où $p$ et $q$ sont étales.
Alors il existe une identification canonique
$\NL_{X/Y}|_{U_\etale} = \NL_{U/V}$ dans $D(\mathcal{O}_U)$.
\end{lemma}

\begin{proof}
La formation du complexe cotangent naïf commute à l'image inverse
(Modules sur les sites, Lemme \ref{sites-modules-lemma-pullback-NL})
et nous avons $p_{small}^{-1}\mathcal{O}_X = \mathcal{O}_U$ ainsi que
$g_{small}^{-1}\mathcal{O}_{V_\etale} =
p_{small}^{-1}f_{small}^{-1}\mathcal{O}_{Y_\etale}$
car $q_{small}^{-1}\mathcal{O}_{Y_\etale} =
\mathcal{O}_{V_\etale}$ d'après Propriétés d'espaces, Lemme
\ref{spaces-properties-lemma-etale-exact-pullback}.
En suivant les définitions, nous concluons que
$\NL_{X/Y}|_{U_\etale} = \NL_{U/V}$.
\end{proof}

\begin{lemma}
\label{lemma-NL-compare-spaces-schemes}
Soit $S$ un schéma. Soit $f : X \to Y$ un morphisme d'espaces algébriques
sur $S$. Supposons $X$ et $Y$ représentés par des schémas $X_0$ et $Y_0$.
Alors il existe une identification canonique
$\NL_{X/Y} = \epsilon^*\NL_{X_0/Y_0}$ dans $D(\mathcal{O}_X)$
où $\epsilon$ est celui de Catégories dérivées d'espaces, Section
\ref{spaces-perfect-section-derived-quasi-coherent-etale}
et $\NL_{X_0/Y_0}$ est défini dans
Compléments sur les morphismes, Définition
\ref{more-morphisms-definition-netherlander}.
\end{lemma}

\begin{proof}
Soit $f_0 : X_0 \to Y_0$ le morphisme de schémas correspondant à $f$.
Il existe un morphisme canonique
$\epsilon^{-1}f_0^{-1}\mathcal{O}_{Y_0} \to f_{small}^{-1}\mathcal{O}_Y$
compatible avec
$\epsilon^\sharp : \epsilon^{-1}\mathcal{O}_{X_0} \to \mathcal{O}_X$
car il existe un diagramme commutatif
$$
\xymatrix{
X_{0, Zar} \ar[d]_{f_0} & X_\etale \ar[l]^\epsilon \ar[d]^f \\
Y_{0, Zar} & Y_\etale \ar[l]_\epsilon
}
$$
voir Catégories dérivées d'espaces, Remarque
\ref{spaces-perfect-remark-match-total-direct-images}.
Nous obtenons ainsi un morphisme canonique
$$
\epsilon^{-1}\NL_{X_0/Y_0} =
\epsilon^{-1}\NL_{\mathcal{O}_{X_0}/f_0^{-1}\mathcal{O}_{Y_0}} =
\NL_{\epsilon^{-1}\mathcal{O}_{X_0}/\epsilon^{-1}f_0^{-1}\mathcal{O}_{Y_0}}
\to
\NL_{\mathcal{O}_X/f^{-1}_{small}\mathcal{O}_Y} = \NL_{X/Y}
$$
par fonctorialité du complexe cotangent naïf.
Pour voir que le morphisme induit $\epsilon^*\NL_{X_0/Y_0} \to \NL_{X/Y}$ est un
isomorphisme dans $D(\mathcal{O}_X)$, il suffit de le vérifier sur les fibres aux points géométriques
(Propriétés d'espaces, Théorème
\ref{spaces-properties-theorem-exactness-stalks}).
Soit $\overline{x} : \Spec(k) \to X_0$
un point géométrique situé au-dessus de $x \in X_0$, et soit
$\overline{y} = f \circ \overline{x}$ situé au-dessus de $y \in Y_0$. Alors
$$
\NL_{X/Y, \overline{x}} =
\NL_{\mathcal{O}_{X, \overline{x}}/\mathcal{O}_{Y, \overline{y}}}
$$
Cela est vrai parce que le passage à la fibre en $\overline{x}$ revient
à prendre l'image inverse par $\overline{x} : \Spec(k) \to X$
et que nous pouvons appliquer Modules sur les sites, Lemme \ref{sites-modules-lemma-pullback-NL}.
D'autre part, nous avons
$$
(\epsilon^*\NL_{X_0/Y_0})_{\overline{x}} =
\NL_{X_0/Y_0, x} \otimes_{\mathcal{O}_{X_0, x}}
\mathcal{O}_{X, \overline{x}} =
\NL_{\mathcal{O}_{X_0, x}/\mathcal{O}_{Y_0, y}}
\otimes_{\mathcal{O}_{X_0, x}} \mathcal{O}_{X, \overline{x}}
$$
Certains détails sont omis (indication : utiliser le fait que la fibre d'une image inverse
est la fibre au point image, voir
Sites, Lemme \ref{sites-lemma-point-morphism-sites},
ainsi que le résultat correspondant pour les modules, voir
Modules sur les sites, Lemme \ref{sites-modules-lemma-pullback-stalk}).
Observons que $\mathcal{O}_{X, \overline{x}}$ est l'hensélisé
strict de $\mathcal{O}_{X_0, x}$, et de même
pour $\mathcal{O}_{Y, \overline{y}}$
(Propriétés d'espaces, Lemme
\ref{spaces-properties-lemma-describe-etale-local-ring}).
Le résultat découle donc de
Compléments sur l'algèbre,
Lemme \ref{more-algebra-lemma-henselization-NL}.
\end{proof}

\begin{lemma}
\label{lemma-netherlander-quasi-coherent}
Soit $S$ un schéma. Soit $f : X \to Y$ un morphisme
d'espaces algébriques sur $S$. Les faisceaux de cohomologie
du complexe $\NL_{X/Y}$ sont quasi-cohérents et nuls en dehors
des degrés $-1$, $0$ ; le faisceau $\Omega_{X/Y}$ est sa cohomologie en degré $0$.
\end{lemma}

\begin{proof}
Par construction du complexe cotangent naïf dans
Modules sur les sites, Section \ref{sites-modules-section-netherlander},
nous savons que $\NL_{X/Y}$ est un complexe nul en dehors des degrés $-1$, $0$
et que sa cohomologie en degré $0$ est $\Omega_{X/Y}$ (d'après notre
construction de $\Omega_{X/Y}$ à la Section \ref{section-sheaf-differentials}).
Le faisceau des différentielles est quasi-cohérent (d'après le
Lemme \ref{lemma-module-differentials-quasi-coherent}).
Pour achever la démonstration, il suffit de montrer que $H^{-1}(\NL_{X/Y})$
est quasi-cohérent. Cela résulte d'une vérification locale pour la topologie étale
(permise par le Lemme \ref{lemma-NL-etale-localization} et
Propriétés d'espaces, Lemme
\ref{spaces-properties-lemma-characterize-quasi-coherent})
qui nous ramène au cas des schémas
(Lemme \ref{lemma-NL-compare-spaces-schemes}),
puis de l'emploi du résultat dans le cas des schémas
(Compléments sur les morphismes, Lemme
\ref{more-morphisms-lemma-netherlander-quasi-coherent}).
\end{proof}

\begin{lemma}
\label{lemma-netherlander-fp}
Soit $S$ un schéma.
Soit $f : X \to Y$ un morphisme d'espaces algébriques sur $S$.
Si $f$ est localement de
présentation finie, alors $\NL_{X/Y}$ est, localement pour la topologie étale sur $X$,
quasi-isomorphe à un complexe
$$
\ldots \to 0 \to \mathcal{F}^{-1} \to \mathcal{F}^0 \to 0 \to \ldots
$$
de $\mathcal{O}_X$-modules quasi-cohérents
où $\mathcal{F}^0$ est de présentation finie
et $\mathcal{F}^{-1}$ de type fini.
\end{lemma}

\begin{proof}
La formation du complexe cotangent naïf commute à la
localisation étale d'après le Lemme \ref{lemma-NL-etale-localization}.
Nous nous ramenons ainsi au cas des schémas par le
Lemme \ref{lemma-NL-compare-spaces-schemes}.
Dans le cas des schémas, le résultat est
Compléments sur les morphismes, Lemme
\ref{more-morphisms-lemma-netherlander-fp}.
\end{proof}

\begin{lemma}
\label{lemma-NL-formally-smooth}
Soit $S$ un schéma.
Soit $f : X \to Y$ un morphisme d'espaces algébriques sur $S$.
Les assertions suivantes sont équivalentes :
\begin{enumerate}
\item $f$ est formellement lisse,
\item $H^{-1}(\NL_{X/Y}) = 0$ et le module $H^0(\NL_{X/Y}) = \Omega_{X/Y}$
est localement projectif.
\end{enumerate}
\end{lemma}

\begin{proof}
Cela résulte du
Lemme \ref{lemma-formally-smooth}, du
Lemme \ref{lemma-NL-etale-localization}, du
Lemme \ref{lemma-NL-compare-spaces-schemes}
et du cas des schémas, traité dans
Compléments sur les morphismes, Lemme \ref{more-morphisms-lemma-NL-formally-smooth}.
\end{proof}

\begin{lemma}
\label{lemma-NL-formally-etale}
Soit $f : X \to Y$ un morphisme de schémas. Les assertions suivantes sont équivalentes :
\begin{enumerate}
\item $f$ est formellement étale,
\item $H^{-1}(\NL_{X/Y}) = H^0(\NL_{X/Y}) = 0$.
\end{enumerate}
\end{lemma}

\begin{proof}
Supposons (1). Tout morphisme formellement étale est formellement lisse.
Ainsi $H^{-1}(\NL_{X/Y}) = 0$ d'après le Lemme \ref{lemma-NL-formally-smooth}.
D'autre part, un morphisme formellement étale est formellement non ramifié ;
par conséquent, $\Omega_{X/Y} = 0$ d'après le
Lemme \ref{lemma-characterize-formally-unramified}.
Réciproquement, si (2) est vérifiée, alors $f$ est formellement lisse d'après le
Lemme \ref{lemma-NL-formally-smooth}
et formellement non ramifié d'après le
Lemme \ref{lemma-characterize-formally-unramified},
et donc formellement étale d'après le
Lemme \ref{lemma-formally-etale-unramified-smooth}.
\end{proof}

\begin{lemma}
\label{lemma-NL-smooth}
Soit $f : X \to Y$ un morphisme de schémas. Les assertions suivantes sont équivalentes :
\begin{enumerate}
\item $f$ est lisse, et
\item $f$ est localement de présentation finie,
$H^{-1}(\NL_{X/Y}) = 0$, et le module $H^0(\NL_{X/Y}) = \Omega_{X/Y}$
est localement libre de type fini.
\end{enumerate}
\end{lemma}

\begin{proof}
Cela résulte du
Lemme \ref{lemma-formally-smooth}, du
Lemme \ref{lemma-NL-etale-localization}, du
Lemme \ref{lemma-NL-compare-spaces-schemes}
et du cas des schémas, traité dans
Compléments sur les morphismes, Lemme \ref{more-morphisms-lemma-NL-smooth}.
\end{proof}











\section{Ouverture du lieu de platitude}
\label{section-open-flat}

\noindent
Cette section est l'analogue de
Compléments sur les morphismes, Section \ref{more-morphisms-section-open-flat}.
Notons que nous avons défini la notion de platitude pour les modules
quasi-cohérents sur les espaces algébriques dans
Morphismes d'espaces, Section \ref{spaces-morphisms-section-flat-modules}.

\begin{theorem}
\label{theorem-openness-flatness}
Soit $S$ un schéma.
Soit $f : X \to Y$ un morphisme d'espaces algébriques sur $S$.
Soit $\mathcal{F}$ un faisceau quasi-cohérent sur $X$.
Supposons que $f$ soit localement de présentation finie et que
$\mathcal{F}$ soit un $\mathcal{O}_X$-module
localement de présentation finie. Alors
$$
\{x \in |X| : \mathcal{F}\text{ is flat over }Y\text{ at }x\}
$$
est ouvert dans $|X|$.
\end{theorem}

\begin{proof}
Choisissons un diagramme commutatif
$$
\xymatrix{
U \ar[d]_p \ar[r]_\alpha &
V \ar[d]^q \\
X \ar[r]^a & Y
}
$$
où $U$, $V$ sont des schémas et $p$, $q$ sont étales et surjectifs comme dans
Espaces, Lemme \ref{spaces-lemma-lift-morphism-presentations}.
D'après
Compléments sur les morphismes, Théorème \ref{more-morphisms-theorem-openness-flatness},
l'ensemble
$U' = \{u \in |U| : p^*\mathcal{F}\text{ is flat over }V\text{ at }u\}$
est ouvert dans $U$. D'après
Morphismes d'espaces, Définition \ref{spaces-morphisms-definition-flat-module},
l'image de $U'$ dans $|X|$ est l'ensemble
du théorème. La démonstration est donc terminée, car l'application $|U| \to |X|$ est
ouverte, voir
Propriétés d'espaces, Lemme \ref{spaces-properties-lemma-topology-points}.
\end{proof}

\begin{lemma}
\label{lemma-flat-locus-base-change}
Soit $S$ un schéma. Soit
$$
\xymatrix{
X' \ar[r]_{g'} \ar[d]_{f'} & X \ar[d]^f \\
Y' \ar[r]^g & Y
}
$$
un diagramme cartésien d'espaces algébriques sur $S$.
Soit $\mathcal{F}$ un $\mathcal{O}_X$-module quasi-cohérent.
Supposons que $g$ soit plat, que $f$ soit localement de présentation finie,
et que $\mathcal{F}$ soit localement de présentation finie.
Alors
$$
\{x' \in |X'| : (g')^*\mathcal{F}\text{ is flat over }Y'\text{ at }x'\}
$$
est l'image inverse de la partie ouverte du
Théorème \ref{theorem-openness-flatness}
par l'application continue $|g'| : |X'| \to |X|$.
\end{lemma}

\begin{proof}
Cela découle de
Morphismes d'espaces, Lemme
\ref{spaces-morphisms-lemma-base-change-module-flat}.
\end{proof}












\section{Crit\`ere de platitude par fibres}
\label{section-criterion-flat-fibres}

\noindent
Soit $S$ un schéma. Considérons un diagramme commutatif d'espaces algébriques
sur $S$
$$
\xymatrix{
X \ar[rr]_f \ar[dr]_g & & Y \ar[dl]^h \\
& Z
}
$$
et un $\mathcal{O}_X$-module quasi-cohérent $\mathcal{F}$.
Étant donné un point $x \in |X|$, nous nous demandons si
$\mathcal{F}$ est plat sur $Y$ en $x$. Si $\mathcal{F}$ est plat
sur $Z$ en $x$, le théorème ci-dessous affirme que cette question est
étroitement liée à celle de savoir si la restriction de
$\mathcal{F}$ à la fibre de $X \to Z$ au-dessus de $g(x)$
est plate sur la fibre de $Y \to Z$ au-dessus de $g(x)$. Pour donner un sens
à cet énoncé, nous proposons le lemme préliminaire suivant.

\begin{lemma}
\label{lemma-flat-on-fibres-at-point}
Dans la situation ci-dessus, les assertions suivantes sont équivalentes :
\begin{enumerate}
\item Choisissons un point géométrique $\overline{x}$ de $X$ situé au-dessus de $x$.
Posons $\overline{y} = f \circ \overline{x}$ et
$\overline{z} = g \circ \overline{x}$. Alors le module
$\mathcal{F}_{\overline{x}}/
\mathfrak m_{\overline{z}}\mathcal{F}_{\overline{x}}$
est plat sur
$\mathcal{O}_{Y, \overline{y}}/
\mathfrak m_{\overline{z}}\mathcal{O}_{Y, \overline{y}}$.
\item Choisissons un morphisme $x : \Spec(K) \to X$ dans la classe d'équivalence de
$x$. Posons $z = g \circ x$, $X_z = \Spec(K) \times_{z, Z} X$,
$Y_z = \Spec(K) \times_{z, Z} Y$, et soit $\mathcal{F}_z$ l'image inverse
de $\mathcal{F}$ sur $X_z$. Alors $\mathcal{F}_z$ est plat en $x$ sur
$Y_z$ (au sens de Morphismes d'espaces,
Définition \ref{spaces-morphisms-definition-flat-module}).
\item Choisissons un diagramme commutatif
$$
\xymatrix{
& & & U \ar[llld]_a \ar[rr] \ar[dr] & & V \ar[llld]_>>>>>>>b \ar[dl] \\
X \ar[rr]_f \ar[dr]_g & & Y \ar[dl]^h &  & W \ar[llld]_c \\
& Z
}
$$
où $U, V, W$ sont des schémas et $a, b, c$ sont étales,
ainsi qu'un point $u \in U$ envoyé sur $x$. Soit $w \in W$ l'image de
$u$. Soit $\mathcal{F}_w$ l'image inverse de $\mathcal{F}$ sur
la fibre $U_w$ de $U \to W$ en $w$. Alors $\mathcal{F}_w$
est plat sur $V_w$ en $u$.
\end{enumerate}
\end{lemma}

\begin{proof}
Notons qu'en (2), le morphisme $x : \Spec(K) \to X$ définit un
point $K$-rationnel de $X_z$ ; l'énoncé a donc bien un sens. En outre,
la condition de (2) est indépendante du choix de $\Spec(K) \to X$
dans la classe d'équivalence de $x$ (les détails sont omis ; cela résultera aussi
des arguments ci-dessous, puisque les autres conditions ne dépendent pas
de ce choix). Notons également que nous pouvons toujours choisir un diagramme comme en
(3) de la manière suivante : choisissons d'abord
un schéma $W$ et un morphisme étale surjectif $W \to Z$, puis
un schéma $V$ et un morphisme étale surjectif $V \to W \times_Z Y$, et
enfin un schéma $U$ et un morphisme étale surjectif
$U \to V \times_Y X$. Ces choix étant faits, nous prenons pour $U \to W$
la composée $U \to V \to W$ et pouvons choisir un point $u \in U$ envoyé
sur $x$, puisque le morphisme $U \to X$ est surjectif.

\medskip\noindent
Supposons donnés à la fois un diagramme comme en (3) et un point géométrique
$\overline{x} : \Spec(k) \to X$ comme en (1). D'après
Propriétés d'espaces, Lemme
\ref{spaces-properties-lemma-geometric-lift-to-usual}
nous pouvons choisir un point géométrique $\overline{u} : \Spec(k) \to U$
au-dessus de $u$ tel que $\overline{x} = a \circ \overline{u}$.
Notons $\overline{v} : \Spec(k) \to V$ et
$\overline{w} : \Spec(k) \to W$ les points géométriques induits de
$V$ et $W$. Dans cette situation, nous savons que
$\mathcal{O}_{X, \overline{x}} = \mathcal{O}_{U, u}^{sh}$
et de même pour $Y$ et $Z$, voir
Propriétés d'espaces,
Lemme \ref{spaces-properties-lemma-describe-etale-local-ring}.
De même, nous avons
$$
\mathcal{F}_{\overline{x}} =
(a^*\mathcal{F})_u \otimes_{\mathcal{O}_{U, u}}
\mathcal{O}_{U, u}^{sh}
$$
voir
Propriétés d'espaces, Lemme \ref{spaces-properties-lemma-stalk-quasi-coherent}.
Notons que la fibre de $\mathcal{F}_w$ en $u$ est donnée par
$$
(\mathcal{F}_w)_u = (a^*\mathcal{F})_u/\mathfrak m_w(a^*\mathcal{F})_u
$$
et que l'anneau local de $V_w$ en $v$ est donné par
$$
\mathcal{O}_{V_w, v} = \mathcal{O}_{V, v}/\mathfrak m_w\mathcal{O}_{V, v}.
$$
Puisque $\mathfrak m_{\overline{z}} =
\mathfrak m_w \mathcal{O}_{Z, \overline{z}} =
\mathfrak m_w \mathcal{O}_{W, w}^{sh}$
nous voyons que
\begin{align*}
\mathcal{F}_{\overline{x}}/
\mathfrak m_{\overline{z}}\mathcal{F}_{\overline{x}} & =
(a^*\mathcal{F})_u \otimes_{\mathcal{O}_{U, u}}
\mathcal{O}_{X, \overline{x}}/
\mathfrak m_{\overline{z}}\mathcal{O}_{X, \overline{x}} \\
& =
(\mathcal{F}_w)_u \otimes_{\mathcal{O}_{U_w, u}}
\mathcal{O}_{U, u}^{sh}/\mathfrak m_w\mathcal{O}_{U, u}^{sh} \\
& = (\mathcal{F}_w)_u \otimes_{\mathcal{O}_{U_w, u}}
\mathcal{O}_{U_w, \overline{u}}^{sh} \\
& = (\mathcal{F}_w)_{\overline{u}}
\end{align*}
l'avant-dernière égalité résultant du
Lemme \ref{algebra-lemma-quotient-strict-henselization} d'Algèbre
et la dernière du
Lemme \ref{spaces-properties-lemma-stalk-quasi-coherent} de Propriétés d'espaces.
Les mêmes arguments appliqués aux faisceaux structuraux de $V$ et $Y$
montrent que
$$
\mathcal{O}_{V_w, \overline{v}}^{sh} =
\mathcal{O}_{V, v}^{sh}/\mathfrak m_w \mathcal{O}_{V, v}^{sh} =
\mathcal{O}_{Y, \overline{y}}/
\mathfrak m_{\overline{z}}\mathcal{O}_{Y, \overline{y}}.
$$
Nous pouvons maintenant appliquer
Morphismes d'espaces, Lemme \ref{spaces-morphisms-lemma-flat-at-point},
pour voir que (1) est équivalente à (3).

\medskip\noindent
Enfin, démontrons l'équivalence de (2) et (3).
Pour cela, choisissons une extension de corps $\tilde K$ de $K$
et un morphisme $\tilde x : \Spec(\tilde K) \to U$ qui
est situé au-dessus de $u$ (c'est possible car $u \times_{X, x} \Spec(K)$
est un schéma non vide). Soit $\tilde z : \Spec(\tilde K) \to U \to W$
la composée. Nous obtenons un diagramme commutatif
$$
\xymatrix{
& & & U_w \times_w \tilde z \ar[llld]_a \ar[rr] \ar[dr] & &
V_w \times_w \tilde z \ar[llld]_>>>>>>>b \ar[dl] \\
X_z \ar[rr]_f \ar[dr]_g & & Y_z \ar[dl]^h &  & \tilde z \ar[llld]_c \\
& z
}
$$
où $z = \Spec(K)$ et $w = \Spec(\kappa(w))$. Il est alors
clair que $\mathcal{F}_w$ et $\mathcal{F}_z$ ont pour images inverses le
même module sur $U_w \times_w \tilde z$. Nous obtenons ainsi un diagramme
commutatif
$$
\xymatrix{
X_z \ar[d] & U_w \times_w \tilde z \ar[l] \ar[d] \ar[r] & U_w \ar[d] \\
Y_z & V_w \times_w \tilde z \ar[l] \ar[r] & V_w
}
$$
dont les deux carrés sont cartésiens et dont les flèches horizontales du bas
sont plates : la flèche horizontale inférieure gauche est la composée
du morphisme $Y \times_Z \tilde z \to Y \times_Z z = Y_z$ (changement de base
d'un morphisme plat), du morphisme étale
$V \times_Z \tilde z \to Y \times_Z \tilde z$, et
du morphisme étale $V \times_W \tilde z \to V \times_Z \tilde z$.
Il résulte donc de
Morphismes d'espaces,
Lemme \ref{spaces-morphisms-lemma-base-change-module-flat},
que
$$
\mathcal{F}_z\text{ flat at }x\text{ over }Y_z
\Leftrightarrow
\mathcal{F}|_{U_w \times_w \tilde z}
\text{ flat at }\tilde x\text{ over }V_w \times_w \tilde z
\Leftrightarrow
\mathcal{F}_w\text{ flat at }u\text{ over }V_w
$$
ce qui conclut.
\end{proof}

\begin{definition}
\label{definition-module-flat-on-fibre}
Soit $S$ un schéma. Soient $X \to Y \to Z$ des morphismes d'espaces
algébriques sur $S$. Soit $\mathcal{F}$ un $\mathcal{O}_X$-module quasi-cohérent.
Soit $x \in |X|$ un point et notons $z \in |Z|$ son image.
\begin{enumerate}
\item Nous disons que {\it la restriction de $\mathcal{F}$ à sa fibre au-dessus de $z$
est plate en $x$ sur la fibre de $Y$ au-dessus de $z$} si les conditions équivalentes du
Lemme \ref{lemma-flat-on-fibres-at-point}
sont vérifiées.
\item Nous disons que {\it la fibre de $X$ au-dessus de $z$ est plate en $x$ sur la fibre de
$Y$ au-dessus de $z$} si les conditions équivalentes du
Lemme \ref{lemma-flat-on-fibres-at-point}
sont vérifiées pour $\mathcal{F} = \mathcal{O}_X$.
\item Nous disons que {\it la fibre de $X$ au-dessus de $z$ est plate sur la fibre de $Y$
au-dessus de $z$} si, pour tout $x \in |X|$ situé au-dessus de $z$, la fibre de $X$ au-dessus de $z$
est plate en $x$ sur la fibre de $Y$ au-dessus de $z$.
\end{enumerate}
\end{definition}

\noindent
Cette définition nous permet d'énoncer la version suivante du critère.
La version noethérienne se trouve à la
Section \ref{section-flat-over-Noetherian}.

\begin{theorem}
\label{theorem-criterion-flatness-fibre}
Soit $S$ un schéma.
Soient $f : X \to Y$ et $Y \to Z$ des morphismes d'espaces algébriques sur $S$.
Soit $\mathcal{F}$ un $\mathcal{O}_X$-module quasi-cohérent.
Supposons que
\begin{enumerate}
\item $X$ soit localement de présentation finie sur $Z$,
\item $\mathcal{F}$ soit un $\mathcal{O}_X$-module de présentation finie, et
\item $Y$ soit localement de type fini sur $Z$.
\end{enumerate}
Soit $x \in |X|$ et soient $y \in |Y|$ et $z \in |Z|$ les images de
$x$. Si $\mathcal{F}_{\overline{x}} \not = 0$, les assertions suivantes sont
équivalentes :
\begin{enumerate}
\item $\mathcal{F}$ est plat sur $Z$ en $x$ et
la restriction de $\mathcal{F}$ à sa fibre au-dessus de $z$
est plate en $x$ sur la fibre de $Y$ au-dessus de $z$, et
\item $Y$ est plat sur $Z$ en $y$ et $\mathcal{F}$ est
plat sur $Y$ en $x$.
\end{enumerate}
En outre, l'ensemble des points $x$ où (1) et (2) sont vérifiées est ouvert dans
$\text{Supp}(\mathcal{F})$.
\end{theorem}

\begin{proof}
Choisissons un diagramme comme dans le
Lemme \ref{lemma-flat-on-fibres-at-point}, partie (3).
Il résulte des définitions que l'on se ramène au
théorème correspondant pour les morphismes de schémas
$U \to V \to W$, le faisceau quasi-cohérent $a^*\mathcal{F}$
et le point $u \in U$. Le théorème découle donc du
résultat correspondant pour les schémas, à savoir
Compléments sur les morphismes,
Théorème \ref{more-morphisms-theorem-criterion-flatness-fibre}.
\end{proof}

\begin{lemma}
\label{lemma-morphism-between-flat}
Soit $S$ un schéma.
Soient $f : X \to Y$ et $Y \to Z$ des morphismes d'espaces algébriques sur $S$.
Supposons que
\begin{enumerate}
\item $X$ soit localement de présentation finie sur $Z$,
\item $X$ soit plat sur $Z$,
\item pour tout $z \in |Z|$, la fibre de $X$ au-dessus de $z$
soit plate sur la fibre de $Y$ au-dessus de $z$, et
\item $Y$ soit localement de type fini sur $Z$.
\end{enumerate}
Alors $f$ est plat. Si $f$ est en outre surjectif, alors $Y$ est plat sur $Z$.
\end{lemma}

\begin{proof}
C'est un cas particulier du
Théorème \ref{theorem-criterion-flatness-fibre}.
\end{proof}

\begin{lemma}
\label{lemma-base-change-criterion-flatness-fibre}
Soit $S$ un schéma. Soient $f : X \to Y$ et $Y \to Z$ des morphismes
d'espaces algébriques sur $S$. Soit $\mathcal{F}$ un
$\mathcal{O}_X$-module quasi-cohérent.
Supposons que
\begin{enumerate}
\item $X$ soit localement de présentation finie sur $Z$,
\item $\mathcal{F}$ soit un $\mathcal{O}_X$-module de présentation finie,
\item $\mathcal{F}$ soit plat sur $Z$, et
\item $Y$ soit localement de type fini sur $Z$.
\end{enumerate}
Alors l'ensemble
$$
A = \{x \in |X| : \mathcal{F} \text{ flat at }x \text{ over }Y\}.
$$
Cet ensemble est ouvert dans $|X|$ et sa formation commute à tout changement de base :
si $Z' \to Z$ est un morphisme d'espaces algébriques et si $A'$ est l'ensemble des
points de $X' = X \times_Z Z'$ où $\mathcal{F}' = \mathcal{F} \times_Z Z'$
est plat sur $Y' = Y \times_Z Z'$, alors $A'$ est l'image inverse de
$A$ par l'application continue $|X'| \to |X|$.
\end{lemma}

\begin{proof}
Une façon de le démontrer consiste à transposer la démonstration donnée dans
Compléments sur les morphismes, Lemme \ref{more-morphisms-lemma-morphism-between-flat},
à la catégorie des espaces algébriques. Nous le démontrerons plutôt
en nous ramenant au cas des schémas. Choisissons donc un diagramme comme dans le
Lemme \ref{lemma-flat-on-fibres-at-point}, partie (3),
tel que $a$, $b$ et $c$ soient surjectifs.
Il résulte des définitions que l'on se ramène au
théorème correspondant pour les morphismes de schémas
$U \to V \to W$, le faisceau quasi-cohérent $a^*\mathcal{F}$
et le point $u \in U$. Le seul point à préciser est que,
étant donné un morphisme d'espaces algébriques $Z' \to Z$, nous choisissons un schéma
$W'$ et un morphisme étale surjectif $W' \to W \times_Z Z'$.
Posons alors $U' = W' \times_W U$ et $V' = W' \times_W V$.
Notons $a', b', c'$ les morphismes de $U', V', W'$ dans
$X', Y', Z'$. Dans ce cas, $A$, resp.\ $A'$ sont les images des parties
ouvertes de $U$, resp.\ $U'$ associées à
$a^*\mathcal{F}$, resp.\ $(a')^*\mathcal{F}'$.
Cela ramène bien le lemme à
Compléments sur les morphismes, Lemme \ref{more-morphisms-lemma-morphism-between-flat}.
\end{proof}

\begin{lemma}
\label{lemma-base-change-flatness-fibres}
Soit $S$ un schéma.
Soient $f : X \to Y$ et $Y \to Z$ des morphismes d'espaces algébriques sur $S$.
Supposons que
\begin{enumerate}
\item $X$ soit localement de présentation finie sur $Z$,
\item $X$ soit plat sur $Z$, et
\item $Y$ soit localement de type fini sur $Z$.
\end{enumerate}
Alors l'ensemble
$$
\{x \in |X| : X\text{ flat at }x \text{ over }Y\}.
$$
Cet ensemble est ouvert dans $|X|$ et sa formation commute à tout changement de base
$Z' \to Z$.
\end{lemma}

\begin{proof}
C'est un cas particulier du
Lemme \ref{lemma-base-change-criterion-flatness-fibre}.
\end{proof}

\begin{lemma}
\label{lemma-flat-and-free-at-point-fibre}
Soit $S$ un schéma. Soit $f : X \to Y$ un morphisme d'espaces algébriques
sur $S$ qui est localement de présentation finie.
Soit $\mathcal{F}$ un $\mathcal{O}_X$-module de présentation finie.
Soit $x \in |X|$, d'image $y \in |Y|$. Si $\mathcal{F}$ est plat en $x$
sur $Y$, les assertions suivantes sont équivalentes :
\begin{enumerate}
\item $(\mathcal{F}_{\overline{y}})_{\overline{x}}$ est un module plat sur
$\mathcal{O}_{X_{\overline{y}}, \overline{x}}$,
\item $(\mathcal{F}_{\overline{y}})_{\overline{x}}$ est un module libre sur
$\mathcal{O}_{X_{\overline{y}}, \overline{x}}$,
\item $\mathcal{F}_{\overline{y}}$ est libre de type fini sur un
voisinage étale de $\overline{x}$ dans $X_{\overline{y}}$, et
\item $\mathcal{F}$ est libre de type fini sur un voisinage étale de $x$ dans $X$.
\end{enumerate}
Ici, $\overline{x}$ est un point géométrique de $X$ situé au-dessus de $x$
et $\overline{y} = f \circ \overline{x}$.
\end{lemma}

\begin{proof}
Choisissons un diagramme commutatif
$$
\xymatrix{
U \ar[d] \ar[r] & V \ar[d] \\
X \ar[r] & Y
}
$$
où $U$ et $V$ sont des schémas et les flèches verticales sont étales,
et tel qu'il existe un point $u \in U$ envoyé sur $x$. Soit $v \in V$
l'image de $u$. En appliquant le Lemme \ref{lemma-flat-on-fibres-at-point}
à $\text{id} : X \to X$ sur $Y$, nous voyons que (1) devient
la condition « $\mathcal{F}|_{U_v}$ est plat sur $U_v$ en $u$ ».
Autrement dit, (1) équivaut à ce que $(\mathcal{F}|_{U_v})_u$
soit un $\mathcal{O}_{U_v, u}$-module plat.
D'après le cas des schémas (Compléments sur les morphismes, Lemme
\ref{more-morphisms-lemma-flat-and-free-at-point-fibre}),
nous voyons que cela implique que
$\mathcal{F}|_U$ est libre de type fini sur un voisinage ouvert
de $u$. Nous voyons ainsi que (1) implique (4).
Les implications (4) $\Rightarrow$ (3) et
(2) $\Rightarrow$ (1) sont immédiates.
Pour l'implication (3) $\Rightarrow$ (2), utilisons
la description des anneaux locaux et des fibres dans
Propriétés d'espaces, Lemmes
\ref{spaces-properties-lemma-describe-etale-local-ring} et
\ref{spaces-properties-lemma-stalk-quasi-coherent}.
\end{proof}

\begin{lemma}
\label{lemma-finite-free-open}
Soit $S$ un schéma. Soit $f : X \to Y$ un morphisme d'espaces algébriques
sur $S$ qui est localement de présentation finie.
Soit $\mathcal{F}$ un $\mathcal{O}_X$-module de présentation finie
plat sur $Y$. Alors l'ensemble
$$
\{x \in |X| : \mathcal{F}\text{ free in an \'etale neighbourhood of }x\}
$$
est ouvert dans $|X|$ et sa formation commute à tout changement de base
$Y' \to Y$.
\end{lemma}

\begin{proof}
L'ouverture est immédiate. Soit $Y' \to Y$ un morphisme d'espaces algébriques,
posons $X' = Y' \times_Y X$,
et soit $x' \in |X'|$ un point situé au-dessus de $x \in |X|$.
D'après le Lemme \ref{lemma-flat-and-free-at-point-fibre},
nous voyons que $x$ appartient à notre ensemble si et seulement si
$(\mathcal{F}_{\overline{y}})_{\overline{x}}$ est un module plat sur
$\mathcal{O}_{X_{\overline{y}}, \overline{x}}$.
De même, $x'$ appartient à l'analogue de notre ensemble pour l'image inverse
$\mathcal{F}'$ de $\mathcal{F}$ sur $X'$ si et seulement si
$(\mathcal{F}'_{\overline{y}'})_{\overline{x}'}$ est un module plat sur
$\mathcal{O}_{X'_{\overline{y}'}, \overline{x}'}$
(avec les notations évidentes). Ces deux assertions sont équivalentes
d'après le Lemme \ref{lemma-flat-on-fibres-at-point} appliqué au
morphisme $\text{id} : X \to X$ sur $Y$.
L'assertion sur le changement de base est donc démontrée.
\end{proof}







\section{Platitude sur une base noethérienne}
\label{section-flat-over-Noetherian}

\noindent
Voici le « critère de platitude par fibres » dans le cas noethérien.

\begin{theorem}
\label{theorem-criterion-flatness-fibre-Noetherian}
Soit $S$ un schéma.
Soient $f : X \to Y$ et $Y \to Z$ des morphismes d'espaces algébriques sur $S$.
Soit $\mathcal{F}$ un $\mathcal{O}_X$-module quasi-cohérent.
Supposons que
\begin{enumerate}
\item $X$, $Y$, $Z$ soient localement noethériens, et
\item $\mathcal{F}$ soit un $\mathcal{O}_X$-module cohérent.
\end{enumerate}
Soit $x \in |X|$ et soient $y \in |Y|$ et $z \in |Z|$ les images de
$x$. Si $\mathcal{F}_{\overline{x}} \not = 0$, les assertions suivantes sont
équivalentes :
\begin{enumerate}
\item $\mathcal{F}$ est plat sur $Z$ en $x$ et
la restriction de $\mathcal{F}$ à sa fibre au-dessus de $z$
est plate en $x$ sur la fibre de $Y$ au-dessus de $z$, et
\item $Y$ est plat sur $Z$ en $y$ et $\mathcal{F}$ est
plat sur $Y$ en $x$.
\end{enumerate}
\end{theorem}

\begin{proof}
Choisissons un diagramme comme dans le
Lemme \ref{lemma-flat-on-fibres-at-point}, partie (3).
Il résulte des définitions que l'on se ramène au
théorème correspondant pour les morphismes de schémas
$U \to V \to W$, le faisceau quasi-cohérent $a^*\mathcal{F}$
et le point $u \in U$. Le théorème découle donc du
résultat correspondant pour les schémas, à savoir
Compléments sur les morphismes,
Théorème \ref{more-morphisms-theorem-criterion-flatness-fibre-Noetherian}.
\end{proof}

\begin{lemma}
\label{lemma-morphism-between-flat-Noetherian}
Soit $S$ un schéma.
Soient $f : X \to Y$ et $Y \to Z$ des morphismes d'espaces algébriques sur $S$.
Supposons que
\begin{enumerate}
\item $X$, $Y$, $Z$ soient localement noethériens,
\item $X$ soit plat sur $Z$,
\item pour tout $z \in |Z|$, la fibre de $X$ au-dessus de $z$
soit plate sur la fibre de $Y$ au-dessus de $z$.
\end{enumerate}
Alors $f$ est plat. Si $f$ est en outre surjectif, alors $Y$ est plat sur $Z$.
\end{lemma}

\begin{proof}
C'est un cas particulier du
Théorème \ref{theorem-criterion-flatness-fibre-Noetherian}.
\end{proof}

\noindent
Tout comme pour vérifier la lissité, lorsque la base est noethérienne, il suffit
de vérifier la platitude sur les anneaux artiniens. Voici un énoncé type.

\begin{lemma}
\label{lemma-flatness-over-Noetherian-ring}
Soit $A$ un anneau noethérien. Soit $I \subset A$ un idéal.
Soit $X$ un espace algébrique localement de présentation finie sur
$S = \Spec(A)$. Pour $n \geq 1$, posons $S_n = \Spec(A/I^n)$ et
$X_n = S_n \times_S X$. Soit $\mathcal{F}$ un $\mathcal{O}_X$-module cohérent.
Si, pour tout $n \geq 1$, l'image inverse $\mathcal{F}_n$ de $\mathcal{F}$ sur $X$
est plate sur $S_n$, alors le lieu (ouvert) où $\mathcal{F}$
est plate sur $X$ contient l'image inverse de $V(I)$ par $X \to S$.
\end{lemma}

\begin{proof}
Le lieu où $\mathcal{F}$ est plat sur $S$ est ouvert dans $|X|$ d'après le
Théorème \ref{theorem-openness-flatness}.
L'énoncé est invariant lorsqu'on remplace $X$ par les membres d'un
recouvrement étale ; nous pouvons donc supposer que $X$ est un schéma affine.
Dans ce cas, le résultat découle immédiatement du
Lemme \ref{algebra-lemma-flat-module-powers} d'Algèbre.
Certains détails sont omis.
\end{proof}







\section{Retour sur la normalisation}
\label{section-normalization}

\noindent
La normalisation commute au changement de base lisse.

\begin{lemma}
\label{lemma-integral-closure-smooth-pullback}
Soit $S$ un schéma. Soit $f : Y \to X$ un morphisme lisse
d'espaces algébriques sur $S$. Soit $\mathcal{A}$ un faisceau quasi-cohérent
de $\mathcal{O}_X$-algèbres. La clôture intégrale
de $\mathcal{O}_Y$ dans $f^*\mathcal{A}$ est égale à $f^*\mathcal{A}'$,
où $\mathcal{A}' \subset \mathcal{A}$ est la clôture intégrale de
$\mathcal{O}_X$ dans $\mathcal{A}$.
\end{lemma}

\begin{proof}
Par notre construction de la clôture intégrale, voir
Morphismes d'espaces, Définition
\ref{spaces-morphisms-definition-integral-closure},
on se ramène immédiatement au cas où $X$ et $Y$ sont affines.
Dans ce cas, le résultat est le
Lemme \ref{algebra-lemma-integral-closure-commutes-smooth} d'Algèbre.
\end{proof}

\begin{lemma}[La normalisation commute au changement de base lisse]
\label{lemma-normalization-smooth-localization}
Soit $S$ un schéma. Soit
$$
\xymatrix{
Y_2 \ar[r] \ar[d] & Y_1 \ar[d]^f \\
X_2 \ar[r]^\varphi & X_1
}
$$
un carré cartésien d'espaces algébriques sur $S$. Supposons que $f$ soit quasi-compact
et quasi-séparé, et que $\varphi$ soit lisse.
Soit $Y_i \to X_i' \to X_i$ la normalisation de $X_i$ dans $Y_i$.
Alors $X_2' \cong X_2 \times_{X_1} X_1'$.
\end{lemma}

\begin{proof}
La factorisation $Y_1 \to X_1' \to X_1$, après changement de base à $X_2$,
devient une factorisation $Y_2 \to X_2 \times_{X_1} X_1' \to X_1$ et le morphisme
$X_2 \times_{X_1} X_1' \to X_1$ est entier
(Morphismes d'espaces, Lemme \ref{spaces-morphisms-lemma-base-change-integral}).
Nous obtenons donc un morphisme
$h : X_2' \to X_2 \times_{X_1} X_1'$ par la propriété universelle de
Morphismes d'espaces, Lemme
\ref{spaces-morphisms-lemma-characterize-normalization}.
Notons que $X_2'$ est le spectre relatif de la clôture intégrale
de $\mathcal{O}_{X_2}$ dans $f_{2, *}\mathcal{O}_{Y_2}$.
Si $\mathcal{A}' \subset f_{1, *}\mathcal{O}_{Y_1}$ désigne la clôture
intégrale de $\mathcal{O}_{X_2}$, alors $X_2 \times_{X_1} X_1'$ est le
spectre relatif de $\varphi^*\mathcal{A}'$, puisque la formation du
spectre relatif commute à tout changement de base. D'après
Cohomologie des espaces, Lemme
\ref{spaces-cohomology-lemma-flat-base-change-cohomology}
nous savons que $f_{2, *}\mathcal{O}_{Y_2} = \varphi^*f_{1, *}\mathcal{O}_{Y_1}$.
Le résultat découle donc du
Lemme \ref{lemma-integral-closure-smooth-pullback}.
\end{proof}









\section{Morphismes de Cohen-Macaulay}
\label{section-CM}

\noindent
Ceci est l'analogue de Compléments sur les morphismes, Section
\ref{more-morphisms-section-CM}.

\begin{lemma}
\label{lemma-CM-local-ring-fibre}
La propriété suivante des morphismes de germes de schémas
\begin{align*}
& \mathcal{P}((X, x) \to (S, s)) = \\
& \text{the local ring }
\mathcal{O}_{X_s, x}
\text{ of the fibre is Noetherian and Cohen-Macaulay}
\end{align*}
est de nature locale pour la topologie étale sur la source et le but (Descente, Définition
\ref{descent-definition-local-source-target-at-point}).
\end{lemma}

\begin{proof}
Étant donné un diagramme comme dans
Descente, Définition \ref{descent-definition-local-source-target-at-point},
nous obtenons un morphisme étale entre les fibres
$U'_{v'} \to U_v$ qui envoie $u'$ sur $u$ ; voir
Descente, Lemme \ref{descent-lemma-etale-on-fiber}.
Les hensélisés stricts des anneaux locaux
$\mathcal{O}_{U'_{v'}, u'}$ et $\mathcal{O}_{U_v, u}$
sont donc les mêmes. Nous concluons par le
Lemme \ref{more-algebra-lemma-henselization-CM} de Compléments sur l'algèbre.
\end{proof}

\begin{definition}
\label{definition-CM}
Soit $S$ un schéma.
Soit $f : X \to Y$ un morphisme d'espaces algébriques sur $S$.
Supposons que les fibres de $f$ soient localement noethériennes
(Diviseurs sur les espaces, Définition
\ref{spaces-divisors-definition-locally-Noetherian-fibre}).
\begin{enumerate}
\item Soient $x \in |X|$ et $y = f(x)$. Nous disons que $f$ est
{\it Cohen-Macaulay en $x$} si $f$ est plat en $x$ et si
les conditions équivalentes de
Morphismes d'espaces, Lemme
\ref{spaces-morphisms-lemma-local-source-target-at-point}
sont satisfaites pour la propriété $\mathcal{P}$ décrite dans le
Lemme \ref{lemma-CM-local-ring-fibre}.
\item Nous disons que $f$ est un {\it morphisme de Cohen-Macaulay} si $f$ est
Cohen-Macaulay en tout point de $X$.
\end{enumerate}
\end{definition}

\noindent
Voici une reformulation.

\begin{lemma}
\label{lemma-CM}
Soit $S$ un schéma. Soit $f : X \to Y$ un morphisme d'espaces algébriques
sur $S$. Supposons que les fibres de $f$ soient localement noethériennes.
Les conditions suivantes sont équivalentes :
\begin{enumerate}
\item $f$ est Cohen-Macaulay,
\item $f$ est plat et il existe un morphisme étale surjectif $V \to Y$
où $V$ est un schéma, tel que les fibres de $X_V \to V$
soient des espaces algébriques de Cohen-Macaulay, et
\item $f$ est plat et, pour tout morphisme étale $V \to Y$
où $V$ est un schéma, les fibres de $X_V \to V$
sont des espaces algébriques de Cohen-Macaulay.
\end{enumerate}
Étant donné $x \in |X|$ d'image $y \in |Y|$, les conditions suivantes sont
équivalentes :
\begin{enumerate}
\item[(a)] $f$ est Cohen-Macaulay en $x$, et
\item[(b)] $\mathcal{O}_{Y, \overline{y}} \to \mathcal{O}_{X, \overline{x}}$
est plat et
$\mathcal{O}_{X, \overline{x}}/
\mathfrak m_{\overline{y}}\mathcal{O}_{X, \overline{x}}$ est Cohen-Macaulay.
\end{enumerate}
\end{lemma}

\begin{proof}
Étant donné un morphisme étale $V \to Y$ où $V$ est un schéma,
choisissons un schéma $U$ et un morphisme étale surjectif $U \to X \times_Y V$.
Considérons le diagramme commutatif
$$
\xymatrix{
U \ar[d] \ar[r] & V \ar[d] \\
X \ar[r] & Y
}
$$
Soit $u \in U$ d'images $x \in |X|$, $y \in |Y|$ et $v \in V$.
Alors $f$ est Cohen-Macaulay en $x$ si et seulement si $U \to V$ l'est
en $u$ (par définition). De plus, le morphisme
$U_v \to X_v = (X_V)_v$ est étale surjectif. Par conséquent, le schéma $U_v$ est
Cohen-Macaulay si et seulement si l'espace algébrique $X_v$ l'est.
L'équivalence de (1), (2) et (3) découle donc de
l'équivalence correspondante pour les morphismes de
schémas ; voir Compléments sur les morphismes, Lemme \ref{more-morphisms-lemma-CM} ;
l'argument est formel.

\medskip\noindent
Démonstration de l'équivalence de (a) et (b). L'équivalence correspondante
pour la platitude est donnée par Morphismes d'espaces, Lemme
\ref{spaces-morphisms-lemma-flat-at-point-etale-local-rings}.
Nous pouvons donc supposer $f$ plat en $x$ pour démontrer l'équivalence.
Considérons un diagramme et $x, y, u, v$ comme ci-dessus. Alors
$\mathcal{O}_{Y, \overline{y}} \to \mathcal{O}_{X, \overline{x}}$
est égal à l'homomorphisme
$\mathcal{O}_{V, v}^{sh} \to \mathcal{O}_{U, u}^{sh}$
entre les hensélisés stricts des anneaux locaux ; voir
Propriétés des espaces, Lemme
\ref{spaces-properties-lemma-describe-etale-local-ring}.
Nous avons donc
$$
\mathcal{O}_{X, \overline{x}}/
\mathfrak m_{\overline{y}}\mathcal{O}_{X, \overline{x}} =
(\mathcal{O}_{U, u}/\mathfrak m_v \mathcal{O}_{U, u})^{sh}
$$
d'après Algèbre, Lemme \ref{algebra-lemma-quotient-strict-henselization}.
Il reste donc à montrer que l'anneau local noethérien
$\mathcal{O}_{U, u}/\mathfrak m_v \mathcal{O}_{U, u}$
est Cohen-Macaulay si et seulement si son hensélisé strict l'est.
C'est Compléments sur l'algèbre, Lemme \ref{more-algebra-lemma-henselization-CM}.
\end{proof}

\begin{lemma}
\label{lemma-composition-CM}
Soit $S$ un schéma.
Soient $f : X \to Y$ et $g : Y \to Z$ des morphismes d'espaces algébriques
sur $S$. Supposons que les
fibres de $f$, $g$ et $g \circ f$ soient localement noethériennes.
Soit $x \in |X|$ d'images $y \in |Y|$ et $z \in |Z|$.
\begin{enumerate}
\item Si $f$ est Cohen-Macaulay en $x$ et $g$ est Cohen-Macaulay
en $f(x)$, alors $g \circ f$ est Cohen-Macaulay en $x$.
\item Si $f$ et $g$ sont Cohen-Macaulay, alors $g \circ f$ l'est.
\item Si $g \circ f$ est Cohen-Macaulay en $x$ et $f$ est plat en $x$,
alors $f$ est Cohen-Macaulay en $x$ et $g$ l'est en $f(x)$.
\item Si $f \circ g$ est Cohen-Macaulay et $f$ est plat, alors
$f$ est Cohen-Macaulay et $g$ l'est en tout point de
l'image de $f$.
\end{enumerate}
\end{lemma}

\begin{proof}
Le résultat correspondant pour les schémas, appliqué localement pour la
topologie étale, donne l'assertion ; voir
Compléments sur les morphismes, Lemme \ref{more-morphisms-lemma-composition-CM}.
On peut aussi utiliser l'équivalence de (a) et (b) du
Lemme \ref{lemma-CM}. Nous considérons alors l'homomorphisme local
d'anneaux locaux noethériens
$$
\mathcal{O}_{Y, \overline{y}}/
\mathfrak m_{\overline{z}}\mathcal{O}_{Y, \overline{y}}
\longrightarrow
\mathcal{O}_{X, \overline{x}}/
\mathfrak m_{\overline{z}}\mathcal{O}_{X, \overline{x}}
$$
dont la fibre est
$$
\mathcal{O}_{X, \overline{x}}/
\mathfrak m_{\overline{y}}\mathcal{O}_{X, \overline{x}}
$$
et nous appliquons Algèbre, Lemme \ref{algebra-lemma-CM-goes-up}.
\end{proof}

\begin{lemma}
\label{lemma-flat-morphism-from-CM}
Soit $S$ un schéma.
Soit $f : X \to Y$ un morphisme plat entre espaces algébriques localement noethériens
sur $S$.
Si $X$ est Cohen-Macaulay, alors $f$ est Cohen-Macaulay et
$\mathcal{O}_{Y, f(\overline{x})}$ est Cohen-Macaulay pour tout $x \in |X|$.
\end{lemma}

\begin{proof}
En traduisant l'énoncé en algèbre à l'aide du Lemme \ref{lemma-CM}
(comparer avec la démonstration du
Lemme \ref{lemma-composition-CM}), le résultat découle du
Lemme \ref{algebra-lemma-CM-goes-up} d'Algèbre.
\end{proof}

\begin{lemma}
\label{lemma-base-change-CM}
Soit $S$ un schéma.
Soit $f : X \to Y$ un morphisme d'espaces algébriques sur $S$.
Supposons que les fibres de $f$ soient localement noethériennes.
Supposons que $Y' \to Y$ soit localement de type fini. Soit $f' : X' \to Y'$
le changement de base de $f$.
Soit $x' \in |X'|$ un point d'image $x \in |X|$.
\begin{enumerate}
\item Si $f$ est Cohen-Macaulay en $x$, alors
$f' : X' \to Y'$ est Cohen-Macaulay en $x'$.
\item Si $f$ est plat en $x$ et $f'$ est Cohen-Macaulay en $x'$, alors $f$
est Cohen-Macaulay en $x$.
\item Si $Y' \to Y$ est plat en $f'(x')$ et $f'$ est Cohen-Macaulay en
$x'$, alors $f$ est Cohen-Macaulay en $x$.
\end{enumerate}
\end{lemma}

\begin{proof}
Notons $y \in |Y|$ et $y' \in |Y'|$ les images de $x'$.
Choisissons un morphisme étale surjectif $V \to Y$ où $V$ est un schéma.
Choisissons un morphisme étale surjectif $U \to X \times_Y V$ où
$U$ est un schéma.
Choisissons un morphisme étale surjectif $V' \to Y' \times_Y V$
où $V'$ est un schéma.
Alors $U' = U \times_V V'$ est un schéma muni
d'un morphisme étale surjectif $U' \to X'$.
Choisissons $u' \in U'$ qui s'envoie sur $x'$. Notons $u \in U$ l'image de $u'$.
Le lemme découle alors du lemme appliqué à
$U \to V$, à son changement de base $U' \to V'$ et aux
points $u'$ et $u$ (cela découle des définitions).
Il découle donc du cas des schémas ; voir
Compléments sur les morphismes, Lemme \ref{more-morphisms-lemma-base-change-CM}.
\end{proof}

\begin{lemma}
\label{lemma-flat-finite-presentation-CM-open}
Soit $S$ un schéma.
Soit $f : X \to Y$ un morphisme d'espaces algébriques sur $S$
qui est plat et localement de présentation finie. Posons
$$
W = \{x \in |X| : f\text{ is Cohen-Macaulay at }x\}
$$
Alors $W$ est ouvert dans $|X|$ et la formation de $W$
commute à tout changement de base de $f$ :
pour tout morphisme $g : Y' \to Y$, considérons
le changement de base $f' : X' \to Y'$ de $f$ et la
projection $g' : X' \to X$. Alors l'ensemble correspondant
$W'$ pour le morphisme $f'$ est donné par $W' = (g')^{-1}(W)$.
\end{lemma}

\begin{proof}
Choisissons un diagramme commutatif
$$
\xymatrix{
U \ar[d] \ar[r] & V \ar[d] \\
X \ar[r] & Y
}
$$
dont les flèches verticales sont étales et où $U$ et $V$ sont des schémas.
Soit $u \in U$ d'image $x \in |X|$.
Alors $f$ est Cohen-Macaulay en $x$ si et seulement si $U \to V$ l'est
en $u$ (par définition).
Nous nous ramenons ainsi au cas du morphisme $U \to V$.
Voir Compléments sur les morphismes, Lemme
\ref{more-morphisms-lemma-flat-finite-presentation-CM-open}.
\end{proof}

\begin{lemma}
\label{lemma-lfp-CM-relative-dimension}
Soient $S$ un schéma et $f : X \to Y$ un
morphisme d'espaces algébriques sur $S$. Supposons que
$f$ soit localement de présentation finie et Cohen-Macaulay.
Il existe alors des sous-schémas ouverts et fermés $X_d \subset X$
tels que $X = \coprod_{d \geq 0} X_d$ et que $f|_{X_d} : X_d \to Y$
soit de dimension relative $d$.
\end{lemma}

\begin{proof}
Choisissons un diagramme commutatif
$$
\xymatrix{
U \ar[d] \ar[r] & V \ar[d] \\
X \ar[r] & Y
}
$$
dont les flèches verticales sont étales et où $U$ et $V$ sont des schémas.
Alors $U \to V$ est localement de présentation finie et Cohen-Macaulay
(cela résulte immédiatement de nos définitions).
Nous avons donc une décomposition $U = \coprod_{d \geq 0} U_d$
en sous-schémas ouverts et fermés telle que $f|_{U_d} : U_d \to V$
soit de dimension relative $d$ ; voir Morphismes, Lemme
\ref{morphisms-lemma-flat-finite-presentation-CM-fibres-relative-dimension}.
Soit $u \in U$ d'image $x \in |X|$. Alors
$f$ est de dimension relative $d$ en $x$ si et seulement si
$U \to V$ est de dimension relative $d$ en $u$
(cela découle de nos définitions).
Nous voyons ainsi que $U_d$ est l'image inverse
d'une partie $X_d \subset |X|$ qui est nécessairement
ouverte et fermée. Notons $X_d$ le sous-espace algébrique
ouvert et fermé correspondant de $X$ ; le lemme est démontré.
\end{proof}







\section{Morphismes de Gorenstein}
\label{section-gorenstein}

\noindent
Ceci est l'analogue de Dualité pour les schémas, Section
\ref{duality-section-gorenstein-morphisms}.

\begin{lemma}
\label{lemma-gorenstein-local-ring-fibre}
La propriété suivante des morphismes de germes de schémas
\begin{align*}
& \mathcal{P}((X, x) \to (S, s)) = \\
& \text{the local ring }
\mathcal{O}_{X_s, x}
\text{ of the fibre is Noetherian and Gorenstein}
\end{align*}
est de nature locale pour la topologie étale sur la source et le but (Descente, Définition
\ref{descent-definition-local-source-target-at-point}).
\end{lemma}

\begin{proof}
Étant donné un diagramme comme dans
Descente, Définition \ref{descent-definition-local-source-target-at-point},
nous obtenons un morphisme étale entre les fibres
$U'_{v'} \to U_v$ qui envoie $u'$ sur $u$ ; voir
Descente, Lemme \ref{descent-lemma-etale-on-fiber}.
Ainsi, $\mathcal{O}_{U_v, u} \to \mathcal{O}_{U'_{v'}, u'}$
est la localisation d'un homomorphisme étale d'anneaux. Par conséquent,
le premier anneau est noethérien si et seulement si le second l'est ; voir
Compléments sur l'algèbre, Lemme \ref{more-algebra-lemma-Noetherian-etale-extension}.
Puisque $\mathcal{O}_{U'_{v'}, u'}/\mathfrak m_u \mathcal{O}_{U'_{v'}, u'}
= \kappa(u')$ (Algèbre, Lemme \ref{algebra-lemma-etale-at-prime})
est un anneau de Gorenstein, nous voyons que
$\mathcal{O}_{U_v, u}$ est de Gorenstein si et seulement si
$\mathcal{O}_{U'_{v'}, u'}$ l'est, d'après
Complexes dualisants, Lemme \ref{dualizing-lemma-flat-under-gorenstein}.
\end{proof}

\begin{definition}
\label{definition-gorenstein}
Soit $S$ un schéma.
Soit $f : X \to Y$ un morphisme d'espaces algébriques sur $S$.
Supposons que les fibres de $f$ soient localement noethériennes
(Diviseurs sur les espaces, Définition
\ref{spaces-divisors-definition-locally-Noetherian-fibre}).
\begin{enumerate}
\item Soient $x \in |X|$ et $y = f(x)$. Nous disons que $f$ est
{\it de Gorenstein en $x$} si $f$ est plat en $x$ et si
les conditions équivalentes de
Morphismes d'espaces, Lemme
\ref{spaces-morphisms-lemma-local-source-target-at-point}
sont satisfaites pour la propriété $\mathcal{P}$ décrite dans le
Lemme \ref{lemma-gorenstein-local-ring-fibre}.
\item Nous disons que $f$ est un {\it morphisme de Gorenstein} si $f$ est
de Gorenstein en tout point de $X$.
\end{enumerate}
\end{definition}

\noindent
Voici une reformulation.

\begin{lemma}
\label{lemma-gorenstein}
Soit $S$ un schéma. Soit $f : X \to Y$ un morphisme d'espaces algébriques
sur $S$. Supposons que les fibres de $f$ soient localement noethériennes.
Les conditions suivantes sont équivalentes :
\begin{enumerate}
\item $f$ est de Gorenstein,
\item $f$ est plat et il existe un morphisme étale surjectif $V \to Y$
où $V$ est un schéma, tel que les fibres de $X_V \to V$
soient des espaces algébriques de Gorenstein, et
\item $f$ est plat et, pour tout morphisme étale $V \to Y$
où $V$ est un schéma, les fibres de $X_V \to V$
sont des espaces algébriques de Gorenstein.
\end{enumerate}
Étant donné $x \in |X|$ d'image $y \in |Y|$, les conditions suivantes sont
équivalentes :
\begin{enumerate}
\item[(a)] $f$ est de Gorenstein en $x$, et
\item[(b)] $\mathcal{O}_{Y, \overline{y}} \to \mathcal{O}_{X, \overline{x}}$
est plat et
$\mathcal{O}_{X, \overline{x}}/
\mathfrak m_{\overline{y}}\mathcal{O}_{X, \overline{x}}$ est de Gorenstein.
\end{enumerate}
\end{lemma}

\begin{proof}
Étant donné un morphisme étale $V \to Y$ où $V$ est un schéma,
choisissons un schéma $U$ et un morphisme étale surjectif $U \to X \times_Y V$.
Considérons le diagramme commutatif
$$
\xymatrix{
U \ar[d] \ar[r] & V \ar[d] \\
X \ar[r] & Y
}
$$
Soit $u \in U$ d'images $x \in |X|$, $y \in |Y|$ et $v \in V$.
Alors $f$ est de Gorenstein en $x$ si et seulement si $U \to V$ l'est
en $u$ (par définition). De plus, le morphisme
$U_v \to X_v = (X_V)_v$ est étale surjectif. Par conséquent, le schéma $U_v$ est
de Gorenstein si et seulement si l'espace algébrique $X_v$ l'est.
L'équivalence de (1), (2) et (3) découle donc de
l'équivalence correspondante pour les morphismes de
schémas ; voir Dualité pour les schémas, Lemme \ref{duality-lemma-gorenstein} ;
l'argument est formel.

\medskip\noindent
Démonstration de l'équivalence de (a) et (b). L'équivalence correspondante
pour la platitude est donnée par Morphismes d'espaces, Lemme
\ref{spaces-morphisms-lemma-flat-at-point-etale-local-rings}.
Nous pouvons donc supposer $f$ plat en $x$ pour démontrer l'équivalence.
Considérons un diagramme et $x, y, u, v$ comme ci-dessus. Alors
$\mathcal{O}_{Y, \overline{y}} \to \mathcal{O}_{X, \overline{x}}$
est égal à l'homomorphisme
$\mathcal{O}_{V, v}^{sh} \to \mathcal{O}_{U, u}^{sh}$
entre les hensélisés stricts des anneaux locaux ; voir
Propriétés des espaces, Lemme
\ref{spaces-properties-lemma-describe-etale-local-ring}.
Nous avons donc
$$
\mathcal{O}_{X, \overline{x}}/
\mathfrak m_{\overline{y}}\mathcal{O}_{X, \overline{x}} =
(\mathcal{O}_{U, u}/\mathfrak m_v \mathcal{O}_{U, u})^{sh}
$$
d'après Algèbre, Lemme \ref{algebra-lemma-quotient-strict-henselization}.
Il reste donc à montrer que l'anneau local noethérien
$\mathcal{O}_{U, u}/\mathfrak m_v \mathcal{O}_{U, u}$
est de Gorenstein si et seulement si son hensélisé strict l'est.
Cela découle immédiatement de
Complexes dualisants, Lemme
\ref{dualizing-lemma-completion-henselization-dualizing}
et de la définition d'un anneau local de Gorenstein comme un
anneau local noethérien qui est un complexe dualisant sur lui-même.
\end{proof}

\begin{lemma}
\label{lemma-composition-gorenstein}
Soit $S$ un schéma.
Soient $f : X \to Y$ et $g : Y \to Z$ des morphismes d'espaces algébriques
sur $S$. Supposons que les
fibres de $f$, $g$ et $g \circ f$ soient localement noethériennes.
Soit $x \in |X|$ d'images $y \in |Y|$ et $z \in |Z|$.
\begin{enumerate}
\item Si $f$ est de Gorenstein en $x$ et $g$ est de Gorenstein
en $f(x)$, alors $g \circ f$ est de Gorenstein en $x$.
\item Si $f$ et $g$ sont de Gorenstein, alors $g \circ f$ l'est.
\item Si $g \circ f$ est de Gorenstein en $x$ et $f$ est plat en $x$,
alors $f$ est de Gorenstein en $x$ et $g$ l'est en $f(x)$.
\item Si $f \circ g$ est de Gorenstein et $f$ est plat, alors
$f$ est de Gorenstein et $g$ l'est en tout point de
l'image de $f$.
\end{enumerate}
\end{lemma}

\begin{proof}
Le résultat correspondant pour les schémas, appliqué localement pour la
topologie étale, donne l'assertion ; voir Dualité pour les schémas, Lemme
\ref{duality-lemma-composition-gorenstein}.
On peut aussi utiliser l'équivalence de (a) et (b) du
Lemme \ref{lemma-gorenstein}. Nous considérons alors l'homomorphisme local
d'anneaux locaux noethériens
$$
\mathcal{O}_{Y, \overline{y}}/
\mathfrak m_{\overline{z}}\mathcal{O}_{Y, \overline{y}}
\longrightarrow
\mathcal{O}_{X, \overline{x}}/
\mathfrak m_{\overline{z}}\mathcal{O}_{X, \overline{x}}
$$
dont la fibre est
$$
\mathcal{O}_{X, \overline{x}}/
\mathfrak m_{\overline{y}}\mathcal{O}_{X, \overline{x}}
$$
et nous appliquons Complexes dualisants, Lemme
\ref{dualizing-lemma-flat-under-gorenstein}.
\end{proof}

\begin{lemma}
\label{lemma-flat-morphism-from-gorenstein}
Soit $S$ un schéma.
Soit $f : X \to Y$ un morphisme plat entre espaces algébriques localement noethériens
sur $S$.
Si $X$ est de Gorenstein, alors $f$ est de Gorenstein et
$\mathcal{O}_{Y, f(\overline{x})}$ est de Gorenstein pour tout $x \in |X|$.
\end{lemma}

\begin{proof}
En traduisant l'énoncé en algèbre à l'aide du Lemme \ref{lemma-gorenstein}
(comparer avec la démonstration du
Lemme \ref{lemma-composition-gorenstein}), le résultat découle du
Lemme \ref{dualizing-lemma-flat-under-gorenstein} de Complexes dualisants.
\end{proof}

\begin{lemma}
\label{lemma-base-change-gorenstein}
Soit $S$ un schéma.
Soit $f : X \to Y$ un morphisme d'espaces algébriques sur $S$.
Supposons que les fibres de $f$ soient localement noethériennes.
Supposons que $Y' \to Y$ soit localement de type fini. Soit $f' : X' \to Y'$
le changement de base de $f$.
Soit $x' \in |X'|$ un point d'image $x \in |X|$.
\begin{enumerate}
\item Si $f$ est de Gorenstein en $x$, alors
$f' : X' \to Y'$ est de Gorenstein en $x'$.
\item Si $f$ est plat en $x$ et $f'$ est de Gorenstein en $x'$, alors $f$
est de Gorenstein en $x$.
\item Si $Y' \to Y$ est plat en $f'(x')$ et $f'$ est de Gorenstein en
$x'$, alors $f$ est de Gorenstein en $x$.
\end{enumerate}
\end{lemma}

\begin{proof}
Notons $y \in |Y|$ et $y' \in |Y'|$ les images de $x'$.
Choisissons un morphisme étale surjectif $V \to Y$ où $V$ est un schéma.
Choisissons un morphisme étale surjectif $U \to X \times_Y V$ où
$U$ est un schéma.
Choisissons un morphisme étale surjectif $V' \to Y' \times_Y V$
où $V'$ est un schéma.
Alors $U' = U \times_V V'$ est un schéma muni
d'un morphisme étale surjectif $U' \to X'$.
Choisissons $u' \in U'$ qui s'envoie sur $x'$. Notons $u \in U$ l'image de $u'$.
Le lemme découle alors du lemme appliqué à
$U \to V$, à son changement de base $U' \to V'$ et aux
points $u'$ et $u$ (cela découle des définitions).
Il découle donc du cas des schémas ; voir
Dualité pour les schémas, Lemme \ref{duality-lemma-base-change-gorenstein}.
\end{proof}

\begin{lemma}
\label{lemma-flat-finite-presentation-gorenstein-open}
Soit $S$ un schéma.
Soit $f : X \to Y$ un morphisme d'espaces algébriques sur $S$
qui est plat et localement de présentation finie. Posons
$$
W = \{x \in |X| : f\text{ is Gorenstein at }x\}
$$
Alors $W$ est ouvert dans $|X|$ et la formation de $W$
commute à tout changement de base de $f$ :
pour tout morphisme $g : Y' \to Y$, considérons
le changement de base $f' : X' \to Y'$ de $f$ et la
projection $g' : X' \to X$. Alors l'ensemble correspondant
$W'$ pour le morphisme $f'$ est donné par $W' = (g')^{-1}(W)$.
\end{lemma}

\begin{proof}
Choisissons un diagramme commutatif
$$
\xymatrix{
U \ar[d] \ar[r] & V \ar[d] \\
X \ar[r] & Y
}
$$
Soit $u \in U$ d'image $x \in |X|$.
Alors $f$ est de Gorenstein en $x$ si et seulement si $U \to V$ l'est
en $u$ (par définition).
Nous nous ramenons ainsi au cas du morphisme $U \to V$
de schémas. L'ouverture est démontrée dans
Dualité pour les schémas, Lemme
\ref{duality-lemma-flat-finite-presentation-characterize-gorenstein}
et la compatibilité au changement de base dans
Dualité pour les schémas, Lemme \ref{duality-lemma-flat-lft-base-change-gorenstein}.
\end{proof}










\section{Tranchage des morphismes de Cohen-Macaulay}
\label{section-slice}

\noindent
Soit $S$ un schéma. Soit $X$ un espace algébrique sur $S$.
Soient $f_1, \ldots, f_r \in \Gamma(X, \mathcal{O}_X)$. Dans ce cas, nous
notons $V(f_1, \ldots, f_r)$ le {\it sous-espace fermé de $X$ défini par
$f_1, \ldots, f_r$}. Plus précisément, on peut définir $V(f_1, \ldots, f_r)$
comme le sous-espace fermé de $X$ correspondant au faisceau quasi-cohérent
d'idéaux engendré par $f_1, \ldots, f_r$ ; voir
Morphismes d'espaces, Lemme
\ref{spaces-morphisms-lemma-closed-immersion-ideals}.
On peut aussi choisir une présentation $X = U/R$ et considérer le
sous-schéma fermé $Z \subset U$ défini par $f_1|U, \ldots, f_r|_U$.
Il est clair que $Z$ est un sous-schéma fermé invariant par $R$ (voir
Groupoïdes, Définition \ref{groupoids-definition-invariant-open})
et l'on peut poser $V(f_1, \ldots, f_r) = Z/R_Z$.

\begin{lemma}
\label{lemma-slice}
Soit $S$ un schéma. Considérons un diagramme cartésien
$$
\xymatrix{
X \ar[d] & F \ar[l]^p \ar[d] \\
Y & \Spec(k) \ar[l]
}
$$
où $X \to Y$ est un morphisme d'espaces algébriques sur $S$
qui est plat et localement de présentation finie, et où
$k$ est un corps sur $S$. Soient $f_1, \ldots, f_r \in \Gamma(X, \mathcal{O}_X)$
et $z \in |F|$ tels que les images de $f_1, \ldots, f_r$ forment une suite régulière
dans l'anneau local $\mathcal{O}_{F, \overline{z}}$.
Alors, quitte à remplacer $X$ par un sous-espace ouvert contenant $p(z)$, le morphisme
$$
V(f_1, \ldots, f_r) \longrightarrow Y
$$
est plat et localement de présentation finie.
\end{lemma}

\begin{proof}
Posons $Z = V(f_1, \ldots, f_r)$. Il est clair que $Z \to X$ est localement de
présentation finie ; la composée $Z \to Y$ est donc localement de
présentation finie ; voir
Morphismes d'espaces,
Lemme \ref{spaces-morphisms-lemma-composition-finite-presentation}.
Il suffit donc de montrer que $Z \to Y$ est plat dans un voisinage de $p(z)$.
Soit $k'/k$ une extension de corps. Alors le morphisme de
$F' = F \times_{\Spec(k)} \Spec(k')$ vers
$F$ est surjectif et plat ; on peut donc trouver un point $z' \in |F'|$ qui s'envoie sur $z$
et l'homomorphisme local
$\mathcal{O}_{F, \overline{z}} \to \mathcal{O}_{F', \overline{z}'}$ est
plat ; voir
Morphismes d'espaces,
Lemme \ref{spaces-morphisms-lemma-flat-at-point-etale-local-rings}.
Par conséquent, les images de $f_1, \ldots, f_r$ dans
$\mathcal{O}_{F', \overline{z}'}$ forment encore une suite régulière ; voir
Algèbre, Lemme \ref{algebra-lemma-flat-increases-depth}.
Ainsi, au cours de la démonstration, nous pouvons remplacer $k$ par une extension de corps.
En particulier, nous pouvons supposer que $z \in |F|$ provient d'une section
$z : \Spec(k) \to F$ du morphisme structural $F \to \Spec(k)$.

\medskip\noindent
Choisissons un schéma $V$ et un morphisme étale surjectif
$V \to Y$. Choisissons un schéma $U$ et un morphisme étale surjectif
$U \to X \times_Y V$. Après avoir éventuellement agrandi encore $k$, nous pouvons
supposer que $\Spec(k) \to F \to X$ se factorise par $U$ (puisque
$U \to X$ est surjectif). Soit
$u : \Spec(k) \to U$ une telle factorisation et notons $v \in V$
l'image de $u$. Notons que les morphismes
$$
U_v \times_{\Spec(\kappa(v))} \Spec(k) =
U \times_V \Spec(k) \to U \times_Y \Spec(k) \to F
$$
sont étales (le premier comme changement de base de $V \to V \times_Y V$ et
le second comme changement de base de $U \to X$). De plus, par construction,
le point $u : \Spec(k) \to U$ définit un point de l'espace tout à gauche
qui s'envoie sur $z$ à droite. Par conséquent, les éléments
$f_1, \ldots, f_r$ ont pour images une suite régulière dans l'anneau local
situé à droite de l'homomorphisme suivant
$$
\mathcal{O}_{U_v, u}
\longrightarrow
\mathcal{O}_{U_v \times_{\Spec(\kappa(v)} \Spec(k), \overline{u}}
=
\mathcal{O}_{U \times_V \Spec(k), \overline{u}}.
$$
Mais, puisque la flèche affichée est plate (en combinant
Compléments sur la platitude, Lemme \ref{flat-lemma-flat-up-down-henselization}
et
Morphismes d'espaces,
Lemme \ref{spaces-morphisms-lemma-flat-at-point-etale-local-rings}),
nous déduisons de
Algèbre, Lemme \ref{algebra-lemma-flat-increases-depth},
que les images de $f_1, \ldots, f_r$ forment une suite régulière dans
$\mathcal{O}_{U_v, u}$. D'après
Compléments sur les morphismes, Lemme \ref{more-morphisms-lemma-slice-given-elements},
nous concluons que le morphisme de schémas
$$
V(f_1, \ldots, f_r) \times_X U = V(f_1|_U, \ldots, f_r|_U) \to V
$$
est plat dans un voisinage ouvert $U'$ de $u$. Soit $X' \subset X$
le sous-espace ouvert correspondant à l'image de
$|U'| \to |X|$ (voir
Propriétés des espaces, Lemmes
\ref{spaces-properties-lemma-topology-points} et
\ref{spaces-properties-lemma-open-subspaces}).
Nous concluons que $V(f_1, \ldots, f_r) \cap X' \to Y$ est plat
(voir
Morphismes d'espaces, Définition \ref{spaces-morphisms-definition-flat}),
car
nous avons le diagramme commutatif
$$
\xymatrix{
V(f_1, \ldots, f_r) \times_X U' \ar[d]_a \ar[r] & V \ar[d]^b \\
V(f_1, \ldots, f_r) \cap X' \ar[r] & Y
}
$$
où $a, b$ sont étales et $a$ est surjectif.
\end{proof}






\section{Reduced fibres}
\label{section-reduced}

\noindent
This section is the analogue of
More on Morphisms, Section \ref{more-morphisms-section-reduced}.

\begin{lemma}
\label{lemma-geometrically-reduced-fibre}
Let $S$ be a scheme. Let $f : X \to Y$ be a
morphism of algebraic spaces over $S$. Let $y \in |Y|$.
The following are equivalent
\begin{enumerate}
\item for some morphism $\Spec(k) \to Y$ in the equivalence class
of $y$ the algebraic space $X_k$ is geometrically reduced over $k$,
\item for every morphism $\Spec(k) \to Y$ in the equivalence class
of $y$ the algebraic space $X_k$ is geometrically reduced over $k$,
\item for every morphism $\Spec(k) \to Y$ in the equivalence class
of $y$ the algebraic space $X_k$ is reduced.
\end{enumerate}
\end{lemma}

\begin{proof}
This follows immediately from Spaces over Fields, Lemma
\ref{spaces-over-fields-lemma-geometrically-reduced-upstairs}
and the definition of the equivalence relation defining $|X|$
given in
Properties of Spaces, Section \ref{spaces-properties-section-points}.
\end{proof}

\begin{definition}
\label{definition-geometrically-reduced-fibre}
Let $S$ be a scheme. Let $f : X \to Y$ be a
morphism of algebraic spaces over $S$. Let $y \in |Y|$.
We say {\it the fibre of $f : X \to Y$ at $y$ is geometrically reduced}
if the equivalent conditions of
Lemma \ref{lemma-geometrically-reduced-fibre} hold.
\end{definition}

\noindent
Here are the obligatory lemmas.

\begin{lemma}
\label{lemma-base-change-fibres-geometrically-reduced}
Let $S$ be a scheme. Let $f : X \to Y$ and $g : Y' \to Y$
be morphisms of algebraic spaces over $S$. Denote
$f' : X' \to Y'$ the base change of $f$ by $g$. Then
\begin{align*}
\{y' \in |Y'| :
\text{the fibre of }f' : X' \to Y'\text{ at }y'
\text{ is geometrically reduced}\} \\
= g^{-1}(\{y \in |Y| :
\text{the fibre of }f : X \to Y\text{ at }y
\text{ is geometrically reduced}\}).
\end{align*}
\end{lemma}

\begin{proof}
For $y' \in |Y'|$ choose a morphism $\Spec(k) \to Y'$
in the equivalence class of $y'$. Then $g(y')$ is
represented by the composition $\Spec(k) \to Y' \to Y$.
Hence $X' \times_{Y'} \Spec(k) = X \times_Y \Spec(k)$
and the result follows from the definition.
\end{proof}

\begin{lemma}
\label{lemma-geometrically-reduced-constructible}
Let $S$ be a scheme. Let $f : X \to Y$ be a morphism of algebraic spaces
over $S$ which is quasi-compact and
locally of finite presentation. Then the set
$$
E = \{y \in |Y| : \text{the fibre of }f : X \to Y\text{ at }y
\text{ is geometrically reduced}\}
$$
is \'etale locally constructible.
\end{lemma}

\begin{proof}
Choose an affine scheme $V$ and an \'etale morphism $V  \to Y$.
The meaning of the statement is that the inverse image of $E$
in $|V|$ is constructible. By
Lemma \ref{lemma-base-change-fibres-geometrically-reduced}
we may replace $Y$ by $V$, i.e., we may assume that $Y$
is an affine scheme. Then $X$ is quasi-compact. Choose an
affine scheme $U$ and a surjective \'etale morphism $U \to X$.
For a morphism $\Spec(k) \to Y$ the morphism between fibres
$U_k \to X_k$ is surjective \'etale. Hence $U_k$ is geometrically
reduced over $k$ if and only if $X_k$ is geometrically reduced
over $k$, see Spaces over Fields, Lemma
\ref{spaces-over-fields-lemma-geometrically-reduced-etale-local}.
Thus the set $E$ for $X \to Y$ is the same as the set $E$
for $U \to Y$.
In this way we see that the lemma follows from the case
of schemes, see More on Morphisms, Lemma
\ref{more-morphisms-lemma-geometrically-reduced-constructible}.
\end{proof}

\begin{lemma}
\label{lemma-proper-flat-over-dvr-reduced-fibre}
Let $X$ be an algebraic space over a discrete valuation ring $R$
whose structure morphism $X \to \Spec(R)$ is proper and flat.
If the special fibre is reduced, then
both $X$ and the generic fibre $X_\eta$ are reduced.
\end{lemma}

\begin{proof}
Choose an \'etale morphism $U \to X$ where $U$ is an affine scheme.
Then $U$ is of finite type over $R$. Let $u \in U$ be in the special fibre.
The local ring $A = \mathcal{O}_{U, u}$ is essentially of finite
type over $R$, hence Noetherian. Let $\pi \in R$ be a uniformizer.
Since $X$ is flat over $R$, we see that $\pi \in \mathfrak m_A$
is a nonzerodivisor on $A$ and since the special fibre of $X$
is reduced, we have that $A/\pi A$ is reduced.
If $a \in A$, $a \not = 0$ then there exists an $n \geq 0$ and an element
$a' \in A$ such that $a = \pi^n a'$ and $a' \not \in \pi A$.
This follows from Krull intersection theorem
(Algebra, Lemma \ref{algebra-lemma-intersect-powers-ideal-module-zero}).
If $a$ is nilpotent, so is $a'$, because $\pi$ is a nonzerodivisor.
But $a'$ maps to a nonzero element of the reduced ring $A/\pi A$
so this is impossible. Hence $A$ is reduced. It follows that
there exists an open neighbourhood of $u$ in $U$ which is reduced
(small detail omitted; use that $U$ is Noetherian).
Thus we can find an \'etale morphism $U \to X$ with $U$ a reduced
scheme, such that every point of the special fibre of $X$ is
in the image. Since $X$ is proper over $R$ it follows that
$U \to X$ is surjective. Hence $X$ is reduced. Since the generic fibre of
$U \to \Spec(R)$ is reduced as well (on affine pieces
it is computed by taking localizations), we conclude the same thing
is true for the generic fibre.
\end{proof}

\begin{lemma}
\label{lemma-geometrically-reduced-open}
Let $S$ be a scheme. Let $f : X \to Y$ be a morphism of algebraic
spaces over $S$. If $f$ is flat, proper, and of finite presentation,
then the set
$$
E = \{y \in |Y| : \text{the fibre of }f : X \to Y\text{ at }y
\text{ is geometrically reduced}\}
$$
is open in $|Y|$.
\end{lemma}

\begin{proof}
By Lemma \ref{lemma-base-change-fibres-geometrically-reduced}
formation of $E$ commutes with base change. To check a subset
of $|Y|$ is open, we may replace $Y$ by the members of an
\'etale covering. Thus we may assume $Y$ is affine.
Then $Y$ is a cofiltered limit of affine
schemes of finite type over $\mathbf{Z}$.
Hence we can assume $X \to Y$ is the
base change of $X_0 \to Y_0$ where $Y_0$ is the spectrum of a finite
type $\mathbf{Z}$-algebra and $X_0 \to Y_0$ is flat and proper.
See Limits of Spaces, Lemma
\ref{spaces-limits-lemma-descend-finite-presentation},
\ref{spaces-limits-lemma-descend-flat}, and
\ref{spaces-limits-lemma-eventually-proper}. Since the formation of
$E$ commutes with base change (see above),
we may assume the base is Noetherian.

\medskip\noindent
Assume $Y$ is Noetherian. The set is constructible by
Lemma \ref{lemma-geometrically-reduced-constructible}.
Hence it suffices to show the set is stable under generalization
(Topology, Lemma \ref{topology-lemma-characterize-closed-Noetherian}). By
Properties, Lemma \ref{properties-lemma-locally-Noetherian-specialization-dvr}
we reduce to the case where $Y = \Spec(R)$, $R$ is a discrete
valuation ring, and the closed fibre $X_y$ is geometrically
reduced. To show: the generic fibre $X_\eta$ is geometrically reduced.

\medskip\noindent
If not then there exists a finite extension $L$ of the fraction
field of $R$ such that $X_L$ is not reduced, see
Spaces over Fields, Lemmas
\ref{spaces-over-fields-lemma-perfect-reduced} (characteristic zero) and
\ref{spaces-over-fields-lemma-geometrically-reduced-positive-characteristic}
(positive characteristic). There exists a discrete valuation ring
$R' \subset L$ with fraction field $L$ dominating $R$, see
Algebra, Lemma \ref{algebra-lemma-integral-closure-Dedekind}.
After replacing $R$ by $R'$ we reduce to
Lemma \ref{lemma-proper-flat-over-dvr-reduced-fibre}.
\end{proof}







\section{Connected components of fibres}
\label{section-connected}

\noindent
This section is the analogue of
More on Morphisms, Section \ref{more-morphisms-section-connected}.

\begin{lemma}
\label{lemma-base-change-fibres-nr-geometrically-connected-components}
Let $S$ be a scheme.
Let $f : X \to Y$ be a morphism of algebraic spaces over $S$. Let
$$
n_{X/Y} : |Y| \to \{0, 1, 2, 3, \ldots, \infty\}
$$
be the function which associates to $y \in Y$ the number of connected
components of $X_k$ where $\Spec(k) \to Y$ is in the equivalence
class of $y$ with $k$ algebraically closed.
This is well defined and if $g : Y' \to Y$ is a morphism then
$$
n_{X'/Y'} = n_{X/Y} \circ g
$$
where $X' \to Y'$ is the base change of $f$.
\end{lemma}

\begin{proof}
Suppose that $y' \in Y'$ has image $y \in Y$. Let $\Spec(k') \to Y'$
be in the equivalence class of $y'$ with $k'$ algebraically closed.
Then we can choose a commutative diagram
$$
\xymatrix{
\Spec(K) \ar[r] \ar[rd] &
\Spec(k') \ar[r] & Y' \ar[d] \\
& \Spec(k) \ar[r] & Y
}
$$
where $K$ is an algebraically closed field.
The result follows as the morphisms of schemes
$$
\xymatrix{
X'_{k'} & (X'_{k'})_K = (X_k)_K \ar[l] \ar[r] & X_k
}
$$
induce bijections between connected components, see
Spaces over Fields, Lemma
\ref{spaces-over-fields-lemma-separably-closed-field-connected-components}.
To use this to prove the function is well defined take $Y' = Y$.
\end{proof}







\section{Dimension of fibres}
\label{section-dimension}

\noindent
This section is the analogue of
More on Morphisms, Section \ref{more-morphisms-section-dimension}.

\begin{lemma}
\label{lemma-dimension-fibre}
Let $S$ be a scheme. Let $f : X \to Y$ be a finite type morphism of
algebraic spaces over $S$. Let $y \in |Y|$. The following quantities
are the same
\begin{enumerate}
\item $d = -\infty$ if $y$ is not in the image of $|f|$ and
otherwise the minimal integer $d$ such that $f$ has relative dimension $\leq d$
at every $x \in |X|$ mapping to $y$,
\item the dimension of the algebraic space $X_k = \Spec(k) \times_Y X$
for any morphism $\Spec(k) \to Y$ in the equivalence class defining $y$.
\end{enumerate}
\end{lemma}

\begin{proof}
To parse this one has to consult
Morphisms of Spaces, Definition
\ref{spaces-morphisms-definition-dimension-fibre},
Properties of Spaces,
Definition \ref{spaces-properties-definition-dimension},
Properties of Spaces,
Definition \ref{spaces-properties-definition-dimension-at-point}.
We will show that the numbers in (1) and (2) are equal for
a fixed morphism $\Spec(k) \to Y$.
Choose an \'etale morphism $V \to Y$ where $V$ is an affine
scheme and a point $v \in V$ mapping to $y$.
Since $V \times_Y \Spec(k) \to \Spec(k)$ is surjective \'etale
(by Properties of Spaces, Lemma \ref{spaces-properties-lemma-points-cartesian})
we can find a finite separable extension $k'/k$
(by Morphisms, Lemma \ref{morphisms-lemma-etale-over-field})
and a commutative diagram
$$
\xymatrix{
\Spec(k') \ar[r] \ar[d] & V \ar[d] \\
\Spec(k) \ar[r] & Y
}
$$
We may replace $X \to Y$ by $V \times_Y X \to V$ and
$X_k$ by $X_{k'} = \Spec(k') \times_V (V \times_Y X)$
because this does not change the dimensions in question by
Properties of Spaces, Lemma
\ref{spaces-properties-lemma-dimension-decent-invariant-under-etale}
and Morphisms of Spaces, Lemma
\ref{spaces-morphisms-lemma-dimension-fibre-after-base-change}.
Thus we may assume that $Y$ is an affine scheme.
In this case we may assume that $k = \kappa(y)$
because the dimension of $X_{\kappa(y)}$ and $X_k$
are the same by the aforementioned Morphisms of Spaces, Lemma
\ref{spaces-morphisms-lemma-dimension-fibre-after-base-change}
and the fact that for an algebraic space $Z$ over a field $K$ the
relative dimension of $Z$ at a point $z \in |Z|$
is the same as $\dim_z(Z)$ by definition.
Assume $Y$ is affine and $k = \kappa(y)$. Then
$X$ is quasi-compact we can choose an affine scheme $U$ and
an surjective \'etale morphism $U \to X$.
Then $\dim(X_k) = \dim(U_k) = \max \dim_u(U_k)$
is equal to the number given in (1) by definition.
\end{proof}

\begin{lemma}
\label{lemma-base-change-dimension-fibres}
Let $S$ be a scheme. Let $f : X \to Y$ be a finite type morphism of
algebraic spaces over $S$. Let
$$
n_{X/Y} : |Y| \to \{-\infty, 0, 1, 2, 3, \ldots\}
$$
be the function which associates to $y \in |Y|$ the
integer discussed in Lemma \ref{lemma-dimension-fibre}.
If $g : Y' \to Y$ is a morphism then
$$
n_{X'/Y'} = n_{X/Y} \circ |g|
$$
where $X' \to Y'$ is the base change of $f$.
\end{lemma}

\begin{proof}
This follows immediately from
Lemma \ref{lemma-dimension-fibre}.
\end{proof}

\begin{lemma}
\label{lemma-dimension-fibres-flat}
Let $S$ be a scheme.
Let $f : X \to Y$ be a flat morphism of finite presentation of
algebraic spaces over $S$. Let
$n_{X/Y}$ be the function on $Y$ giving the dimension of fibres of $f$
introduced in Lemma \ref{lemma-base-change-dimension-fibres}.
Then $n_{X/Y}$ is lower semi-continuous.
\end{lemma}

\begin{proof}
Let $V \to Y$ be a surjective \'etale morphism where $V$ is a scheme.
If we can show that the composition $n_{X/Y} \circ |g|$
is lower semi-continuous, then the lemma follows as $|g|$ is open.
Hence we may assume $Y$ is a scheme.
Working locally we may assume $V$ is an affine scheme.
Then we can choose an affine scheme $U$ and a surjective
\'etale morphism $U \to X$. Then $n_{X/Y} = n_{U/Y}$.
Hence we may assume $X$ and $Y$ are both schemes.
In this case the lemma follows from
More on Morphisms, Lemma \ref{more-morphisms-lemma-dimension-fibres-flat}.
\end{proof}

\begin{lemma}
\label{lemma-dimension-fibres-proper}
Let $S$ be a scheme.
Let $f : X \to Y$ be a proper morphism of algebraic spaces over $S$. Let
$n_{X/Y}$ be the function on $Y$ giving the dimension of fibres of $f$
introduced in Lemma \ref{lemma-base-change-dimension-fibres}.
Then $n_{X/Y}$ is upper semi-continuous.
\end{lemma}

\begin{proof}
Let $Z_d = \{x \in |X| :
\text{the fibre of }f\text{ at }x\text{ has dimension }> d\}$.
Then $Z_d$ is a closed subset of $|X|$ by
Morphisms of Spaces, Lemma
\ref{spaces-morphisms-lemma-openness-bounded-dimension-fibres}.
Since $f$ is proper $f(Z_d)$ is closed in $|Y|$.
Since $y \in f(Z_d) \Leftrightarrow n_{X/Y}(y) > d$
we see that the lemma is true.
\end{proof}

\begin{lemma}
\label{lemma-dimension-fibres-proper-flat}
Let $S$ be a scheme. Let $f : X \to Y$ be a proper, flat, finitely presented
morphism of algebraic spaces over $S$.
Let $n_{X/Y}$ be the function on $Y$ giving the dimension of fibres of $f$
introduced in Lemma \ref{lemma-base-change-dimension-fibres}.
Then $n_{X/Y}$ is locally constant.
\end{lemma}

\begin{proof}
Immediate consequence of
Lemmas \ref{lemma-dimension-fibres-flat} and
\ref{lemma-dimension-fibres-proper}.
\end{proof}





\section{Catenary algebraic spaces}
\label{section-catenary}

\noindent
This section continues the discussion started in
Decent Spaces, Section \ref{decent-spaces-section-catenary}.
The following lemma will be used in the proof of the next one.

\begin{lemma}
\label{lemma-construct-glueing}
Let $S$ be a scheme. Let $f : X \to Y$ be an integral morphism
of algebraic spaces over $S$.
Let $y \in |Y|$ be a point which can be represented by a closed immersion
$y : \Spec(k) \to Y$. Then there exists
a factorization $X \to X' \to Y$ of $f$ such that
\begin{enumerate}
\item $X' \to Y$ is integral,
\item $X \to X'$ is an isomorphism over $X' \setminus X'_y$,
\item $X'_y$ has a unique point $x'$ with $\kappa(x') = k$.
\end{enumerate}
Moreover, if $f$ is finite and $Y$ is locally Noetherian, then
$X' \to Y$ is finite.
\end{lemma}

\begin{proof}
By Morphisms of Spaces, Lemma \ref{spaces-morphisms-lemma-pushforward}
the sheaves $f_*\mathcal{O}_X$, $(X_y \to Y)_*\mathcal{O}_{X_y}$, and
$y_*\mathcal{O}_{\Spec(k)}$ are quasi-coherent sheaves of
$\mathcal{O}_Y$-algebras. Consider the maps
$$
f_*\mathcal{O}_Y \longrightarrow
(X_y \to Y)_*\mathcal{O}_{X_y} \longleftarrow
y_*\mathcal{O}_{\Spec(k)}
$$
The fibre product is a quasi-coherent sheaf of $\mathcal{O}_Y$-algebras
$\mathcal{A}'$ and we can define $X' \to Y$ as the relative spectrum
of $\mathcal{A}'$ over $Y$, see
Morphisms, Lemma \ref{morphisms-lemma-affine-equivalence-algebras}.
This construction commutes with arbitrary change of base.
In particular, it is clear that over the open subspace
$|Y| \setminus \{y\}$ the morphism $X \to X'$ is an isomorphism
and over $|Y| \setminus \{y\}$ the morphism $X' \to Y$ is integral.
It remains to prove the statements in a small neighbourhood of $y$.
Choose an affine scheme $V = \Spec(R)$ and an \'etale
morphism $\varphi : V \to Y$ such that $y$ is in the image of
$\varphi$. Then $V_y$ is a closed subscheme of $V$
\'etale over $k$, whence consists of finitely many points
each with residue field separable over $k$
(see Decent Spaces, Remark \ref{decent-spaces-remark-recall}).
After shrinking $V$
we may assume there is a unique closed point $v = \Spec(l) \to V$
mapping to $y$ with $l/k$ finite separable.
We may write $V \times_Y X = \Spec(C)$ with $R \to C$
an integral ring map. The stated compatibility with
base change gives us that $U \times_X Y' = \Spec(C')$ where
$$
C' = C \times_{C \otimes_R l} l
$$
Since $R \to l$ is surjective, also $C \to C \otimes_R l$
is surjective and we see that this is a fibre product of the
kind studied in More on Algebra, Situation
\ref{more-algebra-situation-module-over-fibre-product}
(with $A, A', B, B'$ corresponding to $C \otimes_R l, C, l, C'$).
Observe that $C'$ is an $R$-subalgebra of $C$ and hence
is integral over $R$; this proves (1).
Finally, More on Algebra, Lemma
\ref{more-algebra-lemma-points-of-fibre-product}
shows that $V \times_X Y' = \Spec(C')$
has a unique point $y''$ lying over $v$ with residue $l$
(this corresponds with the obvious surjective map $C' \to l$).
Thus $X_y \times_{\Spec(k)} \Spec(l)$ has a unique point
with residue field $l$. Since $l/k$ is finite separable,
this implies $X'_y$ has a unique point with residue field $k$, i.e.,
(3) holds.

\medskip\noindent
To prove the final statement, observe that if $Y$ is locally Noetherian,
then $R$ is a Noetherian ring and if $f$ is finite, then $R \to C$ is
finite. Then $C'$ is a finite type $R$-algebra by More on Algebra,
Lemma \ref{more-algebra-lemma-fibre-product-finite-type}.
This proves that $X' \to Y$ is finite.
\end{proof}

\begin{lemma}
\label{lemma-universally-catenary-dimension-function}
Let $S$ be a scheme. Let $B$ be an algebraic space over $S$.
Let $\delta : |B| \to \mathbf{Z}$ be a function.
Assume $B$ is decent, locally Noetherian, and
universally catenary and $\delta$ is a dimension function.
If $X$ is a decent algebraic space over $B$ whose structure morphism
$f : X \to B$ is locally of finite type we define
$\delta_X : |X| \to \mathbf{Z}$ by the rule
$$
\delta_X(x) = \delta(f(x)) + \text{transcendence degreeof }x/f(x)
$$
(Morphisms of Spaces, Definition
\ref{spaces-morphisms-definition-dimension-fibre}).
Then $\delta_X$ is a dimension function.
\end{lemma}

\begin{proof}
The problem is local on $B$. Thus we may assume $B$ is quasi-compact.
By Decent Spaces, Lemma
\ref{decent-spaces-lemma-locally-Noetherian-decent-quasi-separated}
we see $B$ is quasi-separated. By
Limits of Spaces, Proposition
\ref{spaces-limits-proposition-there-is-a-scheme-finite-over}
we can choose a finite surjective morphism $\pi : Y \to X$
where $Y$ is a scheme.
Claim: $\delta_Y$ is a dimension function.

\medskip\noindent
The claim implies the lemma. With $X \to B$ as in the lemma
set $Z = Y \times_B X$ with projections $p : Z \to Y$ and $q : Z \to X$.
Then we have
$$
\delta_Z(z) = \delta_Y(p(z)) + \text{transcendence degreeof }z/p(z)
$$
and $\delta_Z(z) = \delta_X(q(z))$. This follows from
Morphisms of Spaces, Lemma 
\ref{spaces-morphisms-lemma-dimension-fibre-at-a-point-additive}
and the fact that these transcendence degrees are zero
for finite morphisms. By Decent Spaces, Lemma
\ref{decent-spaces-lemma-scheme-with-dimension-function}
and the claim we find that $\delta_Z$ is a dimension function.
Then we find that $\delta_X$ is a dimension function by
Decent Spaces, Lemma
\ref{decent-spaces-lemma-check-dimension-function-finite-cover}.

\medskip\noindent
Proof of the claim. Consider a specialization $y \leadsto y'$,
$y \not = y'$ of points of the Noetherian scheme $Y$.
Then $\delta_Y(y) > \delta_Y(y')$ because there are
no specializations between points in fibres of $Y$
(see Decent Spaces, Lemma
\ref{decent-spaces-lemma-conditions-on-fibre-and-qf}).
Using this for a chain of specializations we find
$$
\delta_Y(y) - \delta_Y(y') \geq
\text{codim}(\overline{\{y'\}}, \overline{\{y\}})
$$
Our task is to show equality. By
Properties, Lemma \ref{properties-lemma-locally-Noetherian-closed-point}
we can choose a specialization $y' \leadsto y_0$.
It suffices to show
$\delta_Y(y) - \delta_Y(y_0) =
\text{codim}(\overline{\{y_0\}}, \overline{\{y\}})$
because this will imply the equality for both
$y \leadsto y'$ and $y' \leadsto y_0$.

\medskip\noindent
Choose a maximal chain
$y = y_c \leadsto y_{c - 1} \leadsto \ldots \leadsto y_0$
of specializations in $Y$.
Set $b = \pi(y)$ and $b_0 = \pi(y_0)$.
Choose a maximal chain
$b = b_e \leadsto b_{e - 1} \leadsto \ldots \leadsto b_0$
of specializations in $|B|$.
We have to show $e = c$.
Since $\pi$ is closed
(Morphisms of Spaces, Lemma \ref{spaces-morphisms-lemma-finite-proper})
we can find a sequence of specializations
$y = y'_e \leadsto y'_{e - 1} \leadsto \ldots \leadsto y'_0$
mapping to
$b = b_e \leadsto b_{e - 1} \leadsto \ldots \leadsto b_0$.
Observe that $y'_e \leadsto y'_{e - 1} \leadsto \ldots \leadsto y'_0$
is a maximal chain as well.
If $y_0 = y'_0$, then because $Y$ is catenary, we conclude
that $e = c$ as desired. In the next paragraph we reduce to
this case by sleight of hand and we conclude in the same manner.

\medskip\noindent
Since $\pi$ is closed we see that $b_0$ is a closed point of $|B|$.
By Decent Spaces, Lemma \ref{decent-spaces-lemma-decent-space-closed-point}
we can represent $b_0$ by a closed immersion $b_0 : \Spec(k) \to B$.
By Lemma \ref{lemma-construct-glueing}
we can find a factorization
$$
Y \to Y' \to X
$$
with $\pi' : Y' \to X$ finite and $Y \to Y'$ a morphism which
map $y_0$ and $y'_0$ to the same point and is an isomorphism
away from the inverse image of $b_0$.
(Of course $Y'$ won't be a scheme but this doesn't matter for the
argument that follows.)
Clearly the maximal chains of specializations
$y_c \leadsto y_{c - 1} \leadsto \ldots \leadsto y_0$ and
$y'_e \leadsto y'_{e - 1} \leadsto \ldots \leadsto y'_0$
map to maximal chains of specializations in $Y'$ having
the same start and end.
Since $B$ is universally catenary, we see that
$|Y'|$ is catenary and we conclude as before.
\end{proof}






\section{\'Etale localization of morphisms}
\label{section-etale-localization}

\noindent
The section is the analogue of
More on Morphisms, Section \ref{more-morphisms-section-etale-localization}.

\begin{lemma}
\label{lemma-etale-splits-off-quasi-finite-part}
Let $S$ be a scheme.
Let $f : X \to Y$ be a morphism of algebraic spaces over $S$.
Let $y \in |Y|$. Let $x_1, \ldots, x_n \in |X|$ mapping to $y$.
Assume that
\begin{enumerate}
\item $f$ is locally of finite type,
\item $f$ is separated,
\item $f$ is quasi-finite at $x_1, \ldots, x_n$, and
\item $f$ is quasi-compact or $Y$ is decent.
\end{enumerate}
Then there exists an \'etale morphism $(U, u) \to (Y, y)$
of pointed algebraic spaces and a decomposition
$$
U \times_Y X = W \amalg V
$$
into open and closed subspaces such that the morphism $V \to U$ is finite,
every point of the fibre of $|V| \to |U|$ over $u$ maps to an $x_i$,
and the fibre of $|W| \to |U|$ over $u$ contains no point mapping to an
$x_i$.
\end{lemma}

\begin{proof}
Let $(U, u) \to (Y, y)$ be an \'etale morphism of algebraic spaces
and consider the set of  $w \in |U \times_Y X|$ mapping to $u \in |U|$
and one of the $x_i \in |X|$. By
Decent Spaces, Lemma \ref{decent-spaces-lemma-qf-and-qc-finite-fibre}
(if $f$ is of finite type) or
Decent Spaces, Lemma \ref{decent-spaces-lemma-decent-finite-fibre}
(if $Y$ is decent) this set is finite.
It follows that we may replace $f$ by the base change
$U \times_Y X \to U$ and $x_1, \ldots, x_n$ by the set of these $w$.
In particular we may and do assume that $Y$ is an affine scheme,
whence $X$ is a separated algebraic space.

\medskip\noindent
Choose an affine scheme $Z$ and an \'etale morphism $Z \to X$ such that
$x_1, \ldots, x_n$ are in the image of $|Z| \to |X|$. The fibres of
$|Z| \to |X|$ are finite, see Properties of Spaces, Lemma
\ref{spaces-properties-lemma-finite-fibres-presentation}
(or the more general discussion in Decent Spaces, Section
\ref{decent-spaces-section-reasonable-decent}).
Let $\{z_1, \ldots, z_m\} \subset |Z|$ be the preimage of
$\{x_1, \ldots, x_n\}$. By More on Morphisms, Lemma
\ref{more-morphisms-lemma-etale-splits-off-quasi-finite-part-technical}
there exists an \'etale morphism $(U, u) \to (Y, y)$ such
that $U \times_Y Z = Z_1 \amalg Z_2$ with $Z_1 \to U$ finite and
$(Z_1)_y = \{z_1, \ldots, z_m\}$. We may assume that $U$ is affine
and hence $Z_1$ is affine too.

\medskip\noindent
Since $f$ is separated, the image $V$ of $Z_1 \to X$ is both open and closed
(Morphisms of Spaces, Lemma
\ref{spaces-morphisms-lemma-universally-closed-permanence}).
Set $W = X \setminus V$ to get a decomposition as in the lemma.
To finish the proof we have to show that $V \to U$ is finite.
As $Z_1 \to V$ is surjective and \'etale, $V$ is the quotient of
$Z_1$ by the \'etale equivalence relation $R = Z_1 \times_V Z_1$, see
Spaces, Lemma \ref{spaces-lemma-space-presentation}.
Since $f$ is separated, $V \to U$ is separated and $R$ is closed in
$Z_1 \times_U Z_1$. Since $Z_1 \to U$ is finite,
the projections $s, t : R \to Z_1$ are finite.
Thus $V$ is an affine scheme by
Groupoids, Proposition \ref{groupoids-proposition-finite-flat-equivalence}.
By Morphisms, Lemma \ref{morphisms-lemma-image-proper-is-proper}
we conclude that $V \to U$ is proper and by
Morphisms, Lemma \ref{morphisms-lemma-finite-proper}
we conclude that $V \to U$ is finite,
thereby finishing the proof.
\end{proof}

\begin{lemma}
\label{lemma-etale-splits-off-just-one-quasi-finite-part}
Let $S$ be a scheme.
Let $f : X \to Y$ be a morphism of algebraic spaces over $S$.
Let $x \in |X|$ with image $y \in |Y|$. Assume that
\begin{enumerate}
\item $f$ is locally of finite type,
\item $f$ is separated, and
\item $f$ is quasi-finite at $x$.
\end{enumerate}
Then there exists an \'etale morphism $(U, u) \to (Y, y)$
of pointed algebraic spaces and a decomposition
$$
U \times_Y X = W \amalg V
$$
into open and closed subspaces such that the morphism $V \to U$ is finite
and there exists a point $v \in |V|$ which maps to $x$ in $|X|$
and $u$ in $|U|$.
\end{lemma}

\begin{proof}
Pick a scheme $U$, a point $u \in U$, and an \'etale morphism
$U \to Y$ mapping $u$ to $y$. There exists a point $x' \in |U \times_Y X|$
mapping to $x$ in $|X|$ and $u$ in $|U|$ (Properties of Spaces,
Lemma \ref{spaces-properties-lemma-points-cartesian}).
To finish, apply Lemma \ref{lemma-etale-splits-off-quasi-finite-part}
to the morphism $U \times_Y X \to U$ and the point $x'$.
It applies because $U$ is a scheme and hence $u$ comes from
the monomorphism $\Spec(\kappa(u)) \to U$.
\end{proof}










\section{Zariski's Main Theorem}
\label{section-structure-quasi-finite}

\noindent
In this section we apply the results of the previous section to
prove Zariski's main theorem for morphisms of algebraic spaces.
This section is the analogue of More on Morphisms, Section
\ref{more-morphisms-section-application-etale-neighbourhoods}.

\begin{lemma}
\label{lemma-finite-type-separated}
Let $S$ be a scheme. Let $f : X \to Y$ be a morphism of algebraic
spaces over $S$ which is of finite type and separated.
Let $Y'$ be the normalization of $Y$ in $X$. Picture:
$$
\xymatrix{
X \ar[rd]_f \ar[rr]_{f'} & & Y' \ar[ld]^\nu \\
& Y &
}
$$
Then there exists an open subspace $U' \subset Y'$ such that
\begin{enumerate}
\item $(f')^{-1}(U') \to U'$ is an isomorphism, and
\item $(f')^{-1}(U') \subset X$ is the set of points at which
$f$ is quasi-finite.
\end{enumerate}
\end{lemma}

\begin{proof}
By Morphisms of Spaces, Lemma
\ref{spaces-morphisms-lemma-locally-finite-type-quasi-finite-part}
there is an open subspace $U \subset X$ corresponding to the points
of $|X|$ where $f$ is quasi-finite. We have to prove
\begin{enumerate}
\item[(a)] the image of $|U| \to |Y'|$ is $|U'|$ for some open subspace
$U'$ of $Y'$,
\item[(b)] $U = f^{-1}(U')$, and
\item[(c)] $U \to U'$ is an isomorphism.
\end{enumerate}
Since formation of $U$ commutes with arbitrary base change
(Morphisms of Spaces, Lemma
\ref{spaces-morphisms-lemma-locally-finite-type-quasi-finite-part}),
since formation of the normalization $Y'$ commutes with smooth
base change (Lemma \ref{lemma-normalization-smooth-localization}),
since \'etale morphisms are open, and
since ``being an isomorphism'' is fpqc local on the base
(Descent on Spaces, Lemma
\ref{spaces-descent-lemma-descending-property-isomorphism}),
it suffices to prove (a), (b), (c) \'etale locally on $Y$
(some details omitted). Thus
we may assume $Y$ is an affine scheme. This implies that $Y'$ is
an (affine) scheme as well.

\medskip\noindent
Let $x \in |U|$. Claim: there exists an open
neighbourhood $f'(x) \in V \subset Y'$ such that $(f')^{-1}V \to V$ is an
isomorphism. We first prove the claim implies the lemma.
Namely, then $(f')^{-1}V \cong V$ is a scheme (as an open
of $Y'$), locally of finite type over $Y$ (as an open subspace of $X$),
and for $v \in V$ the residue field extension
$\kappa(v)/\kappa(\nu(v))$ is algebraic (as
$V \subset Y'$ and $Y'$ is integral over $Y$). Hence the fibres
of $V \to Y$ are discrete (Morphisms, Lemma
\ref{morphisms-lemma-algebraic-residue-field-extension-closed-point-fibre})
and $(f')^{-1}V \to Y$ is locally quasi-finite
(Morphisms, Lemma \ref{morphisms-lemma-locally-quasi-finite-fibres}).
This implies $(f')^{-1}V \subset U$ and $V \subset U'$. Since $x$ was
arbitrary we see that (a), (b), and (c) are true.

\medskip\noindent
Let $y = f(x) \in |Y|$. Let $(T, t) \to (Y, y)$ be an \'etale morphism
of pointed schemes. Denote by a subscript ${}_T$ the base change to $T$.
Let $z \in X_T$ be a point in the fibre $X_t$ lying over $x$.
Note that $U_T \subset X_T$ is the set of points where $f_T$ is
quasi-finite, see Morphisms of Spaces, Lemma
\ref{spaces-morphisms-lemma-locally-finite-type-quasi-finite-part}.
Note that
$$
X_T \xrightarrow{f'_T} Y'_T \xrightarrow{\nu_T} T
$$
is the normalization of $T$ in $X_T$, see
Lemma \ref{lemma-normalization-smooth-localization}.
Suppose that the claim holds for $z \in U_T \subset X_T \to Y'_T \to T$, i.e.,
suppose that we can find an open neighbourhood
$f'_T(z) \in V' \subset Y'_T$ such that $(f'_T)^{-1}V' \to V'$ is an
isomorphism. The morphism $Y'_T \to Y'$ is \'etale hence the image
$V \subset Y'$ of $V'$ is open. Observe that $f'(x) \in V$ as $f'_T(z) \in V'$.
Observe that
$$
\xymatrix{
(f'_T)^{-1}V' \ar[r] \ar[d] & (f')^{-1}(V) \ar[d] \\
V' \ar[r] & V
}
$$
is a fibre square (as $Y'_T \times_{Y'} X = X_T$).
Since the left vertical arrow is an isomorphism
and $\{V' \to V\}$ is a \'etale covering, we conclude that the right vertical
arrow is an isomorphism by
Descent on Spaces, Lemma
\ref{spaces-descent-lemma-descending-property-isomorphism}.
In other words, the claim holds for $x \in U \subset X \to Y' \to Y$.

\medskip\noindent
By the result of the previous paragraph to prove the claim for
$x \in |U|$, we may replace $Y$ by an \'etale neighbourhood $T$ of
$y = f(x)$ and $x$ by any point lying over $x$ in $T \times_Y X$.
Thus we may assume there is a decomposition
$$
X = V \amalg W
$$
into open and closed subspaces where $V \to Y$ is finite and $x \in V$,
see Lemma \ref{lemma-etale-splits-off-quasi-finite-part}.
Since $X$ is a disjoint union of $V$ and $W$ over $Y$ and since
$V \to Y$ is finite we see that the
normalization of $Y$ in $X$ is the morphism
$$
X = V \amalg W \longrightarrow V \amalg W' \longrightarrow S
$$
where $W'$ is the normalization of $Y$ in $W$, see
Morphisms of Spaces, Lemmas
\ref{spaces-morphisms-lemma-normalization-in-disjoint-union},
\ref{spaces-morphisms-lemma-finite-integral}, and
\ref{spaces-morphisms-lemma-normalization-in-integral}.
The claim follows and we win.
\end{proof}

\noindent
The following lemma is a duplicate of
Morphisms of Spaces, Lemma
\ref{spaces-morphisms-lemma-quasi-finite-separated-quasi-affine}.
The reason for having two copies of the same lemma is that the
proofs are somewhat different. The proof given below
rests on Zariski's Main Theorem for nonrepresentable morphisms of
algebraic spaces as presented above, whereas the proof of
Morphisms of Spaces, Lemma
\ref{spaces-morphisms-lemma-quasi-finite-separated-quasi-affine}
rests on Morphisms of Spaces, Proposition
\ref{spaces-morphisms-proposition-locally-quasi-finite-separated-over-scheme}
to reduce to the case of morphisms of schemes.

\begin{lemma}
\label{lemma-quasi-finite-separated-quasi-affine}
Let $S$ be a scheme.
Let $f : X \to Y$ be a morphism of algebraic spaces over $S$.
Assume $f$ is quasi-finite and separated.
Let $Y'$ be the normalization of $Y$ in $X$.
Picture:
$$
\xymatrix{
X \ar[rd]_f \ar[rr]_{f'} & & Y' \ar[ld]^\nu \\
& Y &
}
$$
Then $f'$ is a quasi-compact open immersion and $\nu$ is integral.
In particular $f$ is quasi-affine.
\end{lemma}

\begin{proof}
This follows from Lemma \ref{lemma-finite-type-separated}. Namely, by
that lemma there exists an open subspace $U' \subset Y'$ such that
$(f')^{-1}(U') = X$ (!) and $X \to U'$ is an isomorphism! In other
words, $f'$ is an open immersion. Note that $f'$ is quasi-compact as
$f$ is quasi-compact and $\nu : Y' \to Y$ is separated
(Morphisms of Spaces, Lemma
\ref{spaces-morphisms-lemma-quasi-compact-permanence}).
Hence for every affine scheme $Z$ and morphism $Z \to Y$ the
fibre product $Z \times_Y X$ is a quasi-compact open subscheme
of the affine scheme $Z \times_Y Y'$. Hence $f$ is quasi-affine by
definition.
\end{proof}

\begin{lemma}[Zariski's Main Theorem]
\label{lemma-quasi-finite-separated-pass-through-finite}
Let $S$ be a scheme.
Let $f : X \to Y$ be a morphism of algebraic spaces over $S$.
Assume $f$ is quasi-finite and separated and assume that
$Y$ is quasi-compact and quasi-separated. Then there exists
a factorization
$$
\xymatrix{
X \ar[rd]_f \ar[rr]_j & & T \ar[ld]^\pi \\
& Y &
}
$$
where $j$ is a quasi-compact open immersion and $\pi$ is finite.
\end{lemma}

\begin{proof}
Let $X \to Y' \to Y$ be as in the conclusion of
Lemma \ref{lemma-quasi-finite-separated-quasi-affine}.
By
Limits of Spaces, Lemma
\ref{spaces-limits-lemma-integral-algebra-directed-colimit-finite}
we can write
$\nu_*\mathcal{O}_{Y'} = \colim_{i \in I} \mathcal{A}_i$ as a
directed colimit of finite quasi-coherent $\mathcal{O}_X$-algebras
$\mathcal{A}_i \subset \nu_*\mathcal{O}_{Y'}$. Then
$\pi_i : T_i = \underline{\Spec}_Y(\mathcal{A}_i) \to Y$
is a finite morphism for each $i$.
Note that the transition morphisms $T_{i'} \to T_i$ are affine
and that $Y' = \lim T_i$.

\medskip\noindent
By Limits of Spaces, Lemma \ref{spaces-limits-lemma-descend-opens}
there exists an $i$ and a quasi-compact open
$U_i \subset T_i$ whose inverse image in $Y'$ equals
$f'(X)$. For $i' \geq i$ let $U_{i'}$ be the inverse image
of $U_i$ in $T_{i'}$. Then $X \cong f'(X) = \lim_{i' \geq i} U_{i'}$, see
Limits of Spaces, Lemma
\ref{spaces-limits-lemma-directed-inverse-system-has-limit}.
By
Limits of Spaces, Lemma
\ref{spaces-limits-lemma-finite-type-eventually-closed} we see that
$X \to U_{i'}$ is a closed immersion for some $i' \geq i$.
(In fact $X \cong U_{i'}$ for sufficiently
large $i'$ but we don't need this.) Hence $X \to T_{i'}$ is an immersion. By
Morphisms of Spaces, Lemma
\ref{spaces-morphisms-lemma-factor-the-other-way}
we can factor this as $X \to T \to T_{i'}$ where the first arrow
is an open immersion and the second a closed immersion. Thus we win.
\end{proof}

\begin{lemma}
\label{lemma-quasi-finite-separated-pass-through-finite-addendum}
With notation and hypotheses as in
Lemma \ref{lemma-quasi-finite-separated-pass-through-finite}.
Assume moreover that $f$ is locally of finite presentation. Then we can
choose the factorization such that $T$ is finite and of
finite presentation over $Y$.
\end{lemma}

\begin{proof}
By Limits of Spaces, Lemma
\ref{spaces-limits-lemma-finite-in-finite-and-finite-presentation} we can write
$T = \lim T_i$ where all $T_i$ are finite and of finite presentation
over $Y$ and the transition morphisms $T_{i'} \to T_i$ are closed
immersions. By
Limits of Spaces, Lemma \ref{spaces-limits-lemma-descend-opens}
there exists an $i$ and an open subscheme $U_i \subset T_i$ whose inverse
image in $T$ is $X$. By
Limits of Spaces, Lemma
\ref{spaces-limits-lemma-finite-type-eventually-closed}
we see that $X \cong U_i$ for large enough $i$.
Replacing $T$ by $T_i$ finishes the proof.
\end{proof}












\section{Applications of Zariski's Main Theorem, I}
\label{section-applications-zmt}

\noindent
A first application is the characterization of finite
morphisms as proper morphisms with finite fibres.

\begin{lemma}
\label{lemma-characterize-finite}
Let $S$ be a scheme. Let $f : X \to Y$ be a morphism of algebraic spaces
over $S$. The following are equivalent:
\begin{enumerate}
\item $f$ is finite,
\item $f$ is proper and locally quasi-finite,
\item $f$ is proper and $|X_k|$ is a discrete space for every morphism
$\Spec(k) \to Y$ where $k$ is a field,
\item $f$ is universally closed, separated, locally of finite type
and $|X_k|$ is a discrete space for every morphism $\Spec(k) \to Y$
where $k$ is a field.
\end{enumerate}
\end{lemma}

\begin{proof}
We have (1) $\Rightarrow$ (2) by
Morphisms of Spaces, Lemmas \ref{spaces-morphisms-lemma-finite-proper},
\ref{spaces-morphisms-lemma-finite-quasi-finite}.
We have (2) $\Rightarrow$ (3) by
Morphisms of Spaces, Lemma
\ref{spaces-morphisms-lemma-locally-quasi-finite}.
By definition (3) implies (4).

\medskip\noindent
Assume (4). Since $f$ is universally closed it is quasi-compact
(Morphisms of Spaces, Lemma
\ref{spaces-morphisms-lemma-universally-closed-quasi-compact}).
Pick a point $y$ of $|Y|$. We represent $y$ by a
morphism $\Spec(k) \to Y$. Note that $|X_k|$ is finite discrete
as a quasi-compact discrete space. The map $|X_k| \to |X|$ surjects
onto the fibre of $|X| \to |Y|$ over $y$
(Properties of Spaces, Lemma \ref{spaces-properties-lemma-points-cartesian}).
By
Morphisms of Spaces, Lemma \ref{spaces-morphisms-lemma-quasi-finite-at-point}
we see that $X \to Y$ is quasi-finite at all the points of the fibre
of $|X| \to |Y|$ over $y$.
Choose an elementary \'etale neighbourhood $(U, u) \to (Y, y)$
and decomposition $X_U = V \amalg W$ as in
Lemma \ref{lemma-etale-splits-off-quasi-finite-part}
adapted to all the points of $|X|$ lying over $y$.
Note that $W_u = \emptyset$ because we used all the points
in the fibre of $|X| \to |Y|$ over $y$.
Since $f$ is universally closed we see that
the image of $|W|$ in $|U|$ is a closed set not containing $u$.
After shrinking $U$ we may assume that $W = \emptyset$.
In other words we see that $X_U = V$ is finite over $U$.
Since $y \in |Y|$ was arbitrary
this means there exists a family $\{U_i \to Y\}$
of \'etale morphisms whose images cover $Y$ such that
the base changes $X_{U_i} \to U_i$ are finite.
We conclude that $f$ is finite by
Morphisms of Spaces, Lemma \ref{spaces-morphisms-lemma-integral-local}.
\end{proof}

\noindent
As a consequence we have the following useful result.

\begin{lemma}
\label{lemma-proper-finite-fibre-finite-in-neighbourhood}
Let $S$ be a scheme. Let $f : X \to Y$ be a morphism of algebraic spaces
over $S$. Let $y \in |Y|$. Assume
\begin{enumerate}
\item $f$ is proper, and
\item $f$ is quasi-finite at all $x \in |X|$ lying over $y$
(Decent Spaces, Lemma \ref{decent-spaces-lemma-conditions-on-fibre-and-qf}).
\end{enumerate}
Then there exists an open neighbourhood $V \subset Y$ of $y$
such that $f|_{f^{-1}(V)} : f^{-1}(V) \to V$ is finite.
\end{lemma}

\begin{proof}
By
Morphisms of Spaces, Lemma
\ref{spaces-morphisms-lemma-locally-finite-type-quasi-finite-part} the
set of points at which $f$ is quasi-finite is an open $U \subset X$.
Let $Z = X \setminus U$. Then $y \not \in f(Z)$. Since $f$ is proper
the set $f(Z) \subset Y$ is closed. Choose any open neighbourhood
$V \subset Y$ of $y$ with $Z \cap V = \emptyset$. Then
$f^{-1}(V) \to V$ is locally quasi-finite and proper.
Hence $f^{-1}(V) \to V$ is finite by Lemma \ref{lemma-characterize-finite}.
\end{proof}

\begin{lemma}
\label{lemma-flat-proper-family-cannot-collapse-fibre}
\begin{slogan}
Collapsing a fibre of a proper family forces nearby ones to collapse too.
\end{slogan}
Let $S$ be a scheme. Let
$$
\xymatrix{
X \ar[rr]_h \ar[rd]_f & & Y \ar[ld]^g \\
& B
}
$$
be a commutative diagram of morphism of algebraic spaces over $S$.
Let $b \in B$ and let $\Spec(k) \to B$ be a morphism in the equivalence
class of $b$. Assume
\begin{enumerate}
\item $X \to B$ is a proper morphism,
\item $Y \to B$ is separated and locally of finite type,
\item one of the following is true
\begin{enumerate}
\item the image of $|X_k| \to |Y_k|$ is finite,
\item the image of $|f|^{-1}(\{b\})$ in $|Y|$ is finite
and $B$ is decent.
\end{enumerate}
\end{enumerate}
Then there is an open
subspace $B' \subset B$ containing $b$ such that $X_{B'} \to Y_{B'}$
factors through a closed subspace $Z \subset Y_{B'}$ finite over $B'$.
\end{lemma}

\begin{proof}
Let $Z \subset Y$ be the scheme theoretic image of $h$, see
Morphisms of Spaces, Section
\ref{spaces-morphisms-section-scheme-theoretic-image}.
By Morphisms of Spaces, Lemma
\ref{spaces-morphisms-lemma-scheme-theoretic-image-is-proper}
the morphism $X \to Z$ is surjective and $Z \to B$ is proper.
Thus
$$
\{x \in |X|\text{ lying over }b\} \to
\{z \in |Z|\text{ lying over }b\}
$$
and $|X_k| \to |Z_k|$ are surjective. We see that either
(3)(a) or (3)(b) imply that $Z \to B$ is quasi-finite
all points of $|Z|$ lying over $b$ by
Decent Spaces, Lemma \ref{decent-spaces-lemma-conditions-on-fibre-and-qf}.
Hence $Z \to B$ is finite in an open neighbourhood of $b$ by
Lemma \ref{lemma-proper-finite-fibre-finite-in-neighbourhood}.
\end{proof}







\section{Stein factorization}
\label{section-stein-factorization}

\noindent
Stein factorization is the statement that a proper morphism $f : X \to S$
with $f_*\mathcal{O}_X = \mathcal{O}_S$ has connected fibres.

\begin{lemma}
\label{lemma-stein-universally-closed}
Let $S$ be a scheme. Let $f : X \to Y$ be a universally closed and
quasi-separated morphism of algebraic spaces over $S$.
There exists a factorization
$$
\xymatrix{
X \ar[rr]_{f'} \ar[rd]_f & & Y' \ar[dl]^\pi \\
& Y &
}
$$
with the following properties:
\begin{enumerate}
\item the morphism $f'$ is universally closed, quasi-compact, quasi-separated,
and surjective,
\item the morphism $\pi : Y' \to Y$ is integral,
\item we have $f'_*\mathcal{O}_X = \mathcal{O}_{Y'}$,
\item we have $Y' = \underline{\Spec}_Y(f_*\mathcal{O}_X)$, and
\item $Y'$ is the normalization of $Y$ in $X$ as defined in
Morphisms of Spaces, Definition
\ref{spaces-morphisms-definition-normalization-X-in-Y}.
\end{enumerate}
Formation of the factorization $f = \pi \circ f'$ commutes with flat
base change.
\end{lemma}

\begin{proof}
By Morphisms of Spaces, Lemma
\ref{spaces-morphisms-lemma-universally-closed-quasi-compact}
the morphism $f$ is quasi-compact.
We just define $Y'$ as the normalization of $Y$ in $X$, so (5) and (2) hold
automatically. By
Morphisms of Spaces, Lemma
\ref{spaces-morphisms-lemma-normalization-in-universally-closed}
we see that (4) holds. The morphism $f'$ is universally closed by
Morphisms of Spaces, Lemma
\ref{spaces-morphisms-lemma-universally-closed-permanence}.
It is quasi-compact by
Morphisms of Spaces, Lemma
\ref{spaces-morphisms-lemma-quasi-compact-permanence}
and quasi-separated by
Morphisms of Spaces, Lemma
\ref{spaces-morphisms-lemma-compose-after-separated}.

\medskip\noindent
To show the remaining statements we may assume the base $Y$ is affine
(as taking normalization commutes with \'etale localization).
Say $Y = \Spec(R)$. Then $Y' = \Spec(A)$ with
$A = \Gamma(X, \mathcal{O}_X)$ an integral $R$-algebra.
Thus it is clear that $f'_*\mathcal{O}_X$
is $\mathcal{O}_{Y'}$ (because $f'_*\mathcal{O}_X$ is quasi-coherent,
by Morphisms of Spaces, Lemma
\ref{spaces-morphisms-lemma-pushforward},
and hence equal to $\widetilde{A}$). This proves (3).

\medskip\noindent
Let us show that $f'$ is surjective. As $f'$ is universally closed (see above)
the image of $f'$ is a closed subset
$V(I) \subset Y' = \Spec(A)$. Pick $h \in I$. Then
$h|_X = f^\sharp(h)$ is a global section of the structure sheaf of
$X$ which vanishes at every point. As $X$ is quasi-compact this means
that $h|_X$ is a nilpotent section, i.e., $h^n|X = 0$ for some $n > 0$.
But $A = \Gamma(X, \mathcal{O}_X)$, hence $h^n = 0$.
In other words $I$ is contained in the Jacobson radical of $A$ and we conclude
that $V(I) = Y'$ as desired.
\end{proof}

\begin{lemma}
\label{lemma-stein-universally-closed-residue-fields}
In Lemma \ref{lemma-stein-universally-closed} assume in addition that
$f$ is locally of finite type and $Y$ affine. Then for $y \in Y$ the fibre
$\pi^{-1}(\{y\}) = \{y_1, \ldots, y_n\}$ is finite and the field extensions
$\kappa(y_i)/\kappa(y)$ are finite.
\end{lemma}

\begin{proof}
Recall that there are no specializations among the points of $\pi^{-1}(\{y\})$,
see Algebra, Lemma \ref{algebra-lemma-integral-no-inclusion}.
As $f'$ is surjective, we find that $|X_y| \to \pi^{-1}(\{y\})$ is surjective.
Observe that $X_y$ is a quasi-separated algebraic space of finite type
over a field (quasi-compactness was shown in the proof of the
referenced lemma). Thus $|X_y|$ is a Noetherian topological space
(Morphisms of Spaces, Lemma
\ref{spaces-morphisms-lemma-finite-presentation-noetherian}).
A topological argument (omitted) now shows that $\pi^{-1}(\{y\})$ is finite.
For each $i$ we can pick a finite type point $x_i \in |X_y|$
mapping to $y_i$ (Morphisms of Spaces, Lemma
\ref{spaces-morphisms-lemma-enough-finite-type-points}).
We conclude that $\kappa(y_i)/\kappa(y)$ is finite:
$x_i$ can be represented by a morphism $\Spec(k_i) \to X_y$
of finite type (by our definition of finite type points)
and hence $\Spec(k_i) \to y = \Spec(\kappa(y))$ is of finite type
(as a composition of finite type morphisms),
hence $k_i/\kappa(y)$ is finite (Morphisms, Lemma
\ref{morphisms-lemma-point-finite-type}).
\end{proof}

\noindent
Let $f : X \to Y$ be a morphism of algebraic spaces and let
$\overline{y} : \Spec(k) \to Y$ be a geometric point. Then the
{\it fibre of $f$ over $\overline{y}$} is the algebraic space
$X_{\overline{y}} = X \times_{Y, \overline{y}} \Spec(k)$ over $k$.
If $Y$ is a scheme and $y \in Y$ is a point, then we denote
$X_y = X \times_Y \Spec(\kappa(y))$ the fibre as usual.

\begin{lemma}
\label{lemma-characterize-geometrically-connected-fibres}
Let $S$ be a scheme. Let $f : X \to Y$ be a morphism of algebraic spaces
over $S$. Let $\overline{y}$ be a geometric point of $Y$. Then
$X_{\overline{y}}$ is connected, if and only if for every \'etale
neighbourhood $(V, \overline{v}) \to (Y, \overline{y})$ where $V$
is a scheme the base change $X_V \to V$ has connected fibre $X_v$.
\end{lemma}

\begin{proof}
Since the category of \'etale neighbourhoods of $\overline{y}$ is
cofiltered and contains a cofinal collection of schemes
(Properties of Spaces, Lemma \ref{spaces-properties-lemma-cofinal-etale})
we may replace $Y$ by one of these neighbourhoods and assume that
$Y$ is a scheme. Let $y \in Y$ be the point corresponding to $\overline{y}$.
Then $X_y$ is geometrically connected over $\kappa(y)$ if and only if
$X_{\overline{y}}$ is connected and if and only if $(X_y)_{k'}$
is connected for every finite separable extension $k'$ of $\kappa(y)$.
See Spaces over Fields, Section
\ref{spaces-over-fields-section-geometrically-connected} and especially
Lemma \ref{spaces-over-fields-lemma-characterize-geometrically-disconnected}.
By More on Morphisms, Lemma
\ref{more-morphisms-lemma-realize-prescribed-residue-field-extension-etale}
there exists an affine \'etale neighbourhood $(V, v) \to (Y, y)$ such that
$\kappa(s) \subset \kappa(u)$ is identified with $\kappa(s) \subset k'$
any given finite separable extension. The lemma follows.
\end{proof}

\begin{theorem}[Stein factorization; Noetherian case]
\label{theorem-stein-factorization-Noetherian}
Let $S$ be a scheme. Let $f : X \to Y$ be a proper morphism of algebraic
spaces over $S$ with $Y$ locally Noetherian.
There exists a factorization
$$
\xymatrix{
X \ar[rr]_{f'} \ar[rd]_f & & Y' \ar[dl]^\pi \\
& Y &
}
$$
with the following properties:
\begin{enumerate}
\item the morphism $f'$ is proper with connected geometric fibres,
\item the morphism $\pi : Y' \to Y$ is finite,
\item we have $f'_*\mathcal{O}_X = \mathcal{O}_{Y'}$,
\item we have $Y' = \underline{\Spec}_Y(f_*\mathcal{O}_X)$, and
\item $Y'$ is the normalization of $Y$ in $X$, see
Morphisms, Definition \ref{morphisms-definition-normalization-X-in-Y}.
\end{enumerate}
\end{theorem}

\begin{proof}
Let $f = \pi \circ f'$ be the factorization of
Lemma \ref{lemma-stein-universally-closed}. Note that besides the
conclusions of Lemma \ref{lemma-stein-universally-closed} we
also have that $f'$ is separated
(Morphisms of Spaces, Lemma
\ref{spaces-morphisms-lemma-compose-after-separated})
and finite type
(Morphisms of Spaces, Lemma
\ref{spaces-morphisms-lemma-permanence-finite-type}).
Hence $f'$ is proper. By
Cohomology of Spaces, Lemma
\ref{spaces-cohomology-lemma-proper-pushforward-coherent}
we see that $f_*\mathcal{O}_X$ is a coherent $\mathcal{O}_Y$-module.
Hence we see that $\pi$ is finite, i.e., (2) holds.

\medskip\noindent
This proves all but the most interesting assertion, namely that
the geometric fibres of $f'$ are connected. It is clear from the
discussion above that we may replace $Y$ by $Y'$. 
Then $Y$ is locally Noetherian,
$f : X \to Y$ is proper, and $f_*\mathcal{O}_X = \mathcal{O}_Y$.
Let $\overline{y}$ be a geometric point of $Y$.
At this point we apply the theorem on formal functions,
more precisely Cohomology of Spaces, Lemma
\ref{spaces-cohomology-lemma-formal-functions-stalk}.
It tells us that
$$
\mathcal{O}^\wedge_{Y, \overline{y}} =
\lim_n H^0(X_n, \mathcal{O}_{X_n})
$$
where $X_n =
\Spec(\mathcal{O}_{Y, \overline{y}}/\mathfrak m_{\overline{y}}^n) \times_Y X$.
Note that $X_1 = X_{\overline{y}} \to X_n$ is a (finite order) thickening
and hence the underlying topological space of $X_n$ is equal to that
of $X_{\overline{y}}$. Thus, if $X_{\overline{y}} = T_1 \amalg T_2$
is a disjoint union of nonempty open and closed subspaces, then similarly
$X_n = T_{1, n} \amalg T_{2, n}$ for all $n$. And this in turn means
$H^0(X_n, \mathcal{O}_{X_n})$ contains a nontrivial idempotent $e_{1, n}$,
namely the function which is identically $1$ on $T_{1, n}$ and
identically $0$ on $T_{2, n}$. It is clear that $e_{1, n + 1}$
restricts to $e_{1, n}$ on $X_n$. Hence $e_1 = \lim e_{1, n}$
is a nontrivial idempotent of the limit. This contradicts the fact
that $\mathcal{O}^\wedge_{Y, \overline{y}}$ is a local ring. Thus the
assumption was wrong, i.e., $X_{\overline{y}}$ is connected
as desired.
\end{proof}

\begin{theorem}[Stein factorization; general case]
\label{theorem-stein-factorization-general}
Let $S$ be a scheme. Let $f : X \to Y$ be a proper morphism of algebraic
spaces over $S$. There exists a factorization
$$
\xymatrix{
X \ar[rr]_{f'} \ar[rd]_f & & Y' \ar[dl]^\pi \\
& Y &
}
$$
with the following properties:
\begin{enumerate}
\item the morphism $f'$ is proper with connected geometric fibres,
\item the morphism $\pi : Y' \to Y$ is integral,
\item we have $f'_*\mathcal{O}_X = \mathcal{O}_{Y'}$,
\item we have $Y' = \underline{\Spec}_Y(f_*\mathcal{O}_X)$, and
\item $Y'$ is the normalization of $Y$ in $X$ (Morphisms of Spaces, Definition
\ref{spaces-morphisms-definition-normalization-X-in-Y}).
\end{enumerate}
\end{theorem}

\begin{proof}
We may apply Lemma \ref{lemma-stein-universally-closed} to get the
morphism $f' : X \to Y'$.
Note that besides the
conclusions of Lemma \ref{lemma-stein-universally-closed} we
also have that $f'$ is separated
(Morphisms of Spaces, Lemma
\ref{spaces-morphisms-lemma-compose-after-separated})
and finite type
(Morphisms of Spaces, Lemma
\ref{spaces-morphisms-lemma-permanence-finite-type}).
Hence $f'$ is proper. At this point we have proved all of the
statements except for the statement
that $f'$ has connected geometric fibres.

\medskip\noindent
It is clear from the discussion that we may replace $Y$ by $Y'$. 
Then $f : X \to Y$ is proper and $f_*\mathcal{O}_X = \mathcal{O}_Y$.
Note that these conditions are preserved under flat base change
(Morphisms of Spaces, Lemma \ref{spaces-morphisms-lemma-base-change-proper}
and
Cohomology of Spaces, Lemma
\ref{spaces-cohomology-lemma-flat-base-change-cohomology}).
Let $\overline{y}$ be a geometric point of $Y$. By
Lemma \ref{lemma-characterize-geometrically-connected-fibres}
and the remark just made we reduce to the case where $Y$ is a
scheme, $y \in Y$ is a point, $f : X \to Y$ is a proper algebraic
space over $Y$ with $f_*\mathcal{O}_X = \mathcal{O}_Y$,
and we have to show the fibre $X_y$ is connected.
Replacing $Y$ by an affine neighbourhood of $y$ we may
assume that $Y = \Spec(R)$ is affine. Then $f_*\mathcal{O}_X = \mathcal{O}_Y$
signifies that the ring map
$R \to \Gamma(X, \mathcal{O}_X)$ is bijective.

\medskip\noindent
By Limits of Spaces, Lemma
\ref{spaces-limits-lemma-proper-limit-of-proper-finite-presentation-noetherian}
we can write $(X \to Y) = \lim (X_i \to Y_i)$ with $X_i \to Y_i$
proper and of finite presentation and $Y_i$ Noetherian. For $i$ large
enough $Y_i$ is affine (Limits of Spaces, Lemma
\ref{spaces-limits-lemma-limit-is-affine}).
Say $Y_i = \Spec(R_i)$. Let $R'_i = \Gamma(X_i, \mathcal{O}_{X_i})$.
Observe that we have ring maps $R_i \to R_i' \to R$. Namely, we have
the first because $X_i$ is an algebraic space over $R_i$ and the second because
we have $X \to X_i$ and $R = \Gamma(X, \mathcal{O}_X)$. Note that
$R = \colim R'_i$ by Limits of Spaces, Lemma
\ref{spaces-limits-lemma-descend-section}.
Then 
$$
\xymatrix{
X \ar[d]  \ar[r] & X_i \ar[d] \\
Y \ar[r] & Y'_i \ar[r] & Y_i
}
$$
is commutative with $Y'_i = \Spec(R'_i)$.
Let $y'_i \in Y'_i$ be the image of $y$.
We have $X_y = \lim X_{i, y'_i}$ because $X = \lim X_i$,
$Y = \lim Y'_i$, and $\kappa(y) = \colim \kappa(y'_i)$.
Now let $X_y = U \amalg V$ with $U$ and $V$ open and closed.
Then $U, V$ are the inverse images of opens $U_i, V_i$ in $X_{i, y'_i}$
(Limits of Spaces, Lemma \ref{spaces-limits-lemma-descend-opens}).
By Theorem \ref{theorem-stein-factorization-Noetherian} the fibres
of $X_i \to Y'_i$ are connected, hence either $U$ or $V$ is empty.
This finishes the proof.
\end{proof}

\noindent
Here is an application.

\begin{lemma}
\label{lemma-geometrically-connected-fibres-towards-normal}
Let $S$ be a scheme. Let $f : X \to Y$ be a morphism of algebraic
spaces over $S$. Assume
\begin{enumerate}
\item $f$ is proper,
\item $Y$ is integral (Spaces over Fields, Definition
\ref{spaces-over-fields-definition-integral-algebraic-space})
with generic point $\xi$,
\item $Y$ is normal,
\item $X$ is reduced,
\item every generic point of an irreducible component of $|X|$ maps to $\xi$,
\item we have $H^0(X_\xi, \mathcal{O}) = \kappa(\xi)$.
\end{enumerate}
Then $f_*\mathcal{O}_X = \mathcal{O}_Y$ and $f$
has geometrically connected fibres.
\end{lemma}

\begin{proof}
Apply Theorem \ref{theorem-stein-factorization-general} to get a
factorization $X \to Y' \to Y$. It is enough to show that $Y' = Y$.
It suffices to show that $Y' \times_Y V \to V$ is an isomorphism,
where $V \to Y$ is an \'etale morphism and $V$ an affine integral scheme,
see Spaces over Fields, Lemma
\ref{spaces-over-fields-lemma-normal-integral-cover-by-affines}.
The formation of $Y'$ commutes with \'etale base change, see
Morphisms of Spaces, Lemma
\ref{spaces-morphisms-lemma-properties-normalization}.
The generic points of $X \times_Y V$ lie over the generic points of $X$
(Decent Spaces, Lemma \ref{decent-spaces-lemma-decent-generic-points})
hence map to the generic point of $V$ by assumption (5). Moreover, condition
(6) is preserved under the base change by $V \to Y$, for example by
flat base change (Cohomology of Spaces, Lemma
\ref{spaces-cohomology-lemma-flat-base-change-cohomology}).
Thus it suffices to prove the lemma in case $Y$ is a normal
integral affine scheme.

\medskip\noindent
Assume $Y$ is a normal integral affine scheme. We will show $Y' \to Y$
is an isomorphism by an application of Morphisms, Lemma
\ref{morphisms-lemma-finite-birational-over-normal}.
Namely, $Y'$ is reduced because $X$ is reduced (Morphisms of Spaces, Lemma
\ref{spaces-morphisms-lemma-normalization-in-reduced}).
The morphism $Y' \to Y$ is integral by the theorem cited above.
Since $Y$ is decent and $X \to Y$ is separated, we see that
$X$ is decent too; to see this use
Decent Spaces, Lemmas
\ref{decent-spaces-lemma-properties-trivial-implications} and
\ref{decent-spaces-lemma-property-over-property}. By assumption (5),
Morphisms of Spaces, Lemma \ref{spaces-morphisms-lemma-normalization-generic},
and Decent Spaces, Lemma \ref{decent-spaces-lemma-decent-generic-points}
we see that every generic point of an irreducible component of $|Y'|$
maps to $\xi$. On the other hand, since $Y'$ is the relative
spectrum of $f_*\mathcal{O}_X$ we see that the scheme theoretic fibre
$Y'_\xi$ is the spectrum of $H^0(X_\xi, \mathcal{O})$ which is
equal to $\kappa(\xi)$ by assumption. Hence $Y'$ is an integral
scheme with function field equal to the function field of $Y$.
This finishes the proof.
\end{proof}

\noindent
Here is another application.

\begin{lemma}
\label{lemma-proper-flat-nr-geom-conn-comps-lower-semicontinuous}
Let $S$ be a scheme.
Let $X \to Y$ be a morphism of algebraic spaces over $S$.
If $f$ is proper, flat, and of finite presentation, then the function
$n_{X/Y} : |Y| \to \mathbf{Z}$ counting the number of geometric
connected components of fibres of $f$
(Lemma \ref{lemma-base-change-fibres-nr-geometrically-connected-components})
is lower semi-continuous.
\end{lemma}

\begin{proof}
The question is \'etale local on $Y$, hence we may and do assume $Y$
is an affine scheme. Let $y \in Y$. Set $n = n_{X/S}(y)$.
Note that $n < \infty$ as the geometric fibre of $X \to Y$ at $y$
is a proper algebraic space over a field, hence Noetherian, hence
has a finite number of connected components.
We have to find an open neighbourhood $V$ of $y$ such that $n_{X/S}|_V \geq n$.
Let $X \to Y' \to Y$ be the Stein factorization as in
Theorem \ref{theorem-stein-factorization-general}.
By Lemma \ref{lemma-stein-universally-closed-residue-fields}
there are finitely many points
$y'_1, \ldots, y'_m \in Y'$ lying over $y$
and the extensions $\kappa(y'_i)/\kappa(y)$ are finite.
More on Morphisms, Lemma \ref{more-morphisms-lemma-etale-makes-integral-split}
tells us that after replacing $Y$ by an \'etale neighbourhood
of $y$ we may assume $Y' = V_1 \amalg \ldots \amalg V_m$ as a scheme
with $y'_i \in V_i$ and $\kappa(y'_i)/\kappa(y)$ purely inseparable.
Then the algebraic spaces $X_{y_i'}$ are geometrically connected over
$\kappa(y)$, hence $m = n$. The algebraic spaces
$X_i = (f')^{-1}(V_i)$, $i = 1, \ldots, n$
are flat and of finite presentation over $Y$. Hence the image of $X_i \to Y$
is open (Morphisms of Spaces, Lemma \ref{spaces-morphisms-lemma-fppf-open}).
Thus in a neighbourhood of $y$ we see that $n_{X/Y}$ is
at least $n$.
\end{proof}

\begin{lemma}
\label{lemma-proper-flat-geom-red}
Let $S$ be a scheme.
Let $f : X \to Y$ be a morphism of algebraic spaces over $S$. Assume
\begin{enumerate}
\item $f$ is proper, flat, and of finite presentation, and
\item the geometric fibres of $f$ are reduced.
\end{enumerate}
Then the function $n_{X/S} : |Y| \to \mathbf{Z}$
counting the numbers of geometric connected components
of fibres of $f$
(Lemma \ref{lemma-base-change-fibres-nr-geometrically-connected-components})
is locally constant.
\end{lemma}

\begin{proof}
By Lemma \ref{lemma-proper-flat-nr-geom-conn-comps-lower-semicontinuous}
the function $n_{X/Y}$ is lower semincontinuous.
Thus it suffices to show it is upper semi-continuous.
To do this we may work \'etale locally on $Y$, hence we
may assume $Y$ is an affine scheme.
For $y \in Y$ consider the $\kappa(y)$-algebra
$$
A = H^0(X_y, \mathcal{O}_{X_y})
$$
By Spaces over Fields, Lemma
\ref{spaces-over-fields-lemma-proper-geometrically-reduced-global-sections}
and the fact that $X_y$ is geometrically reduced
$A$ is finite product of finite separable extensions of $\kappa(y)$.
Hence $A \otimes_{\kappa(y)} \kappa(\overline{y})$ is a product
of $\beta_0(y) = \dim_{\kappa(y)} A$ copies of $\kappa(\overline{y})$.
Thus $X_{\overline{y}}$ has $\beta_0(y)$ connected components.
In other words, we have $n_{X/S} = \beta_0$ as functions on $Y$.
Thus $n_{X/Y}$ is upper semi-continuous by
Derived Categories of Spaces, Lemma
\ref{spaces-perfect-lemma-jump-loci-geometric}.
This finishes the proof.
\end{proof}

\begin{lemma}
\label{lemma-stein-factorization-etale}
Let $S$ be a scheme.
Let $f : X \to Y$ be a proper morphism of algebraic spaces over $S$.
Let $X \to Y' \to Y$ be the Stein factorization of $f$
(Theorem \ref{theorem-stein-factorization-general}).
If $f$ is of finite presentation, flat, with geometrically
reduced fibres (Definition \ref{definition-geometrically-reduced-fibre}),
then $Y' \to Y$ is finite \'etale.
\end{lemma}

\begin{proof}
Formation of the Stein factorization commutes with flat base change,
see Lemma \ref{lemma-stein-universally-closed}.
Thus we may work \'etale locally on $Y$ and we may assume $Y$
is an affine scheme. Then $Y'$ is an affine scheme and $Y' \to Y$
is integral.

\medskip\noindent
Let $y \in Y$. Set $n$ be the number of connected components of
the geometric fibre $X_{\overline{y}}$. Note that $n < \infty$
as the geometric fibre of $X \to Y$ at $y$ is a proper
algebraic space over a field, hence Noetherian,
hence has a finite number of connected components.
By Lemma \ref{lemma-stein-universally-closed-residue-fields}
there are finitely many points $y'_1, \ldots, y'_m \in Y'$ lying over $y$
and for each $i$ we can pick a finite type point $x_i \in |X_y|$
mapping to $y'_i$ the extension $\kappa(y'_i)/\kappa(y)$ is finite.
Thus More on Morphisms,
Lemma \ref{more-morphisms-lemma-etale-makes-integral-split}
tells us that after replacing $Y$ by an \'etale neighbourhood
of $y$ we may assume $Y' = V_1 \amalg \ldots \amalg V_m$ as a scheme
with $y'_i \in V_i$ and $\kappa(y'_i)/\kappa(y)$ purely inseparable.
In this case the algebraic spaces $X_{y_i'}$
are geometrically connected over $\kappa(y)$, hence $m = n$.
The algebraic spaces $X_i = (f')^{-1}(V_i)$, $i = 1, \ldots, n$
are proper, flat, of finite presentation, with geometrically
reduced fibres over $Y$. It suffices to prove the lemma
for each of the morphisms $X_i \to Y$. This reduces us to the case where
$X_{\overline{y}}$ is connected.

\medskip\noindent
Assume that $X_{\overline{y}}$ is connected. By
Lemma \ref{lemma-proper-flat-geom-red}
we see that $X \to Y$ has geometrically connected
fibres in a neighbourhood of $y$. Thus
we may assume the fibres of $X \to Y$ are geometrically connected.
Then $f_*\mathcal{O}_X = \mathcal{O}_Y$ by
Derived Categories of Spaces, Lemma
\ref{spaces-perfect-lemma-proper-flat-geom-red-connected}
which finishes the proof.
\end{proof}

\noindent
The proof of the following lemma uses Stein factorization for schemes
which is why it ended up in this section.

\begin{lemma}
\label{lemma-split-off-proper-part-henselian}
Let $(A, I)$ be a henselian pair. Let $X$ be an algebraic space
separated and of finite type over $A$. Set
$X_0 = X \times_{\Spec(A)} \Spec(A/I)$.
Let $Y \subset X_0$ be an open and closed subspace such that
$Y \to \Spec(A/I)$ is proper. Then there exists an open and closed
subspace $W \subset X$ which is proper over $A$ with
$W \times_{\Spec(A)} \Spec(A/I) = Y$.
\end{lemma}

\begin{proof}
We will denote $T \mapsto T_0$ the base change by
$\Spec(A/I) \to \Spec(A)$.
By a weak version of Chow's lemma (in the form of
Cohomology of Spaces, Lemma \ref{spaces-cohomology-lemma-weak-chow})
there exists a surjective proper morphism $\varphi : X' \to X$ such
that $X'$ admits an immersion into $\mathbf{P}^n_A$.
Set $Y' = \varphi^{-1}(Y)$. This is an open and closed subscheme
of $X'_0$. The lemma holds for $(X', Y')$ by
More on Morphisms, Lemma
\ref{more-morphisms-lemma-split-off-proper-part-henselian}.
Let $W' \subset X'$ be the open and closed subscheme proper
over $A$ such that $Y' = W'_0$.
By Morphisms of Spaces, Lemma
\ref{spaces-morphisms-lemma-universally-closed-permanence}
$Q_1 = \varphi(|W'|) \subset |X|$ and
$Q_2 = \varphi(|X' \setminus W'|) \subset |X|$
are closed subsets and by
Morphisms of Spaces, Lemma \ref{spaces-morphisms-lemma-image-proper-is-proper}
any closed subspace structure on $Q_1$ is proper over $A$.
The image of $Q_1 \cap Q_2$ in $\Spec(A)$ is closed.
Since $(A, I)$ is henselian, if $Q_1 \cap Q_2$ is nonempty, then we
find that $Q_1 \cap Q_2$ has a point lying over $\Spec(A/I)$.
This is impossible as $W'_0 = Y' = \varphi^{-1}(Y)$.
We conclude that $Q_1$ is open and closed in $|X|$.
Let $W \subset X$ be the corresponding open and closed
subspace. Then $W$ is proper over $A$ with $W_0 = Y$.
\end{proof}










\section{Extending properties from an open}
\label{section-extending-properties}

\noindent
In this section we collect a number of results of the form: If $f : X \to Y$
is a flat morphism of algebraic spaces and $f$ satisfies some property over
a dense
open of $Y$, then $f$ satisfies the same property over all of $Y$.

\begin{lemma}
\label{lemma-flat-finite-type-finitely-presented-over-dense-open}
Let $S$ be a scheme.
Let $f : X \to Y$ be a morphism of algebraic spaces over $S$.
Let $\mathcal{F}$ be a quasi-coherent $\mathcal{O}_X$-module.
Let $V \subset Y$ be an open subspace. Assume
\begin{enumerate}
\item $f$ is locally of finite presentation,
\item $\mathcal{F}$ is of finite type and flat over $Y$,
\item $V \to Y$ is quasi-compact and scheme theoretically dense,
\item $\mathcal{F}|_{f^{-1}V}$ is of finite presentation.
\end{enumerate}
Then $\mathcal{F}$ is of finite presentation.
\end{lemma}

\begin{proof}
It suffices to prove the pullback of $\mathcal{F}$ to a scheme surjective
and \'etale over $X$ is of finite presentation. Hence we may assume $X$
is a scheme. Similarly, we can replace $Y$ by a scheme surjective and
\'etale and over $Y$ (the inverse image of $V$ in this scheme is scheme
theoretically dense, see
Morphisms of Spaces, Section
\ref{spaces-morphisms-section-scheme-theoretic-closure}).
Thus we reduce to the case of schemes which is
More on Flatness, Lemma
\ref{flat-lemma-flat-finite-type-finitely-presented-over-dense-open}.
\end{proof}

\begin{lemma}
\label{lemma-flat-finite-type-finitely-presented-over-dense-open-X}
Let $S$ be a scheme.
Let $f : X \to Y$ be a morphism of algebraic spaces over $S$.
Let $V \subset Y$ be an open subspace.
Assume
\begin{enumerate}
\item $f$ is locally of finite type and flat,
\item $V \to Y$ is quasi-compact and scheme theoretically dense,
\item $f|_{f^{-1}V} : f^{-1}V \to V$ is locally of finite presentation.
\end{enumerate}
Then $f$ is of locally of finite presentation.
\end{lemma}

\begin{proof}
The proof is identical to the proof of
Lemma \ref{lemma-flat-finite-type-finitely-presented-over-dense-open}
except one uses
More on Flatness, Lemma
\ref{flat-lemma-flat-finite-type-finitely-presented-over-dense-open-X}.
\end{proof}

\begin{lemma}
\label{lemma-flat-fp-dimension-over-dense-open}
Let $S$ be a scheme. Let $f : X \to Y$ be a morphism of algebraic spaces
over $S$ which is flat and locally of finite type. Let $V \subset Y$ be an
open subspace such that $|V| \subset |Y|$ is dense and such that $X_V \to V$
has relative dimension $\leq d$. If also either
\begin{enumerate}
\item $f$ is locally of finite presentation, or
\item $V \to Y$ is quasi-compact,
\end{enumerate}
then $f : X \to Y$ has relative dimension $\leq d$.
\end{lemma}

\begin{proof}
We may replace $Y$ by its reduction, hence we may assume $Y$ is reduced.
Then $V$ is scheme theoretically dense in $Y$, see
Morphisms of Spaces, Lemma
\ref{spaces-morphisms-lemma-quasi-compact-immersion}.
By definition the property of having relative dimension $\leq d$ can
be checked on an \'etale covering, see
Morphisms of Spaces, Sections \ref{spaces-morphisms-section-relative-dimension}.
Thus it suffices to prove $f$ has relative dimension $\leq d$
after replacing $X$ by a scheme surjective and \'etale over $X$.
Similarly, we can replace $Y$ by a scheme surjective and
\'etale and over $Y$. The inverse image of $V$ in this scheme is scheme
theoretically dense, see
Morphisms of Spaces, Section
\ref{spaces-morphisms-section-scheme-theoretic-closure}.
Since a scheme theoretically dense open of a scheme is in particular
dense, we reduce to the case of schemes which is
More on Flatness, Lemma
\ref{flat-lemma-flat-finite-presentation-dimension-over-dense-open}.
\end{proof}

\begin{lemma}
\label{lemma-proper-flat-finite-over-dense-open}
Let $S$ be a scheme. Let $f : X \to Y$ be a morphism of algebraic spaces
over $S$ which is flat and proper. Let $V \to Y$ be an open subspace
with $|V| \subset |Y|$ dense such that $X_V \to V$ is finite. If also
either $f$ is locally of finite presentation or $V \to Y$ is quasi-compact,
then $f$ is finite.
\end{lemma}

\begin{proof}
By Lemma \ref{lemma-flat-fp-dimension-over-dense-open}
the fibres of $f$ have dimension zero.
By Morphisms of Spaces, Lemma
\ref{spaces-morphisms-lemma-locally-quasi-finite-rel-dimension-0}
this implies that $f$ is locally quasi-finite.
By Morphisms of Spaces, Lemma
\ref{spaces-morphisms-lemma-locally-quasi-finite-separated-representable}
this implies that $f$ is representable.
We can check whether $f$ is finite \'etale locally on $Y$,
hence we may assume $Y$ is a scheme. Since  $f$ is representable,
we reduce to the case of schemes which is
More on Flatness, Lemma \ref{flat-lemma-proper-flat-finite-over-dense-open}.
\end{proof}

\begin{lemma}
\label{lemma-zariski}
Let $S$ be a scheme.
Let $f : X \to Y$ be a morphism of algebraic spaces over $S$.
Let $V \subset Y$ be an open subspace. If
\begin{enumerate}
\item $f$ is separated, locally of finite type, and flat,
\item $f^{-1}(V) \to V$ is an isomorphism, and
\item $V \to Y$ is quasi-compact and scheme theoretically dense,
\end{enumerate}
then $f$ is an open immersion.
\end{lemma}

\begin{proof}
Applying
Lemma \ref{lemma-flat-finite-type-finitely-presented-over-dense-open-X}
we see that $f$ is locally of finite presentation. Applying
Lemma \ref{lemma-flat-fp-dimension-over-dense-open}
we see that $f$ has relative dimension $\leq 0$.
By Morphisms of Spaces, Lemma
\ref{spaces-morphisms-lemma-locally-quasi-finite-rel-dimension-0}
this implies that $f$ is locally quasi-finite.
By Morphisms of Spaces, Lemma
\ref{spaces-morphisms-lemma-locally-quasi-finite-separated-representable}
this implies that $f$ is representable.
By Descent on Spaces, Lemma
\ref{spaces-descent-lemma-descending-property-open-immersion}
we can check whether $f$ is an open immersion \'etale locally on $Y$.
Hence we may assume that $Y$ is a scheme. Since $f$ is representable,
we reduce to the case of schemes which is
More on Flatness, Lemma \ref{flat-lemma-zariski}.
\end{proof}







\section{Blowing up and flatness}
\label{section-blowup-flat}

\noindent
Instead of redoing the work in
More on Flatness, Section \ref{flat-section-blowup-flat}
we prove an analogue of More on Flatness, Lemma \ref{flat-lemma-push-ideal}
which tells us that the problem of finding a suitable blowup
is often \'etale local on the base.

\begin{lemma}
\label{lemma-push-ideal}
Let $S$ be a scheme. Let $X$ be a quasi-compact and quasi-separated
algebraic space over $S$. Let $\varphi : W \to X$ be a quasi-compact
separated \'etale morphism. Let $U \subset X$ be a quasi-compact open
subspace. Let $\mathcal{I} \subset \mathcal{O}_W$ be a finite type
quasi-coherent sheaf of ideals such that
$V(\mathcal{I}) \cap \varphi^{-1}(U) = \emptyset$.
Then there exists a finite type quasi-coherent sheaf of ideals
$\mathcal{J} \subset \mathcal{O}_X$ such that
\begin{enumerate}
\item $V(\mathcal{J}) \cap U = \emptyset$, and
\item $\varphi^{-1}(\mathcal{J})\mathcal{O}_W = \mathcal{I} \mathcal{I}'$
for some finite type quasi-coherent ideal
$\mathcal{I}' \subset \mathcal{O}_W$.
\end{enumerate}
\end{lemma}

\begin{proof}
Choose a factorization $W \to Y \to X$ where $j : W \to Y$ is a
quasi-compact open immersion and $\pi : Y \to X$ is a finite
morphism of finite presentation
(Lemma \ref{lemma-quasi-finite-separated-pass-through-finite-addendum}).
Let $V = j(W) \cup \pi^{-1}(U) \subset Y$. Note that
$\mathcal{I}$ on $W \cong j(W)$ and $\mathcal{O}_{\pi^{-1}(U)}$
glue to a finite type quasi-coherent sheaf of ideals
$\mathcal{I}_1 \subset \mathcal{O}_V$. By
Limits of Spaces, Lemma \ref{spaces-limits-lemma-extend}
there exists a finite type quasi-coherent sheaf of ideals
$\mathcal{I}_2 \subset \mathcal{O}_Y$ such that
$\mathcal{I}_2|_V = \mathcal{I}_1$. In other words,
$\mathcal{I}_2 \subset \mathcal{O}_Y$
is a finite type quasi-coherent sheaf of ideals such that
$V(\mathcal{I}_2)$ is disjoint from $\pi^{-1}(U)$ and
$j^{-1}\mathcal{I}_2 = \mathcal{I}$. Denote $i : Z \to Y$
the corresponding closed immersion which is of finite presentation
(Morphisms of Spaces, Lemma
\ref{spaces-morphisms-lemma-closed-immersion-finite-presentation}).
In particular the composition $\tau = \pi \circ i : Z \to X$ is finite
and of finite presentation
(Morphisms of Spaces, Lemmas
\ref{spaces-morphisms-lemma-composition-finite-presentation} and
\ref{spaces-morphisms-lemma-composition-integral}).

\medskip\noindent
Let $\mathcal{F} = \tau_*\mathcal{O}_Z$ which we think of as
a quasi-coherent $\mathcal{O}_X$-module. By
Descent on Spaces, Lemma
\ref{spaces-descent-lemma-finite-finitely-presented-module}
we see that $\mathcal{F}$ is a finitely presented $\mathcal{O}_X$-module.
Let $\mathcal{J} = \text{Fit}_0(\mathcal{F})$. (Insert reference to
fitting modules on ringed topoi here.) This is a finite type quasi-coherent
sheaf of ideals on $X$ (as $\mathcal{F}$ is of finite presentation, see
More on Algebra, Lemma \ref{more-algebra-lemma-fitting-ideal-basics}).
Part (1) of the lemma holds because $|\tau|(|Z|) \cap |U| = \emptyset$
by our choice of $\mathcal{I}_2$ and because the $0$th Fitting ideal
of the trivial module equals the structure sheaf. To prove (2) note that
$\varphi^{-1}(\mathcal{J})\mathcal{O}_W = \text{Fit}_0(\varphi^*\mathcal{F})$
because taking Fitting ideals commutes with base change.
On the other hand, as $\varphi : W \to X$ is separated and \'etale
we see that $(1, j) : W \to W \times_X Y$ is an open and closed immersion.
Hence $W \times_Y Z = V(\mathcal{I}) \amalg Z'$ for some finite and
finitely presented morphism of algebraic spaces $\tau' : Z' \to W$.
Thus we see that
\begin{align*}
\text{Fit}_0(\varphi^*\mathcal{F}) & =
\text{Fit}_0((W \times_Y Z \to W)_*\mathcal{O}_{W \times_Y Z}) \\
& =
\text{Fit}_0(\mathcal{O}_W/\mathcal{I}) \cdot
\text{Fit}_0(\tau'_*\mathcal{O}_{Z'}) \\
& =
\mathcal{I} \cdot \text{Fit}_0(\tau'_*\mathcal{O}_{Z'})
\end{align*}
the second equality by
More on Algebra, Lemma \ref{more-algebra-lemma-fitting-ideal-basics}
translated in sheaves on ringed topoi.
Setting $\mathcal{I}' = \text{Fit}_0(\tau'_*\mathcal{O}_{Z'})$
finishes the proof of the lemma.
\end{proof}

\begin{theorem}
\label{theorem-flatten-module}
Let $S$ be a scheme. Let $B$ be a quasi-compact and quasi-separated
algebraic space over $S$. Let $X$ be an algebraic space over $B$.
Let $\mathcal{F}$ be a quasi-coherent module on $X$.
Let $U \subset B$ be a quasi-compact open subspace. Assume
\begin{enumerate}
\item $X$ is quasi-compact,
\item $X$ is locally of finite presentation over $B$,
\item $\mathcal{F}$ is a module of finite type,
\item $\mathcal{F}_U$ is of finite presentation, and
\item $\mathcal{F}_U$ is flat over $U$.
\end{enumerate}
Then there exists a $U$-admissible blowup $B' \to B$ such that the
strict transform $\mathcal{F}'$ of $\mathcal{F}$ is an
$\mathcal{O}_{X \times_B B'}$-module of finite presentation and
flat over $B'$.
\end{theorem}

\begin{proof}
Choose an affine scheme $V$ and a surjective \'etale morphism $V \to X$.
Because strict transform commutes with \'etale localization
(Divisors on Spaces, Lemma \ref{spaces-divisors-lemma-strict-transform-local})
it suffices to prove the result with $X$ replaced by $V$. Hence we
may assume that $X \to B$ is representable (in addition to the
hypotheses of the lemma).

\medskip\noindent
Assume that $X \to B$ is representable. Choose an affine scheme $W$ and a
surjective \'etale morphism $\varphi : W \to B$. Note that $X \times_B W$
is a scheme. By the case of schemes
(More on Flatness, Theorem \ref{flat-theorem-flatten-module})
we can find a finite type quasi-coherent sheaf of ideals
$\mathcal{I} \subset \mathcal{O}_W$ such that
(a) $|V(\mathcal{I})| \cap |\varphi^{-1}(U)| = \emptyset$ and (b)
the strict transform of $\mathcal{F}|_{X \times_B W}$ with respect
to the blowing up $W' \to W$ in $\mathcal{I}$ becomes flat over $W'$
and is a module of finite presentation. Choose a finite type sheaf of ideals
$\mathcal{J} \subset \mathcal{O}_B$ as in Lemma \ref{lemma-push-ideal}.
Let $B' \to B$ be the blowing up of $\mathcal{J}$. We claim that this
blowup works. Namely, it is clear that $B' \to B$ is $U$-admissible
by our choice of ideal $\mathcal{J}$. Moreover, the base change
$B' \times_B W \to W$ is the blowup of $W$ in
$\varphi^{-1}\mathcal{J} = \mathcal{I}\mathcal{I}'$
(compatibility of blowup with flat base change, see
Divisors on Spaces, Lemma
\ref{spaces-divisors-lemma-flat-base-change-blowing-up}).
Hence there is a factorization
$$
W \times_B B' \to W' \to W
$$
where the first morphism is a blowup as well, see
Divisors on Spaces, Lemma \ref{spaces-divisors-lemma-blowing-up-two-ideals}).
The restriction of $\mathcal{F}'$ (which lives on $B' \times_B X$)
to $W \times_B B' \times_B X$ is the strict transform of
$\mathcal{F}|_{X \times_B W}$
(Divisors on Spaces, Lemma \ref{spaces-divisors-lemma-strict-transform-local})
and hence is the twice repeated strict transform of
$\mathcal{F}|_{X \times_B W}$ by the two blowups displayed above
(Divisors on Spaces, Lemma
\ref{spaces-divisors-lemma-strict-transform-composition-blowups}).
After the first blowup our sheaf is already flat over
the base and of finite presentation (by construction). Whence this holds 
after the second strict transform as well (since this is a
pullback by Divisors on Spaces, Lemma
\ref{spaces-divisors-lemma-strict-transform-flat}).
Thus we see that the restriction of $\mathcal{F}'$
to an \'etale cover of $B' \times_B X$ has the desired properties
and the theorem is proved.
\end{proof}







\section{Applications}
\label{section-applications-flattening-by-blowing-up}

\noindent
In this section we apply the result on flattening by blowing up.

\begin{lemma}
\label{lemma-flat-after-blowing-up}
Let $S$ be a scheme. Let $f : X \to B$ be a morphism of algebraic
spaces over $S$. Let $U \subset B$ be an open subspace. Assume
\begin{enumerate}
\item $B$ is quasi-compact and quasi-separated,
\item $U$ is quasi-compact,
\item $f : X \to B$ is of finite type and quasi-separated, and
\item $f^{-1}(U) \to U$ is flat and locally of finite presentation.
\end{enumerate}
Then there exists a $U$-admissible blowup $B' \to B$ such that
the strict transform $X'$ of $X$ is flat and of finite presentation
over $B'$.
\end{lemma}

\begin{proof}
Let $B' \to B$ be a $U$-admissible blowup. Note that the strict transform
of $X$ is quasi-compact and quasi-separated over $B'$ as $X$ is quasi-compact
and quasi-separated over $B$. Hence we only need to worry about finding
a $U$-admissible blowup such that the strict transform becomes flat and
locally of finite presentation. We cannot directly apply
Theorem \ref{theorem-flatten-module} because $X$ is not locally of finite
presentation over $B$.

\medskip\noindent
Choose an affine scheme $V$ and a surjective \'etale morphism $V \to X$.
(This is possible as $X$ is quasi-compact as a finite type space over
the quasi-compact space $B$.) Then it suffices to show the result for
the morphism $V \to B$ (as strict transform commutes with \'etale
localization, see Divisors on Spaces,
Lemma \ref{spaces-divisors-lemma-strict-transform-local}).
Hence we may assume that $X \to B$ is separated as well as finite type.
In this case we can find a closed immersion $i : X \to Y$ with $Y \to B$
separated and of finite presentation, see
Limits of Spaces, Proposition
\ref{spaces-limits-proposition-separated-closed-in-finite-presentation}.

\medskip\noindent
Apply Theorem \ref{theorem-flatten-module} to $\mathcal{F} = i_*\mathcal{O}_X$
on $Y/B$. We find a $U$-admissible blowup $B' \to B$ such that strict
transform of $\mathcal{F}$ is flat over $B'$ and of finite presentation.
Let $X'$ be the strict transform of $X$ under the blowup $B' \to B$.
Let $i' : X' \to Y \times_B B'$ be the induced morphism.
Since taking strict transform commutes with pushforward along affine
morphisms (Divisors on Spaces, Lemma
\ref{spaces-divisors-lemma-strict-transform-affine}),
we see that $i'_*\mathcal{O}_{X'}$ is flat over $B'$ and of
finite presentation as a $\mathcal{O}_{Y \times_B B'}$-module.
Thus $X' \to B'$ is flat and locally of finite presentation.
This implies the lemma by our earlier remarks.
\end{proof}

\begin{lemma}
\label{lemma-finite-after-blowing-up}
Let $S$ be a scheme. Let $f : X \to B$ be a morphism of algebraic
spaces over $S$. Let $U \subset B$ be an open subspace. Assume
\begin{enumerate}
\item $B$ is quasi-compact and quasi-separated,
\item $U$ is quasi-compact,
\item $f : X \to B$ is proper, and
\item $f^{-1}(U) \to U$ is finite locally free.
\end{enumerate}
Then there exists a $U$-admissible blowup $B' \to B$ such that
the strict transform $X'$ of $X$ is finite locally free over $B'$.
\end{lemma}

\begin{proof}
By Lemma \ref{lemma-flat-after-blowing-up} we may assume that
$X \to B$ is flat and of finite presentation. After replacing
$B$ by a $U$-admissible blowup if necessary, we may assume
that $U \subset B$ is scheme theoretically dense. Then $f$ is
finite by Lemma \ref{lemma-proper-flat-finite-over-dense-open}.
Hence $f$ is finite locally free by
Morphisms of Spaces, Lemma \ref{spaces-morphisms-lemma-finite-flat}.
\end{proof}

\begin{lemma}
\label{lemma-isomorphism-after-blowing-up}
Let $S$ be a scheme. Let $f : X \to B$ be a morphism of algebraic
spaces over $S$. Let $U \subset B$ be an open subspace. Assume
\begin{enumerate}
\item $B$ is quasi-compact and quasi-separated,
\item $U$ is quasi-compact,
\item $f : X \to B$ is proper, and
\item $f^{-1}(U) \to U$ is an isomorphism.
\end{enumerate}
Then there exists a $U$-admissible blowup $B' \to B$ such that
the strict transform $X'$ of $X$ maps isomorphically to $B'$.
\end{lemma}

\begin{proof}
By Lemma \ref{lemma-flat-after-blowing-up} we may assume that
$X \to B$ is flat and of finite presentation. After replacing
$B$ by a $U$-admissible blowup if necessary, we may assume
that $U \subset B$ is scheme theoretically dense. Then $f$ is
finite by Lemma \ref{lemma-proper-flat-finite-over-dense-open}
and an open immersion by Lemma \ref{lemma-zariski}. Hence $f$
is an open immersion whose image is closed and contains the
dense open $U$, whence $f$ is an isomorphism.
\end{proof}

\begin{lemma}
\label{lemma-zariski-after-blowup}
Let $S$ be a scheme. Let $f : X \to B$ be a morphism of algebraic spaces
over $S$. Let $U \subset B$ be an open subspace. Assume
\begin{enumerate}
\item $B$ quasi-compact and quasi-separated,
\item $U$ is quasi-compact,
\item $f$ is of finite type
\item $f^{-1}(U) \to U$ is an isomorphism.
\end{enumerate}
Then there exists a $U$-admissible blowup $B' \to B$ such that $U$
is scheme theoretically dense in $B'$ and such that the strict
transform $X'$ of $X$ maps isomorphically to an open subspace of $B'$.
\end{lemma}

\begin{proof}
This lemma is a generalization of
Lemma \ref{lemma-isomorphism-after-blowing-up}.
As the composition of $U$-admissible blowups is $U$-admissible
(Divisors on Spaces, Lemma
\ref{spaces-divisors-lemma-composition-admissible-blowups})
we can proceed in stages. Pick a finite type quasi-coherent sheaf
of ideals $\mathcal{I} \subset \mathcal{O}_B$ with
$|B| \setminus |U| = |V(\mathcal{I})|$. Replace $B$ by the blowup
of $B$ in $\mathcal{I}$ and $X$ by the strict transform of $X$.
After this replacement $B \setminus U$ is the support of an effective
Cartier divisor $D$ (Divisors on Spaces, Lemma
\ref{spaces-divisors-lemma-blowing-up-gives-effective-Cartier-divisor}).
In particular $U$ is scheme theoretically dense in $B$
(Divisors on Spaces, Lemma
\ref{spaces-divisors-lemma-complement-effective-Cartier-divisor}).
Next, we do another $U$-admissible blowup to get to the situation where
$X \to B$ is flat and of finite presentation, see
Lemma \ref{lemma-flat-after-blowing-up}. Note that $U$ is still scheme
theoretically dense in $B$. Hence $X \to B$ is an open immersion by
Lemma \ref{lemma-zariski}.
\end{proof}

\noindent
The following lemma says that a modification can be dominated
by a blowup.

\begin{lemma}
\label{lemma-dominate-modification-by-blowup}
Let $S$ be a scheme. Let $f : X \to B$ be a morphism of algebraic spaces
over $S$. Let $U \subset B$ be an open subspace. Assume
\begin{enumerate}
\item $B$ is quasi-compact and quasi-separated,
\item $U$ is quasi-compact,
\item $f : X \to B$ is proper,
\item $f^{-1}(U) \to U$ us an isomorphism.
\end{enumerate}
Then there exists a $U$-admissible blowup $B' \to B$ which dominates $X$,
i.e., such that there exists a factorization $B' \to X \to B$
of the blowup morphism.
\end{lemma}

\begin{proof}
By Lemma \ref{lemma-isomorphism-after-blowing-up} we may find a $U$-admissible
blowup $B' \to B$ such that the strict transform $X'$ maps isomorphically to
$B'$. Then we can use $B' = X' \to X$ as the factorization.
\end{proof}

\begin{lemma}
\label{lemma-get-section-after-blowup}
Let $S$ be a scheme. Let $X$, $Y$ be algebraic spaces over $S$.
Let $U \subset W \subset Y$ be open subspaces.
Let $f : X \to W$ and let $s : U \to X$ be morphisms
such that $f \circ s = \text{id}_U$. Assume
\begin{enumerate}
\item $f$ is proper,
\item $Y$ is quasi-compact and quasi-separated, and
\item $U$ and $W$ are quasi-compact.
\end{enumerate}
Then there exists a $U$-admissible blowup $b : Y' \to Y$ and a morphism
$s' : b^{-1}(W) \to X$ extending $s$ with $f \circ s' = b|_{b^{-1}(W)}$.
\end{lemma}

\begin{proof}
We may and do replace $X$ by the scheme theoretic image of $s$.
Then $X \to W$ is an isomorphism over $U$, see
Morphisms of Spaces, Lemma
\ref{spaces-morphisms-lemma-scheme-theoretic-image-of-partial-section}.
By Lemma \ref{lemma-dominate-modification-by-blowup}
there exists a $U$-admissible blowup $W' \to W$ and an
extension $W' \to X$ of $s$.
We finish the proof by applying
Divisors on Spaces, Lemma \ref{spaces-divisors-lemma-extend-admissible-blowups}
to extend $W' \to W$ to a $U$-admissible blowup of $Y$.
\end{proof}















\section{Chow's lemma}
\label{section-chow}

\noindent
In this section we prove Chow's lemma
(Lemma \ref{lemma-chow-noetherian-separated}).
We encourage the reader to take a look at
Cohomology of Spaces, Section \ref{spaces-cohomology-section-weak-chow}
for a weak version of Chow's lemma that is easy to prove and sufficient
for many applications.

\medskip\noindent
Since we have yet to define projective morphisms of algebraic spaces,
the statements of lemmas (see for example
Lemma \ref{lemma-blowup-to-find-embedding}) will
involve representable proper morphisms, rather than projective ones.

\begin{lemma}
\label{lemma-find-common-blowups}
Let $S$ be a scheme. Let $Y$ be a quasi-compact and quasi-separated
algebraic space over $S$. Let $U \to X_1$ and $U \to X_2$ be open immersions
of algebraic spaces over $Y$ and assume $U$, $X_1$, $X_2$ of finite
type and separated over $Y$. Then there exists a commutative diagram
$$
\xymatrix{
X_1' \ar[d] \ar[r] & X & X_2' \ar[l] \ar[d] \\
X_1 & U \ar[l] \ar[lu] \ar[u] \ar[ru] \ar[r] & X_2
}
$$
of algebraic spaces over $Y$ where $X_i' \to X_i$ is a $U$-admissible
blowup, $X_i' \to X$ is an open immersion, and $X$ is separated and finite
type over $Y$.
\end{lemma}

\begin{proof}
Throughout the proof all the algebraic spaces will be separated of finite
type over $Y$. This in particular implies these algebraic spaces are
quasi-compact and quasi-separated and that the morphisms between them
will be quasi-compact and separated. See
Morphisms of Spaces, Sections
\ref{spaces-morphisms-section-separation-axioms} and
\ref{spaces-morphisms-section-quasi-compact}.
We will use that if $U \to W$ is an immersion of such spaces over $Y$,
then the scheme theoretic image $Z$ of $U$ in $W$ is a closed subspace
of $W$ and $U \to Z$ is an open immersion, $U \subset Z$ is scheme
theoretically dense, and $|U| \subset |Z|$ is dense. See
Morphisms of Spaces, Lemma
\ref{spaces-morphisms-lemma-quasi-compact-immersion}.

\medskip\noindent
Let $X_{12} \subset X_1 \times_Y X_2$ be the scheme theoretic image
of $U \to X_1 \times_Y X_2$. The projections $p_i : X_{12} \to X_i$
induce isomorphisms $p_i^{-1}(U) \to U$ by
Morphisms of Spaces, Lemma
\ref{spaces-morphisms-lemma-scheme-theoretic-image-of-partial-section}.
Choose a $U$-admissible blowup $X_i^i \to X_i$ such that
the strict transform $X_{12}^i$ of $X_{12}$ is isomorphic to an
open subspace of $X_i^i$, see
Lemma \ref{lemma-zariski-after-blowup}.
Let $\mathcal{I}_i \subset \mathcal{O}_{X_i}$ be the corresponding
finite type quasi-coherent sheaf of ideals.
Recall that $X_{12}^i \to X_{12}$ is the blowup in
$p_i^{-1}\mathcal{I}_i \mathcal{O}_{X_{12}}$, see
Divisors on Spaces, Lemma \ref{spaces-divisors-lemma-strict-transform}.
Let $X_{12}'$ be the blowup of $X_{12}$ in
$p_1^{-1}\mathcal{I}_1 p_2^{-1}\mathcal{I}_2 \mathcal{O}_{X_{12}}$, see
Divisors on Spaces, Lemma \ref{spaces-divisors-lemma-blowing-up-two-ideals}
for what this entails. We obtain a commutative diagram
$$
\xymatrix{
X_{12}' \ar[d] \ar[r] & X_{12}^2 \ar[d] \\
X_{12}^1 \ar[r] & X_{12}
}
$$
where all the morphisms are $U$-admissible blowing ups. Since
$X_{12}^i \subset X_i^i$ is an open we may choose a $U$-admissible blowup
$X_i' \to X_i^i$ restricting to $X_{12}' \to X_{12}^i$, see
Divisors on Spaces, Lemma
\ref{spaces-divisors-lemma-extend-admissible-blowups}.
Then $X_{12}' \subset X_i'$ is an open subspace and the diagram
$$
\xymatrix{
X_{12}' \ar[d] \ar[r] & X_i' \ar[d] \\
X_{12}^i \ar[r] & X_i^i
}
$$
is commutative with vertical arrows blowing ups and horizontal arrows
open immersions. Note that $X'_{12} \to X_1' \times_Y X_2'$ is
an immersion and proper (use that $X'_{12} \to X_{12}$ is proper
and $X_{12} \to X_1 \times_Y X_2$ is closed and $X_1' \times_Y X_2' \to
X_1 \times_Y X_2$ is separated and apply Morphisms of Spaces, Lemma
\ref{spaces-morphisms-lemma-universally-closed-permanence}).
Thus $X'_{12} \to  X_1' \times_Y X_2'$ is a closed immersion.
If we define $X$ by glueing $X_1'$ and $X_2'$ along the common open
subspace $X_{12}'$, then $X \to Y$ is of finite type and
separated\footnote{Because we may check closedness of the diagonal
$X \to X \times_Y X$ over the four open parts $X'_i \times_Y X'_j$
of $X \times_Y X$ where it is clear.}. As compositions of
$U$-admissible blowups are $U$-admissible blowups
(Divisors on Spaces, Lemma
\ref{spaces-divisors-lemma-composition-admissible-blowups})
the lemma is proved.
\end{proof}

\begin{lemma}
\label{lemma-blowup-to-find-embedding}
Let $S$ be a scheme. Let $f : X \to Y$ be a morphism of algebraic spaces
over $S$. Let $U \subset X$ be an open subspace. Assume
\begin{enumerate}
\item $U$ is quasi-compact,
\item $Y$ is quasi-compact and quasi-separated,
\item there exists an immersion $U \to \mathbf{P}^n_Y$ over $Y$,
\item $f$ is of finite type and separated.
\end{enumerate}
Then there exists a commutative diagram
$$
\xymatrix{
& U \ar[ld] \ar[d] \ar[rd] \ar[rrd] \\
X \ar[rd] & X' \ar[l] \ar[d] \ar[r] & Z' \ar[ld] \ar[r] & Z \ar[ld] \\
& Y & \mathbf{P}^n_Y \ar[l]
}
$$
where
the arrows with source $U$ are open immersions,
$X' \to X$ is a $U$-admissible blowup,
$X' \to Z'$ is an open immersion,
$Z' \to Y$ is a proper and representable morphism of algebraic spaces.
More precisely, $Z' \to Z$ is a $U$-admissible blowup
and $Z \to \mathbf{P}^n_Y$ is a closed immersion.
\end{lemma}

\begin{proof}
Let $Z \subset \mathbf{P}^n_Y$ be the scheme theoretic image of
the immersion $U \to \mathbf{P}^n_Y$. Since $U \to \mathbf{P}^n_Y$
is quasi-compact we see that $U \subset Z$ is a
(scheme theoretically) dense open subspace
(Morphisms of Spaces, Lemma
\ref{spaces-morphisms-lemma-quasi-compact-immersion}).
Apply Lemma \ref{lemma-find-common-blowups} to find a diagram
$$
\xymatrix{
X' \ar[d] \ar[r] & \overline{X}' & Z' \ar[l] \ar[d] \\
X & U \ar[l] \ar[lu] \ar[u] \ar[ru] \ar[r] & Z
}
$$
with properties as listed in the statement of that lemma.
As $X' \to X$ and $Z' \to Z$ are $U$-admissible blowups
we find that $U$ is a scheme theoretically dense open of
both $X'$ and $Z'$ (see Divisors on Spaces, Lemmas
\ref{spaces-divisors-lemma-blowing-up-gives-effective-Cartier-divisor} and
\ref{spaces-divisors-lemma-complement-effective-Cartier-divisor}).
Since $Z' \to Z \to Y$ is proper we see that $Z' \subset \overline{X}'$
is a closed subspace (see Morphisms of Spaces, Lemma
\ref{spaces-morphisms-lemma-universally-closed-permanence}).
It follows that $X' \subset Z'$ (scheme theoretically), hence $X'$
is an open subspace of $Z'$ (small detail omitted) and the lemma is proved.
\end{proof}

\begin{lemma}
\label{lemma-chow-noetherian}
Let $S$ be a scheme. Let $f : X \to Y$ be a morphism of algebraic spaces
over $S$. Assume $f$ separated, of finite type, and $Y$ Noetherian.
Then there exists a dense open subspace $U \subset X$ and
a commutative diagram
$$
\xymatrix{
& U \ar[ld] \ar[d] \ar[rd] \ar[rrd] \\
X \ar[rd] & X' \ar[l] \ar[d] \ar[r] & Z' \ar[ld] \ar[r] & Z \ar[ld] \\
& Y & \mathbf{P}^n_Y \ar[l]
}
$$
where
the arrows with source $U$ are open immersions,
$X' \to X$ is a $U$-admissible blowup,
$X' \to Z'$ is an open immersion,
$Z' \to Y$ is a proper and representable morphism of algebraic spaces.
More precisely, $Z' \to Z$ is a $U$-admissible blowup
and $Z \to \mathbf{P}^n_Y$ is a closed immersion.
\end{lemma}

\begin{proof}
By Limits of Spaces, Lemma
\ref{spaces-limits-lemma-embedding-into-affine-over-qs}
there exists a dense open subspace $U \subset X$ and an immersion
$U \to \mathbf{A}^n_Y$ over $Y$. Composing with the open immersion
$\mathbf{A}^n_Y \to \mathbf{P}^n_Y$ we obtain a situation as in
Lemma \ref{lemma-blowup-to-find-embedding} and the result
follows.
\end{proof}

\begin{remark}
\label{remark-chow-Noetherian}
In Lemmas \ref{lemma-blowup-to-find-embedding} and \ref{lemma-chow-noetherian}
the morphism $g : Z' \to Y$ is a composition of projective morphisms.
Presumably (by the analogue for algebraic spaces of
Morphisms, Lemma \ref{morphisms-lemma-ample-composition})
there exists a $g$-ample invertible sheaf on $Z'$.
If we ever need this, then we will state and prove this here.
\end{remark}

\noindent
The following result is \cite[IV Theorem 3.1]{Kn}. Note that the immersion
$X' \to \mathbf{P}^n_Y$ is quasi-compact, hence can be factored as
$X' \to Z' \to \mathbf{P}^n_Y$ where the first morphism is an
open immersion and the second morphism a closed immersion
(Morphisms of Spaces, Lemma
\ref{spaces-morphisms-lemma-quasi-compact-immersion}).

\begin{lemma}[Chow's lemma]
\label{lemma-chow-noetherian-separated}
\begin{reference}
\cite[IV Theorem 3.1]{Kn}
\end{reference}
Let $S$ be a scheme. Let $f : X \to Y$ be a morphism of algebraic spaces
over $S$. Assume $f$ separated of finite type, and $Y$ separated and
Noetherian. Then there exists a commutative diagram
$$
\xymatrix{
X \ar[rd] & X' \ar[l] \ar[d] \ar[r] & \mathbf{P}^n_Y \ar[ld] \\
& Y
}
$$
where $X' \to X$ is a $U$-admissible blowup for some dense open
$U \subset X$ and the morphism $X' \to \mathbf{P}^n_Y$ is an immersion.
\end{lemma}

\begin{proof}
In this first paragraph of the proof we reduce the lemma to the case
where $Y$ is of finite type over $\Spec(\mathbf{Z})$.
We may and do replace the base scheme $S$ by $\Spec(\mathbf{Z})$.
We can write $Y = \lim Y_i$ as a directed limit of separated
algebraic spaces of finite type over $\Spec(\mathbf{Z})$, see
Limits of Spaces, Proposition \ref{spaces-limits-proposition-approximate} and
Lemma \ref{spaces-limits-lemma-descend-separated}.
For all $i$ sufficiently large we can find a separated finite type morphism
$X_i \to Y_i$ such that $X = Y \times_{Y_i} X_i$, see
Limits of Spaces, Lemmas
\ref{spaces-limits-lemma-descend-finite-presentation} and
\ref{spaces-limits-lemma-descend-separated-morphism}.
Let $\eta_1, \ldots, \eta_n$ be the generic points of the irreducible
components of $|X|$ ($X$ is Noetherian as a finite type separated
algebraic space over the Noetherian algebraic space $Y$ and therefore
$|X|$ is a Noetherian topological space).
By Limits of Spaces, Lemma \ref{spaces-limits-lemma-topology-limit}
we find that the images of $\eta_1, \ldots, \eta_n$ in $|X_i|$
are distinct for $i$ large enough. We may replace
$X_i$ by the scheme theoretic image of the (quasi-compact, in fact affine)
morphism $X \to X_i$.
After this replacement we see that the images
of $\eta_1, \ldots, \eta_n$ in $|X_i|$ are the generic points of the
irreducible components of $|X_i|$, see
Morphisms of Spaces, Lemma
\ref{spaces-morphisms-lemma-quasi-compact-scheme-theoretic-image}.
Having said this, suppose we can find a diagram
$$
\xymatrix{
X_i \ar[rd] & X_i' \ar[l] \ar[d] \ar[r] & \mathbf{P}^n_{Y_i} \ar[ld] \\
& Y
}
$$
where $X_i' \to X_i$ is a $U_i$-admissible blowup for some dense open
$U_i \subset X_i$ and the morphism $X_i' \to \mathbf{P}^n_{Y_i}$
is an immersion. Then the strict transform $X' \to X$ of $X$ relative
to $X_i' \to X_i$ is a $U$-admissible blowing up where $U \subset X$
is the inverse image of $U_i$ in $X$. Because of our carefully chosen
index $i$ it follows that $\eta_1, \ldots, \eta_n \in |U|$ and
$U \subset X$ is dense. Moreover, $X' \to \mathbf{P}^n_Y$ is an
immersion as $X'$ is closed in
$X_i' \times_{X_i} X = X_i' \times_{Y_i} Y$
which comes with an immersion into $\mathbf{P}^n_Y$. Thus we have reduced
to the situation of the following paragraph.

\medskip\noindent
Assume that $Y$ is separated of finite type over $\Spec(\mathbf{Z})$.
Then $X \to \Spec(\mathbf{Z})$ is separated of finite type as well.
We apply Lemma \ref{lemma-chow-noetherian} to $X \to \Spec(\mathbf{Z})$
to find a dense open subspace $U \subset X$ and a commutative diagram
$$
\xymatrix{
& U \ar[ld] \ar[d] \ar[rd] \ar[rrd] \\
X \ar[rd] & X' \ar[l] \ar[d] \ar[r] & Z' \ar[ld] \ar[r] & Z \ar[ld] \\
& \Spec(\mathbf{Z}) & \mathbf{P}^n_\mathbf{Z} \ar[l]
}
$$
with all the properties listed in the lemma.
Note that $Z$ has an ample invertible sheaf, namely
$\mathcal{O}_{\mathbf{P}^n}(1)|_Z$. Hence $Z' \to Z$
is a H-projective morphism by Morphisms, Lemma
\ref{morphisms-lemma-projective-over-quasi-projective-is-H-projective}.
It follows that $Z' \to \Spec(\mathbf{Z})$ is H-projective
by Morphisms, Lemma \ref{morphisms-lemma-H-projective-composition}.
Thus there exists a closed immersion
$Z' \to \mathbf{P}^m_{\Spec(\mathbf{Z})}$ for some $m \geq 0$.
It follows that the diagonal morphism
$$
X' \to Y \times \mathbf{P}^m_\mathbf{Z} = \mathbf{P}^m_Y
$$
is an immersion (because the composition with the projection
to $\mathbf{P}^m_\mathbf{Z}$ is an immersion) and we win.
\end{proof}






\section{Variants of Chow's Lemma}
\label{section-chows-lemma}

\noindent
In this section we prove a number of variants of Chow's lemma dealing
with morphisms between non-Noetherian algebraic spaces.
The Noetherian versions are
Lemma \ref{lemma-chow-noetherian}
and
Lemma \ref{lemma-chow-noetherian-separated}.

\begin{lemma}
\label{lemma-chow-finite-type}
Let $S$ be a  scheme. Let $Y$ be a quasi-compact and quasi-separated
algebraic space over $S$. Let $f : X \to Y$ be a separated morphism of
finite type. Then there exists a commutative diagram
$$
\xymatrix{
X \ar[rd] & X' \ar[l] \ar[d] \ar[r] & \overline{X}' \ar[ld] \\
& Y
}
$$
where $X' \to X$ is proper surjective,
$X' \to \overline{X}'$ is an open immersion, and
$\overline{X}' \to Y$ is proper and representable morphism
of algebraic spaces.
\end{lemma}

\begin{proof}
By
Limits of Spaces, Proposition
\ref{spaces-limits-proposition-separated-closed-in-finite-presentation}
we can find a closed immersion $X \to X_1$ where $X_1$ is separated
and of finite presentation over $Y$. Clearly, if we prove the assertion
for $X_1 \to Y$, then the result follows for $X$. Hence we may assume that
$X$ is of finite presentation over $Y$.

\medskip\noindent
We may and do replace the base scheme $S$ by $\Spec(\mathbf{Z})$.
Write $Y = \lim_i Y_i$ as a directed limit of
quasi-separated algebraic spaces of finite type over $\Spec(\mathbf{Z})$, see
Limits of Spaces,
Proposition \ref{spaces-limits-proposition-approximate}.
By
Limits of Spaces,
Lemma \ref{spaces-limits-lemma-descend-finite-presentation}
we can
find an index $i \in I$ and a scheme $X_i \to Y_i$ of finite presentation
so that $X = Y \times_{Y_i} X_i$.
By
Limits of Spaces,
Lemma \ref{spaces-limits-lemma-descend-separated-morphism}
we may assume that $X_i \to Y_i$ is separated.
Clearly, if we prove the assertion for
$X_i$ over $Y_i$, then the assertion holds for $X$. The case
$X_i \to Y_i$ is treated by
Lemma \ref{lemma-chow-noetherian}.
\end{proof}

\begin{lemma}
\label{lemma-chow-finite-type-separated}
Let $S$ be a scheme. Let $f : X \to Y$ be a morphism of algebraic spaces
over $S$. Assume $f$ separated of finite type, and $Y$ separated and
quasi-compact. Then there exists a commutative diagram
$$
\xymatrix{
X \ar[rd] & X' \ar[l] \ar[d] \ar[r] & \mathbf{P}^n_Y \ar[ld] \\
& Y
}
$$
where $X' \to X$ is proper surjective morphism
and the morphism $X' \to \mathbf{P}^n_Y$ is an immersion.
\end{lemma}

\begin{proof}
By
Limits of Spaces, Proposition
\ref{spaces-limits-proposition-separated-closed-in-finite-presentation}
we can find a closed immersion $X \to X_1$ where $X_1$ is separated
and of finite presentation over $Y$. Clearly, if we prove the assertion
for $X_1 \to Y$, then the result follows for $X$. Hence we may assume that
$X$ is of finite presentation over $Y$.

\medskip\noindent
We may and do replace the base scheme $S$ by $\Spec(\mathbf{Z})$.
Write $Y = \lim_i Y_i$ as a directed limit of
quasi-separated algebraic spaces of finite type over $\Spec(\mathbf{Z})$, see
Limits of Spaces,
Proposition \ref{spaces-limits-proposition-approximate}.
By Limits of Spaces, Lemma \ref{spaces-limits-lemma-descend-separated}
we may assume that $Y_i$ is separated for all $i$.
By
Limits of Spaces,
Lemma \ref{spaces-limits-lemma-descend-finite-presentation}
we can
find an index $i \in I$ and a scheme $X_i \to Y_i$ of finite presentation
so that $X = Y \times_{Y_i} X_i$.
By
Limits of Spaces,
Lemma \ref{spaces-limits-lemma-descend-separated-morphism}
we may assume that $X_i \to Y_i$ is separated.
Clearly, if we prove the assertion for
$X_i$ over $Y_i$, then the assertion holds for $X$. The case
$X_i \to Y_i$ is treated by
Lemma \ref{lemma-chow-noetherian-separated}.
\end{proof}






\section{Grothendieck's existence theorem}
\label{section-existence-proper}

\noindent
In this section we discuss Grothendieck's existence theorem for algebraic
spaces. Instead of developing a theory of ``formal algebraic spaces''
we temporarily develop a bit of language that replaces the notion of a
``coherent module on a Noetherian adic formal space''.

\medskip\noindent
Let $S$ be a scheme. Let $X$ be a Noetherian algebraic space over $S$.
Let $\mathcal{I} \subset \mathcal{O}_X$ be a quasi-coherent sheaf of ideals.
Below we will consider inverse systems $(\mathcal{F}_n)$ of coherent
$\mathcal{O}_X$-modules such that
\begin{enumerate}
\item $\mathcal{F}_n$ is annihilated by $\mathcal{I}^n$, and
\item the transition maps induce isomorphisms
$\mathcal{F}_{n + 1}/\mathcal{I}^n\mathcal{F}_{n + 1} \to \mathcal{F}_n$.
\end{enumerate}
A morphism $\alpha : (\mathcal{F}_n) \to (\mathcal{G}_n)$
of such inverse systems is simply a compatible system of morphisms
$\alpha_n : \mathcal{F}_n \to \mathcal{G}_n$.
Let us denote the category of these inverse systems with
$\textit{Coh}(X, \mathcal{I})$. We will develop some theory regarding
these systems that will parallel to the corresponding
results in the case of schemes, see
Cohomology of Schemes, Sections \ref{coherent-section-existence},
\ref{coherent-section-existence-proper},
\ref{coherent-section-existence-proper-support}, and
\ref{coherent-section-algebraization}.

\medskip\noindent
Functoriality. Let $f : X \to Y$ be a
morphism of Noetherian algebraic spaces over a scheme $S$, and
let $\mathcal{J} \subset \mathcal{O}_Y$ be a quasi-coherent sheaf
of ideals. Set $\mathcal{I} = f^{-1}\mathcal{J}\mathcal{O}_X$.
In this situation there is a functor
$$
f^* : \textit{Coh}(Y, \mathcal{J}) \longrightarrow \textit{Coh}(X, \mathcal{I})
$$
which sends $(\mathcal{G}_n)$ to $(f^*\mathcal{G}_n)$. Compare with
Cohomology of Schemes, Lemma \ref{coherent-lemma-inverse-systems-pullback}.
If $f$ is \'etale, then we may think of this as simply the restriction
of the system to $X$, see Properties of Spaces, 
Equation \ref{spaces-properties-equation-restrict-modules}.

\medskip\noindent
\'Etale descent. Let $S$ be a scheme. Let $U_0 \to X$ be a surjective
\'etale morphism of Noetherian algebraic spaces. Set
$U_1 = U_0 \times_X U_0$ and $U_2 = U_0 \times_X U_0 \times_X U_0$.
Let $\mathcal{I} \subset \mathcal{O}_{X}$ be a quasi-coherent sheaf
of ideals. Set $\mathcal{I}_i = \mathcal{I}|_{U_i}$. In this situation
we obtain a diagram of categories
$$
\xymatrix{
\textit{Coh}(X, \mathcal{I}) \ar[r] &
\textit{Coh}(U_0, \mathcal{I}_0) \ar@<0.5ex>[r] \ar@<-0.5ex>[r] &
\textit{Coh}(U_1, \mathcal{I}_1) \ar@<1ex>[r] \ar[r] \ar@<-1ex>[r] &
\textit{Coh}(U_2, \mathcal{I}_2)
}
$$
an the first arrow presents $\textit{Coh}(X, \mathcal{I})$ as the
homotopy limit of the right part of the diagram. More precisely, given
a {\it descent datum}, i.e., a pair $((\mathcal{G}_n), \varphi)$ where
$(\mathcal{G}_n)$ is an object of $\textit{Coh}(U_0, \mathcal{I}_0)$ and
$\varphi : \text{pr}_0^*(\mathcal{G}_n) \to \text{pr}_1^*(\mathcal{G}_n)$
is an isomorphism in $\textit{Coh}(U_1, \mathcal{I}_1)$
satisfying the cocycle condition in $\textit{Coh}(U_2, \mathcal{I}_2)$,
then there exists a unique object $(\mathcal{F}_n)$ of
$\textit{Coh}(X, \mathcal{I})$ whose associated canonical descent datum
is isomorphic to $((\mathcal{G}_n), \varphi)$. Compare with
Descent on Spaces, Definition
\ref{spaces-descent-definition-descent-datum-effective-quasi-coherent}.
The proof of this statement follows immediately by applying
Descent on Spaces, Proposition
\ref{spaces-descent-proposition-fpqc-descent-quasi-coherent}
to the descent data $(\mathcal{G}_n, \varphi_n)$ for varying $n$.

\begin{lemma}
\label{lemma-inverse-systems-abelian}
Let $S$ be a scheme. Let $X$ be a Noetherian algebraic space over $S$ and
let $\mathcal{I} \subset \mathcal{O}_X$ be a quasi-coherent sheaf of ideals.
\begin{enumerate}
\item The category $\textit{Coh}(X, \mathcal{I})$ is abelian.
\item Exactness in $\textit{Coh}(X, \mathcal{I})$
can be checked \'etale locally.
\item For any flat morphism $f : X' \to X$ of Noetherian algebraic spaces
the functor $f^* : \textit{Coh}(X, \mathcal{I}) \to
\textit{Coh}(X', f^{-1}\mathcal{I}\mathcal{O}_{X'})$ is exact.
\end{enumerate}
\end{lemma}

\begin{proof}
Proof of (1). Choose an affine scheme $U_0$ and a surjective \'etale morphism
$U_0 \to X$. Set $U_1 = U_0 \times_X U_0$ and
$U_2 = U_0 \times_X U_0 \times_X U_0$ as in our discussion of
\'etale descent above. The categories $\textit{Coh}(U_i, \mathcal{I}_i)$
are abelian
(Cohomology of Schemes, Lemma \ref{coherent-lemma-inverse-systems-abelian})
and the pullback functors are exact functors
$\textit{Coh}(U_0, \mathcal{I}_0) \to \textit{Coh}(U_1, \mathcal{I}_1)$
and
$\textit{Coh}(U_1, \mathcal{I}_1) \to \textit{Coh}(U_2, \mathcal{I}_2)$
(Cohomology of Schemes, Lemma \ref{coherent-lemma-inverse-systems-pullback}).
The lemma then follows formally from the description of
$\textit{Coh}(X, \mathcal{I})$ as a category of descent data.
Some details omitted; compare with the proof of
Groupoids, Lemma \ref{groupoids-lemma-abelian}.

\medskip\noindent
Part (2) follows immediately from the discussion in the previous paragraph.
In the situation of (3) choose a commutative diagram
$$
\xymatrix{
U' \ar[d] \ar[r] & U \ar[d] \\
X' \ar[r] & X
}
$$
where $U'$ and $U$ are affine schemes and the vertical morphisms are
surjective \'etale. Then $U' \to U$ is a flat morphism of Noetherian
schemes (Morphisms of Spaces, Lemma \ref{spaces-morphisms-lemma-flat-local})
whence the pullback functor
$\textit{Coh}(U, \mathcal{I}\mathcal{O}_U) \to
\textit{Coh}(U', \mathcal{I}\mathcal{O}_{U'})$
is exact by
Cohomology of Schemes, Lemma \ref{coherent-lemma-inverse-systems-pullback}.
Since we can check exactness in $\textit{Coh}(X, \mathcal{O}_X)$
on $U$ and similarly for $X', U'$ the assertion follows.
\end{proof}

\begin{lemma}
\label{lemma-inverse-systems-surjective}
Let $S$ be a scheme. Let $X$ be a Noetherian algebraic space over $S$
and let $\mathcal{I} \subset \mathcal{O}_X$ be a quasi-coherent sheaf
of ideals. A map $(\mathcal{F}_n) \to (\mathcal{G}_n)$ is surjective in
$\textit{Coh}(X, \mathcal{I})$ if and only if
$\mathcal{F}_1 \to \mathcal{G}_1$ is surjective.
\end{lemma}

\begin{proof}
We can check on an affine \'etale cover of $X$ by
Lemma \ref{lemma-inverse-systems-abelian}.
Thus we reduce to the case of schemes which is
Cohomology of Schemes, Lemma \ref{coherent-lemma-inverse-systems-surjective}.
\end{proof}

\noindent
Let $S$ be a scheme. Let $X$ be a Noetherian algebraic space over $S$ and
let $\mathcal{I} \subset \mathcal{O}_X$ be a quasi-coherent sheaf of ideals.
There is a functor
\begin{equation}
\label{equation-completion-functor}
\textit{Coh}(\mathcal{O}_X) \longrightarrow \textit{Coh}(X, \mathcal{I}), \quad
\mathcal{F} \longmapsto \mathcal{F}^\wedge
\end{equation}
which associates to the coherent $\mathcal{O}_X$-module $\mathcal{F}$
the object $\mathcal{F}^\wedge = (\mathcal{F}/\mathcal{I}^n\mathcal{F})$
of $\textit{Coh}(X, \mathcal{I})$.

\begin{lemma}
\label{lemma-exact}
The functor (\ref{equation-completion-functor}) is exact.
\end{lemma}

\begin{proof}
It suffices to check this \'etale locally on $X$, see
Lemma \ref{lemma-inverse-systems-abelian}.
Thus we reduce to the case of schemes which is
Cohomology of Schemes, Lemma \ref{coherent-lemma-exact}.
\end{proof}

\begin{lemma}
\label{lemma-completion-internal-hom}
Let $S$ be a scheme. Let $X$ be a Noetherian algebraic space over $S$ and
let $\mathcal{I} \subset \mathcal{O}_X$ be a quasi-coherent sheaf of ideals.
Let $\mathcal{F}$, $\mathcal{G}$ be coherent $\mathcal{O}_X$-modules. Set
$\mathcal{H} = \SheafHom_{\mathcal{O}_X}(\mathcal{F}, \mathcal{G})$.
Then
$$
\lim H^0(X, \mathcal{H}/\mathcal{I}^n\mathcal{H}) =
\Mor_{\textit{Coh}(X, \mathcal{I})}
(\mathcal{F}^\wedge, \mathcal{G}^\wedge).
$$
\end{lemma}

\begin{proof}
Since $\mathcal{H}$ is a sheaf on $X_\etale$ and since
we have \'etale descent for objects of $\textit{Coh}(X, \mathcal{I})$
it suffices to prove this \'etale locally.
Thus we reduce to the case of schemes which is
Cohomology of Schemes, Lemma \ref{coherent-lemma-completion-internal-hom}.
\end{proof}

\noindent
We introduce the setting that we will focus on throughout the
rest of this section.

\begin{situation}
\label{situation-existence}
Here $A$ is a Noetherian ring complete with respect to an ideal $I$.
Also $f : X \to \Spec(A)$ is a finite type separated morphism of
algebraic spaces and $\mathcal{I} = I\mathcal{O}_X$.
\end{situation}

\noindent
In this situation we denote
$$
\textit{Coh}_{\text{support proper over } A}(\mathcal{O}_X)
$$
be the full subcategory of $\textit{Coh}(\mathcal{O}_X)$
consisting of those coherent $\mathcal{O}_X$-modules whose
support is proper over $\Spec(A)$, or equivalently whose
scheme theoretic support is proper over $\Spec(A)$, see
Derived Categories of Spaces, Lemma
\ref{spaces-perfect-lemma-module-support-proper-over-base}.
Similarly, we let
$$
\textit{Coh}_{\text{support proper over } A}(X, \mathcal{I})
$$
be the full subcategory of $\textit{Coh}(X, \mathcal{I})$
consisting of those objects $(\mathcal{F}_n)$ such that
the support of $\mathcal{F}_1$ is proper over $\Spec(A)$.
Since the support of a quotient module is contained in the support
of the module, it follows that (\ref{equation-completion-functor})
induces a functor
\begin{equation}
\label{equation-completion-functor-proper-over-A}
\textit{Coh}_{\text{support proper over }A}(\mathcal{O}_X)
\longrightarrow
\textit{Coh}_{\text{support proper over }A}(X, \mathcal{I})
\end{equation}
Our first result is that this functor is fully faithful.

\begin{lemma}
\label{lemma-fully-faithful}
In Situation \ref{situation-existence}.
Let $\mathcal{F}$, $\mathcal{G}$ be coherent $\mathcal{O}_X$-modules.
Assume that the intersection of the supports of
$\mathcal{F}$ and $\mathcal{G}$ is proper over $\Spec(A)$. Then the map
$$
\Mor_{\textit{Coh}(\mathcal{O}_X)}(\mathcal{F}, \mathcal{G})
\longrightarrow
\Mor_{\textit{Coh}(X, \mathcal{I})}
(\mathcal{F}^\wedge, \mathcal{G}^\wedge)
$$
coming from (\ref{equation-completion-functor}) is a bijection.
In particular, (\ref{equation-completion-functor-proper-over-A})
is fully faithful.
\end{lemma}

\begin{proof}
Let $\mathcal{H} = \SheafHom_{\mathcal{O}_X}(\mathcal{G}, \mathcal{F})$.
This is a coherent $\mathcal{O}_X$-module because its restriction
of schemes \'etale over $X$ is coherent by
Modules, Lemma \ref{modules-lemma-internal-hom-locally-kernel-direct-sum}.
By Lemma \ref{lemma-completion-internal-hom} the map
$$
\lim_n H^0(X, \mathcal{H}/\mathcal{I}^n\mathcal{H})
\to
\Mor_{\textit{Coh}(X, \mathcal{I})}
(\mathcal{G}^\wedge, \mathcal{F}^\wedge)
$$
is bijective. Let $i : Z \to X$ be the scheme theoretic support of
$\mathcal{H}$. It is clear that $Z$ is a closed subspace such
that $|Z|$ is contained in the intersection of the supports of $\mathcal{F}$
and $\mathcal{G}$. Hence $Z \to \Spec(A)$ is proper by assumption
(see Derived Categories of Spaces, Section
\ref{spaces-perfect-section-proper-over-base}).
Write $\mathcal{H} = i_*\mathcal{H}'$ for some coherent
$\mathcal{O}_Z$-module $\mathcal{H}'$. We have
$i_*(\mathcal{H}'/I^n\mathcal{H}') = \mathcal{H}/I^n\mathcal{H}$.
Hence we obtain
\begin{align*}
\lim_n H^0(X, \mathcal{H}/\mathcal{I}^n\mathcal{H})
& =
\lim_n H^0(Z, \mathcal{H}'/\mathcal{I}^n\mathcal{H}') \\
& =
H^0(Z, \mathcal{H}') \\
& =
H^0(X, \mathcal{H}) \\
& = 
\Mor_{\textit{Coh}(\mathcal{O}_X)}(\mathcal{F}, \mathcal{G})
\end{align*}
the second equality by the theorem on formal functions
(Cohomology of Spaces, Lemma
\ref{spaces-cohomology-lemma-spell-out-theorem-formal-functions}).
This proves the lemma.
\end{proof}

\begin{remark}
\label{remark-inverse-systems-kernel-cokernel-annihilated-by}
Let $S$ be a scheme. Let $X$ be a Noetherian algebraic space over $S$ and let
$\mathcal{I}, \mathcal{K} \subset \mathcal{O}_X$ be quasi-coherent sheaves of
ideals. Let $\alpha : (\mathcal{F}_n) \to (\mathcal{G}_n)$ be a morphism of
$\textit{Coh}(X, \mathcal{I})$.
Given an affine scheme $U = \Spec(A)$ and a surjective \'etale morphism
$U \to X$ denote $I, K \subset A$ the ideals corresponding to the restrictions
$\mathcal{I}|_U, \mathcal{K}|_U$. Denote $\alpha_U : M \to N$ of finite
$A^\wedge$-modules which corresponds to $\alpha|_U$ via
Cohomology of Schemes, Lemma \ref{coherent-lemma-inverse-systems-affine}.
We claim the following are equivalent
\begin{enumerate}
\item there exists an integer $t \geq 1$ such that
$\Ker(\alpha_n)$ and $\Coker(\alpha_n)$
are annihilated by $\mathcal{K}^t$ for all $n \geq 1$,
\item for any (or some) affine open $\Spec(A) = U \subset X$ as above
the modules $\Ker(\alpha_U)$ and $\Coker(\alpha_U)$
are annihilated by $K^t$ for some integer $t \geq 1$.
\end{enumerate}
If these equivalent conditions hold we will say that $\alpha$ is a
{\it map whose kernel and cokernel are annihilated by a power of
$\mathcal{K}$}. To see the equivalence we refer to
Cohomology of Schemes, Remark
\ref{coherent-remark-inverse-systems-kernel-cokernel-annihilated-by}.
\end{remark}

\begin{lemma}
\label{lemma-existence-easy}
Let $S$ be a scheme. Let $X$ be a Noetherian algebraic space over $S$ and let
$\mathcal{I} \subset \mathcal{O}_X$ be a quasi-coherent sheaf of ideals.
Let $\mathcal{G}$ be a coherent $\mathcal{O}_X$-module, $(\mathcal{F}_n)$
an object of $\textit{Coh}(X, \mathcal{I})$, and
$\alpha : (\mathcal{F}_n) \to \mathcal{G}^\wedge$
a map whose kernel and cokernel are annihilated by a power of $\mathcal{I}$.
Then there exists a unique (up to unique isomorphism) triple
$(\mathcal{F}, a, \beta)$ where
\begin{enumerate}
\item $\mathcal{F}$ is a coherent $\mathcal{O}_X$-module,
\item $a : \mathcal{F} \to \mathcal{G}$ is an $\mathcal{O}_X$-module map
whose kernel and cokernel are annihilated by a power of $\mathcal{I}$,
\item $\beta : (\mathcal{F}_n) \to \mathcal{F}^\wedge$ is an isomorphism, and
\item $\alpha = a^\wedge \circ \beta$.
\end{enumerate}
\end{lemma}

\begin{proof}
The uniqueness and \'etale descent for objects of
$\textit{Coh}(X, \mathcal{I})$ and $\textit{Coh}(\mathcal{O}_X)$
implies it suffices to construct $(\mathcal{F}, a, \beta)$ \'etale
locally on $X$. Thus we reduce to the case of schemes which is
Cohomology of Schemes, Lemma \ref{coherent-lemma-existence-easy}.
\end{proof}

\begin{lemma}
\label{lemma-existence-tricky}
In Situation \ref{situation-existence}. Let $\mathcal{K} \subset \mathcal{O}_X$
be a quasi-coherent sheaf of ideals. Let $X_e \subset X$ be the closed subspace
cut out by $\mathcal{K}^e$. Let $\mathcal{I}_e = \mathcal{I}\mathcal{O}_{X_e}$.
Let $(\mathcal{F}_n)$ be an object of
$\textit{Coh}_{\text{support proper over } A}(X, \mathcal{I})$.
Assume
\begin{enumerate}
\item the functor
$\textit{Coh}_{\text{support proper over } A}(\mathcal{O}_{X_e})
\to \textit{Coh}_{\text{support proper over } A}(X_e, \mathcal{I}_e)$
is an equivalence for all $e \geq 1$, and
\item there exists an object $\mathcal{H}$ of
$\textit{Coh}_{\text{support proper over } A}(\mathcal{O}_X)$ and a map
$\alpha : (\mathcal{F}_n) \to \mathcal{H}^\wedge$ whose
kernel and cokernel are annihilated by a power of $\mathcal{K}$.
\end{enumerate}
Then $(\mathcal{F}_n)$ is in the essential image of
(\ref{equation-completion-functor-proper-over-A}).
\end{lemma}

\begin{proof}
During this proof we will use without further mention that for a closed
immersion $i : Z \to X$ the functor $i_*$ gives an equivalence between the
category of coherent modules on $Z$ and coherent modules on $X$ annihilated
by the ideal sheaf of $Z$, see
Cohomology of Spaces, Lemma \ref{spaces-cohomology-lemma-i-star-equivalence}.
In particular we think of
$$
\textit{Coh}_{\text{support proper over } A}(\mathcal{O}_{X_e})
\subset
\textit{Coh}_{\text{support proper over } A}(\mathcal{O}_X)
$$
as the full subcategory of consisting of modules annihilated by
$\mathcal{K}^e$ and
$$
\textit{Coh}_{\text{support proper over } A}(X_e, \mathcal{I}_e)
\subset
\textit{Coh}_{\text{support proper over } A}(X, \mathcal{I})
$$
as the full subcategory of objects annihilated by $\mathcal{K}^e$.
Moreover (1) tells us these two categories are equivalent under the
completion functor (\ref{equation-completion-functor-proper-over-A}).

\medskip\noindent
Applying this equivalence we get a coherent $\mathcal{O}_X$-module
$\mathcal{G}_e$ annihilated by $\mathcal{K}^e$ corresponding to the system
$(\mathcal{F}_n/\mathcal{K}^e\mathcal{F}_n)$ of
$\textit{Coh}_{\text{support proper over } A}(X, \mathcal{I})$. The maps
$\mathcal{F}_n/\mathcal{K}^{e + 1}\mathcal{F}_n \to
\mathcal{F}_n/\mathcal{K}^e\mathcal{F}_n$ correspond to canonical maps
$\mathcal{G}_{e + 1} \to \mathcal{G}_e$ which induce isomorphisms
$\mathcal{G}_{e + 1}/\mathcal{K}^e\mathcal{G}_{e + 1} \to \mathcal{G}_e$.
We obtain an object $(\mathcal{G}_e)$ of the category
$\textit{Coh}_{\text{support proper over } A}(X, \mathcal{K})$.
The map $\alpha$ induces a system of maps
$$
\mathcal{F}_n/\mathcal{K}^e\mathcal{F}_n
\longrightarrow
\mathcal{H}/(\mathcal{I}^n + \mathcal{K}^e)\mathcal{H}
$$
whence maps $\mathcal{G}_e \to \mathcal{H}/\mathcal{K}^e\mathcal{H}$
(by the equivalence of categories again).
Let $t \geq 1$ be an integer, which exists by assumption (2),
such that $\mathcal{K}^t$ annihilates the kernel and cokernel of all the maps
$\mathcal{F}_n \to \mathcal{H}/\mathcal{I}^n\mathcal{H}$.
Then $\mathcal{K}^{2t}$ annihilates the kernel and cokernel of the maps
$\mathcal{F}_n/\mathcal{K}^e\mathcal{F}_n \to
\mathcal{H}/(\mathcal{I}^n + \mathcal{K}^e)\mathcal{H}$
(details omitted; see Cohomology of Schemes,
Remark \ref{coherent-remark-inverse-systems-kernel-cokernel-annihilated-by}).
Whereupon we conclude that $\mathcal{K}^{4t}$ annihilates the kernel and
the cokernel of the maps
$$
\mathcal{G}_e
\longrightarrow
\mathcal{H}/\mathcal{K}^e\mathcal{H},
$$
(details omitted;  see Cohomology of Schemes,
Remark \ref{coherent-remark-inverse-systems-kernel-cokernel-annihilated-by}).
We apply Lemma \ref{lemma-existence-easy} to obtain a coherent
$\mathcal{O}_X$-module $\mathcal{F}$, a map
$a : \mathcal{F} \to \mathcal{H}$ and an isomorphism
$\beta : (\mathcal{G}_e) \to (\mathcal{F}/\mathcal{K}^e\mathcal{F})$
in $\textit{Coh}(X, \mathcal{K})$. Working backwards, for a given $n$
the triple
$(\mathcal{F}/\mathcal{I}^n\mathcal{F}, a \bmod \mathcal{I}^n, \beta
\bmod \mathcal{I}^n)$ is a triple as in the lemma for the morphism
$\alpha_n \bmod \mathcal{K}^e :
(\mathcal{F}_n/\mathcal{K}^e\mathcal{F}_n) \to
(\mathcal{H}/(\mathcal{I}^n + \mathcal{K}^e)\mathcal{H})$
of $\textit{Coh}(X, \mathcal{K})$. Thus the uniqueness in
Lemma \ref{lemma-existence-easy}
gives a canonical isomorphism
$\mathcal{F}/\mathcal{I}^n\mathcal{F} \to \mathcal{F}_n$
compatible with all the morphisms in sight.

\medskip\noindent
To finish the proof of the lemma we still have to show that the
support of $\mathcal{F}$ is proper over $A$.
By construction the kernel of $a : \mathcal{F} \to \mathcal{H}$
is annihilated by a power of $\mathcal{K}$. Hence the support of
this kernel is contained in the support of $\mathcal{G}_1$. Since
$\mathcal{G}_1$ is an object of
$\textit{Coh}_{\text{support proper over } A}(\mathcal{O}_{X_1})$
we see this is proper over $A$. Combined with the fact that the
support of $\mathcal{H}$ is proper over $A$ we conclude that the
support of $\mathcal{F}$ is proper over $A$ by
Derived Categories of Spaces, Lemma
\ref{spaces-perfect-lemma-union-closed-proper-over-base}.
\end{proof}

\begin{lemma}
\label{lemma-inverse-systems-push-pull}
Let $S$ be a scheme. Let $f : X \to Y$ be a representable
proper morphism of Noetherian algebraic spaces over $S$. Let
$\mathcal{J}, \mathcal{K} \subset \mathcal{O}_Y$
be quasi-coherent sheaves of ideals.
Assume $f$ is an isomorphism over $V = Y \setminus V(\mathcal{K})$.
Set $\mathcal{I} = f^{-1}\mathcal{J} \mathcal{O}_X$.
Let $(\mathcal{G}_n)$ be an object of $\textit{Coh}(Y, \mathcal{J})$,
let $\mathcal{F}$ be a coherent $\mathcal{O}_X$-module, and let
$\beta : (f^*\mathcal{G}_n)  \to \mathcal{F}^\wedge$ be an isomorphism in
$\textit{Coh}(X, \mathcal{I})$. Then there exists a map
$$
\alpha :
(\mathcal{G}_n)
\longrightarrow
(f_*\mathcal{F})^\wedge
$$
in $\textit{Coh}(Y, \mathcal{J})$ whose kernel and cokernel
are annihilated by a power of $\mathcal{K}$.
\end{lemma}

\begin{proof}
Since $f$ is a proper morphism we see that $f_*\mathcal{F}$ is a coherent
$\mathcal{O}_Y$-module (Cohomology of Spaces, Lemma
\ref{spaces-cohomology-lemma-proper-pushforward-coherent}).
Thus the statement of the lemma makes sense. Consider the compositions
$$
\gamma_n : \mathcal{G}_n \to
f_*f^*\mathcal{G}_n \to
f_*(\mathcal{F}/\mathcal{I}^n\mathcal{F}).
$$
Here the first map is the adjunction map and the second is $f_*\beta_n$.
We claim that there exists a unique $\alpha$ as in the lemma
such that the compositions
$$
\mathcal{G}_n \xrightarrow{\alpha_n}
f_*\mathcal{F}/\mathcal{J}^nf_*\mathcal{F} \to
f_*(\mathcal{F}/\mathcal{I}^n\mathcal{F})
$$
equal $\gamma_n$ for all $n$. Because of the uniqueness and \'etale
descent for $\textit{Coh}(Y, \mathcal{J})$ it suffices
to prove this \'etale locally on $Y$. Thus we may assume $Y$
is the spectrum of a Noetherian ring. As $f$ is representable
we see that $X$ is a scheme as well. Thus we reduce to the case of
schemes, see proof of
Cohomology of Schemes, Lemma \ref{coherent-lemma-inverse-systems-push-pull}.
\end{proof}

\begin{theorem}[Grothendieck's existence theorem]
\label{theorem-grothendieck-existence}
In Situation \ref{situation-existence} the functor
(\ref{equation-completion-functor-proper-over-A})
is an equivalence.
\end{theorem}

\begin{proof}
We will use the equivalence of categories of
Cohomology of Spaces, Lemma \ref{spaces-cohomology-lemma-i-star-equivalence}
without further mention in the proof of the theorem.
By Lemma \ref{lemma-fully-faithful} the functor is fully faithful.
Thus we need to prove the functor is essentially surjective.

\medskip\noindent
Consider the collection $\Xi$ of quasi-coherent sheaves of ideals
$\mathcal{K} \subset \mathcal{O}_X$ such that the statement holds
for every object $(\mathcal{F}_n)$ of
$\textit{Coh}_{\text{support proper over }A}(X, \mathcal{I})$
annihilated by $\mathcal{K}$. We want to show $(0)$ is in $\Xi$.
If not, then since $X$ is Noetherian there exists a maximal
quasi-coherent sheaf of ideals $\mathcal{K}$ not in $\Xi$, see
Cohomology of Spaces, Lemma \ref{spaces-cohomology-lemma-acc-coherent}.
After replacing $X$ by the closed subscheme of $X$
corresponding to $\mathcal{K}$ we may assume that every nonzero
$\mathcal{K}$ is in $\Xi$. Let $(\mathcal{F}_n)$ be an object of
$\textit{Coh}_{\text{support proper over }A}(X, \mathcal{I})$.
We will show that this object is in the essential image, thereby
completing the proof of the theorem.

\medskip\noindent
Apply Chow's lemma (Lemma \ref{lemma-chow-noetherian-separated})
to find a proper surjective morphism $f : Y \to X$ which is an isomorphism
over a dense open $U \subset X$ such that $Y$ is H-quasi-projective
over $A$. Note that $Y$ is a scheme and $f$ representable.
Choose an open immersion $j : Y \to Y'$ with $Y'$ projective over $A$, see
Morphisms, Lemma \ref{morphisms-lemma-H-quasi-projective-open-H-projective}.
Let $T_n$ be the scheme theoretic support of $\mathcal{F}_n$.
Note that $|T_n| = |T_1|$, hence $T_n$ is proper over $A$ for all $n$
(Morphisms of Spaces, Lemma
\ref{spaces-morphisms-lemma-image-proper-is-proper}).
Then $f^*\mathcal{F}_n$ is supported on the closed subscheme
$f^{-1}T_n$ which is proper over $A$ (by
Morphisms of Spaces, Lemma \ref{spaces-morphisms-lemma-composition-proper}
and properness of $f$).
In particular, the composition $f^{-1}T_n \to Y \to Y'$ is closed
(Morphisms, Lemma \ref{morphisms-lemma-image-proper-scheme-closed}).
Let $T'_n \subset Y'$ be the corresponding closed subscheme;
it is contained in the open subscheme $Y$ and equal to $f^{-1}T_n$
as a closed subscheme of $Y$. Let $\mathcal{F}_n'$
be the coherent $\mathcal{O}_{Y'}$-module corresponding to
$f^*\mathcal{F}_n$ viewed as a coherent module on $Y'$ via
the closed immersion $f^{-1}T_n = T'_n \subset Y'$.
Then $(\mathcal{F}_n')$
is an object of $\textit{Coh}(Y', I\mathcal{O}_{Y'})$.
By the projective case of Grothendieck's existence theorem
(Cohomology of Schemes, Lemma \ref{coherent-lemma-existence-projective})
there exists a coherent $\mathcal{O}_{Y'}$-module
$\mathcal{F}'$ and an isomorphism
$(\mathcal{F}')^\wedge \cong (\mathcal{F}'_n)$ in
$\textit{Coh}(Y', I\mathcal{O}_{Y'})$.
Let $Z' \subset Y'$ be the scheme theoretic support of $\mathcal{F}'$.
Since $\mathcal{F}'/I\mathcal{F}' = \mathcal{F}'_1$ we see
that $Z' \cap V(I\mathcal{O}_{Y'}) = T'_1$ set-theoretically.
The structure morphism $p' : Y' \to \Spec(A)$ is proper, hence
$p'(Z' \cap (Y' \setminus Y))$ is closed in $\Spec(A)$.
If nonempty, then it would contain a point of $V(I)$
as $I$ is contained in the Jacobson radical of $A$
(Algebra, Lemma \ref{algebra-lemma-radical-completion}).
But we've seen above that $Z' \cap (p')^{-1}V(I) = T'_1 \subset Y$
hence we conclude that $Z' \subset Y$. Thus $\mathcal{F}'|_Y$
is supported on a closed subscheme of $Y$ proper over $A$.

\medskip\noindent
Let $\mathcal{K}$ be the quasi-coherent sheaf of ideals cutting
out the reduced complement $X \setminus U$.
By Cohomology of Spaces, Lemma
\ref{spaces-cohomology-lemma-proper-pushforward-coherent}
the $\mathcal{O}_X$-module $\mathcal{H} = f_*\mathcal{F}'$ is coherent
and by Lemma \ref{lemma-inverse-systems-push-pull}
there exists a morphism $\alpha : (\mathcal{F}_n) \to \mathcal{H}^\wedge$
in the category $\textit{Coh}_{\text{support proper over } A}(X, \mathcal{I})$
whose kernel and cokernel are annihilated by a power of $\mathcal{K}$.
Let $Z_0 \subset X$ be the scheme theoretic support of $\mathcal{H}$.
It is clear that $|Z_0| \subset f(|Z'|)$. Hence
$Z_0 \to \Spec(A)$ is proper
(Morphisms of Spaces, Lemma
\ref{spaces-morphisms-lemma-image-proper-is-proper}).
Thus $\mathcal{H}$ is an object of 
$\textit{Coh}_{\text{support proper over } A}(\mathcal{O}_X)$.
Since each of the sheaves of ideals $\mathcal{K}^e$ is an element of
$\Xi$ we see that the assumptions of Lemma \ref{lemma-existence-tricky}
are satisfied and we conclude.
\end{proof}

\begin{remark}[Unwinding Grothendieck's existence theorem]
\label{remark-reformulate-existence-theorem}
Let $A$ be a Noetherian ring complete with respect to an ideal $I$.
Write $S = \Spec(A)$ and $S_n = \Spec(A/I^n)$.
Let $X \to S$ be a morphism of algebraic spaces that is separated and
of finite type. For $n \geq 1$ we set $X_n = X \times_S S_n$.
Picture:
$$
\xymatrix{
X_1 \ar[r]_{i_1} \ar[d] & X_2 \ar[r]_{i_2} \ar[d] & X_3 \ar[r] \ar[d] &
\ldots & X \ar[d] \\
S_1 \ar[r] & S_2 \ar[r] & S_3 \ar[r] & \ldots & S
}
$$
In this situation we consider systems $(\mathcal{F}_n, \varphi_n)$
where
\begin{enumerate}
\item $\mathcal{F}_n$ is a coherent $\mathcal{O}_{X_n}$-module,
\item $\varphi_n : i_n^*\mathcal{F}_{n + 1} \to \mathcal{F}_n$
is an isomorphism, and
\item $\text{Supp}(\mathcal{F}_1)$ is proper over $S_1$.
\end{enumerate}
Theorem \ref{theorem-grothendieck-existence} says that the
completion functor
$$
\begin{matrix}
\text{coherent }\mathcal{O}_X\text{-modules }\mathcal{F} \\
\text{with support proper over }A
\end{matrix}
\quad
\longrightarrow
\quad
\begin{matrix}
\text{systems }(\mathcal{F}_n) \\
\text{as above}
\end{matrix}
$$
is an equivalence of categories. In the special case that $X$ is
proper over $A$ we can omit the conditions on the supports.
\end{remark}







\section{Grothendieck's algebraization theorem}
\label{section-algebraization}

\noindent
This section is the analogue of
Cohomology of Schemes, Section \ref{coherent-section-algebraization}.
However, this section is missing the result on algebraization of deformations
of proper algebraic spaces endowed with ample invertible sheaves, as a
proper algebraic space which comes with an ample invertible sheaf is
already a scheme. We do have an algebraization result on proper
algebraic spaces of relative dimension $1$.
Our first result is a translation of Grothendieck's existence
theorem in terms of closed subschemes and finite morphisms.

\begin{lemma}
\label{lemma-algebraize-formal-closed-subscheme}
Let $A$ be a Noetherian ring complete with respect to an ideal $I$.
Write $S = \Spec(A)$ and $S_n = \Spec(A/I^n)$.
Let $X \to S$ be a morphism of algebraic spaces that is separated
and of finite type.
For $n \geq 1$ we set $X_n = X \times_S S_n$.
Suppose given a commutative diagram
$$
\xymatrix{
Z_1 \ar[r] \ar[d] & Z_2 \ar[r] \ar[d] & Z_3 \ar[r] \ar[d] & \ldots \\
X_1 \ar[r]^{i_1} & X_2 \ar[r]^{i_2} & X_3 \ar[r] & \ldots
}
$$
of algebraic spaces with cartesian squares. Assume that
\begin{enumerate}
\item $Z_1 \to X_1$ is a closed immersion, and
\item $Z_1 \to S_1$ is proper.
\end{enumerate}
Then there exists a closed immersion of algebraic spaces $Z \to X$ such that
$Z_n = Z \times_S S_n$ for all $n \geq 1$. Moreover, $Z$ is proper over $S$.
\end{lemma}

\begin{proof}
Let's write $j_n : Z_n \to X_n$ for the vertical morphisms.
As the squares in the statement are cartesian
we see that the base change of $j_n$ to $X_1$ is $j_1$.
Thus Limits of Spaces, Lemma
\ref{spaces-limits-lemma-check-closed-infinitesimally}
shows that $j_n$ is a closed immersion.
Set $\mathcal{F}_n = j_{n, *}\mathcal{O}_{Z_n}$, so that
$j_n^\sharp$ is a surjection $\mathcal{O}_{X_n} \to \mathcal{F}_n$.
Again using that the squares are cartesian we see that
the pullback of $\mathcal{F}_{n + 1}$ to $X_n$ is $\mathcal{F}_n$.
Hence Grothendieck's existence theorem, as reformulated in
Remark \ref{remark-reformulate-existence-theorem},
tells us there exists a map
$\mathcal{O}_X \to \mathcal{F}$
of coherent $\mathcal{O}_X$-modules whose restriction to
$X_n$ recovers $\mathcal{O}_{X_n} \to \mathcal{F}_n$.
Moreover, the support of $\mathcal{F}$ is proper over $S$.
As the completion functor is exact (Lemma \ref{lemma-exact})
we see that $\mathcal{O}_X \to \mathcal{F}$
is surjective. Thus $\mathcal{F} = \mathcal{O}_X/\mathcal{J}$
for some quasi-coherent sheaf of ideals $\mathcal{J}$.
Setting $Z = V(\mathcal{J})$ finishes the proof.
\end{proof}

\begin{lemma}
\label{lemma-algebraize-formal-algebraic-space-finite-over-proper}
Let $A$ be a Noetherian ring complete with respect to an ideal $I$.
Write $S = \Spec(A)$ and $S_n = \Spec(A/I^n)$.
Let $X \to S$ be a morphism of algebraic spaces that is separated
and of finite type.
For $n \geq 1$ we set $X_n = X \times_S S_n$.
Suppose given a commutative diagram
$$
\xymatrix{
Y_1 \ar[r] \ar[d] & Y_2 \ar[r] \ar[d] & Y_3 \ar[r] \ar[d] & \ldots \\
X_1 \ar[r]^{i_1} & X_2 \ar[r]^{i_2} & X_3 \ar[r] & \ldots
}
$$
of algebraic spaces with cartesian squares. Assume that
\begin{enumerate}
\item $Y_1 \to X_1$ is a finite morphism, and
\item $Y_1 \to S_1$ is proper.
\end{enumerate}
Then there exists a finite morphism of algebraic spaces $Y \to X$ such that
$Y_n = Y \times_S S_n$ for all $n \geq 1$. Moreover, $Y$ is proper over $S$.
\end{lemma}

\begin{proof}
Let's write $f_n : Y_n \to X_n$ for the vertical morphisms.
As the squares in the statement are cartesian
we see that the base change of $f_n$ to $X_1$ is $f_1$.
Thus Lemma \ref{lemma-thicken-property-morphisms-cartesian}
shows that $f_n$ is a finite morphism.
Set $\mathcal{F}_n = f_{n, *}\mathcal{O}_{Y_n}$.
Using that the squares are cartesian we see that
the pullback of $\mathcal{F}_{n + 1}$ to $X_n$ is $\mathcal{F}_n$.
Hence Grothendieck's existence theorem, as reformulated in
Remark \ref{remark-reformulate-existence-theorem},
tells us there exists a coherent $\mathcal{O}_X$-module $\mathcal{F}$
whose restriction to $X_n$ recovers $\mathcal{F}_n$.
Moreover, the support of $\mathcal{F}$ is proper over $S$.
As the completion functor is fully faithful
(Theorem \ref{theorem-grothendieck-existence})
we see that the multiplication maps
$\mathcal{F}_n \otimes_{\mathcal{O}_{X_n}} \mathcal{F}_n \to
\mathcal{F}_n$ fit together to give an algebra structure on $\mathcal{F}$.
Setting $Y = \underline{\Spec}_X(\mathcal{F})$ finishes the proof.
\end{proof}

\begin{lemma}
\label{lemma-algebraize-morphism}
Let $A$ be a Noetherian ring complete with respect to an ideal $I$.
Write $S = \Spec(A)$ and $S_n = \Spec(A/I^n)$. Let $X$, $Y$ be algebraic
spaces over $S$. For $n \geq 1$ we set $X_n = X \times_S S_n$ and
$Y_n = Y \times_S S_n$. Suppose given a compatible system of
commutative diagrams
$$
\xymatrix{
& & X_{n + 1} \ar[rd] \ar[rr]_{g_{n + 1}} & & Y_{n + 1} \ar[ld] \\
X_n \ar[rru] \ar[rd] \ar[rr]_{g_n} & & Y_n \ar[rru] \ar[ld] & S_{n + 1} \\
& S_n \ar[rru]
}
$$
Assume that
\begin{enumerate}
\item $X \to S$ is proper, and
\item $Y \to S$ is separated of finite type.
\end{enumerate}
Then there exists a unique morphism of algebraic spaces $g : X \to Y$
over $S$ such that $g_n$ is the base change of $g$ to $S_n$.
\end{lemma}

\begin{proof}
The morphisms $(1, g_n) : X_n \to X_n \times_S Y_n$ are closed immersions
because $Y_n \to S_n$ is separated
(Morphisms of Spaces, Lemma \ref{spaces-morphisms-lemma-section-immersion}).
Thus by Lemma \ref{lemma-algebraize-formal-closed-subscheme}
there exists a closed subspace $Z \subset X \times_S Y$
proper over $S$ whose base change to $S_n$ recovers
$X_n \subset X_n \times_S Y_n$. The first projection $p : Z \to X$
is a proper morphism (as $Z$ is proper over $S$, see
Morphisms of Spaces, Lemma
\ref{spaces-morphisms-lemma-universally-closed-permanence})
whose base change to $S_n$ is an isomorphism for all $n$.
In particular, $p : Z \to X$ is quasi-finite on an open subspace
of $Z$ containing every point of $Z_0$ for example by
Morphisms of Spaces, Lemma
\ref{spaces-morphisms-lemma-locally-finite-type-quasi-finite-part}.
As $Z$ is proper over $S$ this open neighbourhood is all of $Z$.
We conclude that $p : Z \to X$ is finite by Zariski's main theorem
(for example apply
Lemma \ref{lemma-quasi-finite-separated-pass-through-finite}
and use properness of $Z$ over $X$ to see that the immersion is
a closed immersion). Applying the equivalence of
Theorem \ref{theorem-grothendieck-existence}
we see that $p_*\mathcal{O}_Z = \mathcal{O}_X$ as this is true
modulo $I^n$ for all $n$. Hence $p$ is an isomorphism and we obtain
the morphism $g$ as the composition $X \cong Z \to Y$.
We omit the proof of uniqueness.
\end{proof}

\begin{remark}
\label{remark-weaken-separation-axioms-question}
We can ask if in Grothendieck's algebraization theorem (in the form
of Lemma \ref{lemma-algebraize-morphism}), we can get
by with weaker separation axioms on the target. Let us be more precise.
Let $A$, $I$, $S$, $S_n$, $X$, $Y$, $X_n$, $Y_n$, and $g_n$ be as in
the statement of Lemma \ref{lemma-algebraize-morphism} and assume that
\begin{enumerate}
\item $X \to S$ is proper, and
\item $Y \to S$ is locally of finite type.
\end{enumerate}
Does there exist a morphism of algebraic spaces $g : X \to Y$
over $S$ such that $g_n$ is the base change of $g$ to $S_n$?
We don't know the answer in general; if you do please email
\href{mailto:stacks.project@gmail.com}{stacks.project@gmail.com}.
If $Y \to S$ is separated, then the result
holds by the lemma (there is an immediate reduction to the case
where $X$ is finite type over $S$, by choosing a quasi-compact
open containing the image of $g_1$). If we only assume $Y \to S$
is quasi-separated, then the result is true as well. First, as before
we may assume $Y$ is quasi-compact as well as quasi-separated.
Then we can use either \cite{Bhatt-Algebraize} or from
\cite{Hall-Rydh-coherent} to algebraize $(g_n)$. Namely, to apply
the first reference, we use
$$
D_{perf}(X) \to \lim D_{perf}(X_n)
\xrightarrow{\lim Lg_n^*}
\lim D_{perf}(Y_n) = D_{perf}(Y)
$$
where the last step uses a Grothendieck existence result for
the derived category of the proper algebraic space $Y$ over $R$
(compare with
Flatness on Spaces, Remark \ref{spaces-flat-remark-correct-generality}).
The paper cited shows that
this arrow determines a morphism $Y \to X$ as desired. To apply the
second reference we use the same argument with coherent modules:
$$
\textit{Coh}(\mathcal{O}_X) \to \lim \textit{Coh}(\mathcal{O}_{X_n})
\xrightarrow{\lim g_n^*}
\lim \textit{Coh}(\mathcal{O}_{Y_n}) = \textit{Coh}(\mathcal{O}_Y)
$$
where the final equality is a consequence of Grothendieck's
existence theorem (Theorem \ref{theorem-grothendieck-existence}).
The second reference tells us that this functor corresponds to a
morphism $Y \to X$ over $R$. If we ever need this generalization
we will precisely state and carefully prove the result here.
\end{remark}

\begin{lemma}
\label{lemma-formal-algebraic-space-proper-reldim-1}
Let $(A, \mathfrak m, \kappa)$ be a complete local Noetherian ring.
Set $S = \Spec(A)$ and $S_n = \Spec(A/\mathfrak m^n)$.
Consider a commutative diagram
$$
\xymatrix{
X_1 \ar[r]_{i_1} \ar[d] & X_2 \ar[r]_{i_2} \ar[d] & X_3 \ar[r] \ar[d] &
\ldots \\
S_1 \ar[r] & S_2 \ar[r] & S_3 \ar[r] & \ldots
}
$$
of algebraic spaces with cartesian squares. If $\dim(X_1) \leq 1$,
then there exists a projective morphism of schemes $X \to S$
and isomorphisms $X_n \cong X \times_S S_n$ compatible with $i_n$.
\end{lemma}

\begin{proof}
By Spaces over Fields, Lemma
\ref{spaces-over-fields-lemma-codim-1-point-in-schematic-locus}
the algebraic space $X_1$ is a scheme. Hence $X_1$
is a proper scheme of dimension $\leq 1$ over $\kappa$.
By Varieties, Lemma \ref{varieties-lemma-dim-1-proper-projective}
we see that $X_1$ is H-projective over $\kappa$.
Let $\mathcal{L}_1$ be an ample invertible sheaf on $X_1$.

\medskip\noindent
We are going to show that $\mathcal{L}_1$ lifts to a compatible system
$\{\mathcal{L}_n\}$ of invertible sheaves on $\{X_n\}$.
Observe that $X_n$ is a scheme too by Lemma \ref{lemma-thickening-scheme}.
Recall that $X_1 \to X_n$ induces homeomorphisms of underlying
topological spaces. In the rest of the proof we do not distinguish
between sheaves on $X_n$ and sheaves on $X_1$.
Suppose, given a lift $\mathcal{L}_n$ to $X_n$. We consider
the exact sequence
$$
1 \to
(1 + \mathfrak m^n\mathcal{O}_{X_{n + 1}})^* \to
\mathcal{O}_{X_{n + 1}}^* \to \mathcal{O}_{X_n}^* \to 1
$$
of sheaves on $X_{n + 1}$. The class of $\mathcal{L}_n$ in
$H^1(X_n, \mathcal{O}_{X_n}^*)$ (see
Cohomology, Lemma \ref{cohomology-lemma-h1-invertible})
can be lifted to an element of $H^1(X_{n + 1}, \mathcal{O}_{X_{n + 1}}^*)$
if and only if the obstruction in
$H^2(X_{n + 1}, (1 + \mathfrak m^n\mathcal{O}_{X_{n + 1}})^*)$
is zero. As $X_1$ is a Noetherian scheme of dimension $\leq 1$
this cohomology group vanishes (Cohomology, Proposition
\ref{cohomology-proposition-vanishing-Noetherian}).

\medskip\noindent
By Grothendieck's algebraization theorem
(Cohomology of Schemes, Theorem \ref{coherent-theorem-algebraization})
we find a projective morphism of schemes $X \to S = \Spec(A)$
and a compatible system of isomorphisms $X_n = S_n \times_S X$.
\end{proof}

\begin{lemma}
\label{lemma-projective-over-complete}
Let $(A, \mathfrak m, \kappa)$ be a complete Noetherian local ring.
Let $X$ be an algebraic space over $\Spec(A)$.
If $X \to \Spec(A)$ is proper and $\dim(X_\kappa) \leq 1$, then
$X$ is a scheme projective over $A$.
\end{lemma}

\begin{proof}
Set $X_n = X \times_{\Spec(A)} \Spec(A/\mathfrak m^n)$.
By Lemma \ref{lemma-formal-algebraic-space-proper-reldim-1}
there exists a projective morphism $Y \to \Spec(A)$
and compatible isomorphisms
$Y \times_{\Spec(A)} \Spec(A/\mathfrak m^n) \cong
X \times_{\Spec(A)} \Spec(A/\mathfrak m^n)$.
By Lemma \ref{lemma-algebraize-morphism}
we see that $X \cong Y$ and the proof is complete.
\end{proof}










\section{Regular immersions}
\label{section-regular-immersions}

\noindent
This section is the analogue of
Divisors, Section \ref{divisors-section-regular-immersions}
for morphisms of algebraic spaces. The reader is encouraged to read up
on regular immersions of schemes in that section first.

\medskip\noindent
In
Divisors, Section \ref{divisors-section-regular-immersions}
we defined four types of regular immersions for morphisms of schemes.
Of these only three are (as far as we know) local on the target for
the \'etale topology; as usual plain old regular immersions aren't.
This is why for morphisms of algebraic spaces we cannot actually define
regular immersions. (These kinds of annoyances prompted Grothendieck
and his school to replace original notion of a regular immersion by a
Koszul-regular immersions, see
\cite[Exposee VII, Definition 1.4]{SGA6}.)
But we can define Koszul-regular, $H_1$-regular, and quasi-regular immersions.
Another remark is that since Koszul-regular immersions are not preserved by
arbitrary base change, we cannot use the strategy of
Morphisms of Spaces, Section \ref{spaces-morphisms-section-representable}
to define them. Similarly, as Koszul-regular immersions are not \'etale local
on the source, we cannot use
Morphisms of Spaces, Lemma \ref{spaces-morphisms-lemma-local-source-target}
to define them either. We replace this lemma instead by the
following.

\begin{lemma}
\label{lemma-representable-etale-local-target}
Let $\mathcal{P}$ be a property of morphisms of schemes which is \'etale
local on the target. Let $S$ be a scheme.
Let $f : X \to Y$ be a representable morphism of algebraic spaces over $S$.
Consider commutative diagrams
$$
\xymatrix{
X \times_Y V \ar[d] \ar[r] & V \ar[d] \\
X \ar[r]^f & Y
}
$$
where $V$ is a scheme and $V \to Y$ is \'etale.
The following are equivalent
\begin{enumerate}
\item for any diagram as above the projection $X \times_Y V \to V$
has property $\mathcal{P}$, and
\item for some diagram as above with $V \to Y$ surjective
the projection $X \times_Y V \to V$ has property $\mathcal{P}$.
\end{enumerate}
If $X$ and $Y$ are representable, then this is also equivalent to
$f$ (as a morphism of schemes) having property $\mathcal{P}$.
\end{lemma}

\begin{proof}
Let us prove the equivalence of (1) and (2).
The implication (1) $\Rightarrow$ (2) is immediate.
Assume
$$
\xymatrix{
X \times_Y V \ar[d] \ar[r] & V \ar[d] \\
X \ar[r]^f & Y
}
\quad\quad
\xymatrix{
X \times_Y V' \ar[d] \ar[r] & V' \ar[d] \\
X \ar[r]^f & Y
}
$$
are two diagrams as in the lemma. Assume $V \to Y$ is
surjective and $X \times_Y V \to V$ has property $\mathcal{P}$.
To show that (2) implies (1) we have to prove that
$X \times_Y V' \to V'$ has $\mathcal{P}$. To do
this consider the diagram
$$
\xymatrix{
X \times_Y V \ar[d] &
(X \times_Y V) \times_X (X \times_Y V') \ar[l] \ar[d] \ar[r] &
X \times_Y V' \ar[d] \\
V &
V \times_Y V' \ar[l] \ar[r] &
V'
}
$$
By our assumption that $\mathcal{P}$ is \'etale local on the source,
we see that $\mathcal{P}$ is preserved under \'etale base change, see
Descent, Lemma \ref{descent-lemma-pullback-property-local-target}.
Hence if the left vertical arrow has $\mathcal{P}$ the so does
the middle vertical arrow. Since $U \times_X U' \to U'$ is surjective
and \'etale (hence defines an \'etale covering of $U'$)
this implies (as $\mathcal{P}$ is assumed local for the \'etale topology
on the target) that the left vertical arrow has $\mathcal{P}$.

\medskip\noindent
If $X$ and $Y$ are representable, then we can take
$\text{id}_Y : Y \to Y$ as our \'etale covering to see the
final statement of the lemma is true.
\end{proof}

\noindent
Note that ``being a Koszul-regular (resp.\ $H_1$-regular, resp.\ quasi-regular)
immersion'' is a property of morphisms of schemes which is fpqc local on the
target, see
Descent, Lemma \ref{descent-lemma-descending-property-regular-immersion}.
Hence the following definition now makes sense.

\begin{definition}
\label{definition-regular-immersion}
Let $S$ be a scheme. Let $i : X \to Y$ be a morphism of algebraic
spaces over $S$.
\begin{enumerate}
\item We say $i$ is a {\it Koszul-regular immersion} if $i$ is representable
and the equivalent conditions of
Lemma \ref{lemma-representable-etale-local-target}
hold with $\mathcal{P}(f) =$``$f$ is a Koszul-regular immersion''.
\item We say $i$ is an {\it $H_1$-regular immersion} if $i$ is representable
and the equivalent conditions of
Lemma \ref{lemma-representable-etale-local-target}
hold with $\mathcal{P}(f) =$``$f$ is an $H_1$-regular immersion''.
\item We say $i$ is a {\it quasi-regular immersion} if $i$ is representable
and the equivalent conditions of
Lemma \ref{lemma-representable-etale-local-target}
hold with $\mathcal{P}(f) =$``$f$ is a quasi-regular immersion''.
\end{enumerate}
\end{definition}

\begin{lemma}
\label{lemma-regular-quasi-regular-immersion}
Let $S$ be a scheme.
Let $i : Z \to X$ be an immersion of algebraic spaces over $S$.
We have the following implications:
$i$ is Koszul-regular $\Rightarrow$
$i$ is $H_1$-regular $\Rightarrow$
$i$ is quasi-regular.
\end{lemma}

\begin{proof}
Via the definition this lemma immediately reduces to
Divisors, Lemma \ref{divisors-lemma-regular-quasi-regular-immersion}.
\end{proof}

\begin{lemma}
\label{lemma-regular-immersion-noetherian}
Let $S$ be a scheme.
Let $i : Z \to X$ be an immersion of algebraic spaces over $S$.
Assume $X$ is locally Noetherian. Then
$i$ is Koszul-regular $\Leftrightarrow$
$i$ is $H_1$-regular $\Leftrightarrow$
$i$ is quasi-regular.
\end{lemma}

\begin{proof}
Via Definition \ref{definition-regular-immersion}
(and the definition of a locally Noetherian algebraic space
in Properties of Spaces, Section
\ref{spaces-properties-section-types-properties})
this immediately translates to the case of schemes which is
Divisors, Lemma \ref{divisors-lemma-regular-immersion-noetherian}.
\end{proof}

\begin{lemma}
\label{lemma-flat-base-change-regular-immersion}
\begin{slogan}
Regular immersions are stable under flat base change.
\end{slogan}
Let $S$ be a scheme. Let $i : Z \to X$ be a Koszul-regular,
$H_1$-regular, or quasi-regular immersion of algebraic spaces over $S$.
Let $X' \to X$ be a flat morphism of algebraic spaces over $S$.
Then the base change $i' : Z \times_X X' \to X'$ is a Koszul-regular,
$H_1$-regular, or quasi-regular immersion.
\end{lemma}

\begin{proof}
Via Definition \ref{definition-regular-immersion}
(and the definition of a flat morphism of algebraic spaces
in Morphisms of Spaces, Section
\ref{spaces-morphisms-section-flat})
this lemma reduces to the case of schemes, see
Divisors, Lemma \ref{divisors-lemma-flat-base-change-regular-immersion}.
\end{proof}

\begin{lemma}
\label{lemma-quasi-regular-immersion}
Let $S$ be a scheme. Let $i : Z \to X$ be an immersion of algebraic spaces
over $S$. Then $i$ is a quasi-regular immersion if and only if the following
conditions are satisfied
\begin{enumerate}
\item $i$ is locally of finite presentation,
\item the conormal sheaf $\mathcal{C}_{Z/X}$ is finite locally free, and
\item the map (\ref{equation-conormal-algebra-quotient}) is an isomorphism.
\end{enumerate}
\end{lemma}

\begin{proof}
Follows from the case of schemes
(Divisors, Lemma \ref{divisors-lemma-quasi-regular-immersion})
via \'etale localization (use Definition \ref{definition-regular-immersion}
and
Lemma \ref{lemma-etale-conormal-algebra}).
\end{proof}

\begin{lemma}
\label{lemma-transitivity-conormal-quasi-regular}
Let $S$ be a scheme. Let $Z \to Y \to X$ be immersions of algebraic spaces
over $S$. Assume that $Z \to Y$ is $H_1$-regular. Then the canonical
sequence of Lemma \ref{lemma-transitivity-conormal}
$$
0 \to i^*\mathcal{C}_{Y/X} \to
\mathcal{C}_{Z/X} \to
\mathcal{C}_{Z/Y} \to 0
$$
is exact and (\'etale) locally split.
\end{lemma}

\begin{proof}
Since $\mathcal{C}_{Z/Y}$ is finite locally free (see
Lemma \ref{lemma-quasi-regular-immersion}
and
Lemma \ref{lemma-regular-quasi-regular-immersion})
it suffices to prove that the sequence is exact.
It suffices to show that the first map is injective
as the sequence is already right exact in general.
After \'etale localization on $X$ this reduces to the case
of schemes, see
Divisors, Lemma \ref{divisors-lemma-transitivity-conormal-quasi-regular}.
\end{proof}

\noindent
A composition of quasi-regular immersions may not be quasi-regular, see
Algebra, Remark \ref{algebra-remark-join-quasi-regular-sequences}.
The other types of regular immersions are preserved under composition.

\begin{lemma}
\label{lemma-composition-regular-immersion}
Let $S$ be a scheme. Let $i : Z \to Y$ and $j : Y \to X$ be immersions of
algebraic spaces over $S$.
\begin{enumerate}
\item If $i$ and $j$ are Koszul-regular immersions, so is $j \circ i$.
\item If $i$ and $j$ are $H_1$-regular immersions, so is $j \circ i$.
\item If $i$ is an $H_1$-regular immersion and $j$ is a quasi-regular
immersion, then $j \circ i$ is a quasi-regular immersion.
\end{enumerate}
\end{lemma}

\begin{proof}
Immediate from the case of schemes, see
Divisors, Lemma \ref{divisors-lemma-composition-regular-immersion}.
\end{proof}

\begin{lemma}
\label{lemma-permanence-regular-immersion}
Let $S$ be a scheme. Let $i : Z \to Y$ and $j : Y \to X$ be immersions of
algebraic spaces over $S$. Assume $j$ is locally of finite presentation
and that the sequence
$$
0 \to i^*\mathcal{C}_{Y/X} \to
\mathcal{C}_{Z/X} \to
\mathcal{C}_{Z/Y} \to 0
$$
of Lemma \ref{lemma-transitivity-conormal} is exact and locally split.
\begin{enumerate}
\item If $j \circ i$ is a quasi-regular immersion, so is $i$.
\item If $j \circ i$ is a $H_1$-regular immersion, so is $i$.
\item If both $j$ and $j \circ i$ are Koszul-regular immersions, so is $i$.
\end{enumerate}
\end{lemma}

\begin{proof}
Immediate from the case of schemes, see
Divisors, Lemma \ref{divisors-lemma-permanence-regular-immersion}.
\end{proof}

\begin{lemma}
\label{lemma-extra-permanence-regular-immersion-noetherian}
Let $S$ be a scheme. Let $i : Z \to Y$ and $j : Y \to X$ be immersions of
algebraic spaces over $S$. Assume $X$ is locally Noetherian.
The following are equivalent
\begin{enumerate}
\item $i$ and $j$ are Koszul regular immersions,
\item $i$ and $j \circ i$ are Koszul regular immersions,
\item $j \circ i$ is a Koszul regular immersion and the conormal sequence
$$
0 \to i^*\mathcal{C}_{Y/X} \to
\mathcal{C}_{Z/X} \to
\mathcal{C}_{Z/Y} \to 0
$$
is exact and locally split.
\end{enumerate}
\end{lemma}

\begin{proof}
Immediate from the case of schemes, see Divisors, Lemma
\ref{divisors-lemma-extra-permanence-regular-immersion-noetherian}.
\end{proof}














\section{Relative pseudo-coherence}
\label{section-relative-pseudo-coherence}

\noindent
This section is the analogue of
More on Morphisms, Section
\ref{more-morphisms-section-relative-pseudo-coherence}.
However, in the treatment of this material for
algebraic spaces we have decided to work exclusively
with objects in the derived category whose cohomology
sheaves are quasi-coherent. There are two reasons for this:
(1) it greatly simplifies the exposition and
(2) we currently have no use for the more general notion.

\begin{remark}
\label{remark-match-relative-pseudo-coherence}
Let $S$ be a scheme. Let $f : X \to Y$ be a morphism of representable
algebraic spaces over $S$ which is locally of finite type. Let
$f_0 : X_0 \to Y_0$ be a morphism of schemes representing $f$
(awkward but temporary notation). Then $f_0$ is locally of finite type.
If $E$ is an object of $D_\QCoh(\mathcal{O}_X)$, then $E$
is the pullback of a unique object $E_0$ in $D_\QCoh(\mathcal{O}_{X_0})$, see
Derived Categories of Spaces, Lemma
\ref{spaces-perfect-lemma-derived-quasi-coherent-small-etale-site}.
In this situation the phrase ``$E$ is $m$-pseudo-coherent relative to $Y$''
will be taken to mean ``$E_0$ is $m$-pseudo-coherent relative to $Y_0$''
as defined in More on Morphisms, Section
\ref{more-morphisms-section-relative-pseudo-coherence}.
\end{remark}

\begin{lemma}
\label{lemma-qcoh-relative-pseudo-coherence-characterize}
Let $S$ be a scheme. Let $f : X \to Y$ be a morphism of algebraic spaces
over $S$ which is locally of finite type. Let $m \in \mathbf{Z}$.
Let $E \in D_\QCoh(\mathcal{O}_X)$. With notation as explained in
Remark \ref{remark-match-relative-pseudo-coherence}
the following are equivalent:
\begin{enumerate}
\item for every commutative diagram
$$
\xymatrix{
U \ar[d] \ar[r] & V \ar[d] \\
X \ar[r] & Y
}
$$
where $U$, $V$ are schemes and the vertical arrows are \'etale, the complex
$E|_U$ is $m$-pseudo-coherent relative to $V$,
\item for some commutative diagram as in (1) with $U \to X$
surjective, the complex $E|_U$ is $m$-pseudo-coherent relative to $V$,
\item for every commutative diagram as in (1) with $U$ and $V$
affine the complex $R\Gamma(U, E)$ of $\mathcal{O}_X(U)$-modules
is $m$-pseudo-coherent relative to $\mathcal{O}_Y(V)$.
\end{enumerate}
\end{lemma}

\begin{proof}
Part (1) implies (3) by More on Morphisms, Lemma
\ref{more-morphisms-lemma-qcoh-relative-pseudo-coherence-characterize}.

\medskip\noindent
Assume (3). Pick any commutative diagram as in (1) with $U \to X$ surjective.
Choose an affine open covering $V = \bigcup V_j$ and affine open coverings
$(U \to V)^{-1}(V_j) = \bigcup U_{ij}$. By (3) and More on Morphisms, Lemma
\ref{more-morphisms-lemma-qcoh-relative-pseudo-coherence-characterize}
we see that $E|_U$ is $m$-pseudo-coherent relative to $V$.
Thus (3) implies (2).

\medskip\noindent
Assume (2). Choose a commutative diagram
$$
\xymatrix{
U \ar[d] \ar[r] & V \ar[d] \\
X \ar[r] & Y
}
$$
where $U$, $V$ are schemes, the vertical arrows are \'etale, the
morphism $U \to X$ is surjective, and $E|_U$ is $m$-pseudo-coherent
relative to $V$. Next, suppose given a second commutative diagram
$$
\xymatrix{
U' \ar[d] \ar[r] & V' \ar[d] \\
X \ar[r] & Y
}
$$
with \'etale vertical arrows and $U', V'$ schemes. We want to show
that $E|_{U'}$ is $m$-pseudo-coherent relative to $V'$.
The morphism $U'' = U \times_X U' \to U'$ is surjective \'etale
and $U'' \to V'$ factors through $V'' = V' \times_Y V$ which
is \'etale over $V'$. Hence it suffices to show that $E|_{U''}$
is $m$-pseudo-coherent relative to $V''$, see
More on Morphisms, Lemmas
\ref{more-morphisms-lemma-relative-pseudo-coherent-descends-fppf} and
\ref{more-morphisms-lemma-relative-pseudo-coherent-post-compose}.
Using the second lemma once more it suffices to show that
$E|_{U''}$ is $m$-pseudo-coherent relative to $V$.
This is true by More on Morphisms, Lemma
\ref{more-morphisms-lemma-pull-relative-pseudo-coherent}
and the fact that an \'etale morphism of schemes is pseudo-coherent by
More on Morphisms, Lemma
\ref{more-morphisms-lemma-flat-finite-presentation-pseudo-coherent}.
\end{proof}

\begin{definition}
\label{definition-relative-pseudo-coherence}
Let $S$ be a scheme. Let $f : X \to Y$ be a morphism of
algebraic spaces over $S$ which is locally of finite type.
Let $E$ be an object of $D_\QCoh(\mathcal{O}_X)$. Let $\mathcal{F}$ be a
quasi-coherent $\mathcal{O}_X$-module. Fix $m \in \mathbf{Z}$.
\begin{enumerate}
\item We say $E$ is {\it $m$-pseudo-coherent relative to $Y$}
if the equivalent conditions of
Lemma \ref{lemma-qcoh-relative-pseudo-coherence-characterize} are satisfied.
\item We say $E$ is {\it pseudo-coherent relative to $Y$}
if $E$ is $m$-pseudo-coherent relative to $Y$ for all $m \in \mathbf{Z}$.
\item We say $\mathcal{F}$ is {\it $m$-pseudo-coherent relative to $Y$} if
$\mathcal{F}$ viewed as an object of $D_\QCoh(\mathcal{O}_X)$ is
$m$-pseudo-coherent relative to $Y$.
\item We say $\mathcal{F}$ is {\it pseudo-coherent relative to $Y$} if
$\mathcal{F}$ viewed as an object of $D_\QCoh(\mathcal{O}_X)$ is
pseudo-coherent relative to $Y$.
\end{enumerate}
\end{definition}

\noindent
Most of the properties of pseudo-coherent complexes relative to a base
will follow immediately from the corresponding properties in the case
of schemes. We will add the relevant lemmas here as needed.

\begin{lemma}
\label{lemma-relative-pseudo-coherent-is-moot}
Let $S$ be a scheme. Let $f : X \to Y$ be a morphism of
algebraic spaces over $S$. Let $E$ in $D_\QCoh(\mathcal{O}_X)$.
If $f$ is flat and locally of finite presentation, then
the following are equivalent
\begin{enumerate}
\item $E$ is pseudo-coherent relative to $Y$, and
\item $E$ is pseudo-coherent on $X$.
\end{enumerate}
\end{lemma}

\begin{proof}
By \'etale localization and the definitions we may assume
$X$ and $Y$ are schemes. For the case of schemes this follows
from More on Morphisms, Lemma
\ref{more-morphisms-lemma-check-relative-pseudo-coherence-on-charts}.
\end{proof}

















\section{Pseudo-coherent morphisms}
\label{section-pseudo-coherent}

\noindent
This section is the analogue of
More on Morphisms, Section \ref{more-morphisms-section-pseudo-coherent}
for morphisms of schemes. The reader is encouraged to read up
on pseudo-coherent morphisms of schemes in that section first.

\medskip\noindent
The property ``pseudo-coherent'' of morphisms of schemes is
\'etale local on the source-and-target. To see this use
More on Morphisms,
Lemmas \ref{more-morphisms-lemma-descending-property-pseudo-coherent} and
\ref{more-morphisms-lemma-pseudo-coherent-fppf-local-source}
and
Descent, Lemma \ref{descent-lemma-etale-local-source-target}.
By
Morphisms of Spaces,
Lemma \ref{spaces-morphisms-lemma-local-source-target}
we may define the notion of a pseudo-coherent morphism of algebraic spaces as
follows and it agrees with the already existing notion defined in
More on Morphisms, Section \ref{more-morphisms-section-pseudo-coherent}
when the algebraic spaces in question are representable.

\begin{definition}
\label{definition-pseudo-coherent}
Let $S$ be a scheme.
Let $f : X \to Y$ be a morphism of algebraic spaces over $S$.
\begin{enumerate}
\item We say $f$ is {\it pseudo-coherent} if the equivalent conditions of
Morphisms of Spaces, Lemma \ref{spaces-morphisms-lemma-local-source-target}
hold with $\mathcal{P} =$``pseudo-coherent''.
\item Let $x \in |X|$. We say $f$ is {\it pseudo-coherent at $x$} if
there exists an open neighbourhood $X' \subset X$ of $x$ such
that $f|_{X'} : X' \to Y$ is pseudo-coherent.
\end{enumerate}
\end{definition}

\noindent
Beware that a base change of a pseudo-coherent morphism is not
pseudo-coherent in general.

\begin{lemma}
\label{lemma-flat-base-change-pseudo-coherent}
A flat base change of a pseudo-coherent morphism is pseudo-coherent.
\end{lemma}

\begin{proof}
Omitted. Hint: Use the schemes version of this lemma, see
More on Morphisms,
Lemma \ref{more-morphisms-lemma-flat-base-change-pseudo-coherent}.
\end{proof}

\begin{lemma}
\label{lemma-composition-pseudo-coherent}
A composition of pseudo-coherent morphisms is pseudo-coherent.
\end{lemma}

\begin{proof}
Omitted. Hint: Use the schemes version of this lemma, see
More on Morphisms,
Lemma \ref{more-morphisms-lemma-composition-pseudo-coherent}.
\end{proof}

\begin{lemma}
\label{lemma-pseudo-coherent-finite-presentation}
A pseudo-coherent morphism is locally of finite presentation.
\end{lemma}

\begin{proof}
Immediate from the definitions.
\end{proof}

\begin{lemma}
\label{lemma-flat-finite-presentation-pseudo-coherent}
A flat morphism which is locally of finite presentation is pseudo-coherent.
\end{lemma}

\begin{proof}
Omitted. Hint: Use the schemes version of this lemma, see
More on Morphisms,
Lemma \ref{more-morphisms-lemma-flat-finite-presentation-pseudo-coherent}.
\end{proof}

\begin{lemma}
\label{lemma-permanence-pseudo-coherent}
Let $f : X \to Y$ be a morphism of algebraic spaces pseudo-coherent
over a base algebraic space $B$. Then $f$ is pseudo-coherent.
\end{lemma}

\begin{proof}
Omitted. Hint: Use the schemes version of this lemma, see
More on Morphisms,
Lemma \ref{more-morphisms-lemma-permanence-pseudo-coherent}.
\end{proof}

\begin{lemma}
\label{lemma-Noetherian-pseudo-coherent}
Let $S$ be a scheme. Let $f : X \to Y$ be a morphism of algebraic spaces
over $S$. If $Y$ is locally Noetherian, then $f$ is pseudo-coherent if
and only if $f$ is locally of finite type.
\end{lemma}

\begin{proof}
Omitted. Hint: Use the schemes version of this lemma, see
More on Morphisms,
Lemma \ref{more-morphisms-lemma-Noetherian-pseudo-coherent}.
\end{proof}








\section{Perfect morphisms}
\label{section-perfect}

\noindent
This section is the analogue of
More on Morphisms, Section \ref{more-morphisms-section-perfect}
for morphisms of schemes. The reader is encouraged to read up
on perfect morphisms of schemes in that section first.

\medskip\noindent
The property ``perfect'' of morphisms of schemes is
\'etale local on the source-and-target. To see this use
More on Morphisms,
Lemmas \ref{more-morphisms-lemma-descending-property-perfect} and
\ref{more-morphisms-lemma-perfect-fppf-local-source}
and
Descent, Lemma \ref{descent-lemma-etale-local-source-target}.
By
Morphisms of Spaces,
Lemma \ref{spaces-morphisms-lemma-local-source-target}
we may define the notion of a perfect morphism of algebraic spaces as
follows and it agrees with the already existing notion defined in
More on Morphisms, Section \ref{more-morphisms-section-perfect}
when the algebraic spaces in question are representable.

\begin{definition}
\label{definition-perfect}
Let $S$ be a scheme.
Let $f : X \to Y$ be a morphism of algebraic spaces over $S$.
\begin{enumerate}
\item We say $f$ is {\it perfect} if the equivalent conditions of
Morphisms of Spaces, Lemma \ref{spaces-morphisms-lemma-local-source-target}
hold with $\mathcal{P} =$``perfect''.
\item Let $x \in |X|$. We say $f$ is {\it perfect at $x$} if
there exists an open neighbourhood $X' \subset X$ of $x$ such
that $f|_{X'} : X' \to Y$ is perfect.
\end{enumerate}
\end{definition}

\noindent
Note that a perfect morphism is pseudo-coherent, hence locally of finite
presentation. Beware that a base change of a perfect morphism is not perfect
in general.

\begin{lemma}
\label{lemma-flat-base-change-perfect}
A flat base change of a perfect morphism is perfect.
\end{lemma}

\begin{proof}
Omitted. Hint: Use the schemes version of this lemma, see
More on Morphisms,
Lemma \ref{more-morphisms-lemma-flat-base-change-perfect}.
\end{proof}

\begin{lemma}
\label{lemma-composition-perfect}
A composition of perfect morphisms is perfect.
\end{lemma}

\begin{proof}
Omitted. Hint: Use the schemes version of this lemma, see
More on Morphisms,
Lemma \ref{more-morphisms-lemma-composition-perfect}.
\end{proof}

\begin{lemma}
\label{lemma-flat-finite-presentation-perfect}
Let $S$ be a scheme. Let $f : X \to Y$ be a morphism of algebraic spaces over
$S$. The following are equivalent
\begin{enumerate}
\item $f$ is flat and perfect, and
\item $f$ is flat and locally of finite presentation.
\end{enumerate}
\end{lemma}

\begin{proof}
Omitted. Hint: Use the schemes version of this lemma, see
More on Morphisms,
Lemma \ref{more-morphisms-lemma-flat-finite-presentation-perfect}.
\end{proof}

\begin{lemma}
\label{lemma-perfect-proper-perfect-direct-image}
Let $S$ be a scheme. Let $Y$ be a Noetherian algebraic space over $S$.
Let $f : X \to Y$ be a perfect proper morphism of algebraic spaces.
Let $E \in D(\mathcal{O}_X)$ be perfect. Then
$Rf_*E$ is a perfect object of $D(\mathcal{O}_Y)$.
\end{lemma}

\begin{proof}
We claim that Derived Categories of Spaces, Lemma
\ref{spaces-perfect-lemma-perfect-direct-image} applies.
Conditions (1) and (2) are immediate. Condition (3) is local
on $X$. Thus we may assume $X$ and $Y$ affine and $E$
represented by a strictly perfect complex of $\mathcal{O}_X$-modules.
Thus it suffices to show that $\mathcal{O}_X$ has finite
tor dimension as a sheaf of $f^{-1}\mathcal{O}_Y$-modules
on the \'etale site. By Derived Categories of Spaces, Lemma
\ref{spaces-perfect-lemma-tor-dimension-rel} it suffices to
check this on the Zariski site. This is equivalent to being perfect
for finite type morphisms of schemes by More on Morphisms,
Lemma \ref{more-morphisms-lemma-check-perfect-stalks}.
\end{proof}











\section{Local complete intersection morphisms}
\label{section-lci}

\noindent
This section is the analogue of
More on Morphisms, Section \ref{more-morphisms-section-lci}
for morphisms of schemes. The reader is encouraged to read up
on local complete intersection morphisms of schemes in that section first.

\medskip\noindent
The property ``being a local complete intersection morphism'' of morphisms
of schemes is \'etale local on the source-and-target. To see this use
More on Morphisms,
Lemmas \ref{more-morphisms-lemma-descending-property-lci} and
\ref{more-morphisms-lemma-lci-syntomic-local-source}
and
Descent, Lemma \ref{descent-lemma-etale-local-source-target}.
By
Morphisms of Spaces,
Lemma \ref{spaces-morphisms-lemma-local-source-target}
we may define the notion of a local complete intersection morphism
of algebraic spaces as follows and it agrees with the already existing
notion defined in
More on Morphisms, Section \ref{more-morphisms-section-lci}
when the algebraic spaces in question are representable.

\begin{definition}
\label{definition-lci}
Let $S$ be a scheme.
Let $f : X \to Y$ be a morphism of algebraic spaces over $S$.
\begin{enumerate}
\item We say $f$ is a {\it Koszul morphism}, or that $f$ is a
{\it local complete intersection morphism} if the equivalent conditions of
Morphisms of Spaces, Lemma \ref{spaces-morphisms-lemma-local-source-target}
hold with $\mathcal{P}(f) =$``$f$ is a local complete intersection morphism''.
\item Let $x \in |X|$. We say $f$ is {\it Koszul at $x$} if
there exists an open neighbourhood $X' \subset X$ of $x$ such
that $f|_{X'} : X' \to Y$ is a local complete intersection morphism.
\end{enumerate}
\end{definition}

\noindent
In some sense the defining property of a local complete intersection
morphism is the result of the following lemma.

\begin{lemma}
\label{lemma-lci}
Let $S$ be a scheme.
Let $f : X \to Y$ be a local complete intersection morphism
of algebraic spaces over $S$.
Let $P$ be an algebraic space smooth over $Y$.
Let $U \to X$ be an \'etale morphism of algebraic spaces
and let $i : U \to P$ an immersion of algebraic spaces over $Y$.
Picture:
$$
\xymatrix{
X \ar[rd] & U \ar[l] \ar[d] \ar[r]_i & P \ar[ld] \\
& Y
}
$$
Then $i$ is a Koszul-regular immersion of algebraic spaces.
\end{lemma}

\begin{proof}
Choose a scheme $V$ and a surjective \'etale morphism $V \to Y$.
Choose a scheme $W$ and a surjective \'etale morphism $W \to P \times_Y V$.
Set $U' = U \times_P W$, which is a scheme \'etale over $U$.
We have to show that $U' \to W$ is a Koszul-regular immersion of
schemes, see
Definition \ref{definition-regular-immersion}.
By
Definition \ref{definition-lci}
above the morphism of schemes $U' \to V$ is a local complete intersection
morphism. Hence the result follows from
More on Morphisms, Lemma \ref{more-morphisms-lemma-lci}.
\end{proof}

\noindent
It seems like a good idea to collect here some properties in common
with all Koszul morphisms.

\begin{lemma}
\label{lemma-lci-properties}
Let $S$ be a scheme. Let $f : X \to Y$ be a local complete intersection
morphism of algebraic spaces over $S$. Then
\begin{enumerate}
\item $f$ is locally of finite presentation,
\item $f$ is pseudo-coherent, and
\item $f$ is perfect.
\end{enumerate}
\end{lemma}

\begin{proof}
Omitted. Hint: Use the schemes version of this lemma, see
More on Morphisms,
Lemma \ref{more-morphisms-lemma-lci-properties}.
\end{proof}

\noindent
Beware that a base change of a Koszul morphism is not Koszul in general.

\begin{lemma}
\label{lemma-flat-base-change-lci}
A flat base change of a local complete intersection morphism is a
local complete intersection morphism.
\end{lemma}

\begin{proof}
Omitted. Hint: Use the schemes version of this lemma, see
More on Morphisms,
Lemma \ref{more-morphisms-lemma-flat-base-change-lci}.
\end{proof}

\begin{lemma}
\label{lemma-composition-lci}
A composition of local complete intersection morphisms is a
local complete intersection morphism.
\end{lemma}

\begin{proof}
Omitted. Hint: Use the schemes version of this lemma, see
More on Morphisms,
Lemma \ref{more-morphisms-lemma-composition-lci}.
\end{proof}

\begin{lemma}
\label{lemma-flat-lci}
\begin{slogan}
Syntomic equals flat plus lci (for algebraic spaces).
\end{slogan}
Let $S$ be a scheme.
Let $f : X \to Y$ be a morphism of algebraic spaces over $S$.
The following are equivalent
\begin{enumerate}
\item $f$ is flat and a local complete intersection morphism, and
\item $f$ is syntomic.
\end{enumerate}
\end{lemma}

\begin{proof}
Omitted. Hint: Use the schemes version of this lemma, see
More on Morphisms,
Lemma \ref{more-morphisms-lemma-flat-lci}.
\end{proof}

\begin{lemma}
\label{lemma-regular-immersion-lci}
Let $S$ be a scheme. A Koszul-regular immersion of algebraic spaces
over $S$ is a local complete intersection morphism.
\end{lemma}

\begin{proof}
Let $i : X \to Y$ be a Koszul-regular immersion of algebraic spaces
over $S$. By definition there exists a surjective \'etale morphism
$V \to Y$ where $V$ is a scheme such that $X \times_Y V$ is a scheme
and the base change $X \times_Y V \to V$ is a Koszul-regular immersion of
schemes. By More on Morphisms, Lemma
\ref{more-morphisms-lemma-regular-immersion-lci} we see that
$X \times_Y V \to V$ is a local complete intersection morphism.
From Definition \ref{definition-lci} we conclude that $i$ is a
local complete intersection morphism of algebraic spaces.
\end{proof}

\begin{lemma}
\label{lemma-lci-permanence}
Let $S$ be a scheme. Let
$$
\xymatrix{
X \ar[rr]_f \ar[rd] & & Y \ar[ld] \\
& Z
}
$$
be a commutative diagram of morphisms of algebraic spaces over $S$.
Assume $Y \to Z$ is smooth and $X \to Z$ is a
local complete intersection morphism.
Then $f : X \to Y$ is a local complete intersection morphism.
\end{lemma}

\begin{proof}
Choose a scheme $W$ and a surjective \'etale morphism $W \to Z$.
Choose a scheme $V$ and a surjective \'etale morphism $V \to W \times_Z Y$.
Choose a scheme $U$ and a surjective \'etale morphism $U \to V \times_Y X$.
Then $U \to W$ is a local complete intersection morphism of schemes and
$V \to W$ is a smooth morphism of schemes. By the result for schemes
(More on Morphisms, Lemma \ref{more-morphisms-lemma-lci-permanence})
we conclude that $U \to V$ is a local complete intersection morphism.
By definition this means that $f$ is a local complete intersection morphism.
\end{proof}

\begin{lemma}
\label{lemma-descending-property-lci}
The property $\mathcal{P}(f) =$``$f$ is a local complete intersection
morphism'' is fpqc local on the base.
\end{lemma}

\begin{proof}
Let $S$ be a scheme. Let $f : X \to Y$ be a
morphism of algebraic spaces over $S$.
Let $\{Y_i \to Y\}$ be an fpqc covering
(Topologies on Spaces, Definition
\ref{spaces-topologies-definition-fpqc-covering}).
Let $f_i : X_i \to Y_i$ be the base change of $f$ by $Y_i \to Y$.
If $f$ is a local complete intersection morphism,
then each $f_i$ is a local complete intersection morphism
by Lemma \ref{lemma-flat-base-change-lci}.

\medskip\noindent
Conversely, assume each $f_i$ is a local complete intersection morphism.
We may replace the covering by a refinement (again because
flat base change preserves the property of being a
local complete intersection morphism). Hence we may assume
$Y_i$ is a scheme for each $i$, see
Topologies on Spaces, Lemma \ref{spaces-topologies-lemma-refine-fpqc-schemes}.
Choose a scheme $V$ and a surjective \'etale morphism $V \to Y$.
Choose a scheme $U$ and a surjective \'etale morphism
$U \to V \times_Y X$. We have to show that $U \to V$ is a
local complete intersection morphism of schemes.
By Topologies on Spaces, Lemma
\ref{spaces-topologies-lemma-recognize-fpqc-covering}
we have that $\{Y_i \times_Y V \to V\}$ is an fpqc covering
of schemes. By the case of schemes
(More on Morphisms, Lemma \ref{more-morphisms-lemma-descending-property-lci})
it suffices to prove the base change
$$
U \times_Y Y_i = U \times_V (V \times_Y Y_i) \longrightarrow V
$$
of $U \to V$ by $V \times_Y Y_i \to V$ is a
local complete intersection morphism. We can write this as the
composition
$$
U \times_Y Y_i \longrightarrow
(V \times_Y X) \times_Y Y_i = V \times_Y X_i \longrightarrow
V \times_Y Y_i
$$
The first arrow is an \'etale morphism of schemes (as a base change of
$U \to V \times_Y X$) and the second arrow is a
local complete intersection morphism of schemes as a flat base change of $f_i$.
The result follows as being a local complete intersection morphism
is syntomic local on the source and since \'etale morphisms
are syntomic (More on Morphisms, Lemma
\ref{more-morphisms-lemma-lci-syntomic-local-source}
and Morphisms, Lemma \ref{morphisms-lemma-etale-syntomic}).
\end{proof}

\begin{lemma}
\label{lemma-lci-syntomic-local-source}
The property $\mathcal{P}(f) =$``$f$ is a local complete intersection
morphism'' is syntomic local on the source.
\end{lemma}

\begin{proof}
This follows from
Descent on Spaces, Lemma \ref{spaces-descent-lemma-transfer-from-schemes} and
More on Morphisms, Lemma \ref{more-morphisms-lemma-lci-syntomic-local-source}.
\end{proof}

\begin{lemma}
\label{lemma-base-change-lci-fibres}
Let $S$ be a scheme. Consider a commutative diagram
$$
\xymatrix{
X \ar[rr]_f \ar[rd]_p & & Y \ar[ld]^q \\
& Z
}
$$
of algebraic spaces over $S$. Assume that both $p$ and $q$
are flat and locally of finite presentation.
Then there exists an open subspace $U(f) \subset X$
such that $|U(f)| \subset |X|$ is the set of points where $f$ is Koszul.
Moreover, for any morphism of algebraic spaces $Z' \to Z$, if
$f' : X' \to Y'$ is the base change of $f$ by $Z' \to Z$, then
$U(f')$ is the inverse image of $U(f)$ under the projection $X' \to X$.
\end{lemma}

\begin{proof}
This lemma is the analogue of
More on Morphisms, Lemma \ref{more-morphisms-lemma-base-change-lci-fibres}
and in fact we will deduce the lemma from it. By
Definition \ref{definition-lci}
the set $\{x \in |X| : f \text{ is Koszul at }x\}$ is
open in $|X|$ hence by
Properties of Spaces, Lemma \ref{spaces-properties-lemma-open-subspaces}
it corresponds to an open subspace $U(f)$ of $X$. Hence we only need to
prove the final statement.

\medskip\noindent
Choose a scheme $W$ and a surjective \'etale morphism $W \to Z$.
Choose a scheme $V$ and a surjective \'etale morphism $V \to W \times_Z Y$.
Choose a scheme $U$ and a surjective \'etale morphism $U \to V \times_Y X$.
Finally, choose a scheme $W'$ and a surjective \'etale morphism
$W' \to W \times_Z Z'$.
Set $V' = W' \times_W V$ and $U' = W' \times_W U$, so that we obtain
surjective \'etale morphisms $V' \to Y'$ and $U' \to X'$.
We will use without further mention an \'etale morphism of algebraic spaces
induces an open map of associated topological spaces (see
Properties of Spaces, Lemma
\ref{spaces-properties-lemma-etale-open}).
Note that by definition $U(f)$ is the image in $|X|$ of the set $T$
of points in $U$ where the morphism of schemes $U \to V$ is Koszul.
Similarly, $U(f')$ is the image in $|X'|$ of the set $T'$ of points in
$U'$ where the morphism of schemes $U' \to V'$ is Koszul. Now, by construction
the diagram
$$
\xymatrix{
U' \ar[r] \ar[d] & U \ar[d] \\
V' \ar[r] & V
}
$$
is cartesian (in the category of schemes). Hence the aforementioned
More on Morphisms, Lemma \ref{more-morphisms-lemma-base-change-lci-fibres}
applies to show that $T'$ is the inverse image of $T$. Since
$|U'| \to |X'|$ is surjective this implies the lemma.
\end{proof}

\begin{lemma}
\label{lemma-unramified-lci}
Let $S$ be a scheme. Let $f : X \to Y$ be a local complete intersection
morphism of algebraic spaces over $S$. Then $f$ is unramified if and only
if $f$ is formally unramified and in this case the conormal sheaf
$\mathcal{C}_{X/Y}$ is finite locally free on $X$.
\end{lemma}

\begin{proof}
This follows from the corresponding result for morphisms of schemes, see
More on Morphisms, Lemma \ref{more-morphisms-lemma-unramified-lci},
by \'etale localization, see
Lemma \ref{lemma-universal-thickening-localize}.
(Note that in the situation of this lemma the morphism $V \to U$
is unramified and a local complete intersection morphism by definition.)
\end{proof}

\begin{lemma}
\label{lemma-transitivity-conormal-lci}
Let $S$ be a scheme. Let $Z \to Y \to X$ be formally unramified morphisms
of algebraic spaces over $S$. Assume that $Z \to Y$ is a local complete
intersection morphism. The exact sequence
$$
0 \to i^*\mathcal{C}_{Y/X} \to
\mathcal{C}_{Z/X} \to
\mathcal{C}_{Z/Y} \to 0
$$
of
Lemma \ref{lemma-transitivity-conormal}
is short exact.
\end{lemma}

\begin{proof}
Choose a scheme $U$ and a surjective \'etale morphism $U \to X$.
Choose a scheme $V$ and a surjective \'etale morphism $V \to U \times_X Y$.
Choose a scheme $W$ and a surjective \'etale morphism $W \to V \times_Y Z$.
By
Lemma \ref{lemma-universal-thickening-localize}
the morphisms $W \to V$ and $V \to U$ are formally unramified.
Moreover the sequence
$i^*\mathcal{C}_{Y/X} \to \mathcal{C}_{Z/X} \to \mathcal{C}_{Z/Y} \to 0$
restricts to the corresponding sequence
$i^*\mathcal{C}_{V/U} \to \mathcal{C}_{W/U} \to \mathcal{C}_{W/V} \to 0$
for $W \to V \to U$. Hence the result follows from the result for schemes
(More on Morphisms, Lemma \ref{more-morphisms-lemma-transitivity-conormal-lci})
as by definition the morphism $W \to V$ is a local complete intersection
morphism.
\end{proof}









\section{When is a morphism an isomorphism?}
\label{section-when-isomorphism}

\noindent
More generally we can ask:
``When does a morphism have property $\mathcal{P}$?''
A more precise question is the following. Suppose given a commutative diagram
$$
\xymatrix{
X \ar[rr]_f \ar[rd]_p & & Y \ar[ld]^q \\
& Z
}
$$
of algebraic spaces. Does there exist a monomorphism of algebraic spaces
$W \to Z$ with the following two properties:
\begin{enumerate}
\item the base change $f_W : X_W \to Y_W$ has property $\mathcal{P}$, and
\item any morphism $Z' \to Z$ of algebraic spaces factors through $W$ if
and only if the base change $f_{Z'} : X_{Z'} \to Y_{Z'}$ has property
$\mathcal{P}$.
\end{enumerate}
In many cases, if $W \to Z$ exists, then it is an immersion, open immersion,
or closed immersion.

\medskip\noindent
The answer to this question may depend on auxiliary properties of the
morphisms $f$, $p$, and $q$. An example is $\mathcal{P}(f) =$``$f$ is flat''
which we have discussed for morphisms of schemes in the case $Y = S$ in
great detail in the chapter ``More on Flatness'', starting with
More on Flatness, Section \ref{flat-section-flattening-functors}.

\begin{lemma}
\label{lemma-where-unramified}
Consider a commutative diagram
$$
\xymatrix{
X \ar[rr]_f \ar[rd]_p & & Y \ar[ld]^q \\
& Z
}
$$
of algebraic spaces. Assume that $p$ is locally of finite type and closed.
Then there exists an open subspace $W \subset Z$
such that a morphism $Z' \to Z$ factors through $W$ if and only if the
base change $f_{Z'} : X_{Z'} \to Y_{Z'}$ is unramified.
\end{lemma}

\begin{proof}
By
Morphisms of Spaces, Lemma \ref{spaces-morphisms-lemma-where-unramified}
there exists an open subspace $U(f) \subset X$ which is the set of
points where $f$ is unramified. Moreover, formation of $U(f)$ commutes
with arbitrary base change. Let $W \subset Z$ be the open subspace
(see
Properties of Spaces, Lemma
\ref{spaces-properties-lemma-open-subspaces})
with underlying set of points
$$
|W| = |Z| \setminus |p|\left(|X| \setminus |U(f)|\right)
$$
i.e., $z \in |Z|$ is a point of $W$ if and only if $f$ is unramified
at every point of $X$ above $z$. Note that this is open because we
assumed that $p$ is closed. Since the formation of $U(f)$
commutes with arbitrary base change we immediately see (using
Properties of Spaces, Lemma
\ref{spaces-properties-lemma-factor-through-open-subspace})
that $W$ has the desired universal property.
\end{proof}

\begin{lemma}
\label{lemma-where-unramified-universally-injective}
Consider a commutative diagram
$$
\xymatrix{
X \ar[rr]_f \ar[rd]_p & & Y \ar[ld]^q \\
& Z
}
$$
of algebraic spaces. Assume that
\begin{enumerate}
\item $p$ is locally of finite type,
\item $p$ is closed, and
\item $p_2 : X \times_Y X \to Z$ is closed.
\end{enumerate}
Then there exists an open subspace $W \subset Z$
such that a morphism $Z' \to Z$ factors through $W$ if and only if the
base change $f_{Z'} : X_{Z'} \to Y_{Z'}$ is unramified and universally
injective.
\end{lemma}

\begin{proof}
After replacing $Z$ by the open subspace found in
Lemma \ref{lemma-where-unramified}
we may assume that $f$ is already unramified; note that this does not
destroy assumption (2) or (3). By
Morphisms of Spaces, Lemma
\ref{spaces-morphisms-lemma-diagonal-unramified-morphism}
we see that $\Delta_{X/Y} : X \to X \times_Y X$ is an open immersion.
This remains true after any base change. Hence by
Morphisms of Spaces, Lemma
\ref{spaces-morphisms-lemma-universally-injective}
we see that $f_{Z'}$ is universally injective if and only if
the base change of the diagonal $X_{Z'} \to (X \times_Y X)_{Z'}$
is an isomorphism. Let $W \subset Z$ be the open subspace
(see
Properties of Spaces, Lemma
\ref{spaces-properties-lemma-open-subspaces})
with underlying set of points
$$
|W| = |Z| \setminus
|p_2|\left(|X \times_Y X| \setminus \Im(|\Delta_{X/Y}|)\right)
$$
i.e., $z \in |Z|$ is a point of $W$ if and only if the fibre of
$|X \times_Y X| \to |Z|$ over $z$ is in the image of
$|X| \to |X \times_Y X|$. Then it is clear from the discussion above
that the restriction $p^{-1}(W) \to q^{-1}(W)$ of $f$ is
unramified and universally injective.

\medskip\noindent
Conversely, suppose that $f_{Z'}$ is unramified and universally injective.
In order to show that $Z' \to Z$ factors through $W$ it suffices to show
that $|Z'| \to |Z|$ has image contained in $|W|$, see
Properties of Spaces, Lemma
\ref{spaces-properties-lemma-factor-through-open-subspace}.
Hence it suffices to prove the result when $Z'$ is the spectrum of a field.
Denote $z \in |Z|$ the image of $|Z'| \to |Z|$. The discussion above shows
that
$$
|X_{Z'}| \longrightarrow |(X \times_Y X)_{Z'}|
$$
is surjective. By
Properties of Spaces,
Lemma \ref{spaces-properties-lemma-points-cartesian}
in the commutative diagram
$$
\xymatrix{
|X_{Z'}| \ar[d] \ar[r] &
|(X \times_Y X)_{Z'}| \ar[d] \\
|p|^{-1}(\{z\}) \ar[r] &
|p_2|^{-1}(\{z\})
}
$$
the vertical arrows are surjective. It follows that $z \in |W|$ as desired.
\end{proof}

\begin{lemma}
\label{lemma-where-closed-immersion}
Consider a commutative diagram
$$
\xymatrix{
X \ar[rr]_f \ar[rd]_p & & Y \ar[ld]^q \\
& Z
}
$$
of algebraic spaces. Assume that
\begin{enumerate}
\item $p$ is locally of finite type,
\item $p$ is universally closed, and
\item $q : Y \to Z$ is separated.
\end{enumerate}
Then there exists an open subspace $W \subset Z$
such that a morphism $Z' \to Z$ factors through $W$ if and only if the
base change $f_{Z'} : X_{Z'} \to Y_{Z'}$ is a closed immersion.
\end{lemma}

\begin{proof}
We will use the characterization of closed immersions as
universally closed, unramified, and universally injective morphisms, see
Lemma \ref{lemma-characterize-closed-immersion}.
First, note that since $p$ is universally closed and $q$ is
separated, we see that $f$ is universally closed, see
Morphisms of Spaces, Lemma
\ref{spaces-morphisms-lemma-universally-closed-permanence}.
It follows that any base change of $f$ is universally closed, see
Morphisms of Spaces, Lemma
\ref{spaces-morphisms-lemma-base-change-universally-closed}.
Thus to finish the proof of the lemma it suffices to prove that
the assumptions of
Lemma \ref{lemma-where-unramified-universally-injective}
are satisfied. The projection $\text{pr}_0 : X \times_Y X \to X$
is universally closed as a base change of $f$, see
Morphisms of Spaces, Lemma
\ref{spaces-morphisms-lemma-base-change-universally-closed}.
Hence $X \times_Y X \to Z$ is universally closed as
a composition of universally closed morphisms (see
Morphisms of Spaces, Lemma
\ref{spaces-morphisms-lemma-composition-universally-closed}).
This finishes the proof of the lemma.
\end{proof}

\begin{lemma}
\label{lemma-where-flat}
Consider a commutative diagram
$$
\xymatrix{
X \ar[rr]_f \ar[rd]_p & & Y \ar[ld]^q \\
& Z
}
$$
of algebraic spaces. Assume that
\begin{enumerate}
\item $p$ is locally of finite presentation,
\item $p$ is flat,
\item $p$ is closed, and
\item $q$ is locally of finite type.
\end{enumerate}
Then there exists an open subspace $W \subset Z$
such that a morphism $Z' \to Z$ factors through $W$ if and only if the
base change $f_{Z'} : X_{Z'} \to Y_{Z'}$ is flat.
\end{lemma}

\begin{proof}
By Lemma \ref{lemma-base-change-flatness-fibres}
the set
$$
A = \{x \in |X| : X\text{ flat at }x \text{ over }Y\}.
$$
is open in $|X|$ and its formation commutes with arbitrary base
change. Let $W \subset Z$ be the open subspace
(see
Properties of Spaces, Lemma
\ref{spaces-properties-lemma-open-subspaces})
with underlying set of points
$$
|W| = |Z| \setminus |p|\left(|X| \setminus A\right)
$$
i.e., $z \in |Z|$ is a point of $W$ if and only if the whole fibre
of $|X| \to |Z|$ over $z$ is contained in $A$. This is open because
$p$ is closed. Since the formation of $A$ commutes with arbitrary
base change it follows that $W$ works.
\end{proof}

\begin{lemma}
\label{lemma-where-surjective-flat}
Consider a commutative diagram
$$
\xymatrix{
X \ar[rr]_f \ar[rd]_p & & Y \ar[ld]^q \\
& Z
}
$$
of algebraic spaces. Assume that
\begin{enumerate}
\item $p$ is locally of finite presentation,
\item $p$ is flat,
\item $p$ is closed,
\item $q$ is locally of finite type, and
\item $q$ is closed.
\end{enumerate}
Then there exists an open subspace $W \subset Z$
such that a morphism $Z' \to Z$ factors through $W$ if and only if the
base change $f_{Z'} : X_{Z'} \to Y_{Z'}$ is surjective and flat.
\end{lemma}

\begin{proof}
By
Lemma \ref{lemma-where-flat}
we may assume that $f$ is flat.
Note that $f$ is locally of finite presentation by
Morphisms of Spaces,
Lemma \ref{spaces-morphisms-lemma-finite-presentation-permanence}.
Hence $f$ is open, see
Morphisms of Spaces, Lemma \ref{spaces-morphisms-lemma-fppf-open}.
Let $W \subset Z$ be the open subspace (see
Properties of Spaces, Lemma
\ref{spaces-properties-lemma-open-subspaces})
with underlying set of points
$$
|W| = |Z| \setminus |q|\left(|Y| \setminus |f|(|X|)\right).
$$
in other words for $z \in |Z|$ we have $z \in |W|$ if and only
if the whole fibre of $|Y| \to |Z|$ over $z$ is in the image of
$|X| \to |Y|$. Since $q$ is closed this set is open in $|Z|$.
The morphism $X_W \to Y_W$ is surjective by construction.
Finally, suppose that $X_{Z'} \to Y_{Z'}$ is surjective.
In order to show that $Z' \to Z$ factors through $W$ it suffices to show
that $|Z'| \to |Z|$ has image contained in $|W|$, see
Properties of Spaces, Lemma
\ref{spaces-properties-lemma-factor-through-open-subspace}.
Hence it suffices to prove the result when $Z'$ is the spectrum of a field.
Denote $z \in |Z|$ the image of $|Z'| \to |Z|$. By
Properties of Spaces,
Lemma \ref{spaces-properties-lemma-points-cartesian}
in the commutative diagram
$$
\xymatrix{
|X_{Z'}| \ar[d] \ar[r] &
|Y_{Z'}| \ar[d] \\
|p|^{-1}(\{z\}) \ar[r] &
|q|^{-1}(\{z\})
}
$$
the vertical arrows are surjective. It follows that $z \in |W|$ as desired.
\end{proof}

\begin{lemma}
\label{lemma-where-isomorphism}
Consider a commutative diagram
$$
\xymatrix{
X \ar[rr]_f \ar[rd]_p & & Y \ar[ld]^q \\
& Z
}
$$
of algebraic spaces. Assume that
\begin{enumerate}
\item $p$ is locally of finite presentation,
\item $p$ is flat,
\item $p$ is universally closed,
\item $q$ is locally of finite type,
\item $q$ is closed, and
\item $q$ is separated.
\end{enumerate}
Then there exists an open subspace $W \subset Z$
such that a morphism $Z' \to Z$ factors through $W$ if and only if the
base change $f_{Z'} : X_{Z'} \to Y_{Z'}$ is an isomorphism.
\end{lemma}

\begin{proof}
By
Lemma \ref{lemma-where-surjective-flat}
there exists an open subspace $W_1 \subset Z$ such that
$f_{Z'}$ is surjective and flat if and only if $Z' \to Z$
factors through $W_1$. By
Lemma \ref{lemma-where-closed-immersion}
there exists an open subspace $W_2 \subset Z$ such that
$f_{Z'}$ is a closed immersion if and only if $Z' \to Z$
factors through $W_2$. We claim that $W = W_1 \cap W_2$ works.
Certainly, if $f_{Z'}$ is an isomorphism, then $Z' \to Z$
factors through $W$. Hence it suffices to show that
$f_W$ is an isomorphism. By construction $f_W$ is a
surjective flat closed immersion. In particular $f_W$ is
representable. Since a surjective flat closed immersion of
schemes is an isomorphism (see
Morphisms, Lemma \ref{morphisms-lemma-characterize-flat-closed-immersions})
we win. (Note that actually $f_W$ is locally of finite presentation,
whence open, so you can avoid the use of this lemma if you like.)
\end{proof}

\begin{lemma}
\label{lemma-where-lci}
Consider a commutative diagram
$$
\xymatrix{
X \ar[rr]_f \ar[rd]_p & & Y \ar[ld]^q \\
& Z
}
$$
of algebraic spaces. Assume that
\begin{enumerate}
\item $p$ is flat and locally of finite presentation,
\item $p$ is closed, and
\item $q$ is flat and locally of finite presentation,
\end{enumerate}
Then there exists an open subspace $W \subset Z$
such that a morphism $Z' \to Z$ factors through $W$ if and only if the
base change $f_{Z'} : X_{Z'} \to Y_{Z'}$ is a local complete intersection
morphism.
\end{lemma}

\begin{proof}
By Lemma \ref{lemma-base-change-lci-fibres}
there exists an open subspace $U(f) \subset X$ which is the set of
points where $f$ is Koszul. Moreover, formation of $U(f)$ commutes
with arbitrary base change. Let $W \subset Z$ be the open subspace
(see
Properties of Spaces, Lemma
\ref{spaces-properties-lemma-open-subspaces})
with underlying set of points
$$
|W| = |Z| \setminus |p|\left(|X| \setminus |U(f)|\right)
$$
i.e., $z \in |Z|$ is a point of $W$ if and only if $f$ is Koszul
at every point of $X$ above $z$. Note that this is open because we
assumed that $p$ is closed. Since the formation of $U(f)$
commutes with arbitrary base change we immediately see (using
Properties of Spaces, Lemma
\ref{spaces-properties-lemma-factor-through-open-subspace})
that $W$ has the desired universal property.
\end{proof}








\section{Exact sequences of differentials and conormal sheaves}
\label{section-exact}

\noindent
In this section we collect some results on exact sequences of conormal
sheaves and sheaves of differentials. In some sense these are all
realizations of the triangle of cotangent complexes associated to
composable morphisms of algebraic spaces.

\medskip\noindent
In the sequences below each of the maps
are as constructed in either
Lemma \ref{lemma-functoriality-differentials} or
Lemma \ref{lemma-universal-thickening-functorial}.
Let $S$ be a scheme. Let $g : Z \to Y$ and $f : Y \to X$ be morphisms of
algebraic spaces over $S$.
\begin{enumerate}
\item There is a canonical exact sequence
$$
g^*\Omega_{Y/X} \to \Omega_{Z/X} \to \Omega_{Z/Y} \to 0,
$$
see
Lemma \ref{lemma-triangle-differentials}.
If $g : Z \to Y$ is formally smooth, then this sequence is a short
exact sequence, see
Lemma \ref{lemma-triangle-differentials-formally-smooth}.
\item If $g$ is formally unramified, then there is a canonical exact sequence
$$
\mathcal{C}_{Z/Y} \to g^*\Omega_{Y/X} \to \Omega_{Z/X} \to 0,
$$
see
Lemma \ref{lemma-universally-unramified-differentials-sequence}.
If $f \circ g : Z \to X$ is formally smooth, then this sequence is a short
exact sequence, see
Lemma \ref{lemma-differentials-formally-unramified-formally-smooth}.
\item if $g$ and $f \circ g$ are formally unramified, then there is a
canonical exact sequence
$$
\mathcal{C}_{Z/X} \to \mathcal{C}_{Z/Y} \to g^*\Omega_{Y/X} \to 0,
$$
see
Lemma \ref{lemma-two-unramified-morphisms}.
If $f : Y \to X$ is formally smooth, then this sequence is a short
exact sequence, see
Lemma \ref{lemma-two-unramified-morphisms-formally-smooth}.
\item if $g$ and $f$ are formally unramified, then there is a canonical
exact sequence
$$
g^*\mathcal{C}_{Y/X} \to \mathcal{C}_{Z/X} \to \mathcal{C}_{Z/Y} \to 0.
$$
see
Lemma \ref{lemma-transitivity-conormal-unramified}.
If $g : Z \to Y$ is a local complete intersection morphism,
then this sequence is a short exact sequence, see
Lemma \ref{lemma-transitivity-conormal-lci}.
\end{enumerate}










\section{Characterizing pseudo-coherent complexes, II}
\label{section-characterize-pseudo-coherent}

\noindent
In this section we discuss a characterization of pseudo-coherent complexes
in terms of cohomology. Earlier material on pseudo-coherent complexes
on algebraic spaces may be found in
Derived Categories of Spaces, Section
\ref{spaces-perfect-section-spell-out}
and in
Derived Categories of Spaces, Section
\ref{spaces-perfect-section-pseudo-coherent-hocolim}.
The analogue of this section for schemes is More on Morphisms, Section
\ref{more-morphisms-section-characterize-pseudo-coherent}.
A basic tool will be to reduce to the case of projective space
using a derived version of Chow's lemma, see
Lemma \ref{lemma-derived-chow}.

\begin{lemma}
\label{lemma-case-of-tor-independence}
Let $S$ be a scheme. Consider a commutative diagram of algebraic spaces
$$
\xymatrix{
Z' \ar[d] \ar[r] & Y' \ar[d] \\
X' \ar[r] & B'
}
$$
over $S$.
Let $B \to B'$ be a morphism. Denote by $X$ and $Y$ the base
changes of $X'$ and $Y'$ to $B$.
Assume $Y' \to B'$ and $Z' \to X'$ are flat.
Then $X \times_B Y$ and $Z'$ are Tor independent over $X' \times_{B'} Y'$.
\end{lemma}

\begin{proof}
By Derived Categories of Spaces, Lemma
\ref{spaces-perfect-lemma-tor-independent}
we may check tor independence \'etale locally on $X \times_B Y$
and $Z'$. This\footnote{Here is the argument in more detail.
Choose a surjective \'etale morphism $W' \to B'$
with $W'$ a scheme. Choose a surjective \'etale morphism
$W \to B \times_{B'} W'$ with $W$ a scheme. Choose a
surjective \'etale morphism
$U' \to X' \times_{B'} W'$ with $U'$ a scheme. Choose a
surjective \'etale morphism $V' \to Y' \times_{B'} W'$ with $V'$ a scheme.
Observe that $U' \times_{W'} V' \to X' \times_{B'} Y'$ is surjective
\'etale. Choose a surjective \'etale morphism
$T' \to Z' \times_{X' \times_{B'} Y'} U' \times_{W'} V'$
with $T'$ a scheme. Denote $U$ and $V$ the base changes of $U'$ and $V'$
to $W$. Then the lemma says that $X \times_B Y$ and $Z'$
are Tor independent over $X' \times_{B'} Y'$ as algebraic spaces
if and only if $U \times_W V$ and $T'$ are Tor independent over
$U' \times_{W'} V'$ as schemes. Thus
it suffices to prove the lemma for
the square with corners $T', U', V', W'$ and base change by $W \to W'$.
The flatness of $Y' \to B'$ and $Z' \to X'$ implies flatness
of $V' \to W'$ and $T' \to U'$.}
reduces the lemma to the case of schemes
which is More on Morphisms, Lemma
\ref{more-morphisms-lemma-case-of-tor-independence}.
\end{proof}

\begin{lemma}[Derived Chow's lemma]
\label{lemma-derived-chow}
Let $A$ be a ring. Let $X$ be a separated algebraic space
of finite presentation over $A$. Let $x \in |X|$. Then there exist
an $n \geq 0$,
a closed subspace $Z \subset X \times_A \mathbf{P}^n_A$,
a point $z \in |Z|$,
an open $V \subset \mathbf{P}^n_A$, and
an object $E$ in $D(\mathcal{O}_{X \times_A \mathbf{P}^n_A})$ such that
\begin{enumerate}
\item $Z \to X \times_A \mathbf{P}^n_A$ is of finite presentation,
\item $c : Z \to \mathbf{P}^n_A$ is a closed immersion over $V$,
set $W = c^{-1}(V)$,
\item the restriction of $b : Z \to X$ to $W$ is \'etale,
$z \in W$, and $b(z) = x$,
\item $E|_{X \times_A V} \cong
(b, c)_*\mathcal{O}_Z|_{X \times_A V}$,
\item $E$ is pseudo-coherent and supported on $Z$.
\end{enumerate}
\end{lemma}

\begin{proof}
We can find a finite type $\mathbf{Z}$-subalgebra $A' \subset A$
and an algebraic space $X'$ separated and of finite presentation over $A'$
whose base change to $A$ is $X$. See
Limits of Spaces, Lemmas
\ref{spaces-limits-lemma-descend-finite-presentation} and
\ref{spaces-limits-lemma-descend-separated-morphism}.
Let $x' \in |X'|$ be the image of $x$.
If we can prove the lemma for $(X'/A', x')$, then
the lemma follows for $(X/A, x)$.
Namely, if $n', Z', z', V', E'$ provide the solution
for $(X'/A', x')$, then we can let
$n = n'$,
let $Z \subset X \times \mathbf{P}^n$ be the inverse image of $Z'$,
let $z \in Z$ be the unique point mapping to $x$,
let $V \subset \mathbf{P}^n_A$ be the inverse image of $V'$, and
let $E$ be the derived pullback of $E'$.
Observe that $E$ is pseudo-coherent by
Cohomology on Sites, Lemma
\ref{sites-cohomology-lemma-pseudo-coherent-pullback}.
It only remains to check (5). To see this
set $W = c^{-1}(V)$ and $W' = (c')^{-1}(V')$
and consider the cartesian square
$$
\xymatrix{
W \ar[d]_{(b, c)} \ar[r] & W' \ar[d]^{(b', c')} \\
X \times_A V \ar[r] & X' \times_{A'} V'
}
$$
By Lemma \ref{lemma-case-of-tor-independence} $X \times_A V$ and $W'$
are tor-independent over $X' \times_{A'} V'$.
Thus the derived pullback of
$(b', c')_*\mathcal{O}_{W'}$ to $X \times_A V$
is $(b, c)_*\mathcal{O}_W$ by
Derived Categories of Spaces,
Lemma \ref{spaces-perfect-lemma-compare-base-change}.
This also uses that $R(b', c')_*\mathcal{O}_{Z'} = (b', c')_*\mathcal{O}_{Z'}$
because $(b', c')$ is a closed immersion and similarly for
$(b, c)_*\mathcal{O}_Z$.
Since $E'|_{U' \times_{A'} V'} =
(b', c')_*\mathcal{O}_{W'}$ we obtain
$E|_{U \times_A V} = (b, c)_*\mathcal{O}_W$
and (5) holds.
This reduces us to the situation described in the next
paragraph.

\medskip\noindent
Assume $A$ is of finite type over $\mathbf{Z}$.
Choose an \'etale morphism $U \to X$ where $U$ is an affine scheme
and a point $u \in U$ mapping to $x$. Then $U$ is of finite type over $A$.
Choose a closed immersion $U \to \mathbf{A}^n_A$ and denote
$j : U \to \mathbf{P}^n_A$ the immersion we get by composing
with the open immersion $\mathbf{A}^n_A \to \mathbf{P}^n_A$.
Let $Z$ be the scheme theoretic closure of
$$
(\text{id}_U, j) : U \longrightarrow X \times_A \mathbf{P}^n_A
$$
Let $z \in Z$ be the image of $u$.
Let $Y \subset \mathbf{P}^n_A$ be the scheme theoretic
closure of $j$. Then it is clear that $Z \subset X \times_A Y$
is the scheme theoretic closure of
$(\text{id}_U, j) : U \to X \times_A Y$.
As $X$ is separated, the morphism
$X \times_A Y \to Y$ is separated as well.
Hence we see that $Z \to Y$ is an isomorphism over
the open subscheme $j(U) \subset Y$ by
Morphisms of Spaces, Lemma
\ref{spaces-morphisms-lemma-scheme-theoretic-image-of-partial-section}.
Choose $V \subset \mathbf{P}^n_A$ open with $V \cap Y = j(U)$.
Then we see that (2) holds, that $W = (\text{id}_U, j)(U)$, and hence
that (3) holds. Part (1) holds because $A$ is Noetherian.

\medskip\noindent
Because $A$ is Noetherian we see that $X$ and $X \times_A \mathbf{P}^n_A$
are Noetherian algebraic spaces. Hence we can take $E = (b, c)_*\mathcal{O}_Z$
in this case: (4) is clear and for (5) see
Derived Categories of Spaces, Lemma
\ref{spaces-perfect-lemma-identify-pseudo-coherent-noetherian}.
This finishes the proof.
\end{proof}

\begin{lemma}
\label{lemma-compute-Fourier-Mukai-for-derived-chow}
Let $X/A$, $x \in |X|$, and
$n, Z, z, V, E$ be as in Lemma \ref{lemma-derived-chow}.
For any $K \in D_\QCoh(\mathcal{O}_X)$ we have
$$
Rq_*(Lp^*K \otimes^\mathbf{L} E)|_V = R(W \to V)_*K|_W
$$
where $p : X \times_A \mathbf{P}^n_A \to X$ and
$q : X \times_A \mathbf{P}^n_A \to \mathbf{P}^n_A$ are
the projections and where the morphism $W \to V$ is
the finitely presented closed immersion $c|_W : W \to V$.
\end{lemma}

\begin{proof}
Since $W = c^{-1}(V)$ and since $c$ is a closed immersion
over $V$, we see that $c|_W$ is a closed immersion.
It is of finite presentation because $W$ and $V$ are of finite
presentation over $A$, see
Morphisms of Spaces, Lemma
\ref{spaces-morphisms-lemma-finite-presentation-permanence}.
First we have
$$
Rq_*(Lp^*K \otimes^\mathbf{L} E)|_V =
Rq'_*\left((Lp^*K \otimes^\mathbf{L} E)|_{X \times_A V}\right)
$$
where $q' : X \times_A V \to V$ is the projection because
formation of total direct image commutes with localization.
Denote $i = (b, c)|_W : W \to X \times_A V$ the given closed immersion.
Then
$$
Rq'_*\left((Lp^*K \otimes^\mathbf{L} E)|_{X \times_A V}\right) =
Rq'_*(Lp^*K|_{X \times_A V} \otimes^\mathbf{L} i_*\mathcal{O}_W)
$$
by property (5). Since $i$ is a closed immersion we have
$i_*\mathcal{O}_W = Ri_*\mathcal{O}_W$.
Using
Derived Categories of Spaces,
Lemma \ref{spaces-perfect-lemma-cohomology-base-change}
we can rewrite this as
$$
Rq'_* Ri_* Li^* Lp^*K|_{X \times_A V} =
R(q' \circ i)_* Lb^*K|_W =
R(W \to V)_* K|_W
$$
which is what we want. (Note that restricting to $W$
and derived pulling back via $W \to X$ is the same thing
as $W$ is \'etale over $X$.)
\end{proof}

\begin{lemma}
\label{lemma-characterize-pseudo-coherent}
Let $A$ be a ring. Let $X$ be an algebraic space separated and
of finite presentation over $A$. Let $K \in D_\QCoh(\mathcal{O}_X)$.
If $R\Gamma(X, E \otimes^\mathbf{L} K)$ is pseudo-coherent
in $D(A)$ for every pseudo-coherent $E$ in $D(\mathcal{O}_X)$,
then $K$ is pseudo-coherent relative to $A$
(Definition \ref{definition-relative-pseudo-coherence}).
\end{lemma}

\begin{proof}
Assume $K \in D_\QCoh(\mathcal{O}_X)$ and
$R\Gamma(X, E \otimes^\mathbf{L} K)$ is pseudo-coherent
in $D(A)$ for every pseudo-coherent $E$ in $D(\mathcal{O}_X)$.
Let $x \in |X|$. We will show that $K$ is pseudo-coherent relative to $A$
in an \'etale neighbourhood of $x$. This will prove the lemma
by our definition of relative pseudo-coherence.

\medskip\noindent
Choose $n, Z, z, V, E$ as in Lemma \ref{lemma-derived-chow}.
Denote $p : X \times \mathbf{P}^n \to X$ and
$q : X \times \mathbf{P}^n \to \mathbf{P}^n_A$
the projections.
Then for any $i \in \mathbf{Z}$ we have
\begin{align*}
& R\Gamma(\mathbf{P}^n_A,
Rq_*(Lp^*K \otimes^\mathbf{L} E)
\otimes^\mathbf{L}
\mathcal{O}_{\mathbf{P}^n_A}(i)) \\
& =
R\Gamma(X \times \mathbf{P}^n,
Lp^*K \otimes^\mathbf{L} E \otimes^\mathbf{L}
Lq^*\mathcal{O}_{\mathbf{P}^n_A}(i)) \\
& =
R\Gamma(X,
K \otimes^\mathbf{L} Rp_*(E \otimes^\mathbf{L}
Lq^*\mathcal{O}_{\mathbf{P}^n_A}(i)))
\end{align*}
by
Derived Categories of Spaces,
Lemma \ref{spaces-perfect-lemma-cohomology-base-change}.
By
Derived Categories of Spaces, Lemma
\ref{spaces-perfect-lemma-flat-proper-pseudo-coherent-direct-image-general}
the complex $Rp_*(E \otimes^\mathbf{L} Lq^*\mathcal{O}_{\mathbf{P}^n_A}(i))$
is pseudo-coherent on $X$. Hence the assumption tells us the expression
in the displayed formula is a pseudo-coherent object of $D(A)$.
By
Derived Categories of Schemes,
Lemma \ref{perfect-lemma-pseudo-coherent-on-projective-space}
we conclude that $Rq_*(Lp^*K \otimes^\mathbf{L} E)$ is
pseudo-coherent on $\mathbf{P}^n_A$.
By Lemma \ref{lemma-compute-Fourier-Mukai-for-derived-chow}
we have
$$
Rq_*(Lp^*K \otimes^\mathbf{L} E)|_{X \times_A V} =
R(W \to V)_*K|_W
$$
Since $W \to V$ is a closed immersion into an open subscheme of
$\mathbf{P}^n_A$ this means $K|_W$ is pseudo-coherent relative to $A$
for example by
More on Morphisms,
Lemma \ref{more-morphisms-lemma-check-relative-pseudo-coherence-on-charts}.
\end{proof}

\begin{lemma}
\label{lemma-characterize-pseudo-coh-improved}
Let $A$ be a ring. Let $X$ be an algebraic space separated and
of finite presentation over $A$. Let $K \in D_\QCoh(\mathcal{O}_X)$. If 
$R \Gamma (X, E \otimes ^{\mathbf{L}} K)$ is
pseudo-coherent in $D(A)$ for every perfect 
$E \in D(\mathcal{O}_X)$, then $K$ is pseudo-coherent
relative to $A$.
\end{lemma}

\begin{proof}
In view of Lemma \ref{lemma-characterize-pseudo-coherent}, it suffices
to show $R \Gamma (X, E \otimes ^{\mathbf{L}} K)$ is
pseudo-coherent in $D(A)$ for every pseudo-coherent 
$E \in D(\mathcal{O}_X)$. By Derived Categories of Spaces,
Proposition \ref{spaces-perfect-proposition-detecting-bounded-above}
it follows that $K \in D^-_\QCoh (\mathcal{O}_X)$. Now the
result follows by Derived Categories of Spaces, Lemma
\ref{spaces-perfect-lemma-perfect-enough}.
\end{proof}












\section{Relatively perfect objects}
\label{section-relatively-perfect}

\noindent
In this section we introduce a notion from \cite{lieblich-complexes}.
This notion has been discussed for morphisms of schemes in
Derived Categories of Schemes, Section
\ref{perfect-section-relatively-perfect}.

\begin{definition}
\label{definition-relatively-perfect}
Let $S$ be a scheme. Let $f : X \to Y$ be a morphism of algebraic spaces
over $S$ which is flat and locally of finite presentation.
An object $E$ of $D(\mathcal{O}_X)$ is {\it perfect relative to $Y$} or
{\it $Y$-perfect} if $E$ is pseudo-coherent
(Cohomology on Sites, Definition
\ref{sites-cohomology-definition-pseudo-coherent}) and
$E$ locally has finite tor dimension as an object of
$D(f^{-1}\mathcal{O}_Y)$
(Cohomology on Sites, Definition
\ref{sites-cohomology-definition-tor-amplitude}).
\end{definition}

\noindent
Please see Derived Categories of Schemes,
Remark \ref{perfect-remark-discuss-rel-perfect} for a discussion;
here we just mention that $E$ being pseudo-coherent is the
same thing as $E$ being pseudo-coherent relative to $Y$ by
Lemma \ref{lemma-relative-pseudo-coherent-is-moot}.
Moreover, pseudo-coherence of $E$ implies $E \in D_\QCoh(\mathcal{O}_X)$, see
Derived Categories of Spaces, Lemma
\ref{spaces-perfect-lemma-pseudo-coherent}.

\begin{example}
\label{example-relatively-perfect-field}
Let $k$ be a field. Let $X$ be an algebraic space of finite presentation
over $k$ (in particular $X$ is quasi-compact). Then an object $E$ of
$D(\mathcal{O}_X)$ is $k$-perfect if and only if it is bounded and
pseudo-coherent (by definition), i.e., if and only if it is in
$D^b_{\textit{Coh}}(X)$ (by
Derived Categories of Spaces, Lemma
\ref{spaces-perfect-lemma-identify-pseudo-coherent-noetherian}).
Thus being relatively perfect does {\bf not} mean ``perfect on the fibres''.
\end{example}

\noindent
The corresponding algebra concept is studied in
More on Algebra, Section \ref{more-algebra-section-relatively-perfect}.
We can link the notion for algebraic spaces with the
algebraic notion as follows.

\begin{lemma}
\label{lemma-affine-locally-rel-perfect}
Let $S$ be a scheme. Let $f : X \to Y$ be a morphism of algebraic spaces
over $S$ which is flat and locally of finite presentation.
Let $E \in D_\QCoh(\mathcal{O}_X)$. The following are equivalent:
\begin{enumerate}
\item $E$ is $Y$-perfect,
\item for every commutative diagram
$$
\xymatrix{
U \ar[d] \ar[r]_g & V \ar[d] \\
X \ar[r]^f & Y
}
$$
where $U$, $V$ are schemes and the vertical arrows are \'etale, the complex
$E|_U$ is $V$-perfect in the sense of Derived Categories of Schemes,
Definition \ref{perfect-definition-relatively-perfect},
\item for some commutative diagram as in (2) with $U \to X$
surjective, the complex $E|_U$ is $V$-perfect in the sense of
Derived Categories of Schemes,
Definition \ref{perfect-definition-relatively-perfect},
\item for every commutative diagram as in (2) with $U$ and $V$
affine the complex $R\Gamma(U, E)$ is $\mathcal{O}_Y(V)$-perfect.
\end{enumerate}
\end{lemma}

\begin{proof}
To make sense of parts (2), (3), (4) of the lemma, observe that
the object $E|_U$ of $D_\QCoh(\mathcal{O}_U)$ corresponds to
an object $E_0$ of $D_\QCoh(\mathcal{O}_{U_0})$ where $U_0$
denotes the scheme underlying $U$, see Derived Categories of Spaces,
Lemma \ref{spaces-perfect-lemma-derived-quasi-coherent-small-etale-site}.
Moreover, in this case $E_0$ is pseudo-coherent if and only if
$E|_U$ is pseudo-coherent, see Derived Categories of Spaces,
Lemma \ref{spaces-perfect-lemma-descend-pseudo-coherent}.
Also, $E|_U$ locally has finite tor dimension over
$f^{-1}\mathcal{O}_Y|_U = g^{-1}\mathcal{O}_V$ if and only if
$E_0$ locally has finite tor dimension over $g_0^{-1}\mathcal{O}_{V_0}$
by Derived Categories of Spaces, Lemma
\ref{spaces-perfect-lemma-tor-dimension-rel}.
Here $g_0 : U_0 \to V_0$ is the morphism of schemes representing
$g : U \to V$ (notation as in Derived Categories of Spaces,
Remark \ref{spaces-perfect-remark-match-total-direct-images}).
Finally, observe that ``being pseudo-coherent'' is \'etale local and
of course ``having locally finite tor dimension'' is \'etale local.
Thus we see that it suffices to check $Y$-perfectness \'etale locally
and by the above discussion we see that (1) implies (2) and
(3) implies (1). Since part (4) is equivalent
to (2) and (3) by
Derived Categories of Schemes, Lemma
\ref{perfect-lemma-affine-locally-rel-perfect}
the proof is complete.
\end{proof}

\begin{lemma}
\label{lemma-triangulated}
Let $S$ be a scheme. Let $f : X \to Y$ be a morphism of
algebraic spaces over $S$ which is flat and locally of finite presentation.
The full subcategory of $D(\mathcal{O}_X)$ consisting of $Y$-perfect objects is
a saturated\footnote{Derived Categories, Definition
\ref{derived-definition-saturated}.} triangulated subcategory.
\end{lemma}

\begin{proof}
This follows from Cohomology on Sites, Lemmas
\ref{sites-cohomology-lemma-cone-pseudo-coherent},
\ref{sites-cohomology-lemma-summands-pseudo-coherent},
\ref{sites-cohomology-lemma-cone-tor-amplitude}, and
\ref{sites-cohomology-lemma-summands-tor-amplitude}.
\end{proof}

\begin{lemma}
\label{lemma-perfect-relatively-perfect}
Let $S$ be a scheme.
Let $f : X \to Y$ be a morphism of algebraic spaces over $S$
which is flat and locally of finite presentation.
A perfect object of $D(\mathcal{O}_X)$ is $Y$-perfect.
If $K, M \in D(\mathcal{O}_X)$, then $K \otimes_{\mathcal{O}_X}^\mathbf{L} M$
is $Y$-perfect if $K$ is perfect and $M$ is $Y$-perfect.
\end{lemma}

\begin{proof}
Reduce to the case of schemes using
Lemma \ref{lemma-affine-locally-rel-perfect}
and then apply
Derived Categories of Schemes, Lemma
\ref{perfect-lemma-perfect-relatively-perfect}.
\end{proof}

\begin{lemma}
\label{lemma-base-change-relatively-perfect}
Let $S$ be a scheme.
Let $f : X \to Y$ be a morphism of algebraic spaces over $S$
which is flat and locally of finite presentation.
Let $g : Y' \to Y$ be a morphism of algebraic spaces over $S$.
Set $X' = Y' \times_Y X$ and denote $g' : X' \to X$ the projection.
If $K \in D(\mathcal{O}_X)$ is $Y$-perfect, then $L(g')^*K$
is $Y'$-perfect.
\end{lemma}

\begin{proof}
Reduce to the case of schemes using
Lemma \ref{lemma-affine-locally-rel-perfect}
and then apply
Derived Categories of Schemes, Lemma
\ref{perfect-lemma-base-change-relatively-perfect}.
\end{proof}

\begin{situation}
\label{situation-relative-descent}
Let $S$ be a scheme.
Let $Y = \lim_{i \in I} Y_i$ be a limit of a directed system of
algebraic spaces over $S$
with affine transition morphisms $g_{i'i} : Y_{i'} \to Y_i$.
We assume that $Y_i$ is quasi-compact and quasi-separated for all $i \in I$.
We denote $g_i : Y \to Y_i$ the projection. We fix an element $0 \in I$
and a flat morphism of finite presentation $X_0 \to Y_0$.
We set $X_i = Y_i \times_{Y_0} X_0$ and $X = Y \times_{Y_0} X_0$
and we denote the transition morphisms $f_{i'i} : X_{i'} \to X_i$
and $f_i : X \to X_i$ the projections.
\end{situation}

\begin{lemma}
\label{lemma-relative-descend-homomorphisms}
In Situation \ref{situation-relative-descent}.
Let $K_0$ and $L_0$ be objects of $D(\mathcal{O}_{X_0})$.
Set $K_i = Lf_{i0}^*K_0$ and $L_i = Lf_{i0}^*L_0$ for $i \geq 0$
and set $K = Lf_0^*K_0$ and $L = Lf_0^*L_0$. Then the map
$$
\colim_{i \geq 0} \Hom_{D(\mathcal{O}_{X_i})}(K_i, L_i)
\longrightarrow
\Hom_{D(\mathcal{O}_X)}(K, L)
$$
is an isomorphism if $K_0$ is pseudo-coherent and
$L_0 \in D_\QCoh(\mathcal{O}_{X_0})$ has (locally)
finite tor dimension as an object of
$D((X_0 \to Y_0)^{-1}\mathcal{O}_{Y_0})$
\end{lemma}

\begin{proof}
For every quasi-compact and quasi-separated object $U_0$ of
$(X_0)_{spaces, \etale}$ consider the condition $P$ that
$$
\colim_{i \geq 0} \Hom_{D(\mathcal{O}_{U_i})}(K_i|_{U_i}, L_i|_{U_i})
\longrightarrow
\Hom_{D(\mathcal{O}_U)}(K|_U, L|_U)
$$
is an isomorphism where $U = X \times_{X_0} U_0$ and
$U_i = X_i \times_{X_0} U_0$. We will prove $P$ holds for each $U_0$.

\medskip\noindent
Suppose that $(U_0 \subset W_0, V_0 \to W_0)$ is an elementary
distinguished square in $(X_0)_{spaces, \etale}$
and $P$ holds for $U_0, V_0, U_0 \times_{W_0} V_0$.
Then $P$ holds for $W_0$ by Mayer-Vietoris
for hom in the derived category, see Derived Categories of Spaces,
Lemma \ref{spaces-perfect-lemma-mayer-vietoris-hom}.

\medskip\noindent
We first consider $U_0 = W_0 \times_{Y_0} X_0$ with $W_0$ a
quasi-compact and quasi-separated object of $(Y_0)_{spaces, \etale}$.
By the induction principle of Derived Categories of Spaces,
Lemma \ref{spaces-perfect-lemma-induction-principle}
applied to these $W_0$ and the previous paragraph,
we find that it is enough to prove
$P$ for $U_0 = W_0 \times_{Y_0} X_0$ with $W_0$ affine.
In other words, we have reduced to the case where $Y_0$ is affine.
Next, we apply the induction principle again, this time to all
quasi-compact and quasi-separated opens of $X_0$, to reduce to the
case where $X_0$ is affine as well.

\medskip\noindent
If $X_0$ and $Y_0$ are affine, then we are back in the case
of schemes which is proved in
Derived Categories of Schemes, Lemma
\ref{perfect-lemma-relative-descend-homomorphisms}.
The reader may use
Derived Categories of Spaces, Lemmas
\ref{spaces-perfect-lemma-pseudo-coherent},
\ref{spaces-perfect-lemma-derived-quasi-coherent-small-etale-site},
\ref{spaces-perfect-lemma-descend-pseudo-coherent}, and
\ref{spaces-perfect-lemma-tor-dimension-rel}
to accomplish the translation of the statement into a statement
involving only schemes and derived categories of modules on schemes.
\end{proof}

\begin{lemma}
\label{lemma-descend-relatively-perfect}
In Situation \ref{situation-relative-descent} the category of
$Y$-perfect objects of $D(\mathcal{O}_X)$ is the colimit of the categories
of $Y_i$-perfect objects of $D(\mathcal{O}_{X_i})$.
\end{lemma}

\begin{proof}
For every quasi-compact and quasi-separated object $U_0$ of
$(X_0)_{spaces, \etale}$ consider the condition $P$
that the functor
$$
\colim_{i \geq 0} D_{Y_i\text{-perfect}}(\mathcal{O}_{U_i})
\longrightarrow
D_{Y\text{-perfect}}(\mathcal{O}_U)
$$
is an equivalence where $U = X \times_{X_0} U_0$ and
$U_i = X_i \times_{X_0} U_0$.
We observe that we already know this functor is fully faithful
by Lemma \ref{lemma-relative-descend-homomorphisms}. Thus it suffices to prove
essential surjectivity.

\medskip\noindent
Suppose that $(U_0 \subset W_0, V_0 \to W_0)$ is an elementary
distinguished square in $(X_0)_{spaces, \etale}$
and $P$ holds for $U_0, V_0, U_0 \times_{W_0} V_0$.
We claim that $P$ holds for $W_0$. We will use the notation
$U_i = X_i \times_{X_0} U_0$, $U = X \times_{X_0} U_0$,
and similarly for $V_0$ and $W_0$. We will abusively use the symbol
$f_i$ for all the morphisms $U \to U_i$, $V \to V_i$,
$U \times_W V \to U_i \times_{W_i} V_i$, and $W \to W_i$.
Suppose $E$ is an $Y$-perfect object of $D(\mathcal{O}_W)$.
Goal: show $E$ is in the essential image of the functor.
By assumption,
we can find $i \geq 0$, an $Y_i$-perfect object $E_{U, i}$ on $U_i$,
an $Y_i$-perfect object $E_{V, i}$ on $V_i$, and
isomorphisms $Lf_i^*E_{U, i} \to E|_U$ and $Lf_i^*E_{V, i} \to E|_V$.
Let
$$
a : E_{U, i} \to (Rf_{i, *}E)|_{U_i}
\quad\text{and}\quad
b : E_{V, i} \to (Rf_{i, *}E)|_{V_i}
$$
the maps adjoint to the isomorphisms $Lf_i^*E_{U, i} \to E|_U$
and $Lf_i^*E_{V, i} \to E|_V$.
By fully faithfulness, after increasing $i$,
we can find an isomorphism
$c : E_{U, i}|_{U_i \times_{W_i} V_i} \to E_{V, i}|_{U_i \times_{W_i} V_i}$
which pulls back to the identifications 
$$
Lf_i^*E_{U, i}|_{U \times_W V} \to E|_{U \times_W V} \to
Lf_i^*E_{V, i}|_{U \times_W V}.
$$
Apply Derived Categories of Spaces, Lemma
\ref{spaces-perfect-lemma-glue}
to get an object $E_i$ on $W_i$ and a map $d : E_i \to Rf_{i, *}E$
which restricts to the maps $a$ and $b$ over $U_i$ and $V_i$.
Then it is clear that $E_i$ is $Y_i$-perfect (because being
relatively perfect is an \'etale local property) and that
$d$ is adjoint to an isomorphism $Lf_i^*E_i \to E$.

\medskip\noindent
By exactly the same argument as used in
the proof of Lemma \ref{lemma-relative-descend-homomorphisms}
using the induction principle
(Derived Categories of Spaces, Lemma
\ref{spaces-perfect-lemma-induction-principle})
we reduce to the case where both $X_0$ and $Y_0$
are affine: first work with quasi-compact and quasi-separated objects
in $(Y_0)_{spaces, \etale}$ to reduce to
$Y_0$ affine, then work with quasi-compact and quasi-separated object
in $(X_0)_{spaces, \etale}$ to reduce to $X_0$ affine.
In the affine case the result follows from the case of schemes which is
Derived Categories of Schemes, Lemma
\ref{perfect-lemma-descend-relatively-perfect}.
The translation into the case for schemes is done by
Lemma \ref{lemma-affine-locally-rel-perfect}.
\end{proof}

\begin{lemma}
\label{lemma-derived-pushforward-rel-perfect}
Let $S$ be a scheme.
Let $f : X \to Y$ be a morphism of algebraic spaces over $S$
which is flat, proper, and
of finite presentation. Let $E \in D(\mathcal{O}_X)$ be $Y$-perfect.
Then $Rf_*E$ is a perfect object of $D(\mathcal{O}_Y)$
and its formation commutes with arbitrary base change.
\end{lemma}

\begin{proof}
The statement on base change is Derived Categories of Spaces,
Lemma \ref{spaces-perfect-lemma-base-change-tensor} (with
$\mathcal{G}^\bullet$ equal to $\mathcal{O}_X$ in degree $0$).
Thus it suffices to show that $Rf_*E$ is a perfect object. We will reduce
to the case where $Y$ is Noetherian affine by a limit argument.

\medskip\noindent
The question is \'etale local on $Y$, hence we may assume $Y$ is affine.
Say $Y = \Spec(R)$. We write $R = \colim R_i$ as a filtered colimit
of Noetherian rings $R_i$. By Limits of Spaces, Lemma
\ref{spaces-limits-lemma-descend-finite-presentation}
there exists an $i$ and an algebraic space
$X_i$ of finite presentation over $R_i$
whose base change to $R$ is $X$. By
Limits of Spaces, Lemmas \ref{spaces-limits-lemma-eventually-proper} and
\ref{spaces-limits-lemma-descend-flat}
we may assume $X_i$ is proper and flat over $R_i$.
By Lemma \ref{lemma-descend-relatively-perfect}
we may assume there exists a $R_i$-perfect object $E_i$ of
$D(\mathcal{O}_{X_i})$ whose pullback to $X$ is $E$.
Applying Derived Categories of Spaces,
Lemma \ref{spaces-perfect-lemma-perfect-direct-image}
to $X_i \to \Spec(R_i)$ and $E_i$ and using the
base change property already shown we obtain the result.
\end{proof}

\begin{lemma}
\label{lemma-compute-ext-rel-perfect}
Let $S$ be a scheme.
Let $f : X \to Y$ be a morphism of algebraic spaces over $S$.
Let $E, K \in D(\mathcal{O}_X)$.
Assume
\begin{enumerate}
\item $Y$ is quasi-compact and quasi-separated,
\item $f$ is proper, flat, and of finite presentation,
\item $E$ is $Y$-perfect,
\item $K$ is pseudo-coherent.
\end{enumerate}
Then there exists a pseudo-coherent $L \in D(\mathcal{O}_Y)$ such that
$$
Rf_*R\SheafHom(K, E) = R\SheafHom(L, \mathcal{O}_Y)
$$
and the same is true after arbitrary base change: given
$$
\vcenter{
\xymatrix{
X' \ar[r]_{g'} \ar[d]_{f'} &
X \ar[d]^f \\
Y' \ar[r]^g &
Y
}
}
\quad\quad
\begin{matrix}
\text{cartesian, then we have } \\
Rf'_*R\SheafHom(L(g')^*K, L(g')^*E) \\
= R\SheafHom(Lg^*L, \mathcal{O}_{Y'})
\end{matrix}
$$
\end{lemma}

\begin{proof}
Since $Y$ is quasi-compact and quasi-separated, the same is true for $X$.
By Derived Categories of Spaces, Lemma
\ref{spaces-perfect-lemma-pseudo-coherent-hocolim} we can write
$K = \text{hocolim} K_n$ with $K_n$ perfect and $K_n \to K$ inducing
an isomorphism on truncations $\tau_{\geq -n}$. Let $K_n^\vee$
be the dual perfect complex
(Cohomology on Sites, Lemma \ref{sites-cohomology-lemma-dual-perfect-complex}).
We obtain an inverse system $\ldots \to K_3^\vee \to K_2^\vee \to K_1^\vee$
of perfect objects. By Lemma \ref{lemma-perfect-relatively-perfect}
we see that $K_n^\vee \otimes_{\mathcal{O}_X} E$ is $Y$-perfect.
Thus we may apply Lemma \ref{lemma-derived-pushforward-rel-perfect}
to $K_n^\vee \otimes_{\mathcal{O}_X} E$ and we obtain an inverse system
$$
\ldots \to M_3 \to M_2 \to M_1
$$
of perfect complexes on $Y$ with
$$
M_n = Rf_*(K_n^\vee \otimes_{\mathcal{O}_X}^\mathbf{L} E) =
Rf_*R\SheafHom(K_n, E)
$$
Moreover, the formation of these complexes commutes with any
base change, namely $Lg^*M_n =
Rf'_*((L(g')^*K_n)^\vee \otimes_{\mathcal{O}_{X'}}^\mathbf{L} L(g')^*E) =
Rf'_*R\SheafHom(L(g')^*K_n, L(g')^*E)$.

\medskip\noindent
As $K_n \to K$ induces an isomorphism on $\tau_{\geq -n}$, we see that
$K_n \to K_{n + 1}$ induces an isomorphism on $\tau_{\geq -n}$.
It follows that $K_{n + 1}^\vee \to K_n^\vee$
induces an isomorphism on $\tau_{\leq n}$ as
$K_n^\vee = R\SheafHom(K_n, \mathcal{O}_X)$.
Suppose that $E$ has tor amplitude in $[a, b]$ as a complex
of $f^{-1}\mathcal{O}_Y$-modules. Then the same is true after
any base change, see
Derived Categories of Spaces, Lemma
\ref{spaces-perfect-lemma-tor-independence-and-tor-amplitude}.
We find that
$K_{n + 1}^\vee \otimes_{\mathcal{O}_X} E \to
K_n^\vee \otimes_{\mathcal{O}_X} E$
induces an isomorphism on $\tau_{\leq n + a}$
and the same is true after any base change.
Applying the right derived functor $Rf_*$
we conclude the maps $M_{n + 1} \to M_n$
induce isomorphisms on $\tau_{\leq n + a}$
and the same is true after any base change.
Choose a distinguished triangle
$$
M_{n + 1} \to M_n \to C_n \to M_{n + 1}[1]
$$
Pick $y \in |Y|$. Choose an elementary \'etale neighbourhood
$(U, u) \to (Y, y)$; this is possible by
Decent Spaces, Lemma
\ref{decent-spaces-lemma-decent-space-elementary-etale-neighbourhood}.
Take $Y'$ equal to the spectrum of the residue field at $u$.
Pull back to see that $C_n|_U \otimes_{\mathcal{O}_U}^\mathbf{L} \kappa(u)$
has nonzero cohomology only in degrees $\geq n + a$. By
More on Algebra, Lemma
\ref{more-algebra-lemma-lift-perfect-from-residue-field}
we see that the perfect complex $C_n|_U$ has tor amplitude in
$[n + a, m_n]$ for some integer $m_n$ and after possibly shrinking $U$.
Thus $C_n$ has tor amplitude in $[n + a, m_n]$ for some integer $m_n$
(because $Y$ is quasi-compact).
In particular, the dual perfect complex $C_n^\vee$ has tor amplitude in
$[-m_n, -n - a]$.

\medskip\noindent
Let $L_n = M_n^\vee$ be the dual perfect complex. The
conclusion from the discussion in the previous paragraph is that
$L_n \to L_{n + 1}$ induces isomorphisms on $\tau_{\geq -n - a}$.
Thus $L = \text{hocolim} L_n$ is pseudo-coherent, see
Derived Categories of Spaces, Lemma
\ref{spaces-perfect-lemma-pseudo-coherent-hocolim}.
Since we have
$$
R\SheafHom(K, E) = R\SheafHom(\text{hocolim} K_n, E) =
R\lim R\SheafHom(K_n, E) = R\lim K_n^\vee \otimes_{\mathcal{O}_X} E
$$
(Cohomology on Sites, Lemma
\ref{sites-cohomology-lemma-colim-and-lim-of-duals})
and since $R\lim$ commutes with $Rf_*$ we find that
$$
Rf_*R\SheafHom(K, E) = R\lim M_n = R\lim R\SheafHom(L_n, \mathcal{O}_Y) =
R\SheafHom(L, \mathcal{O}_Y)
$$
This proves the formula over $Y$. Since the construction of $M_n$ is
compatible with base chance, the formula continues to hold after
any base change.
\end{proof}

\begin{remark}
\label{remark-compare-L}
The reader may have noticed the similarity between
Lemma \ref{lemma-compute-ext-rel-perfect} and
Derived Categories of Spaces, Lemma \ref{spaces-perfect-lemma-compute-ext}.
Indeed, the pseudo-coherent complex $L$ of
Lemma \ref{lemma-compute-ext-rel-perfect}
may be characterized as the unique pseudo-coherent complex
on $Y$ such that there are functorial isomorphisms
$$
\Ext^i_{\mathcal{O}_Y}(L, \mathcal{F}) \longrightarrow
\Ext^i_{\mathcal{O}_X}(K,
E \otimes_{\mathcal{O}_X}^\mathbf{L} Lf^*\mathcal{F})
$$
compatible with boundary maps for $\mathcal{F}$ ranging over
$\QCoh(\mathcal{O}_Y)$. If we ever need this we will
formulate a precise result here and give a detailed proof.
\end{remark}

\begin{lemma}
\label{lemma-bounded-on-fibres}
Let $S$ be a scheme. Let $X$ be an algebraic space over $S$ such that
the structure morphism $f : X \to S$ is flat and
locally of finite presentation. Let $E$ be a pseudo-coherent
object of $D(\mathcal{O}_X)$. The following are equivalent
\begin{enumerate}
\item $E$ is $S$-perfect, and
\item $E$ is locally bounded below and for every point $s \in S$
the object $L(X_s \to X)^*E$ of $D(\mathcal{O}_{X_s})$
is locally bounded below.
\end{enumerate}
\end{lemma}

\begin{proof}
Since everything is local we immediately reduce to the
case that $X$ and $S$ are affine, see Lemma
\ref{lemma-affine-locally-rel-perfect}.
This case is handled by
Derived Categories of Schemes, Lemma \ref{perfect-lemma-bounded-on-fibres}.
\end{proof}

\begin{lemma}
\label{lemma-characterize-relatively-perfect}
Let $A$ be a ring. Let $X$ be an algebraic space separated, of
finite presentation, and flat over $A$. Let $K \in D_\QCoh(\mathcal{O}_X)$.
If  $R \Gamma (X, E \otimes^\mathbf{L} K)$ is perfect in
$D(A)$ for every perfect $E \in D(\mathcal{O}_X)$, then $K$ is
$\Spec(A)$-perfect.
\end{lemma}

\begin{proof}
By Lemma \ref{lemma-characterize-pseudo-coh-improved},
$K$ is pseudo-coherent relative to $A$. By Lemma
\ref{lemma-relative-pseudo-coherent-is-moot},
$K$ is pseudo-coherent in $D(\mathcal{O}_X)$.
By Derived Categories of Spaces, Proposition
\ref{spaces-perfect-proposition-detecting-bounded-below}
we see that $K$ is in $D^-(\mathcal{O}_X)$.
Let $\mathfrak{p}$ be a prime ideal of $A$ and denote
$i : Y \to X$ the inclusion of the scheme theoretic
fibre over $\mathfrak{p}$, i.e., $Y$ is a scheme over $\kappa(\mathfrak p)$.
By Lemma \ref{lemma-bounded-on-fibres},
we will be done if we can show  $Li^*(K)$ is bounded below. 
Let $G \in D_{perf} (\mathcal{O}_X)$ be a perfect
complex which generates $D_\QCoh (\mathcal{O}_X)$,
see Derived Categories of Spaces, Theorem
\ref{spaces-perfect-theorem-bondal-van-den-Bergh}.
We have
\begin{align*}
R\Hom _{\mathcal{O}_Y}(Li^*(G), Li^*(K))
& =
R\Gamma(Y, Li^*(G ^\vee \otimes ^\mathbf{L} K)) \\
& =
R\Gamma(X, G^\vee \otimes ^{\mathbf{L}} K)
\otimes^\mathbf{L}_A \kappa(\mathfrak{p})
\end{align*}
The first equality uses that $Li^*$ preserves perfect objects and duals
and Cohomology on Sites, Lemma
\ref{sites-cohomology-lemma-dual-perfect-complex}; we omit
some details. The second equality follows from
Derived Categories of Spaces, Lemma
\ref{spaces-perfect-lemma-compare-base-change}
as $X$ is flat over $A$. It follows from our hypothesis that this is a
perfect object of $D(\kappa(\mathfrak{p}))$. The object
$Li^*(G) \in D_{perf}(\mathcal{O}_Y)$  generates $D_\QCoh(\mathcal{O}_Y)$ by
Derived Categories of Spaces, Remark
\ref{spaces-perfect-remark-pullback-generator}.
Hence Derived Categories of Spaces, Proposition
\ref{spaces-perfect-proposition-detecting-bounded-below}
now implies that $Li^*(K)$ is bounded below and we win.
\end{proof}














\section{Theorem of the cube}
\label{section-theorem-cube}

\noindent
This section is the analogue of More on Morphisms, Section
\ref{more-morphisms-section-theorem-cube}.
The following lemma tells us that the diagonal of the Picard
functor is representable by locally closed immersions under
the assumptions made in the lemma.

\begin{lemma}
\label{lemma-diagonal-picard-flat-proper}
Let $S$ be a scheme.
Let $f : X \to Y$ be a flat, proper morphism of finite presentation
of algebraic spaces over $S$.
Let $\mathcal{E}$ be a finite locally free $\mathcal{O}_X$-module.
For a morphism $g : Y' \to Y$ consider the base change diagram
$$
\xymatrix{
X' \ar[d]_{f'} \ar[r]_{g'} & X \ar[d]^f \\
Y' \ar[r]^g & Y
}
$$
Assume $\mathcal{O}_{Y'} \to f'_*\mathcal{O}_{X'}$ is an
isomorphism for all $g : Y' \to Y$.
Then there exists an immersion $j : Z \to Y$ of finite presentation
such that a morphism $g : Y' \to Y$ factors through $Z$ if and only if
there exists a finite locally free $\mathcal{O}_{Y'}$-module $\mathcal{N}$
with $(f')^*\mathcal{N} \cong (g')^*\mathcal{L}$.
\end{lemma}

\begin{proof}
Let $y : \Spec(k) \to Y$ be a field valued point.
Then the fibre $X_y$ of $f$ at $y$ is connected by our assumption
that $H^0(X_y, \mathcal{O}_{X_y}) = k$. Thus the rank of
$\mathcal{E}$ is constant on the fibres. Since $f$ is open
(Morphisms of Spaces, Lemma \ref{spaces-morphisms-lemma-fppf-open})
and closed we conclude that there is a decomposition $Y = \coprod Y_r$
of $Y$ into open and closed subspaces such that $\mathcal{E}$
has constant rank $r$ on the inverse image of $Y_r$.
Thus we may assume $\mathcal{E}$ has constant rank $r$.
We will denote $\mathcal{E}^\vee = \SheafHom(\mathcal{E}, \mathcal{O}_X)$
the dual rank $r$ module.

\medskip\noindent
By cohomology and base change (more precisely by
Derived Categories of Spaces, Lemma
\ref{spaces-perfect-lemma-flat-proper-perfect-direct-image-general})
we see that $E = Rf_*\mathcal{E}$ is a perfect object of the
derived category of $Y$ and that its formation commutes with
arbitrary change of base. Similarly for $E' = Rf_*\mathcal{E}^\vee$.
Since there is never any cohomology in degrees $< 0$, we see that
$E$ and $E'$ have (locally) tor-amplitude in $[0, b]$ for some $b$.
Observe that for any $g : Y' \to Y$ we have
$f'_*((g')^*\mathcal{E}) = H^0(Lg^*E)$ and
$f'_*((g')^*\mathcal{E}^\vee) = H^0(Lg^*E')$.
Let $j : Z \to Y$ and $j' : Z' \to Y$ be the locally closed
immersions constructed in Derived Categories of Spaces, Lemma
\ref{spaces-perfect-lemma-locally-closed-where-H0-locally-free}
for $E$ and $E'$ with $a = 0$ and $r = r$; these are characterized
by the property that $H^0(Lj^*E)$ and $H^0((j')^*E')$
are locally free modules of rank $r$ compatible with pullback.

\medskip\noindent
Let $g : Y' \to Y$ be a morphism. If there exists an $\mathcal{N}$
as in the lemma, then, using the projection formula
Cohomology on Sites, Lemma \ref{sites-cohomology-lemma-projection-formula},
we see that the modules
$$
f'_*((g')^*\mathcal{E}) \cong
f'_*((f')^*\mathcal{N}) \cong
\mathcal{N} \otimes_{\mathcal{O}_{Y'}} f'_*\mathcal{O}_{X'} \cong
\mathcal{N}\quad\text{and similarly }\quad
f'_*((g')^*\mathcal{E}^\vee) \cong \mathcal{N}^\vee
$$
are locally free of rank $r$ and remain locally free of rank $r$
after any further base change $Y'' \to Y'$.
Hence in this case $g : Y' \to Y$ factors through $j$ and through $j'$.
Thus we may replace $Y$ by $Z \times_Y Z'$ and assume that
$f_*\mathcal{E}$ and $f_*\mathcal{E}^\vee$ are locally free
$\mathcal{O}_Y$-modules of rank $r$
whose formation commutes with arbitrary change of base.

\medskip\noindent
In this situation if $g : Y' \to Y$ is a morphism and there exists an
$\mathcal{N}$ as in the lemma, then the map (cup product in degree $0$)
$$
f'_*((g')^*\mathcal{E})
\otimes_{\mathcal{O}_{Y'}}
f'_*((g')^*\mathcal{E}^\vee)
\longrightarrow \mathcal{O}_{Y'}
$$
is a perfect pairing. Conversely, if this cup product map is a
perfect pairing, then we see that locally on $Y'$ we have
a basis of sections
$\sigma_1, \ldots, \sigma_r$ in $f'_*((g')^*\mathcal{L})$ and
$\tau_1, \ldots, \tau_r$ in
$f'_*((g')^*\mathcal{E}^\vee)$ whose products satisfy
$\sigma_i \tau_j = \delta_{ij}$. Thinking of $\sigma_i$
as a section of $(g')^*\mathcal{L}$ on $X'$
and $\tau_j$ as a section of $(g')^*\mathcal{L}^\vee$ on $X'$,
we conclude that
$$
\sigma_1, \ldots, \sigma_r :
\mathcal{O}_{X'}^{\oplus r}
\longrightarrow
(g')^*\mathcal{E}
$$
is an isomorphism with inverse given by
$$
\tau_1, \ldots, \tau_r :
(g')^*\mathcal{E}
\longrightarrow
\mathcal{O}_{X'}^{\oplus r}
$$
In other words, we see that
$(f')^*f'_*(g')^*\mathcal{E} \cong (g')^*\mathcal{E}$.
But the condition that the cup product is nondegenerate
picks out a retrocompact open subscheme (namely, the locus where a suitable
determinant is nonzero) and the proof is complete.
\end{proof}









\section{Descent of finiteness properties of complexes}
\label{section-descent-finiteness}

\noindent
This section is the analogue of More on Morphisms,
Section \ref{more-morphisms-section-descent-finiteness}
and
Derived Categories of Schemes, Section
\ref{perfect-section-descent-finiteness}.

\begin{lemma}
\label{lemma-pseudo-coherent-descends-fpqc}
Let $S$ be a scheme. Let $\{f_i : X_i \to X\}$ be an fpqc covering of
algebraic spaces over $S$. Let $E \in D_\QCoh(\mathcal{O}_X)$.
Let $m \in \mathbf{Z}$. Then $E$ is $m$-pseudo-coherent if and only if each
$Lf_i^*E$ is $m$-pseudo-coherent.
\end{lemma}

\begin{proof}
Pullback always preserves $m$-pseudo-coherence, see
Cohomology on Sites, Lemma
\ref{sites-cohomology-lemma-pseudo-coherent-pullback}.
Thus it suffices to assume $Lf_i^*E$ is $m$-pseudo-coherent
and to prove that $E$ is $m$-pseudo-coherent.
Then first we may assume $X_i$ is a scheme for all $i$, see
Topologies on Spaces, Lemma \ref{spaces-topologies-lemma-refine-fpqc-schemes}.
Next, choose a  surjective \'etale morphism $U \to X$ where $U$ is a scheme.
Then $U_i = U \times_X X_i$ is a scheme and we obtain an fpqc covering
$\{U_i \to U\}$ of schemes, see
Topologies on Spaces, Lemma
\ref{spaces-topologies-lemma-recognize-fpqc-covering}.
We know the result is true for
$\{U_i \to U\}_{i \in I}$ by the case for schemes, see
Derived Categories of Schemes, Lemma
\ref{perfect-lemma-pseudo-coherent-descends-fpqc}.
On the other hand, the restriction $E|_U$ comes from
an object of $D_\QCoh(\mathcal{O}_U)$ (defined using the Zariski
topology and the ``usual'' structure sheaf of $U$), see
Derived Categories of Spaces, Lemma
\ref{spaces-perfect-lemma-derived-quasi-coherent-small-etale-site}.
The lemma follows as the two notions of pseudo-coherent
(\'etale and Zariski) agree by
Derived Categories of Spaces,
Lemma \ref{spaces-perfect-lemma-descend-pseudo-coherent}.
\end{proof}

\begin{lemma}
\label{lemma-tor-amplitude-descends-fppf}
Let $S$ be a scheme. Let $\{g_i : Y_i \to Y\}$ be an fpqc covering of
algebraic spaces over $S$. Let $f : X \to Y$ be a morphism of algebraic
spaces and set $X_i = Y_i \times_Y X$ with projections $f_i : X_i \to Y_i$
and $g'_i : X_i \to X$. Let $E \in D_\QCoh(\mathcal{O}_X)$.
Let $a, b \in \mathbf{Z}$. Then the following are equivalent
\begin{enumerate}
\item $E$ has tor amplitude in $[a, b]$ as an object of
$D(f^{-1}\mathcal{O}_Y)$, and
\item $L(g'_i)^*E$ has tor amplitude in $[a, b]$ as a object of
$D(f_i^{-1}\mathcal{O}_{Y_i})$ for all $i$.
\end{enumerate}
Also true if ``tor amplitude in $[a, b]$'' is replaced by
``locally finite tor dimension''.
\end{lemma}

\begin{proof}
Pullback preserves ``tor amplitude in $[a, b]$'' by
Derived Categories of Spaces, Lemma
\ref{spaces-perfect-lemma-tor-independence-and-tor-amplitude}
Observe that $Y_i$ and $X$ are tor independent over $Y$
as $Y_i \to Y$ is flat. Let us assume (2) and prove (1).
We can compute tor dimension at stalks, see
Cohomology on Sites, Lemma \ref{sites-cohomology-lemma-tor-amplitude-stalk}
and Properties of Spaces, Theorem
\ref{spaces-properties-theorem-exactness-stalks}.
Let $\overline{x}$ be a geometric point of $X$. Choose
an $i$ and a geometric point $\overline{x}_i$ in $X_i$ with image
$\overline{x}$ in $X$. Then
$$
(L(g_i')^*E)_{\overline{x}_i} =
E_{\overline{x}}
\otimes_{\mathcal{O}_{X, \overline{x}}}^\mathbf{L}
\mathcal{O}_{X_, \overline{x}_i}
$$
Let $\overline{y}_i$ in $Y_i$ and $\overline{y}$ in $Y$
be the image of $\overline{x}_i$ and $\overline{x}$.
Since $X$ and $Y_i$ are tor independent over $Y$, we can apply
More on Algebra, Lemma \ref{more-algebra-lemma-base-change-comparison}
to see that the right hand side of the displayed formula is equal to
$E_{\overline{x}}
\otimes_{\mathcal{O}_{Y, \overline{y}}}^\mathbf{L}
\mathcal{O}_{Y_i, \overline{y}_i}$
in $D(\mathcal{O}_{Y_i, \overline{y}_i})$.
Since we have assume the tor amplitude of this is in
$[a, b]$, we conclude that the tor amplitude of
$E_{\overline{x}}$ in $D(\mathcal{O}_{Y, \overline{y}})$
is in $[a, b]$ by More on Algebra, Lemma
\ref{more-algebra-lemma-flat-descent-tor-amplitude}.
Thus (1) follows.

\medskip\noindent
Using some elementary topology the case
``locally finite tor dimension'' follows too.
\end{proof}

\noindent
The following lemmas do not really belong in this section.

\begin{lemma}
\label{lemma-thickening-pseudo-coherent}
Let $S$ be a scheme. Let $i : X \to X'$ be a finite order thickening of
algebraic spaces. Let $K' \in D(\mathcal{O}_{X'})$ be an object such that
$K = Li^*K'$ is pseudo-coherent. Then $K'$ is pseudo-coherent.
\end{lemma}

\begin{proof}
We first prove $K'$ has quasi-coherent cohomology sheaves; we urge
the reader to skip this part.
To do this, we may reduce to the case of a first order thickening, see
Section \ref{section-thickenings}. Let $\mathcal{I} \subset \mathcal{O}_{X'}$
be the quasi-coherent sheaf of ideals cutting out $X$.
Tensoring the short exact sequence
$$
0 \to \mathcal{I} \to \mathcal{O}_{X'} \to i_*\mathcal{O}_X \to 0
$$
with $K'$ we obtain a distinguished triangle
$$
K' \otimes_{\mathcal{O}_{X'}}^\mathbf{L} \mathcal{I}
\to K' \to
K' \otimes_{\mathcal{O}_{X'}}^\mathbf{L} i_*\mathcal{O}_X
\to
(K' \otimes_{\mathcal{O}_{X'}}^\mathbf{L} \mathcal{I})[1]
$$
Since $i_* = Ri_*$ and since we may view $\mathcal{I}$
as a quasi-coherent $\mathcal{O}_X$-module (as we have a first
order thickening) we may rewrite this as
$$
i_*(K \otimes_{\mathcal{O}_X}^\mathbf{L} \mathcal{I})
\to K' \to
i_*K \to
i_*(K \otimes_{\mathcal{O}_X}^\mathbf{L} \mathcal{I})[1]
$$
Please use Cohomology of Spaces, Lemma
\ref{spaces-cohomology-lemma-projection-formula-finite}
to identify the terms. Since $K$ is in
$D_\QCoh(\mathcal{O}_X)$ we conclude that
$K'$ is in $D_\QCoh(\mathcal{O}_{X'})$; this uses
Derived Categories of Spaces, Lemmas
\ref{spaces-perfect-lemma-pseudo-coherent},
\ref{spaces-perfect-lemma-quasi-coherence-tensor-product}, and
\ref{spaces-perfect-lemma-quasi-coherence-direct-image}.

\medskip\noindent
Assume $K'$ is in $D_\QCoh(\mathcal{O}_{X'})$.
The question is \'etale local on $X'$
hence we may assume $X'$ is affine.
In this case the result follows from the case of schemes
(More on Morphisms, Lemma
\ref{more-morphisms-lemma-thickening-pseudo-coherent}).
The translation into the language of schemes uses
Derived Categories of Spaces, Lemmas
\ref{spaces-perfect-lemma-derived-quasi-coherent-small-etale-site} and
\ref{spaces-perfect-lemma-descend-pseudo-coherent} and
Remark \ref{spaces-perfect-remark-match-total-direct-images}.
\end{proof}

\begin{lemma}
\label{lemma-thickening-relatively-perfect}
Let $S$ be a scheme. Consider a cartesian diagram
$$
\xymatrix{
X \ar[r]_i \ar[d]_f & X' \ar[d]^{f'} \\
Y \ar[r]^j & Y'
}
$$
of algebraic spaces over $S$. Assume $X' \to Y'$ is flat and locally
of finite presentation and $Y \to Y'$ is a finite order thickening.
Let $E' \in D(\mathcal{O}_{X'})$. If $E = Li^*(E')$ is $Y$-perfect,
then $E'$ is $Y'$-perfect.
\end{lemma}

\begin{proof}
Recall that being $Y$-perfect for $E$ means $E$ is
pseudo-coherent and locally has finite tor dimension as a complex
of $f^{-1}\mathcal{O}_Y$-modules
(Definition \ref{definition-relatively-perfect}).
By Lemma \ref{lemma-thickening-pseudo-coherent}
we find that $E'$ is pseudo-coherent.
In particular, $E'$ is in $D_\QCoh(\mathcal{O}_{X'})$, see
Derived Categories of Spaces, Lemma
\ref{spaces-perfect-lemma-pseudo-coherent}.
By Lemma \ref{lemma-affine-locally-rel-perfect}
this reduces us to the case of schemes.
The case of schemes is
More on Morphisms, Lemma
\ref{more-morphisms-lemma-thickening-relatively-perfect}.
\end{proof}

\begin{lemma}
\label{lemma-henselian-relatively-perfect}
Let $(R, I)$ be a pair consisting of a ring and an ideal $I$
contained in the Jacobson radical. Set $S = \Spec(R)$ and $S_0 = \Spec(R/I)$.
Let $X$ be an algebraic space over $R$ whose structure morphism
$f : X \to S$ is proper, flat, and of finite presentation.
Denote $X_0 = S_0 \times_S X$. Let $E \in D(\mathcal{O}_X)$
be pseudo-coherent. If the derived restriction $E_0$ of $E$
to $X_0$ is $S_0$-perfect, then $E$ is $S$-perfect.
\end{lemma}

\begin{proof}
Choose a surjective \'etale morphism $U \to X$ with $U$ affine.
Choose a closed immersion $U \to \mathbf{A}^d_S$.
Set $U_0 = S_0 \times_S U$.
The complex $E_0|_{U_0}$ has tor amplitude
in $[a, b]$ for some $a, b \in \mathbf{Z}$.
Let $\overline{x}$ be a geometric point of $X$.
We will show that the tor amplitude of
$E_{\overline{x}}$ over $R$ is in $[a - d, b]$.
This will finish the proof as the tor amplitude can be
read off from the stalks by
Cohomology on Sites, Lemma \ref{sites-cohomology-lemma-tor-amplitude-stalk}
and Properties of Spaces, Theorem
\ref{spaces-properties-theorem-exactness-stalks}.

\medskip\noindent
Let $x \in |X|$ be the point determined by $\overline{x}$.
Recall that $|X| \to |S|$ is closed (by definition of proper morphisms).
Since $I$ is contained in the Jacobson radical, any nonempty closed
subset of $S$ contains a point of the closed subscheme $S_0$.
Hence we can find a specialization $x \leadsto x_0$ in $|X|$
with $x_0 \in |X_0|$. Choose $u_0 \in U_0$ mapping to $x_0$.
By Decent Spaces, Lemma \ref{decent-spaces-lemma-generalizations-lift-flat}
(or by Decent Spaces, Lemma \ref{decent-spaces-lemma-specialization}
which applies directly to \'etale morphisms)
we find a specialization $u \leadsto u_0$ in $U$
such that $u$ maps to $x$. We may lift $\overline{x}$
to a geometric point $\overline{u}$ of $U$ lying over $u$.
Then we have $E_{\overline{x}} = (E|_U)_{\overline{u}}$.

\medskip\noindent
Write $U = \Spec(A)$. Then $A$ is a flat, finitely presented
$R$-algebra which is a quotient of a polynomial $R$-algebra in
$d$-variables. The restriction $E|_U$ corresponds
(by Derived Categories of Spaces, Lemmas
\ref{spaces-perfect-lemma-pseudo-coherent},
\ref{spaces-perfect-lemma-derived-quasi-coherent-small-etale-site}, and
\ref{spaces-perfect-lemma-descend-pseudo-coherent}
and
Derived Categories of Schemes, Lemma
\ref{perfect-lemma-affine-compare-bounded} and
\ref{perfect-lemma-pseudo-coherent-affine})
to a pseudo-coherent object $K$ of $D(A)$.
Observe that $E_0$ corresponds to $K \otimes_A^\mathbf{L} A/IA$.
Let $\mathfrak q \subset \mathfrak q_0 \subset A$ be the prime
ideals corresponding to $u \leadsto u_0$.
Then
$$
E_{\overline{x}} =
(E|_U)_{\overline{u}} =
E_u \otimes_{\mathcal{O}_{U, u}}^\mathbf{L} \mathcal{O}_{U, \overline{u}} =
K_{\mathfrak q} \otimes_{A_\mathfrak q}^\mathbf{L} A_{\mathfrak q}^{sh}
$$
(some details omitted). Since $A_\mathfrak q \to A_\mathfrak q^{sh}$
is flat, the tor amplitude of this as an $R$-module is the same as
the tor amplitude of $K_\mathfrak q$ as an $R$-module
(More on Algebra, Lemma \ref{more-algebra-lemma-no-change-tor-amplitude}).
Also, $K_{\mathfrak q}$ is a localization of $K_{\mathfrak q_0}$.
Hence it suffices to show that $K_{\mathfrak q_0}$ has tor amplitude in
$[a - d, b]$ as a complex of $R$-modules.

\medskip\noindent
Let $I \subset \mathfrak p_0 \subset R$ be the prime
ideal corresponding to $f(x_0)$. Then we have
\begin{align*}
K \otimes_R^\mathbf{L} \kappa(\mathfrak p_0)
& =
(K \otimes_R^\mathbf{L} R/I) \otimes_{R/I}^\mathbf{L}
\kappa(\mathfrak p_0) \\
& =
(K \otimes_A^\mathbf{L} A/IA) \otimes_{R/I}^\mathbf{L} \kappa(\mathfrak p_0)
\end{align*}
the second equality because $R \to A$ is flat.
By our choice of $a, b$ this complex has cohomology
only in degrees in the interval $[a, b]$.
Thus we may finally apply
More on Algebra, Lemma
\ref{more-algebra-lemma-lift-from-fibre-relatively-perfect}
to $R \to A$, $\mathfrak q_0$, $\mathfrak p_0$ and $K$
to conclude.
\end{proof}







\section{Families of nodal curves}
\label{section-families-nodal}

\noindent
This section is the continuation of
Algebraic Curves, Section \ref{curves-section-families-nodal}.
Please also see that section for our choice
of terminology.

\medskip\noindent
The property ``at-worst-nodal of relative dimension $1$''
of morphisms of schemes is \'etale local on the source-and-target, see
Descent, Lemma \ref{descent-lemma-etale-local-source-target}
and
Algebraic Curves, Lemmas \ref{curves-lemma-nodal-family-etale-local-source},
\ref{curves-lemma-nodal-family-fpqc-local-target}, and
\ref{curves-lemma-nodal-family-postcompose-etale}.
It is also stable under base change and fpqc local on the target, see
Algebraic Curves, Lemmas \ref{curves-lemma-base-change-nodal-family} and
\ref{curves-lemma-nodal-family-fpqc-local-target}. Hence, by
Morphisms of Spaces, Lemma \ref{spaces-morphisms-lemma-local-source-target}
we may define the notion of an at-worst-nodal morphism of relative
dimension $1$ for algebraic spaces as follows and it agrees with the
already existing notion defined in
Morphisms of Spaces, Section \ref{spaces-morphisms-section-representable}
when the morphism is representable.

\begin{definition}
\label{definition-nodal-family}
Let $S$ be a scheme.
Let $f : X \to Y$ be a morphism of algebraic spaces over $S$.
We say $f$ is {\it at-worst-nodal of relative dimension $1$}
if the equivalent conditions of
Morphisms of Spaces, Lemma \ref{spaces-morphisms-lemma-local-source-target}
hold with $\mathcal{P} =$``at-worst-nodal of relative dimension $1$''.
\end{definition}

\begin{lemma}
\label{lemma-base-change-nodal}
The property of being at-worst-nodal of relative dimension $1$
is preserved under base change.
\end{lemma}

\begin{proof}
See
Morphisms of Spaces, Remark \ref{spaces-morphisms-remark-base-change-P}
and
Algebraic Curves, Lemma \ref{curves-lemma-base-change-nodal-family}.
\end{proof}

\begin{lemma}
\label{lemma-nodal-local}
Let $S$ be a scheme.
Let $f : X \to Y$ be a morphism of algebraic spaces over $S$.
The following are equivalent:
\begin{enumerate}
\item $f$ is at-worst-nodal of relative dimension $1$,
\item for every scheme $Z$ and any morphism $Z \to Y$ the morphism
$Z \times_Y X \to Z$ is at-worst-nodal of relative dimension $1$,
\item for every affine scheme $Z$ and any morphism
$Z \to Y$ the morphism $Z \times_Y X \to Z$ is
at-worst-nodal of relative dimension $1$,
\item there exists a scheme $V$ and a surjective \'etale morphism
$V \to Y$ such that $V \times_Y X \to V$ is
at-worst-nodal of relative dimension $1$,
\item there exists a scheme $U$ and a surjective \'etale morphism
$\varphi : U \to X$ such that the composition $f \circ \varphi$
is at-worst-nodal of relative dimension $1$,
\item for every commutative diagram
$$
\xymatrix{
U \ar[d] \ar[r] & V \ar[d] \\
X \ar[r] & Y
}
$$
where $U$, $V$ are schemes and the vertical arrows are \'etale
the top horizontal arrow is at-worst-nodal of relative dimension $1$,
\item there exists a commutative diagram
$$
\xymatrix{
U \ar[d] \ar[r] & V \ar[d] \\
X \ar[r] & Y
}
$$
where $U$, $V$ are schemes, the vertical arrows are \'etale, and
$U \to X$ is surjective such that the top horizontal arrow is
at-worst-nodal of relative dimension $1$, and
\item there exist Zariski coverings $Y = \bigcup_{i \in I} Y_i$,
and $f^{-1}(Y_i) = \bigcup X_{ij}$ such that
each morphism $X_{ij} \to Y_i$ is
at-worst-nodal of relative dimension $1$.
\end{enumerate}
\end{lemma}

\begin{proof}
Omitted.
\end{proof}

\noindent
The following lemma tells us that we can check whether a morphism
is at-worst-nodal of relative dimension $1$ on the fibres.

\begin{lemma}
\label{lemma-locus-where-nodal}
Let $S$ be a scheme.
Let $f : X \to Y$ be a morphism of algebraic spaces over $S$
which is flat and locally of finite presentation. Then there
is a maximal open subspace $X' \subset X$ such that $f|_{X'} : X' \to Y$
is at-worst-nodal of relative dimension $1$. Moreover, formation
of $X'$ commutes with arbitrary base change.
\end{lemma}

\begin{proof}
Choose a commutative diagram
$$
\xymatrix{
U \ar[d] \ar[r]_h & V \ar[d] \\
X \ar[r]^f & Y
}
$$
where $U$, $V$ are schemes, the vertical arrows are \'etale, and
$U \to X$ is surjective. By the lemma for the case of schemes
(Algebraic Curves, Lemma \ref{curves-lemma-locus-where-nodal})
we find a maximal open subscheme $U' \subset U$
such that $h|_{U'} : U' \to V$ is at-worst-nodal of relative dimension $1$
and such that formation of $U'$ commutes with base change.
Let $X' \subset X$ be the open subspace whose points correspond
to the open subset $\Im(|U'| \to |X|)$.
By Lemma \ref{lemma-nodal-local} we see that $X' \to Y$ is
at-worst-nodal of relative dimension $1$ and that $X'$ is the
largest open subspace with this property (this also implies
that $U'$ is the inverse image of $X'$ in $U$, but we do
not need this). Since the same is true after base change
the proof is complete.
\end{proof}




\section{The resolution property}
\label{section-resolution-property}

\noindent
We continue the discussion in Derived Categories of Spaces, Section
\ref{spaces-perfect-section-resolution-property}.

\begin{situation}
\label{situation-descend-resolution-property}
Let $S$ be a scheme. Let $X$ be a quasi-compact and quasi-separated
algebraic space over $S$. Let $V \to X$ be a surjective \'etale
morphism where $V$ is an affine scheme (such a thing exists by
Properties of Spaces, Lemma
\ref{spaces-properties-lemma-quasi-compact-affine-cover}).
Choose a commutative diagram
$$
\xymatrix{
V \ar[rd]_\varphi \ar[rr]_j & & Y \ar[ld]^\pi \\
& X
}
$$
where $j$ is an open immersion and $\pi$ is a finite morphism
of algebraic spaces (such a diagram exists by
Lemma \ref{lemma-quasi-finite-separated-pass-through-finite}).
Let $\mathcal{I} \subset \mathcal{O}_Y$ be a finite type
quasi-coherent sheaf of ideals on $Y$ with
$V(\mathcal{I}) = Y \setminus j(V)$ (such a sheaf of ideals exists
by Limits of Spaces, Lemma
\ref{spaces-limits-lemma-quasi-coherent-finite-type-ideals}).
\end{situation}

\begin{lemma}
\label{lemma-surjection-in-Noetherian-case}
In Situation \ref{situation-descend-resolution-property}, assume
$X$ is Noetherian. Then for any coherent $\mathcal{O}_X$-module
$\mathcal{F}$ there exist $r \geq 0$,
integers $n_1, \ldots, n_r \geq 0$, and a surjection
$$
\bigoplus\nolimits_{i = 1, \ldots, r} \pi_*(\mathcal{I}^{n_i})
\longrightarrow \mathcal{F}
$$
of $\mathcal{O}_X$-modules.
\end{lemma}

\begin{proof}
Denote $\omega_{Y/X}$ the coherent $\mathcal{O}_Y$-module such that there
is an isomorphism
$$
\pi_*\omega_{Y/X} \cong
\SheafHom_{\mathcal{O}_X}(\pi_*\mathcal{O}_Y, \mathcal{O}_X)
$$
of $\pi_*\mathcal{O}_Y$-modules, see Morphisms of Spaces, Lemma
\ref{spaces-morphisms-lemma-affine-equivalence-modules} and
Descent on Spaces, Lemma \ref{spaces-descent-lemma-finite-over-finite-module}.
The canonical map $\mathcal{O}_X \to \pi_*\mathcal{O}_Y$
produces a canonical map
$$
\text{Tr}_\pi : \pi_*\omega_{Y/X} \longrightarrow \mathcal{O}_X
$$
Since $V$ is Noetherian affine we may choose sections
$$
s_1, \ldots, s_r \in
\Gamma(V, \pi^*\mathcal{F} \otimes_{\mathcal{O}_Y} \omega_{Y/X})
$$
generating the coherent module
$\pi^*\mathcal{F} \otimes_{\mathcal{O}_X} \omega_{Y/X}$
over $V$. By
Cohomology of Spaces, Lemma \ref{spaces-cohomology-lemma-homs-over-open}
we can choose integers $n_i \geq 0$ such that $s_i$
extends to a map
$s_i' : \mathcal{I}^{n_i} \to
\pi^*\mathcal{F} \otimes_{\mathcal{O}_Y} \omega_{Y/X}$.
Pushing to $X$ we obtain maps
$$
\sigma_i :
\pi_*\mathcal{I}^{n_i}
\xrightarrow{\pi_*s'_i}
\pi_*(\pi^*\mathcal{F} \otimes_{\mathcal{O}_Y} \omega_{Y/X}) =
\mathcal{F} \otimes_{\mathcal{O}_X} \pi_*\omega_{Y/X}
\xrightarrow{\text{Tr}_\pi} \mathcal{F}
$$
where the equality sign is
Cohomology of Spaces, Lemma \ref{spaces-cohomology-lemma-finite-rings}.
To finish the proof we will show that the sum of these maps is surjective.

\medskip\noindent
Let $x \in |X|$ be a point of $X$. Let $v \in |V|$ be a point mapping to $x$.
We may choose an \'etale neighbourhood $(U, u) \to (X, x)$ such that
$$
U \times_X Y = W \coprod W'
$$
(disjoint union of algebraic spaces) such that $W \to U$ is an
isomorphism and such that the unique point $w \in W$ lying over
$u$ maps to $v$ in $V \subset Y$. To see this is true use
Lemma \ref{lemma-etale-splits-off-just-one-quasi-finite-part} and
\'Etale Morphisms, Lemma \ref{etale-lemma-etale-etale-local}.
After shrinking $U$ further if necessary we may assume $W$ maps into
$V \subset Y$ by the projection. Since the formation of $\omega_{Y/X}$
commutes with \'etale localization we see that
$$
\pi_*\omega_{Y/X}|_U = (\pi|_W)_*\omega_{W/U} \oplus (\pi|_{W'})_*\omega_{W'/U}
$$
We have $(\pi|_W)_*\omega_{W/U} = \mathcal{O}_U$ and this isomorphism
is given by the trace map $\text{Tr}_\pi|_U$ restricted
to the first summand in the decomposition above. Since $W$ maps
into $V$ we see that $\mathcal{I}^{n_i}|_W = \mathcal{O}_W$.
Hence
$$
\pi_*(\mathcal{I}^{n_i})|_U = \mathcal{O}_U \oplus
(W' \to U)_*(\mathcal{I}^{n_i}|_{W'})
$$
Chasing diagrams the reader sees (details omitted) that $\sigma_i|_U$
on the summand $\mathcal{O}_U$ is the map $\mathcal{O}_U \to \mathcal{F}$
corresponding to the section
$$
s_i|_W  \in
\Gamma(W,\pi^*\mathcal{F} \otimes_{\mathcal{O}_Y} \omega_{Y/X})
=
\Gamma(W, \mathcal{F}|_W \otimes_{\mathcal{O}_W} \omega_{W/U})
=
\Gamma(U, \mathcal{F})
$$
Since the sections $s_i$ generate the module
$\pi^*\mathcal{F} \otimes_{\mathcal{O}_Y} \omega_{Y/X}$
over $V$ and since $W$ maps into $V$ we conclude
that the restriction of $\bigoplus \sigma_i$ to $U$ is surjective.
Since $x$ was an arbitrary point the proof is complete.
\end{proof}

\begin{lemma}
\label{lemma-resolution-property-in-Noetherian-case}
In Situation \ref{situation-descend-resolution-property}, assume
$X$ is Noetherian. Then $X$ has the resolution property if and only
if $\pi_*\mathcal{I}$ is the quotient of a finite locally free
$\mathcal{O}_X$-module.
\end{lemma}

\begin{proof}
The module $\pi_*\mathcal{I}$ is coherent by
Cohomology of Spaces, Lemma
\ref{spaces-cohomology-lemma-finite-pushforward-coherent}.
Hence if $X$ has the resolution property then $\pi_*\mathcal{I}$
is the quotient of a finite locally free $\mathcal{O}_X$-module.
Conversely, assume given a surjection $\mathcal{E} \to \pi_*\mathcal{I}$
for some finite locally free $\mathcal{O}_X$-module $\mathcal{E}$.
Observe that for all $n \geq 1$ there is a surjection
$$
\pi_*\mathcal{I} \otimes_{\mathcal{O}_X} \pi_*\mathcal{I}^n
\longrightarrow
\pi_*\mathcal{I}^{n + 1}
$$
Hence $\mathcal{E}^{\otimes n}$ surjects onto
$\pi_*\mathcal{I}^n$ for all $n \geq 1$. We conclude that $X$
has the resolution property if we combine this with the
result of Lemma \ref{lemma-surjection-in-Noetherian-case}.
\end{proof}

\begin{lemma}
\label{lemma-resolution-property-one-sheaf}
In Situation \ref{situation-descend-resolution-property},
the algebraic space $X$ has the resolution property
if and only if $\pi_*\mathcal{I}$ is the quotient of a
finite locally free $\mathcal{O}_X$-module.
\end{lemma}

\begin{proof}
The pushforward $\pi_*\mathcal{G}$ of a finite type
quasi-coherent $\mathcal{O}_Y$-module $\mathcal{G}$ is a finite
type quasi-coherent $\mathcal{O}_X$-module by
Descent on Spaces, Lemma
\ref{spaces-descent-lemma-finite-over-finite-module}.
In particular, if $X$ has the resolution property, then $\pi_*\mathcal{I}$
is the quotient of a finite locally free $\mathcal{O}_X$-module by
Derived Categories of Spaces, Definition
\ref{spaces-perfect-definition-resolution-property}.

\medskip\noindent
Assume that we have a surjection $\mathcal{E} \to \pi_*\mathcal{I}$
for some finite locally free $\mathcal{O}_X$-module $\mathcal{E}$.
In the rest of the proof we show that $X$ has the resolution property
by reducing to the Noetherian case handled in
Lemma \ref{lemma-resolution-property-in-Noetherian-case}.
We suggest the reader skip the rest of the proof.

\medskip\noindent
A first reduction is that we may view $X$ as an algebraic space over
$\Spec(\mathbf{Z})$, see
Spaces, Definition \ref{spaces-definition-base-change}.
(This doesn't affect the conditions nor the conclusion of the lemma.)

\medskip\noindent
By Limits of Spaces, Lemma
\ref{spaces-limits-lemma-finite-in-finite-and-finite-presentation}
we can write $Y = \lim Y_i$ with $Y_i$ finite and of finite
presentation over $X$ and where the transition maps are closed
immersions. Consider the closed subspace $Z = V(\mathcal{I})$ of $Y$.
Since $\mathcal{I}$ is of finite type, the morphism $Z \to Y$ is of
finite presentation. Hence we can find an $i$ and a morphism
$Z_i \to Y_i$ of finite presentation whose base change to $Y$
is $Z \to Y$, see Limits of Spaces, Lemma
\ref{spaces-limits-lemma-descend-finite-presentation}.
For $i' \geq i$ denote $Z_{i'} = Z_i \times_{Y_i} Y_{i'}$.
After increasing $i$ we may assume $Z_i \to Y_i$ is a
closed immersion (of finite presentation), see Limits of Spaces, Lemma
\ref{spaces-limits-lemma-descend-closed-immersion}. Denote
$\mathcal{I}_i \subset \mathcal{O}_{Y_i}$ the ideal
sheaf of $Z_i$ and denote $\pi_i : Y_i \to X$ the structure morphism.
Similarly for $i' \geq i$. Since $Z = \lim_{i' \geq i} Z_{i'}$ we have
$$
\pi_*\mathcal{I} = \colim \pi_{i', *}\mathcal{I}_{i'}
$$
The transition maps in the system are all surjective
as follows from the surjectivity of the maps
$\pi_{i, *}\mathcal{O}_{Y_i} \to \pi_{i', *}\mathcal{O}_{Y_{i'}}$
and the fact that $Z_{i'} = Z_i \times_{Y_i} Y_{i'}$.
By Cohomology of Spaces, Lemma
\ref{spaces-cohomology-lemma-finite-presentation-quasi-compact-colimit}
for some $i' \geq i$ the map $\mathcal{E} \to \pi_*\mathcal{I}$
lifts to a map $\mathcal{E} \to \pi_{i', *}\mathcal{I}_{i'}$.
After increasing $i'$ this map
$\mathcal{E} \to \pi_{i', *}\mathcal{I}_{i'}$
becomes surjective (since if not the colimit of the cokernels,
having surjective transition maps, is nonzero).
This reduces us to the case discussed in the next paragraph.

\medskip\noindent
Assume $X$ is an algebraic space over $\mathbf{Z}$ and that
$Y \to X$ is of finite presentation. By absolute Noetherian
approximation we can write $X = \lim X_i$ as a directed limit,
where each $X_i$ is a quasi-separated algebraic space of finite type
over $\mathbf{Z}$ and the transition morphisms are affine, see
Limits of Spaces, Proposition \ref{spaces-limits-proposition-approximate}.
Since $\pi : Y \to X$ is of finite presentation we can find an $i$
and a morphism $\pi_i : Y_i \to X_i$ of finite presentation
whose base change to $X$ is $\pi$, see
Limits of Spaces, Lemma
\ref{spaces-limits-lemma-descend-finite-presentation}.
After increasing $i$ we may assume $\pi_i$ is finite, see
Limits of Spaces, Lemma
\ref{spaces-limits-lemma-descend-finite}.
Next, we may assume there exists a finite locally
free $\mathcal{O}_{X_i}$-module $\mathcal{E}_i$
whose pullback to $X$ is $\mathcal{E}$, see
Limits of Spaces, Lemma \ref{spaces-limits-lemma-descend-invertible-modules}.
We may also assume there is a map
$\mathcal{E}_i \to \pi_{i, *}\mathcal{O}_{Y_i}$
whose pullback to $X$ is the composition
$\mathcal{E} \to \pi_*\mathcal{I} \to \pi_*\mathcal{O}_Y$, see
Limits of Spaces, Lemma
\ref{spaces-limits-lemma-descend-modules-finite-presentation}.
The cokernel
$$
\mathcal{E}_i \to \pi_{i, *}\mathcal{O}_{Y_i} \to \mathcal{Q}_i \to 0
$$
is a coherent $\mathcal{O}_{Y_i}$-module whose pullback to $X$
is the (finitely presented) cokernel $\mathcal{Q}$ of the map
$\mathcal{E} \to \pi_*\mathcal{O}_Y$. In other words, we have
$\mathcal{Q} = \pi_*(\mathcal{O}_Y/\mathcal{I})$. Consider
the map
$$
\mathcal{E}_i
\otimes_{\mathcal{O}_{X_i}}
\pi_{i, *} \mathcal{O}_{Y_i}
\longrightarrow
\pi_{i, *}\mathcal{O}_{Y_i}
\otimes_{\mathcal{O}_{X_i}}
\pi_{i, *} \mathcal{O}_{Y_i}
\to
\pi_{i, *} \mathcal{O}_{Y_i}
\to
\mathcal{Q}_i
$$
where the second arrow is given by the algebra structure
on $\pi_{i, *}\mathcal{O}_{Y_i}$. The pullback of this map
to $Y$ is zero because the image of $\mathcal{E} \to \pi_*\mathcal{O}_Y$
is the ideal $\pi_*\mathcal{I}$. Hence by
Limits of Spaces, Lemma
\ref{spaces-limits-lemma-descend-modules-finite-presentation}
after increasing $i$ we may assume the displayed composition
is zero. This exactly means that the imag of
$\mathcal{E}_i \to \pi_{i, *}\mathcal{O}_{Y_i}$
is of the form $\pi_{i, *}\mathcal{I}_i$
for some coherent ideal sheaf $\mathcal{I}_i \subset \mathcal{O}_{Y_i}$.
Since $\mathcal{E}_i \to \pi_{i, *}\mathcal{O}_{Y_i}$
pulls back to $\mathcal{E} \to \pi_*\mathcal{O}_Y$ we
see that the pullback of $\mathcal{I}_i$ to $Y$ generates $\mathcal{I}$.
Denote $V_i \subset Y_i$ the open subspace whose complement
is $V(\mathcal{I}_i) \subset Y_i$. Then $V$ is the inverse image
of $V_i$ by the comments above. After increasing $i$ we may assume
that $V_i$ is affine and that $\pi_i|_{V_i} : V_i \to X_i$ is 
\'etale, see
Limits of Spaces, Lemmas
\ref{spaces-limits-lemma-limit-is-affine} and
\ref{spaces-limits-lemma-descend-etale}.
Having said all of this, we may apply
Lemma \ref{lemma-resolution-property-in-Noetherian-case}
to conclude that $X_i$ has the
resolution property. Since $X \to X_i$ is affine
we conclude that $X$ has the resolution property too
by Derived Categories of Spaces, Lemma
\ref{spaces-perfect-lemma-resolution-property-goes-up-affine}.
\end{proof}

\begin{lemma}
\label{lemma-resolution-property-descends}
Let $S$ be a scheme.
Let $X = \lim X_i$ be a limit of a direct system of quasi-compact
and quasi-separated algebraic spaces over $S$ with affine transition morphisms.
Then $X$ has the resolution property if and only if $X_i$ has
the resolution properties for some $i$.
\end{lemma}

\begin{proof}
If $X_i$ has the resolution property, then $X$ does by
Derived Categories of Spaces, Lemma
\ref{spaces-perfect-lemma-resolution-property-goes-up-affine}.
Assume $X$ has the resolution property.
Choose $i \in I$. We may choose an affine scheme $V_i$ and a surjective
\'etale morphism $V_i \to X_i$ 
(Properties of Spaces, Lemma
\ref{spaces-properties-lemma-quasi-compact-affine-cover}).
We may choose an embedding $j : V_i \to Y_i$
with $Y_i$ finite and finitely presented over $X_i$
(Lemma \ref{lemma-quasi-finite-separated-pass-through-finite-addendum}).
We may choose a finite type
quasi-coherent ideal $\mathcal{I}_i \subset \mathcal{O}_{Y_i}$
such that $V_i = Y_i \setminus V(\mathcal{I}_i)$
(Limits of Spaces, Lemma
\ref{spaces-limits-lemma-quasi-coherent-finite-type-ideals}).
Denote $V \to Y \to X$ the base changes of $V_i \to Y_i \to X_i$
to $X$. Denote $\mathcal{I} \subset \mathcal{O}_Y$ the pullback
of the ideal $\mathcal{I}_i$. By the easy direction of
Lemma \ref{lemma-resolution-property-one-sheaf}
there exists a finite locally free $\mathcal{O}_X$-module
$\mathcal{E}$ and a surjection $\mathcal{E} \to \pi_*\mathcal{I}$.
Note that since $\pi_i : Y_i \to X_i$ is finite and of finite presentation
we also have that $\pi : Y \to X$ is finite and of finite presentation
and that the $\mathcal{O}_{X_i}$-modules
$\pi_{i, *}\mathcal{O}_{Y_i}$ and
$\pi_{i, *}(\mathcal{O}_{Y_i}/\mathcal{I}_i)$
are of finite presentation and pullback to $X$ to give
$\pi_*\mathcal{O}_Y$ and
$\pi_*(\mathcal{O}_Y/\mathcal{I})$. Thus by Limits of Spaces, Lemma
\ref{spaces-limits-lemma-descend-modules-finite-presentation}
after increasing $i$ we can find a finite locally free
$\mathcal{O}_{X_i}$-module $\mathcal{E}_i$
and a map $\mathcal{E}_i \to \pi_{i, *}\mathcal{O}_{Y_i}$
whose base change to $X$ recovers the composition
$\mathcal{E} \to \pi_*\mathcal{I} \to \pi_*\mathcal{O}_Y$.
The pullbacks of the finitely presented $\mathcal{O}_{X_i}$-modules
$\Coker(\mathcal{E}_i \to \pi_{i, *}\mathcal{O}_{Y_i})$
and $\pi_{i, *}(\mathcal{O}_{Y_i}/\mathcal{I}_i)$
to $X$ agree as quotients of $\pi_*\mathcal{O}_Y$.
Hence by Limits of Spaces, Lemma
\ref{spaces-limits-lemma-descend-modules-finite-presentation}
we may assume that these agree, in other words that
the image of $\mathcal{E}_i \to \pi_{i, *}\mathcal{O}_{X_i}$
is equal to $\pi_{i, *}\mathcal{I}_i$.
Then we conclude that $X_i$ has the resolution property by
Lemma \ref{lemma-resolution-property-one-sheaf}.
\end{proof}

\begin{lemma}
\label{lemma-resolution-property-affine-diagonal}
Let $S$ be a scheme.
Let $X$ be a quasi-compact and quasi-separated algebraic space with
the resolution property. Then $X$ has affine diagonal over $\mathbf{Z}$
(as in Properties of Spaces, Definition
\ref{spaces-properties-definition-separated}).
\end{lemma}

\begin{proof}
We could prove this as in the case of schemes, but instead we will
deduce the lemma from the case of schemes.
First, we may and do assume $S = \Spec(\mathbf{Z})$.
Next, we choose a scheme $Y$ and a surjective integral morphism
$f : Y \to X$, see
Decent Spaces, Lemma \ref{decent-spaces-lemma-there-is-a-scheme-integral-over}.
Then $f$ is affine, hence $Y$ has the resolution property by
Derived Categories of Spaces, Lemma
\ref{spaces-perfect-lemma-resolution-property-goes-up-affine}.
Hence by the case of schemes, the scheme $Y$ has affine diagonal, see
Derived Categories of Schemes, Lemma
\ref{perfect-lemma-resolution-property-affine-diagonal}.
Next, we consider the commutative diagram
$$
\xymatrix{
Y \ar[d] \ar[rr]_{\Delta_Y} & & Y \times_{\mathbf{Z}} Y \ar[d] \\
X \ar[rr]^{\Delta_X} & & X \times_{\mathbf{Z}} X
}
$$
Observe that the right vertical arrow is integral, in particular
affine. Let $W \to X \times_{\mathbf{Z}} X$ be a morphism with $W$ affine.
Then we see that
$$
Y \times_{X \times_{\mathbf{Z}} X} W =
Y \times_{\Delta_Y, Y \times_{\mathbf{Z}} Y}
(Y \times_{\mathbf{Z}} Y) \times_{X \times_{\mathbf{Z}} X} W
$$
is affine. On the other hand, $Y \to X$ is integral and surjective
hence
$$
Y \times_{X \times_{\mathbf{Z}} X} W
\longrightarrow
X \times_{X \times_{\mathbf{Z}} X} W
$$
is integral surjective as the base change of $Y \to X$ to $W$.
We conclude that the target of this arrow is affine by
Limits of Spaces, Proposition \ref{spaces-limits-proposition-affine}.
It follows that $\Delta_X$ is affine as desired.
\end{proof}







\section{Blowing up and the resolution property}
\label{section-resolution-property-by-blowing-up}

\noindent
We prove that the resolution property is satisfied after a blowing up.

\begin{lemma}
\label{lemma-blow-up-resolution-property}
Let $S$ be a scheme. Let $X$ be a quasi-compact and quasi-separated
algebraic space over $S$. Assume that $|X|$ has finitely many irreducible
components. There exists a dense quasi-compact open $U \subset X$
and a $U$-admissible blowing up $X' \to X$ such that the algebraic space
$X'$ has the resolution property.
\end{lemma}

\begin{proof}
By Limits of Spaces, Lemma
\ref{spaces-limits-lemma-there-is-a-scheme-finite-over-refined}
there exists a surjective, finite, and finitely presented morphism
$f : Y \to X$ where $Y$ is a scheme and a quasi-compact dense open
$U \subset X$ such that $f^{-1}(U) \to U$ is finite \'etale.
By More on Morphisms, Lemma \ref{more-morphisms-lemma-blow-up-ample-family}
there is a quasi-compact dense open $V \subset Y$ and a $V$-admissible
blowing up $Y' \to Y$ such that $Y'$ has an ample family of invertible
modules. After shrinking $U$ we may assume that $f^{-1}(U) \subset V$
(details omitted). Hence $f' : Y' \to X$ is finite \'etale over $U$
and in particular, the morphism $(f')^{-1}(U) \to U$ is finite locally free.
By Lemma \ref{lemma-finite-after-blowing-up} there is a $U$-admissible
blowing up $X' \to X$ such that the strict transform $Y''$ of $Y'$
is finite locally free over $X'$. Picture
$$
\xymatrix{
Y'' \ar[d] \ar[r]_g &
Y' \ar[r] &
Y \ar[d] \\
X' \ar[rr] & & X
}
$$
Since $g : Y'' \to Y'$ is a blowing up
(Divisors on Spaces, Lemma \ref{spaces-divisors-lemma-strict-transform})
in the inverse image of the center of $X' \to X$, we see that
$g : Y'' \to Y'$ is projective and that there exists some 
$g$-ample invertible module on $Y''$. Hence by
More on Morphisms, Lemma \ref{more-morphisms-lemma-ample-family-ample-relative}
we see that $Y''$ has an ample family of invertible modules.
Hence $Y''$ has the resolution property, see
Derived Categories of Schemes, Lemma
\ref{perfect-lemma-resolution-property-ample-family}.
We conclude that $X'$ has the resolution property
by
Derived Categories of Spaces, Lemma
\ref{spaces-perfect-lemma-resolution-property-goes-down-finite-flat}.
\end{proof}

\begin{lemma}
\label{lemma-blow-up-resolution-property-filtered}
Let $S$ be a scheme. Let $X$ be a quasi-compact and quasi-separated
algebraic space over $S$. There exists a $t \geq 0$ and closed
subspaces
$$
X \supset Z_0 \supset Z_1 \supset \ldots \supset Z_t = \emptyset
$$
such that $Z_i \to X$ is of finite presentation,
$Z_0 \subset X$ is a thickening, and for each $i = 0, \ldots t - 1$
there exists a $(Z_i \setminus Z_{i - 1})$-admissible
blowing up $Z'_i \to Z_i$ such that $Z'_i$ has the resolution property.
\end{lemma}

\begin{proof}
In this paragraph we use absolute Noetherian approximation to reduce to the
case of algebraic spaces of finite presentation over $\Spec(\mathbf{Z})$.
We may view $X$ as an algebraic space over $\Spec(\mathbf{Z})$, see
Spaces, Definition \ref{spaces-definition-base-change} and
Properties of Spaces, Definition \ref{spaces-properties-definition-separated}.
Thus we may apply Limits of Spaces, Proposition
\ref{spaces-limits-proposition-approximate}.
It follows that we can find an affine morphism $X \to X_0$
with $X_0$ of finite presentation over $\mathbf{Z}$.
If we can prove the lemma for $X_0$, then we can pull back
the stratification and the centers of the blowing ups to $X$
and get the result for $X$; this uses that the resolution property
goes up along affine morphisms (Derived Categories of Spaces,
Lemma \ref{spaces-perfect-lemma-resolution-property-goes-up-affine}) and
that the strict transform of an affine morphism is affine -- details omitted.
This reduces us to the case discussed in the next paragraph.

\medskip\noindent
Assume $X$ is of finite presentation over $\mathbf{Z}$.
Then $X$ is Noetherian and $|X|$ is a Noetherian topological
space (with finitely many irreducible components) of finite dimension.
Hence we may use induction on $\dim(|X|)$.
By Lemma \ref{lemma-blow-up-resolution-property}
there exists a dense open $U \subset X$ and a $U$-admissible
blowing up $X' \to X$ such that $X'$ has the resolution property.
Set $Z_0 = X$ and let $Z_1 \subset X$ be the reduced closed subspace
with $|Z_1| = |X| \setminus |U|$.
By induction we find an integer $t \geq 0$ and a filtration
$$
Z_1 \supset Z_{1, 0} \supset Z_{1, 1} \supset \ldots
\supset Z_{1, t} = \emptyset
$$
by closed subspaces, where $Z_{1, 0} \to Z_1$ is a thickening
and there exist $(Z_{1, i} \setminus Z_{1, i + 1})$-admissible
blowing ups $Z'_{1, i} \to Z_{1, i}$ such that $Z'_{1, i}$
has the resolution property. Since $Z_1$ is reduced, we have $Z_1 = Z_{1, 0}$.
Hence we can set $Z_i = Z_{1, i - 1}$ and $Z'_i = Z'_{1, i - 1}$
for $i \geq 1$ and the lemma is proved.
\end{proof}










\begin{multicols}{2}[\section{Autres chapitres}]
\noindent
Préliminaires
\begin{enumerate}
\item \hyperref[introduction-section-phantom]{Introduction}
\item \hyperref[conventions-section-phantom]{Conventions}
\item \hyperref[sets-section-phantom]{Théorie des ensembles}
\item \hyperref[categories-section-phantom]{Catégories}
\item \hyperref[topology-section-phantom]{Topologie}
\item \hyperref[sheaves-section-phantom]{Faisceaux sur les espaces}
\item \hyperref[sites-section-phantom]{Sites et faisceaux}
\item \hyperref[stacks-section-phantom]{Champs}
\item \hyperref[fields-section-phantom]{Corps}
\item \hyperref[algebra-section-phantom]{Algèbre commutative}
\item \hyperref[brauer-section-phantom]{Groupes de Brauer}
\item \hyperref[homology-section-phantom]{Algèbre homologique}
\item \hyperref[derived-section-phantom]{Catégories dérivées}
\item \hyperref[simplicial-section-phantom]{Méthodes simpliciales}
\item \hyperref[more-algebra-section-phantom]{Compléments d'algèbre}
\item \hyperref[smoothing-section-phantom]{Lissage des morphismes d'anneaux}
\item \hyperref[modules-section-phantom]{Faisceaux de modules}
\item \hyperref[sites-modules-section-phantom]{Modules sur les sites}
\item \hyperref[injectives-section-phantom]{Objets injectifs}
\item \hyperref[cohomology-section-phantom]{Cohomologie des faisceaux}
\item \hyperref[sites-cohomology-section-phantom]{Cohomologie sur les sites}
\item \hyperref[dga-section-phantom]{Algèbre différentielle graduée}
\item \hyperref[dpa-section-phantom]{Algèbre à puissances divisées}
\item \hyperref[sdga-section-phantom]{Faisceaux différentiels gradués}
\item \hyperref[hypercovering-section-phantom]{Hyperrecouvrements}
\end{enumerate}
Schémas
\begin{enumerate}
\setcounter{enumi}{25}
\item \hyperref[schemes-section-phantom]{Schémas}
\item \hyperref[constructions-section-phantom]{Constructions de schémas}
\item \hyperref[properties-section-phantom]{Propriétés des schémas}
\item \hyperref[morphisms-section-phantom]{Morphismes de schémas}
\item \hyperref[coherent-section-phantom]{Cohomologie des schémas}
\item \hyperref[divisors-section-phantom]{Diviseurs}
\item \hyperref[limits-section-phantom]{Limites de schémas}
\item \hyperref[varieties-section-phantom]{Variétés}
\item \hyperref[topologies-section-phantom]{Topologies sur les schémas}
\item \hyperref[descent-section-phantom]{Descente}
\item \hyperref[perfect-section-phantom]{Catégories dérivées de schémas}
\item \hyperref[more-morphisms-section-phantom]{Compléments sur les morphismes}
\item \hyperref[flat-section-phantom]{Compléments sur la platitude}
\item \hyperref[groupoids-section-phantom]{Schémas en groupoïdes}
\item \hyperref[more-groupoids-section-phantom]{Compléments sur les schémas en groupoïdes}
\item \hyperref[etale-section-phantom]{Morphismes étales de schémas}
\end{enumerate}
Thèmes de théorie des schémas
\begin{enumerate}
\setcounter{enumi}{41}
\item \hyperref[chow-section-phantom]{Homologie de Chow}
\item \hyperref[intersection-section-phantom]{Théorie de l'intersection}
\item \hyperref[pic-section-phantom]{Schémas de Picard des courbes}
\item \hyperref[weil-section-phantom]{Théories cohomologiques de Weil}
\item \hyperref[adequate-section-phantom]{Modules adéquats}
\item \hyperref[dualizing-section-phantom]{Complexes dualisants}
\item \hyperref[duality-section-phantom]{Dualité pour les schémas}
\item \hyperref[discriminant-section-phantom]{Discriminants et différentes}
\item \hyperref[derham-section-phantom]{Cohomologie de de Rham}
\item \hyperref[local-cohomology-section-phantom]{Cohomologie locale}
\item \hyperref[algebraization-section-phantom]{Géométrie algébrique et formelle}
\item \hyperref[curves-section-phantom]{Courbes algébriques}
\item \hyperref[resolve-section-phantom]{Résolution des surfaces}
\item \hyperref[models-section-phantom]{Réduction semi-stable}
\item \hyperref[functors-section-phantom]{Foncteurs et morphismes}
\item \hyperref[equiv-section-phantom]{Catégories dérivées de variétés}
\item \hyperref[pione-section-phantom]{Groupes fondamentaux de schémas}
\item \hyperref[etale-cohomology-section-phantom]{Cohomologie étale}
\item \hyperref[crystalline-section-phantom]{Cohomologie cristalline}
\item \hyperref[proetale-section-phantom]{Cohomologie pro-étale}
\item \hyperref[relative-cycles-section-phantom]{Cycles relatifs}
\item \hyperref[more-etale-section-phantom]{Compléments de cohomologie étale}
\item \hyperref[trace-section-phantom]{La formule des traces}
\end{enumerate}
Espaces algébriques
\begin{enumerate}
\setcounter{enumi}{64}
\item \hyperref[spaces-section-phantom]{Espaces algébriques}
\item \hyperref[spaces-properties-section-phantom]{Propriétés des espaces algébriques}
\item \hyperref[spaces-morphisms-section-phantom]{Morphismes d'espaces algébriques}
\item \hyperref[decent-spaces-section-phantom]{Espaces algébriques décents}
\item \hyperref[spaces-cohomology-section-phantom]{Cohomologie des espaces algébriques}
\item \hyperref[spaces-limits-section-phantom]{Limites d'espaces algébriques}
\item \hyperref[spaces-divisors-section-phantom]{Diviseurs sur les espaces algébriques}
\item \hyperref[spaces-over-fields-section-phantom]{Espaces algébriques sur des corps}
\item \hyperref[spaces-topologies-section-phantom]{Topologies sur les espaces algébriques}
\item \hyperref[spaces-descent-section-phantom]{Descente et espaces algébriques}
\item \hyperref[spaces-perfect-section-phantom]{Catégories dérivées d'espaces}
\item \hyperref[spaces-more-morphisms-section-phantom]{Compléments sur les morphismes d'espaces}
\item \hyperref[spaces-flat-section-phantom]{Platitude sur les espaces algébriques}
\item \hyperref[spaces-groupoids-section-phantom]{Groupoïdes dans les espaces algébriques}
\item \hyperref[spaces-more-groupoids-section-phantom]{Compléments sur les groupoïdes dans les espaces}
\item \hyperref[bootstrap-section-phantom]{Amorçage}
\item \hyperref[spaces-pushouts-section-phantom]{Sommes amalgamées d'espaces algébriques}
\end{enumerate}
Thèmes de géométrie
\begin{enumerate}
\setcounter{enumi}{81}
\item \hyperref[spaces-chow-section-phantom]{Groupes de Chow des espaces}
\item \hyperref[groupoids-quotients-section-phantom]{Quotients de groupoïdes}
\item \hyperref[spaces-more-cohomology-section-phantom]{Compléments sur la cohomologie des espaces}
\item \hyperref[spaces-simplicial-section-phantom]{Espaces simpliciaux}
\item \hyperref[spaces-duality-section-phantom]{Dualité pour les espaces}
\item \hyperref[formal-spaces-section-phantom]{Espaces algébriques formels}
\item \hyperref[restricted-section-phantom]{Algébrisation des espaces formels}
\item \hyperref[spaces-resolve-section-phantom]{Résolution des surfaces revisitée}
\end{enumerate}
Théorie des déformations
\begin{enumerate}
\setcounter{enumi}{89}
\item \hyperref[formal-defos-section-phantom]{Théorie formelle des déformations}
\item \hyperref[defos-section-phantom]{Théorie des déformations}
\item \hyperref[cotangent-section-phantom]{Le complexe cotangent}
\item \hyperref[examples-defos-section-phantom]{Problèmes de déformation}
\end{enumerate}
Champs algébriques
\begin{enumerate}
\setcounter{enumi}{93}
\item \hyperref[algebraic-section-phantom]{Champs algébriques}
\item \hyperref[examples-stacks-section-phantom]{Exemples de champs}
\item \hyperref[stacks-sheaves-section-phantom]{Faisceaux sur les champs algébriques}
\item \hyperref[criteria-section-phantom]{Critères de représentabilité}
\item \hyperref[artin-section-phantom]{Axiomes d'Artin}
\item \hyperref[quot-section-phantom]{Espaces de Quot et de Hilbert}
\item \hyperref[stacks-properties-section-phantom]{Propriétés des champs algébriques}
\item \hyperref[stacks-morphisms-section-phantom]{Morphismes de champs algébriques}
\item \hyperref[stacks-limits-section-phantom]{Limites de champs algébriques}
\item \hyperref[stacks-cohomology-section-phantom]{Cohomologie des champs algébriques}
\item \hyperref[stacks-perfect-section-phantom]{Catégories dérivées de champs}
\item \hyperref[stacks-introduction-section-phantom]{Introduction aux champs algébriques}
\item \hyperref[stacks-more-morphisms-section-phantom]{Compléments sur les morphismes de champs}
\item \hyperref[stacks-geometry-section-phantom]{La géométrie des champs}
\end{enumerate}
Thèmes de théorie des espaces de modules
\begin{enumerate}
\setcounter{enumi}{107}
\item \hyperref[moduli-section-phantom]{Champs de modules}
\item \hyperref[moduli-curves-section-phantom]{Espaces de modules de courbes}
\end{enumerate}
Divers
\begin{enumerate}
\setcounter{enumi}{109}
\item \hyperref[examples-section-phantom]{Exemples}
\item \hyperref[exercises-section-phantom]{Exercices}
\item \hyperref[guide-section-phantom]{Guide bibliographique}
\item \hyperref[desirables-section-phantom]{Desiderata}
\item \hyperref[coding-section-phantom]{Style de programmation}
\item \hyperref[obsolete-section-phantom]{Obsolète}
\item \hyperref[fdl-section-phantom]{Licence de documentation libre GNU}
\item \hyperref[index-section-phantom]{Index généré automatiquement}
\end{enumerate}
\end{multicols}


\bibliography{my}
\bibliographystyle{amsalpha}

\end{document}
